% 逐条呈现账本事件；不把共享年序坐标当作具体行动。
\section{辽、西夏与北方边地的逐年入口}
每一项均以其出处、行动者、空间线索和可追踪后果展开。材料没有说明地点、人数或地方反应时，正文保留这一限度。
\begin{textwindow}{本节材料怎样阅读}
《辽史》卷一以“太祖姓耶律氏，讳阿保机”界定本纪人物；黑水城文书保留河西的行政与宗教痕迹。前者适于校核先后，后者限定局部实践；两者都不等于整个北方边地的经验。
\end{textwindow}

\subsection{907年}
\noindent\eventanchor{LIAO-0907-EVENT-026}{耶律阿保机取得汗位，开始重组部族权力}\textbf{契丹。}
契丹 在 907年 的材料痕迹是〔耶律阿保机取得汗位，开始重组部族权力〕，出处为 《辽史》相关本纪及志；《宋史》辽国传、同年本纪；宋使行录、契丹／汉文碑刻与考古材料（据事核）。材料未将地点细化到城址，不能用中枢所在地替它补足。

前一层压力可由以下线索辨认：“耶律阿保机取得汗位，开始重组部族权力”发生在辽初部族与多都城整合的具体压力与机会之中，相关行动者的选择不能脱离此前事件链解释。 这一处置留下的后续条件是：“耶律阿保机取得汗位，开始重组部族权力”为后续政治、边防或资源配置留下条件；其社会影响与实际覆盖范围仍须以异类材料限定。

《辽史》为元末官修且契丹本方连续文书有限；宋使与碑刻的观察位置不同，战果、族称及制度覆盖范围不得互推。 因此，辽代契丹国家、都城、部族与跨境关系研究 只能扩展问题，不能代替未保存的细节。

\subsection{916年}
\noindent\eventanchor{LIAO-0916-EVENT-027}{阿保机称帝并建国号契丹}\textbf{契丹。}
从 《辽史》相关本纪及志；《宋史》辽国传、同年本纪；宋使行录、契丹／汉文碑刻与考古材料（据事核） 的 916年 条目看，〔阿保机称帝并建国号契丹〕使 契丹 的一个具体行动进入叙事；材料未将地点细化到城址，不能用中枢所在地替它补足。

辽初部族与多都城整合中前任统治的结束与宫廷、部族或官僚的权力安排相互交织，促成“阿保机称帝并建国号契丹”这一继承节点。 记录所见的结果则是：继承并不自动消除既有矛盾；“阿保机称帝并建国号契丹”后的任官、军权和对外策略需由后续条目追踪。 日期相邻不等于因果已经证实。

关于可说范围，须保留：《辽史》为元末官修且契丹本方连续文书有限；宋使与碑刻的观察位置不同，战果、族称及制度覆盖范围不得互推。 进一步讨论可参照 辽代契丹国家、都城、部族与跨境关系研究

\subsection{926年}
\noindent\eventanchor{LIAO-0926-EVENT-028}{辽军攻灭渤海，设置东丹等安排}\textbf{辽与渤海。}
926年 的记载将 辽与渤海 与〔辽军攻灭渤海，设置东丹等安排〕相连。此处可确认的是这项被记录的行动；题名保留的城、路、水道或机构名称，是目前可以落地的空间线索。

把本项放回事件链，能够看到的前提为：辽初部族与多都城整合中既有的联盟、交通与补给压力，为“辽军攻灭渤海，设置东丹等安排”提供了直接的军事或动员背景。 而它所打开或收紧的下一步是：“辽军攻灭渤海，设置东丹等安排”改变了相关城镇、道路或政权间的军事选择；伤亡、俘获和控制范围仅可按各方材料分别确认。

证据的边界是：《辽史》为元末官修且契丹本方连续文书有限；宋使与碑刻的观察位置不同，战果、族称及制度覆盖范围不得互推。 辽代契丹国家、都城、部族与跨境关系研究 可用于继续检验这一结论。

\subsection{927年}
\noindent\eventanchor{LIAO-0927-EVENT-029}{耶律德光即位为辽太宗}\textbf{辽。}
〔耶律德光即位为辽太宗〕见于 《辽史》相关本纪及志；《宋史》辽国传、同年本纪；宋使行录、契丹／汉文碑刻与考古材料（据事核） 的 927年 记录。辽 的选择因此有了可回查的时间位置；材料未将地点细化到城址，不能用中枢所在地替它补足。

它承接的条件是：辽初部族与多都城整合中前任统治的结束与宫廷、部族或官僚的权力安排相互交织，促成“耶律德光即位为辽太宗”这一继承节点。 随后可追踪的变化为：继承并不自动消除既有矛盾；“耶律德光即位为辽太宗”后的任官、军权和对外策略需由后续条目追踪。

《辽史》为元末官修且契丹本方连续文书有限；宋使与碑刻的观察位置不同，战果、族称及制度覆盖范围不得互推。 因此，辽代契丹国家、都城、部族与跨境关系研究 只能扩展问题，不能代替未保存的细节。

\subsection{936年}
\noindent\eventanchor{LIAO-0936-EVENT-030}{辽支持石敬瑭建立后晋，十六州问题进入北方政治结构}\textbf{辽与后晋。}
936年，《辽史》相关本纪及志；《宋史》辽国传、同年本纪；宋使行录、契丹／汉文碑刻与考古材料（据事核） 所见的〔辽支持石敬瑭建立后晋，十六州问题进入北方政治结构〕把 辽与后晋 的处置留在可复核的年序中；题名保留的城、路、水道或机构名称，是目前可以落地的空间线索。

前一层压力可由以下线索辨认：“辽支持石敬瑭建立后晋，十六州问题进入北方政治结构”发生在辽初部族与多都城整合的具体压力与机会之中，相关行动者的选择不能脱离此前事件链解释。 这一处置留下的后续条件是：“辽支持石敬瑭建立后晋，十六州问题进入北方政治结构”为后续政治、边防或资源配置留下条件；其社会影响与实际覆盖范围仍须以异类材料限定。

关于可说范围，须保留：《辽史》为元末官修且契丹本方连续文书有限；宋使与碑刻的观察位置不同，战果、族称及制度覆盖范围不得互推。 进一步讨论可参照 辽代契丹国家、都城、部族与跨境关系研究

\subsection{938年}
\noindent\eventanchor{LIAO-0938-EVENT-031}{辽以幽州为南京，经营燕云城市与交通节点}\textbf{辽。}
辽 在 938年 的材料痕迹是〔辽以幽州为南京，经营燕云城市与交通节点〕，出处为 《辽史》相关本纪及志；《宋史》辽国传、同年本纪；宋使行录、契丹／汉文碑刻与考古材料（据事核）。题名保留的城、路、水道或机构名称，是目前可以落地的空间线索。

“辽以幽州为南京，经营燕云城市与交通节点”发生在辽初部族与多都城整合的具体压力与机会之中，相关行动者的选择不能脱离此前事件链解释。 记录所见的结果则是：“辽以幽州为南京，经营燕云城市与交通节点”为后续政治、边防或资源配置留下条件；其社会影响与实际覆盖范围仍须以异类材料限定。 日期相邻不等于因果已经证实。

证据的边界是：《辽史》为元末官修且契丹本方连续文书有限；宋使与碑刻的观察位置不同，战果、族称及制度覆盖范围不得互推。 辽代契丹国家、都城、部族与跨境关系研究 可用于继续检验这一结论。

\subsection{947年}
\noindent\eventanchor{LIAO-0947-EVENT-032}{辽军进入开封，后因补给与抵抗压力撤离}\textbf{辽与后晋。}
从 《辽史》相关本纪及志；《宋史》辽国传、同年本纪；宋使行录、契丹／汉文碑刻与考古材料（据事核） 的 947年 条目看，〔辽军进入开封，后因补给与抵抗压力撤离〕使 辽与后晋 的一个具体行动进入叙事；材料未将地点细化到城址，不能用中枢所在地替它补足。

把本项放回事件链，能够看到的前提为：“辽军进入开封，后因补给与抵抗压力撤离”发生在辽初部族与多都城整合的具体压力与机会之中，相关行动者的选择不能脱离此前事件链解释。 而它所打开或收紧的下一步是：“辽军进入开封，后因补给与抵抗压力撤离”为后续政治、边防或资源配置留下条件；其社会影响与实际覆盖范围仍须以异类材料限定。

《辽史》为元末官修且契丹本方连续文书有限；宋使与碑刻的观察位置不同，战果、族称及制度覆盖范围不得互推。 因此，辽代契丹国家、都城、部族与跨境关系研究 只能扩展问题，不能代替未保存的细节。

\subsection{969年}
\noindent\eventanchor{LIAO-0969-EVENT-033}{辽景宗即位，萧绰逐步取得政治作用}\textbf{辽。}
969年 的记载将 辽 与〔辽景宗即位，萧绰逐步取得政治作用〕相连。此处可确认的是这项被记录的行动；材料未将地点细化到城址，不能用中枢所在地替它补足。

它承接的条件是：辽初部族与多都城整合中前任统治的结束与宫廷、部族或官僚的权力安排相互交织，促成“辽景宗即位，萧绰逐步取得政治作用”这一继承节点。 随后可追踪的变化为：继承并不自动消除既有矛盾；“辽景宗即位，萧绰逐步取得政治作用”后的任官、军权和对外策略需由后续条目追踪。

关于可说范围，须保留：《辽史》为元末官修且契丹本方连续文书有限；宋使与碑刻的观察位置不同，战果、族称及制度覆盖范围不得互推。 进一步讨论可参照 辽代契丹国家、都城、部族与跨境关系研究

\subsection{979年}
\noindent\eventanchor{LIAO-0979-EVENT-034}{辽军在高梁河击退北伐宋军}\textbf{辽与宋。}
〔辽军在高梁河击退北伐宋军〕见于 《辽史》相关本纪及志；《宋史》辽国传、同年本纪；宋使行录、契丹／汉文碑刻与考古材料（据事核） 的 979年 记录。辽与宋 的选择因此有了可回查的时间位置；题名保留的城、路、水道或机构名称，是目前可以落地的空间线索。

前一层压力可由以下线索辨认：辽初部族与多都城整合中既有的联盟、交通与补给压力，为“辽军在高梁河击退北伐宋军”提供了直接的军事或动员背景。 这一处置留下的后续条件是：“辽军在高梁河击退北伐宋军”改变了相关城镇、道路或政权间的军事选择；伤亡、俘获和控制范围仅可按各方材料分别确认。

证据的边界是：《辽史》为元末官修且契丹本方连续文书有限；宋使与碑刻的观察位置不同，战果、族称及制度覆盖范围不得互推。 辽代契丹国家、都城、部族与跨境关系研究 可用于继续检验这一结论。

\subsection{982年}
\noindent\eventanchor{LIAO-0982-EVENT-035}{辽圣宗即位，萧太后摄政}\textbf{辽。}
982年，《辽史》相关本纪及志；《宋史》辽国传、同年本纪；宋使行录、契丹／汉文碑刻与考古材料（据事核） 所见的〔辽圣宗即位，萧太后摄政〕把 辽 的处置留在可复核的年序中；材料未将地点细化到城址，不能用中枢所在地替它补足。

辽初部族与多都城整合中前任统治的结束与宫廷、部族或官僚的权力安排相互交织，促成“辽圣宗即位，萧太后摄政”这一继承节点。 记录所见的结果则是：继承并不自动消除既有矛盾；“辽圣宗即位，萧太后摄政”后的任官、军权和对外策略需由后续条目追踪。 日期相邻不等于因果已经证实。

《辽史》为元末官修且契丹本方连续文书有限；宋使与碑刻的观察位置不同，战果、族称及制度覆盖范围不得互推。 因此，辽代契丹国家、都城、部族与跨境关系研究 只能扩展问题，不能代替未保存的细节。

\subsection{986年}
\noindent\eventanchor{LIAO-0986-EVENT-036}{辽在雍熙北伐中应对宋军多路进攻并取胜}\textbf{辽与宋。}
辽与宋 在 986年 的材料痕迹是〔辽在雍熙北伐中应对宋军多路进攻并取胜〕，出处为 《辽史》相关本纪及志；《宋史》辽国传、同年本纪；宋使行录、契丹／汉文碑刻与考古材料（据事核）。题名保留的城、路、水道或机构名称，是目前可以落地的空间线索。

把本项放回事件链，能够看到的前提为：辽初部族与多都城整合中既有的联盟、交通与补给压力，为“辽在雍熙北伐中应对宋军多路进攻并取胜”提供了直接的军事或动员背景。 而它所打开或收紧的下一步是：“辽在雍熙北伐中应对宋军多路进攻并取胜”改变了相关城镇、道路或政权间的军事选择；伤亡、俘获和控制范围仅可按各方材料分别确认。

关于可说范围，须保留：《辽史》为元末官修且契丹本方连续文书有限；宋使与碑刻的观察位置不同，战果、族称及制度覆盖范围不得互推。 进一步讨论可参照 辽代契丹国家、都城、部族与跨境关系研究

\subsection{960 年（共享年序桥接）}
\noindent\eventanchor{SY-960-CHRONOLOGY-BRIDGE}{【C级年序桥接】保留960 年的多政权并行检索入口；本条不主张所列政权在当年发生直接互动，具体事件须另以单项材料入账。}\textbf{宋、辽与北汉／高丽（年序桥接）。}
从 《宋史》；《辽史》；《续资治通鉴长编》相关年条（据事核） 的 960 年（共享年序桥接） 条目看，〔【C级年序桥接】保留960 年的多政权并行检索入口；本条不主张所列政权在当年发生直接互动，具体事件须另以单项材料入账。〕使 宋、辽与北汉／高丽（年序桥接） 的一个具体行动进入叙事；题名保留的城、路、水道或机构名称，是目前可以落地的空间线索。

它承接的条件是：相邻已入账事件之间缺少该年度的共享年序入口，易使跨政权材料被单一政权叙事遮蔽。 随后可追踪的变化为：保留该年度的检索与补账位置；后续仅在获得可核单项材料时拆分为具体政治、财政、军事、社会或文化事件。

证据的边界是：本条只确认该年位于可追踪的共享年序中；正史、地方文书和外部材料的保存密度不一，不能由桥接标签推断行动、统属、疆界或社会影响。 北宋、辽与东北亚并行年序及边境网络研究 可用于继续检验这一结论。

\subsection{961 年（共享年序桥接）}
\noindent\eventanchor{SY-961-CHRONOLOGY-BRIDGE}{【C级年序桥接】保留961 年的多政权并行检索入口；本条不主张所列政权在当年发生直接互动，具体事件须另以单项材料入账。}\textbf{宋、辽与北汉／高丽（年序桥接）。}
961 年（共享年序桥接） 的记载将 宋、辽与北汉／高丽（年序桥接） 与〔【C级年序桥接】保留961 年的多政权并行检索入口；本条不主张所列政权在当年发生直接互动，具体事件须另以单项材料入账。〕相连。此处可确认的是这项被记录的行动；题名保留的城、路、水道或机构名称，是目前可以落地的空间线索。

前一层压力可由以下线索辨认：相邻已入账事件之间缺少该年度的共享年序入口，易使跨政权材料被单一政权叙事遮蔽。 这一处置留下的后续条件是：保留该年度的检索与补账位置；后续仅在获得可核单项材料时拆分为具体政治、财政、军事、社会或文化事件。

本条只确认该年位于可追踪的共享年序中；正史、地方文书和外部材料的保存密度不一，不能由桥接标签推断行动、统属、疆界或社会影响。 因此，北宋、辽与东北亚并行年序及边境网络研究 只能扩展问题，不能代替未保存的细节。

\subsection{962 年（共享年序桥接）}
\noindent\eventanchor{SY-962-CHRONOLOGY-BRIDGE}{【C级年序桥接】保留962 年的多政权并行检索入口；本条不主张所列政权在当年发生直接互动，具体事件须另以单项材料入账。}\textbf{宋、辽与北汉／高丽（年序桥接）。}
〔【C级年序桥接】保留962 年的多政权并行检索入口；本条不主张所列政权在当年发生直接互动，具体事件须另以单项材料入账。〕见于 《宋史》；《辽史》；《续资治通鉴长编》相关年条（据事核） 的 962 年（共享年序桥接） 记录。宋、辽与北汉／高丽（年序桥接） 的选择因此有了可回查的时间位置；题名保留的城、路、水道或机构名称，是目前可以落地的空间线索。

相邻已入账事件之间缺少该年度的共享年序入口，易使跨政权材料被单一政权叙事遮蔽。 记录所见的结果则是：保留该年度的检索与补账位置；后续仅在获得可核单项材料时拆分为具体政治、财政、军事、社会或文化事件。 日期相邻不等于因果已经证实。

关于可说范围，须保留：本条只确认该年位于可追踪的共享年序中；正史、地方文书和外部材料的保存密度不一，不能由桥接标签推断行动、统属、疆界或社会影响。 进一步讨论可参照 北宋、辽与东北亚并行年序及边境网络研究

\subsection{963 年（共享年序桥接）}
\noindent\eventanchor{SY-963-CHRONOLOGY-BRIDGE}{【C级年序桥接】保留963 年的多政权并行检索入口；本条不主张所列政权在当年发生直接互动，具体事件须另以单项材料入账。}\textbf{宋、辽与北汉／高丽（年序桥接）。}
963 年（共享年序桥接），《宋史》；《辽史》；《续资治通鉴长编》相关年条（据事核） 所见的〔【C级年序桥接】保留963 年的多政权并行检索入口；本条不主张所列政权在当年发生直接互动，具体事件须另以单项材料入账。〕把 宋、辽与北汉／高丽（年序桥接） 的处置留在可复核的年序中；题名保留的城、路、水道或机构名称，是目前可以落地的空间线索。

把本项放回事件链，能够看到的前提为：相邻已入账事件之间缺少该年度的共享年序入口，易使跨政权材料被单一政权叙事遮蔽。 而它所打开或收紧的下一步是：保留该年度的检索与补账位置；后续仅在获得可核单项材料时拆分为具体政治、财政、军事、社会或文化事件。

证据的边界是：本条只确认该年位于可追踪的共享年序中；正史、地方文书和外部材料的保存密度不一，不能由桥接标签推断行动、统属、疆界或社会影响。 北宋、辽与东北亚并行年序及边境网络研究 可用于继续检验这一结论。

\subsection{964 年（共享年序桥接）}
\noindent\eventanchor{SY-964-CHRONOLOGY-BRIDGE}{【C级年序桥接】保留964 年的多政权并行检索入口；本条不主张所列政权在当年发生直接互动，具体事件须另以单项材料入账。}\textbf{宋、辽与北汉／高丽（年序桥接）。}
宋、辽与北汉／高丽（年序桥接） 在 964 年（共享年序桥接） 的材料痕迹是〔【C级年序桥接】保留964 年的多政权并行检索入口；本条不主张所列政权在当年发生直接互动，具体事件须另以单项材料入账。〕，出处为 《宋史》；《辽史》；《续资治通鉴长编》相关年条（据事核）。题名保留的城、路、水道或机构名称，是目前可以落地的空间线索。

它承接的条件是：相邻已入账事件之间缺少该年度的共享年序入口，易使跨政权材料被单一政权叙事遮蔽。 随后可追踪的变化为：保留该年度的检索与补账位置；后续仅在获得可核单项材料时拆分为具体政治、财政、军事、社会或文化事件。

本条只确认该年位于可追踪的共享年序中；正史、地方文书和外部材料的保存密度不一，不能由桥接标签推断行动、统属、疆界或社会影响。 因此，北宋、辽与东北亚并行年序及边境网络研究 只能扩展问题，不能代替未保存的细节。

\subsection{965 年（共享年序桥接）}
\noindent\eventanchor{SY-965-CHRONOLOGY-BRIDGE}{【C级年序桥接】保留965 年的多政权并行检索入口；本条不主张所列政权在当年发生直接互动，具体事件须另以单项材料入账。}\textbf{宋、辽与北汉／高丽（年序桥接）。}
从 《宋史》；《辽史》；《续资治通鉴长编》相关年条（据事核） 的 965 年（共享年序桥接） 条目看，〔【C级年序桥接】保留965 年的多政权并行检索入口；本条不主张所列政权在当年发生直接互动，具体事件须另以单项材料入账。〕使 宋、辽与北汉／高丽（年序桥接） 的一个具体行动进入叙事；题名保留的城、路、水道或机构名称，是目前可以落地的空间线索。

前一层压力可由以下线索辨认：相邻已入账事件之间缺少该年度的共享年序入口，易使跨政权材料被单一政权叙事遮蔽。 这一处置留下的后续条件是：保留该年度的检索与补账位置；后续仅在获得可核单项材料时拆分为具体政治、财政、军事、社会或文化事件。

关于可说范围，须保留：本条只确认该年位于可追踪的共享年序中；正史、地方文书和外部材料的保存密度不一，不能由桥接标签推断行动、统属、疆界或社会影响。 进一步讨论可参照 北宋、辽与东北亚并行年序及边境网络研究

\subsection{966 年（共享年序桥接）}
\noindent\eventanchor{SY-966-CHRONOLOGY-BRIDGE}{【C级年序桥接】保留966 年的多政权并行检索入口；本条不主张所列政权在当年发生直接互动，具体事件须另以单项材料入账。}\textbf{宋、辽与北汉／高丽（年序桥接）。}
966 年（共享年序桥接） 的记载将 宋、辽与北汉／高丽（年序桥接） 与〔【C级年序桥接】保留966 年的多政权并行检索入口；本条不主张所列政权在当年发生直接互动，具体事件须另以单项材料入账。〕相连。此处可确认的是这项被记录的行动；题名保留的城、路、水道或机构名称，是目前可以落地的空间线索。

相邻已入账事件之间缺少该年度的共享年序入口，易使跨政权材料被单一政权叙事遮蔽。 记录所见的结果则是：保留该年度的检索与补账位置；后续仅在获得可核单项材料时拆分为具体政治、财政、军事、社会或文化事件。 日期相邻不等于因果已经证实。

证据的边界是：本条只确认该年位于可追踪的共享年序中；正史、地方文书和外部材料的保存密度不一，不能由桥接标签推断行动、统属、疆界或社会影响。 北宋、辽与东北亚并行年序及边境网络研究 可用于继续检验这一结论。

\subsection{967 年（共享年序桥接）}
\noindent\eventanchor{SY-967-CHRONOLOGY-BRIDGE}{【C级年序桥接】保留967 年的多政权并行检索入口；本条不主张所列政权在当年发生直接互动，具体事件须另以单项材料入账。}\textbf{宋、辽与北汉／高丽（年序桥接）。}
〔【C级年序桥接】保留967 年的多政权并行检索入口；本条不主张所列政权在当年发生直接互动，具体事件须另以单项材料入账。〕见于 《宋史》；《辽史》；《续资治通鉴长编》相关年条（据事核） 的 967 年（共享年序桥接） 记录。宋、辽与北汉／高丽（年序桥接） 的选择因此有了可回查的时间位置；题名保留的城、路、水道或机构名称，是目前可以落地的空间线索。

把本项放回事件链，能够看到的前提为：相邻已入账事件之间缺少该年度的共享年序入口，易使跨政权材料被单一政权叙事遮蔽。 而它所打开或收紧的下一步是：保留该年度的检索与补账位置；后续仅在获得可核单项材料时拆分为具体政治、财政、军事、社会或文化事件。

本条只确认该年位于可追踪的共享年序中；正史、地方文书和外部材料的保存密度不一，不能由桥接标签推断行动、统属、疆界或社会影响。 因此，北宋、辽与东北亚并行年序及边境网络研究 只能扩展问题，不能代替未保存的细节。

\subsection{968 年（共享年序桥接）}
\noindent\eventanchor{SY-968-CHRONOLOGY-BRIDGE}{【C级年序桥接】保留968 年的多政权并行检索入口；本条不主张所列政权在当年发生直接互动，具体事件须另以单项材料入账。}\textbf{宋、辽与北汉／高丽（年序桥接）。}
968 年（共享年序桥接），《宋史》；《辽史》；《续资治通鉴长编》相关年条（据事核） 所见的〔【C级年序桥接】保留968 年的多政权并行检索入口；本条不主张所列政权在当年发生直接互动，具体事件须另以单项材料入账。〕把 宋、辽与北汉／高丽（年序桥接） 的处置留在可复核的年序中；题名保留的城、路、水道或机构名称，是目前可以落地的空间线索。

它承接的条件是：相邻已入账事件之间缺少该年度的共享年序入口，易使跨政权材料被单一政权叙事遮蔽。 随后可追踪的变化为：保留该年度的检索与补账位置；后续仅在获得可核单项材料时拆分为具体政治、财政、军事、社会或文化事件。

关于可说范围，须保留：本条只确认该年位于可追踪的共享年序中；正史、地方文书和外部材料的保存密度不一，不能由桥接标签推断行动、统属、疆界或社会影响。 进一步讨论可参照 北宋、辽与东北亚并行年序及边境网络研究

\subsection{969 年（共享年序桥接）}
\noindent\eventanchor{SY-969-CHRONOLOGY-BRIDGE}{【C级年序桥接】保留969 年的多政权并行检索入口；本条不主张所列政权在当年发生直接互动，具体事件须另以单项材料入账。}\textbf{宋、辽与北汉／高丽（年序桥接）。}
宋、辽与北汉／高丽（年序桥接） 在 969 年（共享年序桥接） 的材料痕迹是〔【C级年序桥接】保留969 年的多政权并行检索入口；本条不主张所列政权在当年发生直接互动，具体事件须另以单项材料入账。〕，出处为 《宋史》；《辽史》；《续资治通鉴长编》相关年条（据事核）。题名保留的城、路、水道或机构名称，是目前可以落地的空间线索。

前一层压力可由以下线索辨认：相邻已入账事件之间缺少该年度的共享年序入口，易使跨政权材料被单一政权叙事遮蔽。 这一处置留下的后续条件是：保留该年度的检索与补账位置；后续仅在获得可核单项材料时拆分为具体政治、财政、军事、社会或文化事件。

证据的边界是：本条只确认该年位于可追踪的共享年序中；正史、地方文书和外部材料的保存密度不一，不能由桥接标签推断行动、统属、疆界或社会影响。 北宋、辽与东北亚并行年序及边境网络研究 可用于继续检验这一结论。

\subsection{970 年（共享年序桥接）}
\noindent\eventanchor{SY-970-CHRONOLOGY-BRIDGE}{【C级年序桥接】保留970 年的多政权并行检索入口；本条不主张所列政权在当年发生直接互动，具体事件须另以单项材料入账。}\textbf{宋、辽与北汉／高丽（年序桥接）。}
从 《宋史》；《辽史》；《续资治通鉴长编》相关年条（据事核） 的 970 年（共享年序桥接） 条目看，〔【C级年序桥接】保留970 年的多政权并行检索入口；本条不主张所列政权在当年发生直接互动，具体事件须另以单项材料入账。〕使 宋、辽与北汉／高丽（年序桥接） 的一个具体行动进入叙事；题名保留的城、路、水道或机构名称，是目前可以落地的空间线索。

相邻已入账事件之间缺少该年度的共享年序入口，易使跨政权材料被单一政权叙事遮蔽。 记录所见的结果则是：保留该年度的检索与补账位置；后续仅在获得可核单项材料时拆分为具体政治、财政、军事、社会或文化事件。 日期相邻不等于因果已经证实。

本条只确认该年位于可追踪的共享年序中；正史、地方文书和外部材料的保存密度不一，不能由桥接标签推断行动、统属、疆界或社会影响。 因此，北宋、辽与东北亚并行年序及边境网络研究 只能扩展问题，不能代替未保存的细节。

\subsection{971 年（共享年序桥接）}
\noindent\eventanchor{SY-971-CHRONOLOGY-BRIDGE}{【C级年序桥接】保留971 年的多政权并行检索入口；本条不主张所列政权在当年发生直接互动，具体事件须另以单项材料入账。}\textbf{宋、辽与北汉／高丽（年序桥接）。}
971 年（共享年序桥接） 的记载将 宋、辽与北汉／高丽（年序桥接） 与〔【C级年序桥接】保留971 年的多政权并行检索入口；本条不主张所列政权在当年发生直接互动，具体事件须另以单项材料入账。〕相连。此处可确认的是这项被记录的行动；题名保留的城、路、水道或机构名称，是目前可以落地的空间线索。

把本项放回事件链，能够看到的前提为：相邻已入账事件之间缺少该年度的共享年序入口，易使跨政权材料被单一政权叙事遮蔽。 而它所打开或收紧的下一步是：保留该年度的检索与补账位置；后续仅在获得可核单项材料时拆分为具体政治、财政、军事、社会或文化事件。

关于可说范围，须保留：本条只确认该年位于可追踪的共享年序中；正史、地方文书和外部材料的保存密度不一，不能由桥接标签推断行动、统属、疆界或社会影响。 进一步讨论可参照 北宋、辽与东北亚并行年序及边境网络研究

\subsection{972 年（共享年序桥接）}
\noindent\eventanchor{SY-972-CHRONOLOGY-BRIDGE}{【C级年序桥接】保留972 年的多政权并行检索入口；本条不主张所列政权在当年发生直接互动，具体事件须另以单项材料入账。}\textbf{宋、辽与北汉／高丽（年序桥接）。}
〔【C级年序桥接】保留972 年的多政权并行检索入口；本条不主张所列政权在当年发生直接互动，具体事件须另以单项材料入账。〕见于 《宋史》；《辽史》；《续资治通鉴长编》相关年条（据事核） 的 972 年（共享年序桥接） 记录。宋、辽与北汉／高丽（年序桥接） 的选择因此有了可回查的时间位置；题名保留的城、路、水道或机构名称，是目前可以落地的空间线索。

它承接的条件是：相邻已入账事件之间缺少该年度的共享年序入口，易使跨政权材料被单一政权叙事遮蔽。 随后可追踪的变化为：保留该年度的检索与补账位置；后续仅在获得可核单项材料时拆分为具体政治、财政、军事、社会或文化事件。

证据的边界是：本条只确认该年位于可追踪的共享年序中；正史、地方文书和外部材料的保存密度不一，不能由桥接标签推断行动、统属、疆界或社会影响。 北宋、辽与东北亚并行年序及边境网络研究 可用于继续检验这一结论。

\subsection{973 年（共享年序桥接）}
\noindent\eventanchor{SY-973-CHRONOLOGY-BRIDGE}{【C级年序桥接】保留973 年的多政权并行检索入口；本条不主张所列政权在当年发生直接互动，具体事件须另以单项材料入账。}\textbf{宋、辽与北汉／高丽（年序桥接）。}
973 年（共享年序桥接），《宋史》；《辽史》；《续资治通鉴长编》相关年条（据事核） 所见的〔【C级年序桥接】保留973 年的多政权并行检索入口；本条不主张所列政权在当年发生直接互动，具体事件须另以单项材料入账。〕把 宋、辽与北汉／高丽（年序桥接） 的处置留在可复核的年序中；题名保留的城、路、水道或机构名称，是目前可以落地的空间线索。

前一层压力可由以下线索辨认：相邻已入账事件之间缺少该年度的共享年序入口，易使跨政权材料被单一政权叙事遮蔽。 这一处置留下的后续条件是：保留该年度的检索与补账位置；后续仅在获得可核单项材料时拆分为具体政治、财政、军事、社会或文化事件。

本条只确认该年位于可追踪的共享年序中；正史、地方文书和外部材料的保存密度不一，不能由桥接标签推断行动、统属、疆界或社会影响。 因此，北宋、辽与东北亚并行年序及边境网络研究 只能扩展问题，不能代替未保存的细节。

\subsection{974 年（共享年序桥接）}
\noindent\eventanchor{SY-974-CHRONOLOGY-BRIDGE}{【C级年序桥接】保留974 年的多政权并行检索入口；本条不主张所列政权在当年发生直接互动，具体事件须另以单项材料入账。}\textbf{宋、辽与北汉／高丽（年序桥接）。}
宋、辽与北汉／高丽（年序桥接） 在 974 年（共享年序桥接） 的材料痕迹是〔【C级年序桥接】保留974 年的多政权并行检索入口；本条不主张所列政权在当年发生直接互动，具体事件须另以单项材料入账。〕，出处为 《宋史》；《辽史》；《续资治通鉴长编》相关年条（据事核）。题名保留的城、路、水道或机构名称，是目前可以落地的空间线索。

相邻已入账事件之间缺少该年度的共享年序入口，易使跨政权材料被单一政权叙事遮蔽。 记录所见的结果则是：保留该年度的检索与补账位置；后续仅在获得可核单项材料时拆分为具体政治、财政、军事、社会或文化事件。 日期相邻不等于因果已经证实。

关于可说范围，须保留：本条只确认该年位于可追踪的共享年序中；正史、地方文书和外部材料的保存密度不一，不能由桥接标签推断行动、统属、疆界或社会影响。 进一步讨论可参照 北宋、辽与东北亚并行年序及边境网络研究

\subsection{975 年（共享年序桥接）}
\noindent\eventanchor{SY-975-CHRONOLOGY-BRIDGE}{【C级年序桥接】保留975 年的多政权并行检索入口；本条不主张所列政权在当年发生直接互动，具体事件须另以单项材料入账。}\textbf{宋、辽与北汉／高丽（年序桥接）。}
从 《宋史》；《辽史》；《续资治通鉴长编》相关年条（据事核） 的 975 年（共享年序桥接） 条目看，〔【C级年序桥接】保留975 年的多政权并行检索入口；本条不主张所列政权在当年发生直接互动，具体事件须另以单项材料入账。〕使 宋、辽与北汉／高丽（年序桥接） 的一个具体行动进入叙事；题名保留的城、路、水道或机构名称，是目前可以落地的空间线索。

把本项放回事件链，能够看到的前提为：相邻已入账事件之间缺少该年度的共享年序入口，易使跨政权材料被单一政权叙事遮蔽。 而它所打开或收紧的下一步是：保留该年度的检索与补账位置；后续仅在获得可核单项材料时拆分为具体政治、财政、军事、社会或文化事件。

证据的边界是：本条只确认该年位于可追踪的共享年序中；正史、地方文书和外部材料的保存密度不一，不能由桥接标签推断行动、统属、疆界或社会影响。 北宋、辽与东北亚并行年序及边境网络研究 可用于继续检验这一结论。

\subsection{976 年（共享年序桥接）}
\noindent\eventanchor{SY-976-CHRONOLOGY-BRIDGE}{【C级年序桥接】保留976 年的多政权并行检索入口；本条不主张所列政权在当年发生直接互动，具体事件须另以单项材料入账。}\textbf{宋、辽与北汉／高丽（年序桥接）。}
976 年（共享年序桥接） 的记载将 宋、辽与北汉／高丽（年序桥接） 与〔【C级年序桥接】保留976 年的多政权并行检索入口；本条不主张所列政权在当年发生直接互动，具体事件须另以单项材料入账。〕相连。此处可确认的是这项被记录的行动；题名保留的城、路、水道或机构名称，是目前可以落地的空间线索。

它承接的条件是：相邻已入账事件之间缺少该年度的共享年序入口，易使跨政权材料被单一政权叙事遮蔽。 随后可追踪的变化为：保留该年度的检索与补账位置；后续仅在获得可核单项材料时拆分为具体政治、财政、军事、社会或文化事件。

本条只确认该年位于可追踪的共享年序中；正史、地方文书和外部材料的保存密度不一，不能由桥接标签推断行动、统属、疆界或社会影响。 因此，北宋、辽与东北亚并行年序及边境网络研究 只能扩展问题，不能代替未保存的细节。

\subsection{977 年（共享年序桥接）}
\noindent\eventanchor{SY-977-CHRONOLOGY-BRIDGE}{【C级年序桥接】保留977 年的多政权并行检索入口；本条不主张所列政权在当年发生直接互动，具体事件须另以单项材料入账。}\textbf{宋、辽与北汉／高丽（年序桥接）。}
〔【C级年序桥接】保留977 年的多政权并行检索入口；本条不主张所列政权在当年发生直接互动，具体事件须另以单项材料入账。〕见于 《宋史》；《辽史》；《续资治通鉴长编》相关年条（据事核） 的 977 年（共享年序桥接） 记录。宋、辽与北汉／高丽（年序桥接） 的选择因此有了可回查的时间位置；题名保留的城、路、水道或机构名称，是目前可以落地的空间线索。

前一层压力可由以下线索辨认：相邻已入账事件之间缺少该年度的共享年序入口，易使跨政权材料被单一政权叙事遮蔽。 这一处置留下的后续条件是：保留该年度的检索与补账位置；后续仅在获得可核单项材料时拆分为具体政治、财政、军事、社会或文化事件。

关于可说范围，须保留：本条只确认该年位于可追踪的共享年序中；正史、地方文书和外部材料的保存密度不一，不能由桥接标签推断行动、统属、疆界或社会影响。 进一步讨论可参照 北宋、辽与东北亚并行年序及边境网络研究

\subsection{978 年（共享年序桥接）}
\noindent\eventanchor{SY-978-CHRONOLOGY-BRIDGE}{【C级年序桥接】保留978 年的多政权并行检索入口；本条不主张所列政权在当年发生直接互动，具体事件须另以单项材料入账。}\textbf{宋、辽与北汉／高丽（年序桥接）。}
978 年（共享年序桥接），《宋史》；《辽史》；《续资治通鉴长编》相关年条（据事核） 所见的〔【C级年序桥接】保留978 年的多政权并行检索入口；本条不主张所列政权在当年发生直接互动，具体事件须另以单项材料入账。〕把 宋、辽与北汉／高丽（年序桥接） 的处置留在可复核的年序中；题名保留的城、路、水道或机构名称，是目前可以落地的空间线索。

相邻已入账事件之间缺少该年度的共享年序入口，易使跨政权材料被单一政权叙事遮蔽。 记录所见的结果则是：保留该年度的检索与补账位置；后续仅在获得可核单项材料时拆分为具体政治、财政、军事、社会或文化事件。 日期相邻不等于因果已经证实。

证据的边界是：本条只确认该年位于可追踪的共享年序中；正史、地方文书和外部材料的保存密度不一，不能由桥接标签推断行动、统属、疆界或社会影响。 北宋、辽与东北亚并行年序及边境网络研究 可用于继续检验这一结论。

\subsection{979 年（共享年序桥接）}
\noindent\eventanchor{SY-979-CHRONOLOGY-BRIDGE}{【C级年序桥接】保留979 年的多政权并行检索入口；本条不主张所列政权在当年发生直接互动，具体事件须另以单项材料入账。}\textbf{宋、辽与北汉／高丽（年序桥接）。}
宋、辽与北汉／高丽（年序桥接） 在 979 年（共享年序桥接） 的材料痕迹是〔【C级年序桥接】保留979 年的多政权并行检索入口；本条不主张所列政权在当年发生直接互动，具体事件须另以单项材料入账。〕，出处为 《宋史》；《辽史》；《续资治通鉴长编》相关年条（据事核）。题名保留的城、路、水道或机构名称，是目前可以落地的空间线索。

把本项放回事件链，能够看到的前提为：相邻已入账事件之间缺少该年度的共享年序入口，易使跨政权材料被单一政权叙事遮蔽。 而它所打开或收紧的下一步是：保留该年度的检索与补账位置；后续仅在获得可核单项材料时拆分为具体政治、财政、军事、社会或文化事件。

本条只确认该年位于可追踪的共享年序中；正史、地方文书和外部材料的保存密度不一，不能由桥接标签推断行动、统属、疆界或社会影响。 因此，北宋、辽与东北亚并行年序及边境网络研究 只能扩展问题，不能代替未保存的细节。

\subsection{980 年（共享年序桥接）}
\noindent\eventanchor{SY-980-CHRONOLOGY-BRIDGE}{【C级年序桥接】保留980 年的多政权并行检索入口；本条不主张所列政权在当年发生直接互动，具体事件须另以单项材料入账。}\textbf{宋、辽与北汉／高丽（年序桥接）。}
从 《宋史》；《辽史》；《续资治通鉴长编》相关年条（据事核） 的 980 年（共享年序桥接） 条目看，〔【C级年序桥接】保留980 年的多政权并行检索入口；本条不主张所列政权在当年发生直接互动，具体事件须另以单项材料入账。〕使 宋、辽与北汉／高丽（年序桥接） 的一个具体行动进入叙事；题名保留的城、路、水道或机构名称，是目前可以落地的空间线索。

它承接的条件是：相邻已入账事件之间缺少该年度的共享年序入口，易使跨政权材料被单一政权叙事遮蔽。 随后可追踪的变化为：保留该年度的检索与补账位置；后续仅在获得可核单项材料时拆分为具体政治、财政、军事、社会或文化事件。

关于可说范围，须保留：本条只确认该年位于可追踪的共享年序中；正史、地方文书和外部材料的保存密度不一，不能由桥接标签推断行动、统属、疆界或社会影响。 进一步讨论可参照 北宋、辽与东北亚并行年序及边境网络研究

\subsection{981 年（共享年序桥接）}
\noindent\eventanchor{SY-981-CHRONOLOGY-BRIDGE}{【C级年序桥接】保留981 年的多政权并行检索入口；本条不主张所列政权在当年发生直接互动，具体事件须另以单项材料入账。}\textbf{宋、辽与北汉／高丽（年序桥接）。}
981 年（共享年序桥接） 的记载将 宋、辽与北汉／高丽（年序桥接） 与〔【C级年序桥接】保留981 年的多政权并行检索入口；本条不主张所列政权在当年发生直接互动，具体事件须另以单项材料入账。〕相连。此处可确认的是这项被记录的行动；题名保留的城、路、水道或机构名称，是目前可以落地的空间线索。

前一层压力可由以下线索辨认：相邻已入账事件之间缺少该年度的共享年序入口，易使跨政权材料被单一政权叙事遮蔽。 这一处置留下的后续条件是：保留该年度的检索与补账位置；后续仅在获得可核单项材料时拆分为具体政治、财政、军事、社会或文化事件。

证据的边界是：本条只确认该年位于可追踪的共享年序中；正史、地方文书和外部材料的保存密度不一，不能由桥接标签推断行动、统属、疆界或社会影响。 北宋、辽与东北亚并行年序及边境网络研究 可用于继续检验这一结论。

\subsection{982 年（共享年序桥接）}
\noindent\eventanchor{SY-982-CHRONOLOGY-BRIDGE}{【C级年序桥接】保留982 年的多政权并行检索入口；本条不主张所列政权在当年发生直接互动，具体事件须另以单项材料入账。}\textbf{宋、辽与北汉／高丽（年序桥接）。}
〔【C级年序桥接】保留982 年的多政权并行检索入口；本条不主张所列政权在当年发生直接互动，具体事件须另以单项材料入账。〕见于 《宋史》；《辽史》；《续资治通鉴长编》相关年条（据事核） 的 982 年（共享年序桥接） 记录。宋、辽与北汉／高丽（年序桥接） 的选择因此有了可回查的时间位置；题名保留的城、路、水道或机构名称，是目前可以落地的空间线索。

相邻已入账事件之间缺少该年度的共享年序入口，易使跨政权材料被单一政权叙事遮蔽。 记录所见的结果则是：保留该年度的检索与补账位置；后续仅在获得可核单项材料时拆分为具体政治、财政、军事、社会或文化事件。 日期相邻不等于因果已经证实。

本条只确认该年位于可追踪的共享年序中；正史、地方文书和外部材料的保存密度不一，不能由桥接标签推断行动、统属、疆界或社会影响。 因此，北宋、辽与东北亚并行年序及边境网络研究 只能扩展问题，不能代替未保存的细节。

\subsection{983 年（共享年序桥接）}
\noindent\eventanchor{SY-983-CHRONOLOGY-BRIDGE}{【C级年序桥接】保留983 年的多政权并行检索入口；本条不主张所列政权在当年发生直接互动，具体事件须另以单项材料入账。}\textbf{宋、辽与北汉／高丽（年序桥接）。}
983 年（共享年序桥接），《宋史》；《辽史》；《续资治通鉴长编》相关年条（据事核） 所见的〔【C级年序桥接】保留983 年的多政权并行检索入口；本条不主张所列政权在当年发生直接互动，具体事件须另以单项材料入账。〕把 宋、辽与北汉／高丽（年序桥接） 的处置留在可复核的年序中；题名保留的城、路、水道或机构名称，是目前可以落地的空间线索。

把本项放回事件链，能够看到的前提为：相邻已入账事件之间缺少该年度的共享年序入口，易使跨政权材料被单一政权叙事遮蔽。 而它所打开或收紧的下一步是：保留该年度的检索与补账位置；后续仅在获得可核单项材料时拆分为具体政治、财政、军事、社会或文化事件。

关于可说范围，须保留：本条只确认该年位于可追踪的共享年序中；正史、地方文书和外部材料的保存密度不一，不能由桥接标签推断行动、统属、疆界或社会影响。 进一步讨论可参照 北宋、辽与东北亚并行年序及边境网络研究

\subsection{984 年（共享年序桥接）}
\noindent\eventanchor{SY-984-CHRONOLOGY-BRIDGE}{【C级年序桥接】保留984 年的多政权并行检索入口；本条不主张所列政权在当年发生直接互动，具体事件须另以单项材料入账。}\textbf{宋、辽与北汉／高丽（年序桥接）。}
宋、辽与北汉／高丽（年序桥接） 在 984 年（共享年序桥接） 的材料痕迹是〔【C级年序桥接】保留984 年的多政权并行检索入口；本条不主张所列政权在当年发生直接互动，具体事件须另以单项材料入账。〕，出处为 《宋史》；《辽史》；《续资治通鉴长编》相关年条（据事核）。题名保留的城、路、水道或机构名称，是目前可以落地的空间线索。

它承接的条件是：相邻已入账事件之间缺少该年度的共享年序入口，易使跨政权材料被单一政权叙事遮蔽。 随后可追踪的变化为：保留该年度的检索与补账位置；后续仅在获得可核单项材料时拆分为具体政治、财政、军事、社会或文化事件。

证据的边界是：本条只确认该年位于可追踪的共享年序中；正史、地方文书和外部材料的保存密度不一，不能由桥接标签推断行动、统属、疆界或社会影响。 北宋、辽与东北亚并行年序及边境网络研究 可用于继续检验这一结论。

\subsection{985 年（共享年序桥接）}
\noindent\eventanchor{SY-985-CHRONOLOGY-BRIDGE}{【C级年序桥接】保留985 年的多政权并行检索入口；本条不主张所列政权在当年发生直接互动，具体事件须另以单项材料入账。}\textbf{宋、辽与北汉／高丽（年序桥接）。}
从 《宋史》；《辽史》；《续资治通鉴长编》相关年条（据事核） 的 985 年（共享年序桥接） 条目看，〔【C级年序桥接】保留985 年的多政权并行检索入口；本条不主张所列政权在当年发生直接互动，具体事件须另以单项材料入账。〕使 宋、辽与北汉／高丽（年序桥接） 的一个具体行动进入叙事；题名保留的城、路、水道或机构名称，是目前可以落地的空间线索。

前一层压力可由以下线索辨认：相邻已入账事件之间缺少该年度的共享年序入口，易使跨政权材料被单一政权叙事遮蔽。 这一处置留下的后续条件是：保留该年度的检索与补账位置；后续仅在获得可核单项材料时拆分为具体政治、财政、军事、社会或文化事件。

本条只确认该年位于可追踪的共享年序中；正史、地方文书和外部材料的保存密度不一，不能由桥接标签推断行动、统属、疆界或社会影响。 因此，北宋、辽与东北亚并行年序及边境网络研究 只能扩展问题，不能代替未保存的细节。

\subsection{986 年（共享年序桥接）}
\noindent\eventanchor{SY-986-CHRONOLOGY-BRIDGE}{【C级年序桥接】保留986 年的多政权并行检索入口；本条不主张所列政权在当年发生直接互动，具体事件须另以单项材料入账。}\textbf{宋、辽与北汉／高丽（年序桥接）。}
986 年（共享年序桥接） 的记载将 宋、辽与北汉／高丽（年序桥接） 与〔【C级年序桥接】保留986 年的多政权并行检索入口；本条不主张所列政权在当年发生直接互动，具体事件须另以单项材料入账。〕相连。此处可确认的是这项被记录的行动；题名保留的城、路、水道或机构名称，是目前可以落地的空间线索。

相邻已入账事件之间缺少该年度的共享年序入口，易使跨政权材料被单一政权叙事遮蔽。 记录所见的结果则是：保留该年度的检索与补账位置；后续仅在获得可核单项材料时拆分为具体政治、财政、军事、社会或文化事件。 日期相邻不等于因果已经证实。

关于可说范围，须保留：本条只确认该年位于可追踪的共享年序中；正史、地方文书和外部材料的保存密度不一，不能由桥接标签推断行动、统属、疆界或社会影响。 进一步讨论可参照 北宋、辽与东北亚并行年序及边境网络研究

\subsection{987 年（共享年序桥接）}
\noindent\eventanchor{SY-987-CHRONOLOGY-BRIDGE}{【C级年序桥接】保留987 年的多政权并行检索入口；本条不主张所列政权在当年发生直接互动，具体事件须另以单项材料入账。}\textbf{宋、辽与北汉／高丽（年序桥接）。}
〔【C级年序桥接】保留987 年的多政权并行检索入口；本条不主张所列政权在当年发生直接互动，具体事件须另以单项材料入账。〕见于 《宋史》；《辽史》；《续资治通鉴长编》相关年条（据事核） 的 987 年（共享年序桥接） 记录。宋、辽与北汉／高丽（年序桥接） 的选择因此有了可回查的时间位置；题名保留的城、路、水道或机构名称，是目前可以落地的空间线索。

把本项放回事件链，能够看到的前提为：相邻已入账事件之间缺少该年度的共享年序入口，易使跨政权材料被单一政权叙事遮蔽。 而它所打开或收紧的下一步是：保留该年度的检索与补账位置；后续仅在获得可核单项材料时拆分为具体政治、财政、军事、社会或文化事件。

证据的边界是：本条只确认该年位于可追踪的共享年序中；正史、地方文书和外部材料的保存密度不一，不能由桥接标签推断行动、统属、疆界或社会影响。 北宋、辽与东北亚并行年序及边境网络研究 可用于继续检验这一结论。

\subsection{988 年（共享年序桥接）}
\noindent\eventanchor{SY-988-CHRONOLOGY-BRIDGE}{【C级年序桥接】保留988 年的多政权并行检索入口；本条不主张所列政权在当年发生直接互动，具体事件须另以单项材料入账。}\textbf{宋、辽与北汉／高丽（年序桥接）。}
988 年（共享年序桥接），《宋史》；《辽史》；《续资治通鉴长编》相关年条（据事核） 所见的〔【C级年序桥接】保留988 年的多政权并行检索入口；本条不主张所列政权在当年发生直接互动，具体事件须另以单项材料入账。〕把 宋、辽与北汉／高丽（年序桥接） 的处置留在可复核的年序中；题名保留的城、路、水道或机构名称，是目前可以落地的空间线索。

它承接的条件是：相邻已入账事件之间缺少该年度的共享年序入口，易使跨政权材料被单一政权叙事遮蔽。 随后可追踪的变化为：保留该年度的检索与补账位置；后续仅在获得可核单项材料时拆分为具体政治、财政、军事、社会或文化事件。

本条只确认该年位于可追踪的共享年序中；正史、地方文书和外部材料的保存密度不一，不能由桥接标签推断行动、统属、疆界或社会影响。 因此，北宋、辽与东北亚并行年序及边境网络研究 只能扩展问题，不能代替未保存的细节。

\subsection{989 年（共享年序桥接）}
\noindent\eventanchor{SY-989-CHRONOLOGY-BRIDGE}{【C级年序桥接】保留989 年的多政权并行检索入口；本条不主张所列政权在当年发生直接互动，具体事件须另以单项材料入账。}\textbf{宋、辽与北汉／高丽（年序桥接）。}
宋、辽与北汉／高丽（年序桥接） 在 989 年（共享年序桥接） 的材料痕迹是〔【C级年序桥接】保留989 年的多政权并行检索入口；本条不主张所列政权在当年发生直接互动，具体事件须另以单项材料入账。〕，出处为 《宋史》；《辽史》；《续资治通鉴长编》相关年条（据事核）。题名保留的城、路、水道或机构名称，是目前可以落地的空间线索。

前一层压力可由以下线索辨认：相邻已入账事件之间缺少该年度的共享年序入口，易使跨政权材料被单一政权叙事遮蔽。 这一处置留下的后续条件是：保留该年度的检索与补账位置；后续仅在获得可核单项材料时拆分为具体政治、财政、军事、社会或文化事件。

关于可说范围，须保留：本条只确认该年位于可追踪的共享年序中；正史、地方文书和外部材料的保存密度不一，不能由桥接标签推断行动、统属、疆界或社会影响。 进一步讨论可参照 北宋、辽与东北亚并行年序及边境网络研究

\subsection{990 年（共享年序桥接）}
\noindent\eventanchor{SY-990-CHRONOLOGY-BRIDGE}{【C级年序桥接】保留990 年的多政权并行检索入口；本条不主张所列政权在当年发生直接互动，具体事件须另以单项材料入账。}\textbf{宋、辽与北汉／高丽（年序桥接）。}
从 《宋史》；《辽史》；《续资治通鉴长编》相关年条（据事核） 的 990 年（共享年序桥接） 条目看，〔【C级年序桥接】保留990 年的多政权并行检索入口；本条不主张所列政权在当年发生直接互动，具体事件须另以单项材料入账。〕使 宋、辽与北汉／高丽（年序桥接） 的一个具体行动进入叙事；题名保留的城、路、水道或机构名称，是目前可以落地的空间线索。

相邻已入账事件之间缺少该年度的共享年序入口，易使跨政权材料被单一政权叙事遮蔽。 记录所见的结果则是：保留该年度的检索与补账位置；后续仅在获得可核单项材料时拆分为具体政治、财政、军事、社会或文化事件。 日期相邻不等于因果已经证实。

证据的边界是：本条只确认该年位于可追踪的共享年序中；正史、地方文书和外部材料的保存密度不一，不能由桥接标签推断行动、统属、疆界或社会影响。 北宋、辽与东北亚并行年序及边境网络研究 可用于继续检验这一结论。

\subsection{991 年（共享年序桥接）}
\noindent\eventanchor{SY-991-CHRONOLOGY-BRIDGE}{【C级年序桥接】保留991 年的多政权并行检索入口；本条不主张所列政权在当年发生直接互动，具体事件须另以单项材料入账。}\textbf{宋、辽与北汉／高丽（年序桥接）。}
991 年（共享年序桥接） 的记载将 宋、辽与北汉／高丽（年序桥接） 与〔【C级年序桥接】保留991 年的多政权并行检索入口；本条不主张所列政权在当年发生直接互动，具体事件须另以单项材料入账。〕相连。此处可确认的是这项被记录的行动；题名保留的城、路、水道或机构名称，是目前可以落地的空间线索。

把本项放回事件链，能够看到的前提为：相邻已入账事件之间缺少该年度的共享年序入口，易使跨政权材料被单一政权叙事遮蔽。 而它所打开或收紧的下一步是：保留该年度的检索与补账位置；后续仅在获得可核单项材料时拆分为具体政治、财政、军事、社会或文化事件。

本条只确认该年位于可追踪的共享年序中；正史、地方文书和外部材料的保存密度不一，不能由桥接标签推断行动、统属、疆界或社会影响。 因此，北宋、辽与东北亚并行年序及边境网络研究 只能扩展问题，不能代替未保存的细节。

\subsection{992 年（共享年序桥接）}
\noindent\eventanchor{SY-992-CHRONOLOGY-BRIDGE}{【C级年序桥接】保留992 年的多政权并行检索入口；本条不主张所列政权在当年发生直接互动，具体事件须另以单项材料入账。}\textbf{宋、辽与北汉／高丽（年序桥接）。}
〔【C级年序桥接】保留992 年的多政权并行检索入口；本条不主张所列政权在当年发生直接互动，具体事件须另以单项材料入账。〕见于 《宋史》；《辽史》；《续资治通鉴长编》相关年条（据事核） 的 992 年（共享年序桥接） 记录。宋、辽与北汉／高丽（年序桥接） 的选择因此有了可回查的时间位置；题名保留的城、路、水道或机构名称，是目前可以落地的空间线索。

它承接的条件是：相邻已入账事件之间缺少该年度的共享年序入口，易使跨政权材料被单一政权叙事遮蔽。 随后可追踪的变化为：保留该年度的检索与补账位置；后续仅在获得可核单项材料时拆分为具体政治、财政、军事、社会或文化事件。

关于可说范围，须保留：本条只确认该年位于可追踪的共享年序中；正史、地方文书和外部材料的保存密度不一，不能由桥接标签推断行动、统属、疆界或社会影响。 进一步讨论可参照 北宋、辽与东北亚并行年序及边境网络研究

\subsection{993 年（共享年序桥接）}
\noindent\eventanchor{SY-993-CHRONOLOGY-BRIDGE}{【C级年序桥接】保留993 年的多政权并行检索入口；本条不主张所列政权在当年发生直接互动，具体事件须另以单项材料入账。}\textbf{宋、辽与北汉／高丽（年序桥接）。}
993 年（共享年序桥接），《宋史》；《辽史》；《续资治通鉴长编》相关年条（据事核） 所见的〔【C级年序桥接】保留993 年的多政权并行检索入口；本条不主张所列政权在当年发生直接互动，具体事件须另以单项材料入账。〕把 宋、辽与北汉／高丽（年序桥接） 的处置留在可复核的年序中；题名保留的城、路、水道或机构名称，是目前可以落地的空间线索。

前一层压力可由以下线索辨认：相邻已入账事件之间缺少该年度的共享年序入口，易使跨政权材料被单一政权叙事遮蔽。 这一处置留下的后续条件是：保留该年度的检索与补账位置；后续仅在获得可核单项材料时拆分为具体政治、财政、军事、社会或文化事件。

证据的边界是：本条只确认该年位于可追踪的共享年序中；正史、地方文书和外部材料的保存密度不一，不能由桥接标签推断行动、统属、疆界或社会影响。 北宋、辽与东北亚并行年序及边境网络研究 可用于继续检验这一结论。

\subsection{994 年（共享年序桥接）}
\noindent\eventanchor{SY-994-CHRONOLOGY-BRIDGE}{【C级年序桥接】保留994 年的多政权并行检索入口；本条不主张所列政权在当年发生直接互动，具体事件须另以单项材料入账。}\textbf{宋、辽与北汉／高丽（年序桥接）。}
宋、辽与北汉／高丽（年序桥接） 在 994 年（共享年序桥接） 的材料痕迹是〔【C级年序桥接】保留994 年的多政权并行检索入口；本条不主张所列政权在当年发生直接互动，具体事件须另以单项材料入账。〕，出处为 《宋史》；《辽史》；《续资治通鉴长编》相关年条（据事核）。题名保留的城、路、水道或机构名称，是目前可以落地的空间线索。

相邻已入账事件之间缺少该年度的共享年序入口，易使跨政权材料被单一政权叙事遮蔽。 记录所见的结果则是：保留该年度的检索与补账位置；后续仅在获得可核单项材料时拆分为具体政治、财政、军事、社会或文化事件。 日期相邻不等于因果已经证实。

本条只确认该年位于可追踪的共享年序中；正史、地方文书和外部材料的保存密度不一，不能由桥接标签推断行动、统属、疆界或社会影响。 因此，北宋、辽与东北亚并行年序及边境网络研究 只能扩展问题，不能代替未保存的细节。

\subsection{995 年（共享年序桥接）}
\noindent\eventanchor{SY-995-CHRONOLOGY-BRIDGE}{【C级年序桥接】保留995 年的多政权并行检索入口；本条不主张所列政权在当年发生直接互动，具体事件须另以单项材料入账。}\textbf{宋、辽与北汉／高丽（年序桥接）。}
从 《宋史》；《辽史》；《续资治通鉴长编》相关年条（据事核） 的 995 年（共享年序桥接） 条目看，〔【C级年序桥接】保留995 年的多政权并行检索入口；本条不主张所列政权在当年发生直接互动，具体事件须另以单项材料入账。〕使 宋、辽与北汉／高丽（年序桥接） 的一个具体行动进入叙事；题名保留的城、路、水道或机构名称，是目前可以落地的空间线索。

把本项放回事件链，能够看到的前提为：相邻已入账事件之间缺少该年度的共享年序入口，易使跨政权材料被单一政权叙事遮蔽。 而它所打开或收紧的下一步是：保留该年度的检索与补账位置；后续仅在获得可核单项材料时拆分为具体政治、财政、军事、社会或文化事件。

关于可说范围，须保留：本条只确认该年位于可追踪的共享年序中；正史、地方文书和外部材料的保存密度不一，不能由桥接标签推断行动、统属、疆界或社会影响。 进一步讨论可参照 北宋、辽与东北亚并行年序及边境网络研究

\subsection{996 年（共享年序桥接）}
\noindent\eventanchor{SY-996-CHRONOLOGY-BRIDGE}{【C级年序桥接】保留996 年的多政权并行检索入口；本条不主张所列政权在当年发生直接互动，具体事件须另以单项材料入账。}\textbf{宋、辽与北汉／高丽（年序桥接）。}
996 年（共享年序桥接） 的记载将 宋、辽与北汉／高丽（年序桥接） 与〔【C级年序桥接】保留996 年的多政权并行检索入口；本条不主张所列政权在当年发生直接互动，具体事件须另以单项材料入账。〕相连。此处可确认的是这项被记录的行动；题名保留的城、路、水道或机构名称，是目前可以落地的空间线索。

它承接的条件是：相邻已入账事件之间缺少该年度的共享年序入口，易使跨政权材料被单一政权叙事遮蔽。 随后可追踪的变化为：保留该年度的检索与补账位置；后续仅在获得可核单项材料时拆分为具体政治、财政、军事、社会或文化事件。

证据的边界是：本条只确认该年位于可追踪的共享年序中；正史、地方文书和外部材料的保存密度不一，不能由桥接标签推断行动、统属、疆界或社会影响。 北宋、辽与东北亚并行年序及边境网络研究 可用于继续检验这一结论。

\subsection{997 年（共享年序桥接）}
\noindent\eventanchor{SY-997-CHRONOLOGY-BRIDGE}{【C级年序桥接】保留997 年的多政权并行检索入口；本条不主张所列政权在当年发生直接互动，具体事件须另以单项材料入账。}\textbf{宋、辽与北汉／高丽（年序桥接）。}
〔【C级年序桥接】保留997 年的多政权并行检索入口；本条不主张所列政权在当年发生直接互动，具体事件须另以单项材料入账。〕见于 《宋史》；《辽史》；《续资治通鉴长编》相关年条（据事核） 的 997 年（共享年序桥接） 记录。宋、辽与北汉／高丽（年序桥接） 的选择因此有了可回查的时间位置；题名保留的城、路、水道或机构名称，是目前可以落地的空间线索。

前一层压力可由以下线索辨认：相邻已入账事件之间缺少该年度的共享年序入口，易使跨政权材料被单一政权叙事遮蔽。 这一处置留下的后续条件是：保留该年度的检索与补账位置；后续仅在获得可核单项材料时拆分为具体政治、财政、军事、社会或文化事件。

本条只确认该年位于可追踪的共享年序中；正史、地方文书和外部材料的保存密度不一，不能由桥接标签推断行动、统属、疆界或社会影响。 因此，北宋、辽与东北亚并行年序及边境网络研究 只能扩展问题，不能代替未保存的细节。

\subsection{998 年（共享年序桥接）}
\noindent\eventanchor{SY-998-CHRONOLOGY-BRIDGE}{【C级年序桥接】保留998 年的多政权并行检索入口；本条不主张所列政权在当年发生直接互动，具体事件须另以单项材料入账。}\textbf{宋、辽与北汉／高丽（年序桥接）。}
998 年（共享年序桥接），《宋史》；《辽史》；《续资治通鉴长编》相关年条（据事核） 所见的〔【C级年序桥接】保留998 年的多政权并行检索入口；本条不主张所列政权在当年发生直接互动，具体事件须另以单项材料入账。〕把 宋、辽与北汉／高丽（年序桥接） 的处置留在可复核的年序中；题名保留的城、路、水道或机构名称，是目前可以落地的空间线索。

相邻已入账事件之间缺少该年度的共享年序入口，易使跨政权材料被单一政权叙事遮蔽。 记录所见的结果则是：保留该年度的检索与补账位置；后续仅在获得可核单项材料时拆分为具体政治、财政、军事、社会或文化事件。 日期相邻不等于因果已经证实。

关于可说范围，须保留：本条只确认该年位于可追踪的共享年序中；正史、地方文书和外部材料的保存密度不一，不能由桥接标签推断行动、统属、疆界或社会影响。 进一步讨论可参照 北宋、辽与东北亚并行年序及边境网络研究

\subsection{999 年（共享年序桥接）}
\noindent\eventanchor{SY-999-CHRONOLOGY-BRIDGE}{【C级年序桥接】保留999 年的多政权并行检索入口；本条不主张所列政权在当年发生直接互动，具体事件须另以单项材料入账。}\textbf{宋、辽与北汉／高丽（年序桥接）。}
宋、辽与北汉／高丽（年序桥接） 在 999 年（共享年序桥接） 的材料痕迹是〔【C级年序桥接】保留999 年的多政权并行检索入口；本条不主张所列政权在当年发生直接互动，具体事件须另以单项材料入账。〕，出处为 《宋史》；《辽史》；《续资治通鉴长编》相关年条（据事核）。题名保留的城、路、水道或机构名称，是目前可以落地的空间线索。

把本项放回事件链，能够看到的前提为：相邻已入账事件之间缺少该年度的共享年序入口，易使跨政权材料被单一政权叙事遮蔽。 而它所打开或收紧的下一步是：保留该年度的检索与补账位置；后续仅在获得可核单项材料时拆分为具体政治、财政、军事、社会或文化事件。

证据的边界是：本条只确认该年位于可追踪的共享年序中；正史、地方文书和外部材料的保存密度不一，不能由桥接标签推断行动、统属、疆界或社会影响。 北宋、辽与东北亚并行年序及边境网络研究 可用于继续检验这一结论。

\subsection{961年（宋建隆二年、辽应历11年；公元换算据两国年号对照）}
\noindent\eventanchor{SYD-961-LIAO-001}{乙未，李繼勳敗北漢軍，俘遼州刺史傅廷彥、弟勳來獻}\textbf{宋与辽。}
从 《宋史》卷001〈建隆二年〉；《辽史》本纪同年条（互校） 的 961年（宋建隆二年、辽应历11年；公元换算据两国年号对照） 条目看，〔乙未，李繼勳敗北漢軍，俘遼州刺史傅廷彥、弟勳來獻〕使 宋与辽 的一个具体行动进入叙事；题名保留的城、路、水道或机构名称，是目前可以落地的空间线索。

它承接的条件是：本纪从朝廷视角记录军政行动；出兵与战果需分辨。 随后可追踪的变化为：后续路线、控制与损失应与对方记录和遗址材料复核。

《宋史》为元末官修，本纪保存宋廷纪年与处置；战果、数字、地方执行及对方动机须以编年、会要、对方史料或地方材料互校。 因此，宋辽边政、使节记录与契丹碑刻研究 只能扩展问题，不能代替未保存的细节。

\subsection{963年（宋乾德元年、辽应历13年；公元换算据两国年号对照）}
\noindent\eventanchor{SYD-963-LIAO-002}{己亥，契丹幽州岐溝關使柴廷翰等來降}\textbf{宋与辽。}
963年（宋乾德元年、辽应历13年；公元换算据两国年号对照） 的记载将 宋与辽 与〔己亥，契丹幽州岐溝關使柴廷翰等來降〕相连。此处可确认的是这项被记录的行动；题名保留的城、路、水道或机构名称，是目前可以落地的空间线索。

前一层压力可由以下线索辨认：本纪从朝廷视角记录军政行动；出兵与战果需分辨。 这一处置留下的后续条件是：后续路线、控制与损失应与对方记录和遗址材料复核。

关于可说范围，须保留：《宋史》为元末官修，本纪保存宋廷纪年与处置；战果、数字、地方执行及对方动机须以编年、会要、对方史料或地方材料互校。 进一步讨论可参照 宋辽边政、使节记录与契丹碑刻研究

\noindent\eventanchor{SYD-963-LIAO-086}{戊寅，北漢引契丹兵攻平晉，遣洺州防禦使郭進等救之}\textbf{宋与辽。}
〔戊寅，北漢引契丹兵攻平晉，遣洺州防禦使郭進等救之〕见于 《宋史》卷001〈乾德元年〉；《辽史》本纪同年条（互校） 的 963年（宋乾德元年、辽应历13年；公元换算据两国年号对照） 记录。宋与辽 的选择因此有了可回查的时间位置；题名保留的城、路、水道或机构名称，是目前可以落地的空间线索。

本纪从朝廷视角记录军政行动；出兵与战果需分辨。 记录所见的结果则是：后续路线、控制与损失应与对方记录和遗址材料复核。 日期相邻不等于因果已经证实。

证据的边界是：《宋史》为元末官修，本纪保存宋廷纪年与处置；战果、数字、地方执行及对方动机须以编年、会要、对方史料或地方材料互校。 宋辽边政、使节记录与契丹碑刻研究 可用于继续检验这一结论。

\subsection{964年（宋乾德二年、辽应历14年；公元换算据两国年号对照）}
\noindent\eventanchor{SYD-964-LIAO-003}{二月戊申朔，北漢遼州刺史杜延韜以城來降}\textbf{宋与辽。}
964年（宋乾德二年、辽应历14年；公元换算据两国年号对照），《宋史》卷001〈乾德二年〉；《辽史》本纪同年条（互校） 所见的〔二月戊申朔，北漢遼州刺史杜延韜以城來降〕把 宋与辽 的处置留在可复核的年序中；题名保留的城、路、水道或机构名称，是目前可以落地的空间线索。

把本项放回事件链，能够看到的前提为：本纪从朝廷视角记录军政行动；出兵与战果需分辨。 而它所打开或收紧的下一步是：后续路线、控制与损失应与对方记录和遗址材料复核。

《宋史》为元末官修，本纪保存宋廷纪年与处置；战果、数字、地方执行及对方动机须以编年、会要、对方史料或地方材料互校。 因此，宋辽边政、使节记录与契丹碑刻研究 只能扩展问题，不能代替未保存的细节。

\subsection{966年（宋乾德四年、辽应历16年；公元换算据两国年号对照）}
\noindent\eventanchor{SYD-966-LIAO-004}{丁巳，契丹天德軍節度使于延超與其子來降}\textbf{宋与辽。}
宋与辽 在 966年（宋乾德四年、辽应历16年；公元换算据两国年号对照） 的材料痕迹是〔丁巳，契丹天德軍節度使于延超與其子來降〕，出处为 《宋史》卷002〈乾德四年〉；《辽史》本纪同年条（互校）。材料未将地点细化到城址，不能用中枢所在地替它补足。

它承接的条件是：本纪从朝廷视角记录军政行动；出兵与战果需分辨。 随后可追踪的变化为：后续路线、控制与损失应与对方记录和遗址材料复核。

关于可说范围，须保留：《宋史》为元末官修，本纪保存宋廷纪年与处置；战果、数字、地方执行及对方动机须以编年、会要、对方史料或地方材料互校。 进一步讨论可参照 宋辽边政、使节记录与契丹碑刻研究

\subsection{969年（宋開寶二年、辽保宁1年；公元换算据两国年号对照）}
\noindent\eventanchor{SYD-969-LIAO-005}{己未，何繼筠敗契丹於陽曲，斬首數千級，俘武州刺史王彥符以獻，命陳示所獲首級、鎧甲於城下}\textbf{宋与辽。}
从 《宋史》卷002〈開寶二年〉；《辽史》本纪同年条（互校） 的 969年（宋開寶二年、辽保宁1年；公元换算据两国年号对照） 条目看，〔己未，何繼筠敗契丹於陽曲，斬首數千級，俘武州刺史王彥符以獻，命陳示所獲首級、鎧甲於城下〕使 宋与辽 的一个具体行动进入叙事；题名保留的城、路、水道或机构名称，是目前可以落地的空间线索。

前一层压力可由以下线索辨认：本纪年条记录政令或机构处置；实施背景须参照制度材料。 这一处置留下的后续条件是：可在同年及后续政令中追踪；条文不单独证明地方执行。

证据的边界是：《宋史》为元末官修，本纪保存宋廷纪年与处置；战果、数字、地方执行及对方动机须以编年、会要、对方史料或地方材料互校。 宋辽边政、使节记录与契丹碑刻研究 可用于继续检验这一结论。

\noindent\eventanchor{SYD-969-LIAO-087}{命彰德軍節度使韓仲贇為北面都部署，彰義軍節度使郭延義副之，以防契丹}\textbf{宋与辽。}
969年（宋開寶二年、辽保宁1年；公元换算据两国年号对照） 的记载将 宋与辽 与〔命彰德軍節度使韓仲贇為北面都部署，彰義軍節度使郭延義副之，以防契丹〕相连。此处可确认的是这项被记录的行动；题名保留的城、路、水道或机构名称，是目前可以落地的空间线索。

本纪年条记录政令或机构处置；实施背景须参照制度材料。 记录所见的结果则是：可在同年及后续政令中追踪；条文不单独证明地方执行。 日期相邻不等于因果已经证实。

《宋史》为元末官修，本纪保存宋廷纪年与处置；战果、数字、地方执行及对方动机须以编年、会要、对方史料或地方材料互校。 因此，宋辽边政、使节记录与契丹碑刻研究 只能扩展问题，不能代替未保存的细节。

\subsection{970年（宋開寶三年、辽保宁2年；公元换算据两国年号对照）}
\noindent\eventanchor{SYD-970-LIAO-006}{癸亥，定州駐泊都監田欽祚敗契丹於遂城}\textbf{宋与辽。}
〔癸亥，定州駐泊都監田欽祚敗契丹於遂城〕见于 《宋史》卷002〈開寶三年〉；《辽史》本纪同年条（互校） 的 970年（宋開寶三年、辽保宁2年；公元换算据两国年号对照） 记录。宋与辽 的选择因此有了可回查的时间位置；题名保留的城、路、水道或机构名称，是目前可以落地的空间线索。

把本项放回事件链，能够看到的前提为：本纪从朝廷视角记录军政行动；出兵与战果需分辨。 而它所打开或收紧的下一步是：后续路线、控制与损失应与对方记录和遗址材料复核。

关于可说范围，须保留：《宋史》为元末官修，本纪保存宋廷纪年与处置；战果、数字、地方执行及对方动机须以编年、会要、对方史料或地方材料互校。 进一步讨论可参照 宋辽边政、使节记录与契丹碑刻研究

\subsection{975年（宋開寶八年、辽保宁7年；公元换算据两国年号对照）}
\noindent\eventanchor{SYD-975-LIAO-007}{庚辰，遣閣門使郝崇信、太常丞呂端使契丹}\textbf{宋与辽。}
975年（宋開寶八年、辽保宁7年；公元换算据两国年号对照），《宋史》卷003〈開寶八年〉；《辽史》本纪同年条（互校） 所见的〔庚辰，遣閣門使郝崇信、太常丞呂端使契丹〕把 宋与辽 的处置留在可复核的年序中；材料未将地点细化到城址，不能用中枢所在地替它补足。

它承接的条件是：本纪记录使节、贡赐或往来，不等于稳定统属关系。 随后可追踪的变化为：可与对方编年和交通、文书材料核对其后续落实。

证据的边界是：《宋史》为元末官修，本纪保存宋廷纪年与处置；战果、数字、地方执行及对方动机须以编年、会要、对方史料或地方材料互校。 宋辽边政、使节记录与契丹碑刻研究 可用于继续检验这一结论。

\noindent\eventanchor{SYD-975-LIAO-088}{己亥，契丹遣使克沙骨慎思以書來講和}\textbf{宋与辽。}
宋与辽 在 975年（宋開寶八年、辽保宁7年；公元换算据两国年号对照） 的材料痕迹是〔己亥，契丹遣使克沙骨慎思以書來講和〕，出处为 《宋史》卷003〈開寶八年〉；《辽史》本纪同年条（互校）。材料未将地点细化到城址，不能用中枢所在地替它补足。

前一层压力可由以下线索辨认：本纪记录使节、贡赐或往来，不等于稳定统属关系。 这一处置留下的后续条件是：可与对方编年和交通、文书材料核对其后续落实。

《宋史》为元末官修，本纪保存宋廷纪年与处置；战果、数字、地方执行及对方动机须以编年、会要、对方史料或地方材料互校。 因此，宋辽边政、使节记录与契丹碑刻研究 只能扩展问题，不能代替未保存的细节。

\subsection{976年（宋開寶九年、辽保宁8年；公元换算据两国年号对照）}
\noindent\eventanchor{SYD-976-LIAO-008}{契丹遣使耶律延寧頁以御衣、玉帶、名馬、散馬、白鶻來賀長春節}\textbf{宋与辽。}
从 《宋史》卷003〈開寶九年〉；《辽史》本纪同年条（互校） 的 976年（宋開寶九年、辽保宁8年；公元换算据两国年号对照） 条目看，〔契丹遣使耶律延寧頁以御衣、玉帶、名馬、散馬、白鶻來賀長春節〕使 宋与辽 的一个具体行动进入叙事；材料未将地点细化到城址，不能用中枢所在地替它补足。

本纪记录使节、贡赐或往来，不等于稳定统属关系。 记录所见的结果则是：可与对方编年和交通、文书材料核对其后续落实。 日期相邻不等于因果已经证实。

关于可说范围，须保留：《宋史》为元末官修，本纪保存宋廷纪年与处置；战果、数字、地方执行及对方动机须以编年、会要、对方史料或地方材料互校。 进一步讨论可参照 宋辽边政、使节记录与契丹碑刻研究

\noindent\eventanchor{SYD-976-LIAO-089}{己丑，遣著作郎馮正、佐郎張玘使契丹告哀}\textbf{宋与辽。}
976年（宋開寶九年、辽保宁8年；公元换算据两国年号对照） 的记载将 宋与辽 与〔己丑，遣著作郎馮正、佐郎張玘使契丹告哀〕相连。此处可确认的是这项被记录的行动；材料未将地点细化到城址，不能用中枢所在地替它补足。

把本项放回事件链，能够看到的前提为：本纪记录使节、贡赐或往来，不等于稳定统属关系。 而它所打开或收紧的下一步是：可与对方编年和交通、文书材料核对其后续落实。

证据的边界是：《宋史》为元末官修，本纪保存宋廷纪年与处置；战果、数字、地方执行及对方动机须以编年、会要、对方史料或地方材料互校。 宋辽边政、使节记录与契丹碑刻研究 可用于继续检验这一结论。

\subsection{977年（宋太平興國二年、辽保宁9年；公元换算据两国年号对照）}
\noindent\eventanchor{SYD-977-LIAO-009}{丁酉，契丹遣使來會葬}\textbf{宋与辽。}
〔丁酉，契丹遣使來會葬〕见于 《宋史》卷004〈太平興國二年〉；《辽史》本纪同年条（互校） 的 977年（宋太平興國二年、辽保宁9年；公元换算据两国年号对照） 记录。宋与辽 的选择因此有了可回查的时间位置；材料未将地点细化到城址，不能用中枢所在地替它补足。

它承接的条件是：本纪记录使节、贡赐或往来，不等于稳定统属关系。 随后可追踪的变化为：可与对方编年和交通、文书材料核对其后续落实。

《宋史》为元末官修，本纪保存宋廷纪年与处置；战果、数字、地方执行及对方动机须以编年、会要、对方史料或地方材料互校。 因此，宋辽边政、使节记录与契丹碑刻研究 只能扩展问题，不能代替未保存的细节。

\noindent\eventanchor{SYD-977-LIAO-090}{二月甲午，契丹遣使來賀即位及正旦}\textbf{宋与辽。}
977年（宋太平興國二年、辽保宁9年；公元换算据两国年号对照），《宋史》卷004〈太平興國二年〉；《辽史》本纪同年条（互校） 所见的〔二月甲午，契丹遣使來賀即位及正旦〕把 宋与辽 的处置留在可复核的年序中；材料未将地点细化到城址，不能用中枢所在地替它补足。

前一层压力可由以下线索辨认：本纪记录使节、贡赐或往来，不等于稳定统属关系。 这一处置留下的后续条件是：可与对方编年和交通、文书材料核对其后续落实。

关于可说范围，须保留：《宋史》为元末官修，本纪保存宋廷纪年与处置；战果、数字、地方执行及对方动机须以编年、会要、对方史料或地方材料互校。 进一步讨论可参照 宋辽边政、使节记录与契丹碑刻研究

\subsection{978年（宋太平興國三年、辽保宁10年；公元换算据两国年号对照）}
\noindent\eventanchor{SYD-978-LIAO-010}{冬十月癸丑朔，契丹遣使來賀乾明節}\textbf{宋与辽。}
宋与辽 在 978年（宋太平興國三年、辽保宁10年；公元换算据两国年号对照） 的材料痕迹是〔冬十月癸丑朔，契丹遣使來賀乾明節〕，出处为 《宋史》卷004〈太平興國三年〉；《辽史》本纪同年条（互校）。材料未将地点细化到城址，不能用中枢所在地替它补足。

本纪记录使节、贡赐或往来，不等于稳定统属关系。 记录所见的结果则是：可与对方编年和交通、文书材料核对其后续落实。 日期相邻不等于因果已经证实。

证据的边界是：《宋史》为元末官修，本纪保存宋廷纪年与处置；战果、数字、地方执行及对方动机须以编年、会要、对方史料或地方材料互校。 宋辽边政、使节记录与契丹碑刻研究 可用于继续检验这一结论。

\noindent\eventanchor{SYD-978-LIAO-091}{癸巳，遣李從吉等使契丹}\textbf{宋与辽。}
从 《宋史》卷004〈太平興國三年〉；《辽史》本纪同年条（互校） 的 978年（宋太平興國三年、辽保宁10年；公元换算据两国年号对照） 条目看，〔癸巳，遣李從吉等使契丹〕使 宋与辽 的一个具体行动进入叙事；材料未将地点细化到城址，不能用中枢所在地替它补足。

把本项放回事件链，能够看到的前提为：本纪记录使节、贡赐或往来，不等于稳定统属关系。 而它所打开或收紧的下一步是：可与对方编年和交通、文书材料核对其后续落实。

《宋史》为元末官修，本纪保存宋廷纪年与处置；战果、数字、地方执行及对方动机须以编年、会要、对方史料或地方材料互校。 因此，宋辽边政、使节记录与契丹碑刻研究 只能扩展问题，不能代替未保存的细节。

\subsection{979年（宋太平興國四年、辽乾亨1年；公元换算据两国年号对照）}
\noindent\eventanchor{SYD-979-LIAO-011}{丙午，鎮州都鈐轄劉廷翰及契丹戰于遂城西，大敗之，斬首萬三百級，獲三將、馬萬匹}\textbf{宋与辽。}
979年（宋太平興國四年、辽乾亨1年；公元换算据两国年号对照） 的记载将 宋与辽 与〔丙午，鎮州都鈐轄劉廷翰及契丹戰于遂城西，大敗之，斬首萬三百級，獲三將、馬萬匹〕相连。此处可确认的是这项被记录的行动；题名保留的城、路、水道或机构名称，是目前可以落地的空间线索。

它承接的条件是：本纪从朝廷视角记录军政行动；出兵与战果需分辨。 随后可追踪的变化为：后续路线、控制与损失应与对方记录和遗址材料复核。

关于可说范围，须保留：《宋史》为元末官修，本纪保存宋廷纪年与处置；战果、数字、地方执行及对方动机须以编年、会要、对方史料或地方材料互校。 进一步讨论可参照 宋辽边政、使节记录与契丹碑刻研究

\noindent\eventanchor{SYD-979-LIAO-092}{庚申，帝復自將伐契丹}\textbf{宋与辽。}
〔庚申，帝復自將伐契丹〕见于 《宋史》卷004〈太平興國四年〉；《辽史》本纪同年条（互校） 的 979年（宋太平興國四年、辽乾亨1年；公元换算据两国年号对照） 记录。宋与辽 的选择因此有了可回查的时间位置；材料未将地点细化到城址，不能用中枢所在地替它补足。

前一层压力可由以下线索辨认：本纪年条记录政令或机构处置；实施背景须参照制度材料。 这一处置留下的后续条件是：可在同年及后续政令中追踪；条文不单独证明地方执行。

证据的边界是：《宋史》为元末官修，本纪保存宋廷纪年与处置；战果、数字、地方执行及对方动机须以编年、会要、对方史料或地方材料互校。 宋辽边政、使节记录与契丹碑刻研究 可用于继续检验这一结论。

\subsection{980年（宋太平興國五年、辽乾亨2年；公元换算据两国年号对照）}
\noindent\eventanchor{SYD-980-LIAO-012}{關南與契丹戰，大破之}\textbf{宋与辽。}
980年（宋太平興國五年、辽乾亨2年；公元换算据两国年号对照），《宋史》卷004〈太平興國五年〉；《辽史》本纪同年条（互校） 所见的〔關南與契丹戰，大破之〕把 宋与辽 的处置留在可复核的年序中；题名保留的城、路、水道或机构名称，是目前可以落地的空间线索。

本纪从朝廷视角记录军政行动；出兵与战果需分辨。 记录所见的结果则是：后续路线、控制与损失应与对方记录和遗址材料复核。 日期相邻不等于因果已经证实。

《宋史》为元末官修，本纪保存宋廷纪年与处置；战果、数字、地方执行及对方动机须以编年、会要、对方史料或地方材料互校。 因此，宋辽边政、使节记录与契丹碑刻研究 只能扩展问题，不能代替未保存的细节。

\noindent\eventanchor{SYD-980-LIAO-093}{癸巳，代州言，宣徽南院使潘美敗契丹之師於雁門，殺其駙馬侍中蕭咄李，獲都指揮使李重誨}\textbf{宋与辽。}
宋与辽 在 980年（宋太平興國五年、辽乾亨2年；公元换算据两国年号对照） 的材料痕迹是〔癸巳，代州言，宣徽南院使潘美敗契丹之師於雁門，殺其駙馬侍中蕭咄李，獲都指揮使李重誨〕，出处为 《宋史》卷004〈太平興國五年〉；《辽史》本纪同年条（互校）。题名保留的城、路、水道或机构名称，是目前可以落地的空间线索。

把本项放回事件链，能够看到的前提为：本纪从朝廷视角记录军政行动；出兵与战果需分辨。 而它所打开或收紧的下一步是：后续路线、控制与损失应与对方记录和遗址材料复核。

关于可说范围，须保留：《宋史》为元末官修，本纪保存宋廷纪年与处置；战果、数字、地方执行及对方动机须以编年、会要、对方史料或地方材料互校。 进一步讨论可参照 宋辽边政、使节记录与契丹碑刻研究

\subsection{981年（宋太平興國六年、辽乾亨3年；公元换算据两国年号对照）}
\noindent\eventanchor{SYD-981-LIAO-013}{平塞軍與契丹戰，破之}\textbf{宋与辽。}
从 《宋史》卷004〈太平興國六年〉；《辽史》本纪同年条（互校） 的 981年（宋太平興國六年、辽乾亨3年；公元换算据两国年号对照） 条目看，〔平塞軍與契丹戰，破之〕使 宋与辽 的一个具体行动进入叙事；材料未将地点细化到城址，不能用中枢所在地替它补足。

它承接的条件是：本纪从朝廷视角记录军政行动；出兵与战果需分辨。 随后可追踪的变化为：后续路线、控制与损失应与对方记录和遗址材料复核。

证据的边界是：《宋史》为元末官修，本纪保存宋廷纪年与处置；战果、数字、地方执行及对方动机须以编年、会要、对方史料或地方材料互校。 宋辽边政、使节记录与契丹碑刻研究 可用于继续检验这一结论。

\noindent\eventanchor{SYD-981-LIAO-094}{七月丙午，詔渤海琰府王助討契丹}\textbf{宋与辽。}
981年（宋太平興國六年、辽乾亨3年；公元换算据两国年号对照） 的记载将 宋与辽 与〔七月丙午，詔渤海琰府王助討契丹〕相连。此处可确认的是这项被记录的行动；题名保留的城、路、水道或机构名称，是目前可以落地的空间线索。

前一层压力可由以下线索辨认：本纪年条记录政令或机构处置；实施背景须参照制度材料。 这一处置留下的后续条件是：可在同年及后续政令中追踪；条文不单独证明地方执行。

《宋史》为元末官修，本纪保存宋廷纪年与处置；战果、数字、地方执行及对方动机须以编年、会要、对方史料或地方材料互校。 因此，宋辽边政、使节记录与契丹碑刻研究 只能扩展问题，不能代替未保存的细节。

\subsection{982年（宋太平興國七年、辽乾亨4年；公元换算据两国年号对照）}
\noindent\eventanchor{SYD-982-LIAO-014}{己未，府州破契丹於新澤砦，獲其將校以下百人}\textbf{宋与辽。}
〔己未，府州破契丹於新澤砦，獲其將校以下百人〕见于 《宋史》卷004〈太平興國七年〉；《辽史》本纪同年条（互校） 的 982年（宋太平興國七年、辽乾亨4年；公元换算据两国年号对照） 记录。宋与辽 的选择因此有了可回查的时间位置；题名保留的城、路、水道或机构名称，是目前可以落地的空间线索。

本纪从朝廷视角记录军政行动；出兵与战果需分辨。 记录所见的结果则是：后续路线、控制与损失应与对方记录和遗址材料复核。 日期相邻不等于因果已经证实。

关于可说范围，须保留：《宋史》为元末官修，本纪保存宋廷纪年与处置；战果、数字、地方执行及对方动机须以编年、会要、对方史料或地方材料互校。 进一步讨论可参照 宋辽边政、使节记录与契丹碑刻研究

\noindent\eventanchor{SYD-982-LIAO-095}{閏月戊子朔，豐州與契丹戰，破之，獲其天德軍節度使蕭太}\textbf{宋与辽。}
982年（宋太平興國七年、辽乾亨4年；公元换算据两国年号对照），《宋史》卷004〈太平興國七年〉；《辽史》本纪同年条（互校） 所见的〔閏月戊子朔，豐州與契丹戰，破之，獲其天德軍節度使蕭太〕把 宋与辽 的处置留在可复核的年序中；题名保留的城、路、水道或机构名称，是目前可以落地的空间线索。

把本项放回事件链，能够看到的前提为：本纪从朝廷视角记录军政行动；出兵与战果需分辨。 而它所打开或收紧的下一步是：后续路线、控制与损失应与对方记录和遗址材料复核。

证据的边界是：《宋史》为元末官修，本纪保存宋廷纪年与处置；战果、数字、地方执行及对方动机须以编年、会要、对方史料或地方材料互校。 宋辽边政、使节记录与契丹碑刻研究 可用于继续检验这一结论。

\subsection{983年（宋太平興國八年、辽统和1年；公元换算据两国年号对照）}
\noindent\eventanchor{SYD-983-LIAO-015}{豐州破契丹兵，降三千餘帳}\textbf{宋与辽。}
宋与辽 在 983年（宋太平興國八年、辽统和1年；公元换算据两国年号对照） 的材料痕迹是〔豐州破契丹兵，降三千餘帳〕，出处为 《宋史》卷004〈太平興國八年〉；《辽史》本纪同年条（互校）。题名保留的城、路、水道或机构名称，是目前可以落地的空间线索。

它承接的条件是：本纪从朝廷视角记录军政行动；出兵与战果需分辨。 随后可追踪的变化为：后续路线、控制与损失应与对方记录和遗址材料复核。

《宋史》为元末官修，本纪保存宋廷纪年与处置；战果、数字、地方执行及对方动机须以编年、会要、对方史料或地方材料互校。 因此，宋辽边政、使节记录与契丹碑刻研究 只能扩展问题，不能代替未保存的细节。

\subsection{984年（宋雍熙元年、西夏天授礼法延祚-53年；公元换算据两国年号对照）}
\noindent\eventanchor{SYD-984-XIA-001}{夏州言，掩擊李繼遷，獲其母妻，俘千四百餘帳，繼遷走}\textbf{宋与西夏。}
从 《宋史》卷004〈雍熙元年〉；西夏文书、碑刻及《辽史》或《金史》同年条（互校） 的 984年（宋雍熙元年、西夏天授礼法延祚-53年；公元换算据两国年号对照） 条目看，〔夏州言，掩擊李繼遷，獲其母妻，俘千四百餘帳，繼遷走〕使 宋与西夏 的一个具体行动进入叙事；题名保留的城、路、水道或机构名称，是目前可以落地的空间线索。

前一层压力可由以下线索辨认：本纪从朝廷视角记录军政行动；出兵与战果需分辨。 这一处置留下的后续条件是：后续路线、控制与损失应与对方记录和遗址材料复核。

关于可说范围，须保留：《宋史》为元末官修，本纪保存宋廷纪年与处置；战果、数字、地方执行及对方动机须以编年、会要、对方史料或地方材料互校。 进一步讨论可参照 西夏文书、河西—宁夏考古与宋夏关系研究

\subsection{985年（宋雍熙二年、西夏天授礼法延祚-52年；公元换算据两国年号对照）}
\noindent\eventanchor{SYD-985-XIA-002}{辛丑，夏州行營破西蕃息利族，斬其代州刺史折羅遇并弟埋乞，又破保、洗兩族，降五十餘族}\textbf{宋与西夏。}
985年（宋雍熙二年、西夏天授礼法延祚-52年；公元换算据两国年号对照） 的记载将 宋与西夏 与〔辛丑，夏州行營破西蕃息利族，斬其代州刺史折羅遇并弟埋乞，又破保、洗兩族，降五十餘族〕相连。此处可确认的是这项被记录的行动；题名保留的城、路、水道或机构名称，是目前可以落地的空间线索。

本纪从朝廷视角记录军政行动；出兵与战果需分辨。 记录所见的结果则是：后续路线、控制与损失应与对方记录和遗址材料复核。 日期相邻不等于因果已经证实。

证据的边界是：《宋史》为元末官修，本纪保存宋廷纪年与处置；战果、数字、地方执行及对方动机须以编年、会要、对方史料或地方材料互校。 西夏文书、河西—宁夏考古与宋夏关系研究 可用于继续检验这一结论。

\noindent\eventanchor{SYD-985-XIA-053}{乙未，夏州李繼遷誘殺汝州團練使曹光實}\textbf{宋与西夏。}
〔乙未，夏州李繼遷誘殺汝州團練使曹光實〕见于 《宋史》卷005〈雍熙二年〉；西夏文书、碑刻及《辽史》或《金史》同年条（互校） 的 985年（宋雍熙二年、西夏天授礼法延祚-52年；公元换算据两国年号对照） 记录。宋与西夏 的选择因此有了可回查的时间位置；题名保留的城、路、水道或机构名称，是目前可以落地的空间线索。

把本项放回事件链，能够看到的前提为：本纪从朝廷视角记录军政行动；出兵与战果需分辨。 而它所打开或收紧的下一步是：后续路线、控制与损失应与对方记录和遗址材料复核。

《宋史》为元末官修，本纪保存宋廷纪年与处置；战果、数字、地方执行及对方动机须以编年、会要、对方史料或地方材料互校。 因此，西夏文书、河西—宁夏考古与宋夏关系研究 只能扩展问题，不能代替未保存的细节。

\subsection{986年（宋雍熙三年、辽统和4年；公元换算据两国年号对照）}
\noindent\eventanchor{SYD-986-LIAO-016}{代州副部署盧漢贇敗契丹於土鐙堡，斬獲甚眾，殺監軍舍利二人}\textbf{宋与辽。}
986年（宋雍熙三年、辽统和4年；公元换算据两国年号对照），《宋史》卷005〈雍熙三年〉；《辽史》本纪同年条（互校） 所见的〔代州副部署盧漢贇敗契丹於土鐙堡，斬獲甚眾，殺監軍舍利二人〕把 宋与辽 的处置留在可复核的年序中；题名保留的城、路、水道或机构名称，是目前可以落地的空间线索。

它承接的条件是：本纪从朝廷视角记录军政行动；出兵与战果需分辨。 随后可追踪的变化为：后续路线、控制与损失应与对方记录和遗址材料复核。

关于可说范围，须保留：《宋史》为元末官修，本纪保存宋廷纪年与处置；战果、数字、地方执行及对方动机须以编年、会要、对方史料或地方材料互校。 进一步讨论可参照 宋辽边政、使节记录与契丹碑刻研究

\noindent\eventanchor{SYD-986-LIAO-096}{會契丹十萬眾復陷寰州，楊業護送遷民遇之，苦戰力盡，為所禽，守節而死}\textbf{宋与辽。}
宋与辽 在 986年（宋雍熙三年、辽统和4年；公元换算据两国年号对照） 的材料痕迹是〔會契丹十萬眾復陷寰州，楊業護送遷民遇之，苦戰力盡，為所禽，守節而死〕，出处为 《宋史》卷005〈雍熙三年〉；《辽史》本纪同年条（互校）。题名保留的城、路、水道或机构名称，是目前可以落地的空间线索。

前一层压力可由以下线索辨认：本纪年条记录政令或机构处置；实施背景须参照制度材料。 这一处置留下的后续条件是：可在同年及后续政令中追踪；条文不单独证明地方执行。

证据的边界是：《宋史》为元末官修，本纪保存宋廷纪年与处置；战果、数字、地方执行及对方动机须以编年、会要、对方史料或地方材料互校。 宋辽边政、使节记录与契丹碑刻研究 可用于继续检验这一结论。

\subsection{988年（宋端拱元年、辽统和6年；公元换算据两国年号对照）}
\noindent\eventanchor{SYD-988-LIAO-017}{己丑，郭守文破契丹于唐河}\textbf{宋与辽。}
从 《宋史》卷005〈端拱元年〉；《辽史》本纪同年条（互校） 的 988年（宋端拱元年、辽统和6年；公元换算据两国年号对照） 条目看，〔己丑，郭守文破契丹于唐河〕使 宋与辽 的一个具体行动进入叙事；题名保留的城、路、水道或机构名称，是目前可以落地的空间线索。

本纪保留灾害或工程的报告地点与朝廷处置；灾情范围需另证。 记录所见的结果则是：可与地方文书、河工和环境材料比较，不将单点记录外推为全域因果。 日期相邻不等于因果已经证实。

《宋史》为元末官修，本纪保存宋廷纪年与处置；战果、数字、地方执行及对方动机须以编年、会要、对方史料或地方材料互校。 因此，宋辽边政、使节记录与契丹碑刻研究 只能扩展问题，不能代替未保存的细节。

\subsection{989年（宋端拱二年、辽统和7年；公元换算据两国年号对照）}
\noindent\eventanchor{SYD-989-LIAO-018}{辛丑，契丹犯威虜軍，崇儀使尹繼倫擊破之，殺其相皮室，大將于越遁去}\textbf{宋与辽。}
989年（宋端拱二年、辽统和7年；公元换算据两国年号对照） 的记载将 宋与辽 与〔辛丑，契丹犯威虜軍，崇儀使尹繼倫擊破之，殺其相皮室，大將于越遁去〕相连。此处可确认的是这项被记录的行动；材料未将地点细化到城址，不能用中枢所在地替它补足。

把本项放回事件链，能够看到的前提为：本纪从朝廷视角记录军政行动；出兵与战果需分辨。 而它所打开或收紧的下一步是：后续路线、控制与损失应与对方记录和遗址材料复核。

关于可说范围，须保留：《宋史》为元末官修，本纪保存宋廷纪年与处置；战果、数字、地方执行及对方动机须以编年、会要、对方史料或地方材料互校。 进一步讨论可参照 宋辽边政、使节记录与契丹碑刻研究

\subsection{991年（宋淳化二年、辽统和9年；公元换算据两国年号对照）}
\noindent\eventanchor{SYD-991-LIAO-019}{是歲，女真表請伐契丹，詔不許，自是遂屬契丹}\textbf{宋与辽。}
〔是歲，女真表請伐契丹，詔不許，自是遂屬契丹〕见于 《宋史》卷005〈淳化二年〉；《辽史》本纪同年条（互校） 的 991年（宋淳化二年、辽统和9年；公元换算据两国年号对照） 记录。宋与辽 的选择因此有了可回查的时间位置；材料未将地点细化到城址，不能用中枢所在地替它补足。

它承接的条件是：本纪年条记录政令或机构处置；实施背景须参照制度材料。 随后可追踪的变化为：可在同年及后续政令中追踪；条文不单独证明地方执行。

证据的边界是：《宋史》为元末官修，本纪保存宋廷纪年与处置；战果、数字、地方执行及对方动机须以编年、会要、对方史料或地方材料互校。 宋辽边政、使节记录与契丹碑刻研究 可用于继续检验这一结论。

\subsection{991年（宋淳化二年、西夏天授礼法延祚-46年；公元换算据两国年号对照）}
\noindent\eventanchor{SYD-991-XIA-003}{丙子，遣商州團練使翟守素帥兵援趙保忠于夏州}\textbf{宋与西夏。}
991年（宋淳化二年、西夏天授礼法延祚-46年；公元换算据两国年号对照），《宋史》卷005〈淳化二年〉；西夏文书、碑刻及《辽史》或《金史》同年条（互校） 所见的〔丙子，遣商州團練使翟守素帥兵援趙保忠于夏州〕把 宋与西夏 的处置留在可复核的年序中；题名保留的城、路、水道或机构名称，是目前可以落地的空间线索。

前一层压力可由以下线索辨认：本纪记录使节、贡赐或往来，不等于稳定统属关系。 这一处置留下的后续条件是：可与对方编年和交通、文书材料核对其后续落实。

《宋史》为元末官修，本纪保存宋廷纪年与处置；战果、数字、地方执行及对方动机须以编年、会要、对方史料或地方材料互校。 因此，西夏文书、河西—宁夏考古与宋夏关系研究 只能扩展问题，不能代替未保存的细节。

\subsection{994年（宋淳化五年、辽统和12年；公元换算据两国年号对照）}
\noindent\eventanchor{SYD-994-LIAO-020}{高麗遣使，以契丹來侵乞師}\textbf{宋与辽。}
宋与辽 在 994年（宋淳化五年、辽统和12年；公元换算据两国年号对照） 的材料痕迹是〔高麗遣使，以契丹來侵乞師〕，出处为 《宋史》卷005〈淳化五年〉；《辽史》本纪同年条（互校）。材料未将地点细化到城址，不能用中枢所在地替它补足。

本纪记录使节、贡赐或往来，不等于稳定统属关系。 记录所见的结果则是：可与对方编年和交通、文书材料核对其后续落实。 日期相邻不等于因果已经证实。

关于可说范围，须保留：《宋史》为元末官修，本纪保存宋廷纪年与处置；战果、数字、地方执行及对方动机须以编年、会要、对方史料或地方材料互校。 进一步讨论可参照 宋辽边政、使节记录与契丹碑刻研究

\subsection{994年（宋淳化五年、西夏天授礼法延祚-43年；公元换算据两国年号对照）}
\noindent\eventanchor{SYD-994-XIA-004}{三月乙亥，趙保忠為趙保吉所襲，奔還夏州，指揮使趙光嗣執之以獻，李繼隆帥師入夏州}\textbf{宋与西夏。}
从 《宋史》卷005〈淳化五年〉；西夏文书、碑刻及《辽史》或《金史》同年条（互校） 的 994年（宋淳化五年、西夏天授礼法延祚-43年；公元换算据两国年号对照） 条目看，〔三月乙亥，趙保忠為趙保吉所襲，奔還夏州，指揮使趙光嗣執之以獻，李繼隆帥師入夏州〕使 宋与西夏 的一个具体行动进入叙事；题名保留的城、路、水道或机构名称，是目前可以落地的空间线索。

把本项放回事件链，能够看到的前提为：本纪年条提供日期和中枢可见行动；前因不据单条补定。 而它所打开或收紧的下一步是：可回查同年及后续条目，地方影响仍须另证。

证据的边界是：《宋史》为元末官修，本纪保存宋廷纪年与处置；战果、数字、地方执行及对方动机须以编年、会要、对方史料或地方材料互校。 西夏文书、河西—宁夏考古与宋夏关系研究 可用于继续检验这一结论。

\subsection{995年（宋至道元年、辽统和13年；公元换算据两国年号对照）}
\noindent\eventanchor{SYD-995-LIAO-021}{癸亥，契丹大將韓德威誘党項勒浪、嵬族自振武犯邊，永安節度使折御卿邀擊，敗之于子河汊，勒浪等乘亂反擊德威，遂殺其將突厥大尉、司徒、舍利等，獲吐渾首領一人，德威僅以身免}\textbf{宋与辽。}
995年（宋至道元年、辽统和13年；公元换算据两国年号对照） 的记载将 宋与辽 与〔癸亥，契丹大將韓德威誘党項勒浪、嵬族自振武犯邊，永安節度使折御卿邀擊，敗之于子河汊，勒浪等乘亂反擊德威，遂殺其將突厥大尉、司徒、舍利等，獲吐渾首領一人，德威僅以身免〕相连。此处可确认的是这项被记录的行动；题名保留的城、路、水道或机构名称，是目前可以落地的空间线索。

它承接的条件是：本纪保留灾害或工程的报告地点与朝廷处置；灾情范围需另证。 随后可追踪的变化为：可与地方文书、河工和环境材料比较，不将单点记录外推为全域因果。

《宋史》为元末官修，本纪保存宋廷纪年与处置；战果、数字、地方执行及对方动机须以编年、会要、对方史料或地方材料互校。 因此，宋辽边政、使节记录与契丹碑刻研究 只能扩展问题，不能代替未保存的细节。

\noindent\eventanchor{SYD-995-LIAO-097}{乙酉，契丹犯雄州，知州何承矩擊敗之，斬其鐵林大將一人}\textbf{宋与辽。}
〔乙酉，契丹犯雄州，知州何承矩擊敗之，斬其鐵林大將一人〕见于 《宋史》卷005〈至道元年〉；《辽史》本纪同年条（互校） 的 995年（宋至道元年、辽统和13年；公元换算据两国年号对照） 记录。宋与辽 的选择因此有了可回查的时间位置；题名保留的城、路、水道或机构名称，是目前可以落地的空间线索。

前一层压力可由以下线索辨认：本纪从朝廷视角记录军政行动；出兵与战果需分辨。 这一处置留下的后续条件是：后续路线、控制与损失应与对方记录和遗址材料复核。

关于可说范围，须保留：《宋史》为元末官修，本纪保存宋廷纪年与处置；战果、数字、地方执行及对方动机须以编年、会要、对方史料或地方材料互校。 进一步讨论可参照 宋辽边政、使节记录与契丹碑刻研究

\subsection{996年（宋至道二年、西夏天授礼法延祚-41年；公元换算据两国年号对照）}
\noindent\eventanchor{SYD-996-XIA-005}{己卯，夏州、延州行營言破李繼遷於烏白池，獲未幕軍主、吃囉指揮使等二十七人，繼遷遁}\textbf{宋与西夏。}
996年（宋至道二年、西夏天授礼法延祚-41年；公元换算据两国年号对照），《宋史》卷005〈至道二年〉；西夏文书、碑刻及《辽史》或《金史》同年条（互校） 所见的〔己卯，夏州、延州行營言破李繼遷於烏白池，獲未幕軍主、吃囉指揮使等二十七人，繼遷遁〕把 宋与西夏 的处置留在可复核的年序中；题名保留的城、路、水道或机构名称，是目前可以落地的空间线索。

本纪从朝廷视角记录军政行动；出兵与战果需分辨。 记录所见的结果则是：后续路线、控制与损失应与对方记录和遗址材料复核。 日期相邻不等于因果已经证实。

证据的边界是：《宋史》为元末官修，本纪保存宋廷纪年与处置；战果、数字、地方执行及对方动机须以编年、会要、对方史料或地方材料互校。 西夏文书、河西—宁夏考古与宋夏关系研究 可用于继续检验这一结论。

\subsection{997年（宋至道三年、西夏天授礼法延祚-40年；公元换算据两国年号对照）}
\noindent\eventanchor{SYD-997-XIA-006}{冬十月，夏人寇靈州，合河都部署楊瓊擊走之}\textbf{宋与西夏。}
宋与西夏 在 997年（宋至道三年、西夏天授礼法延祚-40年；公元换算据两国年号对照） 的材料痕迹是〔冬十月，夏人寇靈州，合河都部署楊瓊擊走之〕，出处为 《宋史》卷006〈至道三年〉；西夏文书、碑刻及《辽史》或《金史》同年条（互校）。题名保留的城、路、水道或机构名称，是目前可以落地的空间线索。

把本项放回事件链，能够看到的前提为：本纪保留灾害或工程的报告地点与朝廷处置；灾情范围需另证。 而它所打开或收紧的下一步是：可与地方文书、河工和环境材料比较，不将单点记录外推为全域因果。

《宋史》为元末官修，本纪保存宋廷纪年与处置；战果、数字、地方执行及对方动机须以编年、会要、对方史料或地方材料互校。 因此，西夏文书、河西—宁夏考古与宋夏关系研究 只能扩展问题，不能代替未保存的细节。

\subsection{999年（宋咸平二年、辽统和17年；公元换算据两国年号对照）}
\noindent\eventanchor{SYD-999-LIAO-022}{府州言官軍入契丹五合川，拔黃太尉砦，殲其眾，焚其車帳，獲馬牛萬計}\textbf{宋与辽。}
从 《宋史》卷006〈咸平二年〉；《辽史》本纪同年条（互校） 的 999年（宋咸平二年、辽统和17年；公元换算据两国年号对照） 条目看，〔府州言官軍入契丹五合川，拔黃太尉砦，殲其眾，焚其車帳，獲馬牛萬計〕使 宋与辽 的一个具体行动进入叙事；题名保留的城、路、水道或机构名称，是目前可以落地的空间线索。

它承接的条件是：本纪年条提供日期和中枢可见行动；前因不据单条补定。 随后可追踪的变化为：可回查同年及后续条目，地方影响仍须另证。

关于可说范围，须保留：《宋史》为元末官修，本纪保存宋廷纪年与处置；战果、数字、地方执行及对方动机须以编年、会要、对方史料或地方材料互校。 进一步讨论可参照 宋辽边政、使节记录与契丹碑刻研究

\noindent\eventanchor{SYD-999-LIAO-098}{冀州言敗契丹兵於城南，殺千餘人，奪馬百餘匹}\textbf{宋与辽。}
999年（宋咸平二年、辽统和17年；公元换算据两国年号对照） 的记载将 宋与辽 与〔冀州言敗契丹兵於城南，殺千餘人，奪馬百餘匹〕相连。此处可确认的是这项被记录的行动；题名保留的城、路、水道或机构名称，是目前可以落地的空间线索。

前一层压力可由以下线索辨认：本纪从朝廷视角记录军政行动；出兵与战果需分辨。 这一处置留下的后续条件是：后续路线、控制与损失应与对方记录和遗址材料复核。

证据的边界是：《宋史》为元末官修，本纪保存宋廷纪年与处置；战果、数字、地方执行及对方动机须以编年、会要、对方史料或地方材料互校。 宋辽边政、使节记录与契丹碑刻研究 可用于继续检验这一结论。

\subsection{982年}
\noindent\eventanchor{XIA-0982-EVENT-049}{李继迁脱离宋朝羁縻控制，在西北边地集结力量}\textbf{党项李氏与宋。}
〔李继迁脱离宋朝羁縻控制，在西北边地集结力量〕见于 《宋史》〈夏国传〉及同年本纪；《辽史》《金史》或《元史》相关条；西夏文文书、碑刻与城址材料（据事核） 的 982年 记录。党项李氏与宋 的选择因此有了可回查的时间位置；材料未将地点细化到城址，不能用中枢所在地替它补足。

“李继迁脱离宋朝羁縻控制，在西北边地集结力量”发生在党项李氏在宋朝羁縻边地累积军政资源的具体压力与机会之中，相关行动者的选择不能脱离此前事件链解释。 记录所见的结果则是：“李继迁脱离宋朝羁縻控制，在西北边地集结力量”为后续政治、边防或资源配置留下条件；其社会影响与实际覆盖范围仍须以异类材料限定。 日期相邻不等于因果已经证实。

西夏无存世连续国史，汉文叙事多出自对手；西夏文文书的释读、定年和地域代表性须逐项说明。 因此，西夏文文书、河西—宁夏考古与宋夏金蒙关系研究 只能扩展问题，不能代替未保存的细节。

\subsection{990年}
\noindent\eventanchor{XIA-0990-EVENT-050}{李继迁在战与和之间接受宋廷封授并保持独立军事行动}\textbf{党项李氏与宋。}
990年，《宋史》〈夏国传〉及同年本纪；《辽史》《金史》或《元史》相关条；西夏文文书、碑刻与城址材料（据事核） 所见的〔李继迁在战与和之间接受宋廷封授并保持独立军事行动〕把 党项李氏与宋 的处置留在可复核的年序中；材料未将地点细化到城址，不能用中枢所在地替它补足。

把本项放回事件链，能够看到的前提为：党项李氏在宋朝羁縻边地累积军政资源中既有的联盟、交通与补给压力，为“李继迁在战与和之间接受宋廷封授并保持独立军事行动”提供了直接的军事或动员背景。 而它所打开或收紧的下一步是：“李继迁在战与和之间接受宋廷封授并保持独立军事行动”改变了相关城镇、道路或政权间的军事选择；伤亡、俘获和控制范围仅可按各方材料分别确认。

关于可说范围，须保留：西夏无存世连续国史，汉文叙事多出自对手；西夏文文书的释读、定年和地域代表性须逐项说明。 进一步讨论可参照 西夏文文书、河西—宁夏考古与宋夏金蒙关系研究

\subsection{1000 年（共享年序桥接）}
\noindent\eventanchor{SY-1000-CHRONOLOGY-BRIDGE}{【C级年序桥接】保留1000 年的多政权并行检索入口；本条不主张所列政权在当年发生直接互动，具体事件须另以单项材料入账。}\textbf{宋、辽与北汉／高丽（年序桥接）。}
宋、辽与北汉／高丽（年序桥接） 在 1000 年（共享年序桥接） 的材料痕迹是〔【C级年序桥接】保留1000 年的多政权并行检索入口；本条不主张所列政权在当年发生直接互动，具体事件须另以单项材料入账。〕，出处为 《宋史》；《辽史》；《续资治通鉴长编》相关年条（据事核）。题名保留的城、路、水道或机构名称，是目前可以落地的空间线索。

它承接的条件是：相邻已入账事件之间缺少该年度的共享年序入口，易使跨政权材料被单一政权叙事遮蔽。 随后可追踪的变化为：保留该年度的检索与补账位置；后续仅在获得可核单项材料时拆分为具体政治、财政、军事、社会或文化事件。

证据的边界是：本条只确认该年位于可追踪的共享年序中；正史、地方文书和外部材料的保存密度不一，不能由桥接标签推断行动、统属、疆界或社会影响。 北宋、辽与东北亚并行年序及边境网络研究 可用于继续检验这一结论。

\subsection{1000年（宋咸平三年、辽统和18年；公元换算据两国年号对照）}
\noindent\eventanchor{SYD-1000-LIAO-023}{甲子，契丹稅木監使黃顒等率屬內附，賜冠帶}\textbf{宋与辽。}
从 《宋史》卷006〈咸平三年〉；《辽史》本纪同年条（互校） 的 1000年（宋咸平三年、辽统和18年；公元换算据两国年号对照） 条目看，〔甲子，契丹稅木監使黃顒等率屬內附，賜冠帶〕使 宋与辽 的一个具体行动进入叙事；材料未将地点细化到城址，不能用中枢所在地替它补足。

前一层压力可由以下线索辨认：本纪所记财用、赈济或征发属于特定报告场景。 这一处置留下的后续条件是：可与食货志、文书和物质材料比较，不外推为全境总量。

《宋史》为元末官修，本纪保存宋廷纪年与处置；战果、数字、地方执行及对方动机须以编年、会要、对方史料或地方材料互校。 因此，宋辽边政、使节记录与契丹碑刻研究 只能扩展问题，不能代替未保存的细节。

\noindent\eventanchor{SYD-1000-LIAO-099}{九月庚辰，賜契丹降人蕭肯頭名懷忠，為右領軍衛將軍、嚴州刺史}\textbf{宋与辽。}
1000年（宋咸平三年、辽统和18年；公元换算据两国年号对照） 的记载将 宋与辽 与〔九月庚辰，賜契丹降人蕭肯頭名懷忠，為右領軍衛將軍、嚴州刺史〕相连。此处可确认的是这项被记录的行动；题名保留的城、路、水道或机构名称，是目前可以落地的空间线索。

本纪从朝廷视角记录军政行动；出兵与战果需分辨。 记录所见的结果则是：后续路线、控制与损失应与对方记录和遗址材料复核。 日期相邻不等于因果已经证实。

关于可说范围，须保留：《宋史》为元末官修，本纪保存宋廷纪年与处置；战果、数字、地方执行及对方动机须以编年、会要、对方史料或地方材料互校。 进一步讨论可参照 宋辽边政、使节记录与契丹碑刻研究

\subsection{1001 年（共享年序桥接）}
\noindent\eventanchor{SY-1001-CHRONOLOGY-BRIDGE}{【C级年序桥接】保留1001 年的多政权并行检索入口；本条不主张所列政权在当年发生直接互动，具体事件须另以单项材料入账。}\textbf{宋、辽与北汉／高丽（年序桥接）。}
〔【C级年序桥接】保留1001 年的多政权并行检索入口；本条不主张所列政权在当年发生直接互动，具体事件须另以单项材料入账。〕见于 《宋史》；《辽史》；《续资治通鉴长编》相关年条（据事核） 的 1001 年（共享年序桥接） 记录。宋、辽与北汉／高丽（年序桥接） 的选择因此有了可回查的时间位置；题名保留的城、路、水道或机构名称，是目前可以落地的空间线索。

把本项放回事件链，能够看到的前提为：相邻已入账事件之间缺少该年度的共享年序入口，易使跨政权材料被单一政权叙事遮蔽。 而它所打开或收紧的下一步是：保留该年度的检索与补账位置；后续仅在获得可核单项材料时拆分为具体政治、财政、军事、社会或文化事件。

证据的边界是：本条只确认该年位于可追踪的共享年序中；正史、地方文书和外部材料的保存密度不一，不能由桥接标签推断行动、统属、疆界或社会影响。 北宋、辽与东北亚并行年序及边境网络研究 可用于继续检验这一结论。

\subsection{1001年（宋咸平四年、辽统和19年；公元换算据两国年号对照）}
\noindent\eventanchor{SYD-1001-LIAO-024}{己卯，邊臣言契丹謀入寇}\textbf{宋与辽。}
1001年（宋咸平四年、辽统和19年；公元换算据两国年号对照），《宋史》卷006〈咸平四年〉；《辽史》本纪同年条（互校） 所见的〔己卯，邊臣言契丹謀入寇〕把 宋与辽 的处置留在可复核的年序中；材料未将地点细化到城址，不能用中枢所在地替它补足。

它承接的条件是：本纪从朝廷视角记录军政行动；出兵与战果需分辨。 随后可追踪的变化为：后续路线、控制与损失应与对方记录和遗址材料复核。

《宋史》为元末官修，本纪保存宋廷纪年与处置；战果、数字、地方执行及对方动机须以编年、会要、对方史料或地方材料互校。 因此，宋辽边政、使节记录与契丹碑刻研究 只能扩展问题，不能代替未保存的细节。

\noindent\eventanchor{SYD-1001-LIAO-100}{己未，張斌破契丹于長城口}\textbf{宋与辽。}
宋与辽 在 1001年（宋咸平四年、辽统和19年；公元换算据两国年号对照） 的材料痕迹是〔己未，張斌破契丹于長城口〕，出处为 《宋史》卷006〈咸平四年〉；《辽史》本纪同年条（互校）。题名保留的城、路、水道或机构名称，是目前可以落地的空间线索。

前一层压力可由以下线索辨认：本纪从朝廷视角记录军政行动；出兵与战果需分辨。 这一处置留下的后续条件是：后续路线、控制与损失应与对方记录和遗址材料复核。

关于可说范围，须保留：《宋史》为元末官修，本纪保存宋廷纪年与处置；战果、数字、地方执行及对方动机须以编年、会要、对方史料或地方材料互校。 进一步讨论可参照 宋辽边政、使节记录与契丹碑刻研究

\subsection{1001年（宋咸平四年、西夏天授礼法延祚-36年；公元换算据两国年号对照）}
\noindent\eventanchor{SYD-1001-XIA-007}{壬午，靈州言河外砦主李瓊等以城降西夏}\textbf{宋与西夏。}
从 《宋史》卷006〈咸平四年〉；西夏文书、碑刻及《辽史》或《金史》同年条（互校） 的 1001年（宋咸平四年、西夏天授礼法延祚-36年；公元换算据两国年号对照） 条目看，〔壬午，靈州言河外砦主李瓊等以城降西夏〕使 宋与西夏 的一个具体行动进入叙事；题名保留的城、路、水道或机构名称，是目前可以落地的空间线索。

本纪保留灾害或工程的报告地点与朝廷处置；灾情范围需另证。 记录所见的结果则是：可与地方文书、河工和环境材料比较，不将单点记录外推为全域因果。 日期相邻不等于因果已经证实。

证据的边界是：《宋史》为元末官修，本纪保存宋廷纪年与处置；战果、数字、地方执行及对方动机须以编年、会要、对方史料或地方材料互校。 西夏文书、河西—宁夏考古与宋夏关系研究 可用于继续检验这一结论。

\subsection{1002 年（共享年序桥接）}
\noindent\eventanchor{SY-1002-CHRONOLOGY-BRIDGE}{【C级年序桥接】保留1002 年的多政权并行检索入口；本条不主张所列政权在当年发生直接互动，具体事件须另以单项材料入账。}\textbf{宋、辽与北汉／高丽（年序桥接）。}
1002 年（共享年序桥接） 的记载将 宋、辽与北汉／高丽（年序桥接） 与〔【C级年序桥接】保留1002 年的多政权并行检索入口；本条不主张所列政权在当年发生直接互动，具体事件须另以单项材料入账。〕相连。此处可确认的是这项被记录的行动；题名保留的城、路、水道或机构名称，是目前可以落地的空间线索。

把本项放回事件链，能够看到的前提为：相邻已入账事件之间缺少该年度的共享年序入口，易使跨政权材料被单一政权叙事遮蔽。 而它所打开或收紧的下一步是：保留该年度的检索与补账位置；后续仅在获得可核单项材料时拆分为具体政治、财政、军事、社会或文化事件。

本条只确认该年位于可追踪的共享年序中；正史、地方文书和外部材料的保存密度不一，不能由桥接标签推断行动、统属、疆界或社会影响。 因此，北宋、辽与东北亚并行年序及边境网络研究 只能扩展问题，不能代替未保存的细节。

\subsection{1002年（宋咸平五年、辽统和20年；公元换算据两国年号对照）}
\noindent\eventanchor{SYD-1002-LIAO-025}{壬戌，契丹大林砦使王昭敏等來降}\textbf{宋与辽。}
〔壬戌，契丹大林砦使王昭敏等來降〕见于 《宋史》卷006〈咸平五年〉；《辽史》本纪同年条（互校） 的 1002年（宋咸平五年、辽统和20年；公元换算据两国年号对照） 记录。宋与辽 的选择因此有了可回查的时间位置；材料未将地点细化到城址，不能用中枢所在地替它补足。

它承接的条件是：本纪从朝廷视角记录军政行动；出兵与战果需分辨。 随后可追踪的变化为：后续路线、控制与损失应与对方记录和遗址材料复核。

关于可说范围，须保留：《宋史》为元末官修，本纪保存宋廷纪年与处置；战果、数字、地方执行及对方动机须以编年、会要、对方史料或地方材料互校。 进一步讨论可参照 宋辽边政、使节记录与契丹碑刻研究

\subsection{1003 年（共享年序桥接）}
\noindent\eventanchor{SY-1003-CHRONOLOGY-BRIDGE}{【C级年序桥接】保留1003 年的多政权并行检索入口；本条不主张所列政权在当年发生直接互动，具体事件须另以单项材料入账。}\textbf{宋、辽与北汉／高丽（年序桥接）。}
1003 年（共享年序桥接），《宋史》；《辽史》；《续资治通鉴长编》相关年条（据事核） 所见的〔【C级年序桥接】保留1003 年的多政权并行检索入口；本条不主张所列政权在当年发生直接互动，具体事件须另以单项材料入账。〕把 宋、辽与北汉／高丽（年序桥接） 的处置留在可复核的年序中；题名保留的城、路、水道或机构名称，是目前可以落地的空间线索。

前一层压力可由以下线索辨认：相邻已入账事件之间缺少该年度的共享年序入口，易使跨政权材料被单一政权叙事遮蔽。 这一处置留下的后续条件是：保留该年度的检索与补账位置；后续仅在获得可核单项材料时拆分为具体政治、财政、军事、社会或文化事件。

证据的边界是：本条只确认该年位于可追踪的共享年序中；正史、地方文书和外部材料的保存密度不一，不能由桥接标签推断行动、统属、疆界或社会影响。 北宋、辽与东北亚并行年序及边境网络研究 可用于继续检验这一结论。

\subsection{1003年（宋咸平六年、辽统和21年；公元换算据两国年号对照）}
\noindent\eventanchor{SYD-1003-LIAO-026}{契丹來侵，戰望都縣，副都部署王繼忠陷於敵，發河東廣銳兵赴援}\textbf{宋与辽。}
宋与辽 在 1003年（宋咸平六年、辽统和21年；公元换算据两国年号对照） 的材料痕迹是〔契丹來侵，戰望都縣，副都部署王繼忠陷於敵，發河東廣銳兵赴援〕，出处为 《宋史》卷007〈咸平六年〉；《辽史》本纪同年条（互校）。题名保留的城、路、水道或机构名称，是目前可以落地的空间线索。

本纪保留灾害或工程的报告地点与朝廷处置；灾情范围需另证。 记录所见的结果则是：可与地方文书、河工和环境材料比较，不将单点记录外推为全域因果。 日期相邻不等于因果已经证实。

《宋史》为元末官修，本纪保存宋廷纪年与处置；战果、数字、地方执行及对方动机须以编年、会要、对方史料或地方材料互校。 因此，宋辽边政、使节记录与契丹碑刻研究 只能扩展问题，不能代替未保存的细节。

\subsection{1004年（辽统和二十二年、宋景德元年；干支、月日待核；公元换算据两朝本纪年号对照）}
\noindent\eventanchor{LIAO-1004-EVENT-037}{辽圣宗、萧太后率军南下至澶州一线}\textbf{辽与宋。}
从 《辽史》相关本纪及志；《宋史》辽国传、同年本纪；宋使行录、契丹／汉文碑刻与考古材料（据事核） 的 1004年（辽统和二十二年、宋景德元年；干支、月日待核；公元换算据两朝本纪年号对照） 条目看，〔辽圣宗、萧太后率军南下至澶州一线〕使 辽与宋 的一个具体行动进入叙事；题名保留的城、路、水道或机构名称，是目前可以落地的空间线索。

把本项放回事件链，能够看到的前提为：“辽圣宗、萧太后率军南下至澶州一线”发生在燕云城镇与辽宋边境秩序的具体压力与机会之中，相关行动者的选择不能脱离此前事件链解释。 而它所打开或收紧的下一步是：“辽圣宗、萧太后率军南下至澶州一线”为后续政治、边防或资源配置留下条件；其社会影响与实际覆盖范围仍须以异类材料限定。

关于可说范围，须保留：《辽史》为元末官修且契丹本方连续文书有限；宋使与碑刻的观察位置不同，战果、族称及制度覆盖范围不得互推。 进一步讨论可参照 辽代契丹国家、都城、部族与跨境关系研究

\subsection{1004年（宋景德元年、辽统和22年；公元换算据两国年号对照）}
\noindent\eventanchor{SYD-1004-LIAO-027}{契丹兵至澶州北，直犯前軍西陣，其大帥撻覽耀兵出陣，俄中伏弩死}\textbf{宋与辽。}
1004年（宋景德元年、辽统和22年；公元换算据两国年号对照） 的记载将 宋与辽 与〔契丹兵至澶州北，直犯前軍西陣，其大帥撻覽耀兵出陣，俄中伏弩死〕相连。此处可确认的是这项被记录的行动；题名保留的城、路、水道或机构名称，是目前可以落地的空间线索。

它承接的条件是：本纪年条提供日期和中枢可见行动；前因不据单条补定。 随后可追踪的变化为：可回查同年及后续条目，地方影响仍须另证。

证据的边界是：《宋史》为元末官修，本纪保存宋廷纪年与处置；战果、数字、地方执行及对方动机须以编年、会要、对方史料或地方材料互校。 宋辽边政、使节记录与契丹碑刻研究 可用于继续检验这一结论。

\noindent\eventanchor{SYD-1004-LIAO-101}{保、莫州、威虜、岢嵐軍及北平砦皆擊敗契丹}\textbf{宋与辽。}
〔保、莫州、威虜、岢嵐軍及北平砦皆擊敗契丹〕见于 《宋史》卷007〈景德元年〉；《辽史》本纪同年条（互校） 的 1004年（宋景德元年、辽统和22年；公元换算据两国年号对照） 记录。宋与辽 的选择因此有了可回查的时间位置；题名保留的城、路、水道或机构名称，是目前可以落地的空间线索。

前一层压力可由以下线索辨认：本纪从朝廷视角记录军政行动；出兵与战果需分辨。 这一处置留下的后续条件是：后续路线、控制与损失应与对方记录和遗址材料复核。

《宋史》为元末官修，本纪保存宋廷纪年与处置；战果、数字、地方执行及对方动机须以编年、会要、对方史料或地方材料互校。 因此，宋辽边政、使节记录与契丹碑刻研究 只能扩展问题，不能代替未保存的细节。

\subsection{1004年}
\noindent\eventanchor{XIA-1004-EVENT-051}{李继迁去世，李德明继承部众}\textbf{党项李氏。}
1004年，《宋史》〈夏国传〉及同年本纪；《辽史》《金史》或《元史》相关条；西夏文文书、碑刻与城址材料（据事核） 所见的〔李继迁去世，李德明继承部众〕把 党项李氏 的处置留在可复核的年序中；材料未将地点细化到城址，不能用中枢所在地替它补足。

党项李氏在宋朝羁縻边地累积军政资源中前任统治的结束与宫廷、部族或官僚的权力安排相互交织，促成“李继迁去世，李德明继承部众”这一继承节点。 记录所见的结果则是：继承并不自动消除既有矛盾；“李继迁去世，李德明继承部众”后的任官、军权和对外策略需由后续条目追踪。 日期相邻不等于因果已经证实。

关于可说范围，须保留：西夏无存世连续国史，汉文叙事多出自对手；西夏文文书的释读、定年和地域代表性须逐项说明。 进一步讨论可参照 西夏文文书、河西—宁夏考古与宋夏金蒙关系研究

\subsection{1005年（辽统和二十三年、宋景德二年；干支、月日待核；公元换算据两朝本纪年号对照）}
\noindent\eventanchor{LIAO-1005-EVENT-038}{辽宋以澶渊盟约建立岁币、使节和边境秩序}\textbf{辽与宋。}
辽与宋 在 1005年（辽统和二十三年、宋景德二年；干支、月日待核；公元换算据两朝本纪年号对照） 的材料痕迹是〔辽宋以澶渊盟约建立岁币、使节和边境秩序〕，出处为 《辽史》相关本纪及志；《宋史》辽国传、同年本纪；宋使行录、契丹／汉文碑刻与考古材料（据事核）。材料未将地点细化到城址，不能用中枢所在地替它补足。

把本项放回事件链，能够看到的前提为：燕云城镇与辽宋边境秩序里的军事消耗、边境供给与使节往来构成谈判条件，促成“辽宋以澶渊盟约建立岁币、使节和边境秩序”所涉的协议或名义关系。 而它所打开或收紧的下一步是：“辽宋以澶渊盟约建立岁币、使节和边境秩序”重排了贡赐、使节或停战的制度接口；条款是否执行及边地实际控制不能由盟约文字直接推定。

证据的边界是：《辽史》为元末官修且契丹本方连续文书有限；宋使与碑刻的观察位置不同，战果、族称及制度覆盖范围不得互推。 辽代契丹国家、都城、部族与跨境关系研究 可用于继续检验这一结论。

\subsection{1005年（宋景德二年、辽统和23年；公元换算据两国年号对照）}
\noindent\eventanchor{SYD-1005-LIAO-028}{丙戌，遣職方郎中韓國華等使契丹}\textbf{宋与辽。}
从 《宋史》卷007〈景德二年〉；《辽史》本纪同年条（互校） 的 1005年（宋景德二年、辽统和23年；公元换算据两国年号对照） 条目看，〔丙戌，遣職方郎中韓國華等使契丹〕使 宋与辽 的一个具体行动进入叙事；材料未将地点细化到城址，不能用中枢所在地替它补足。

它承接的条件是：本纪记录使节、贡赐或往来，不等于稳定统属关系。 随后可追踪的变化为：可与对方编年和交通、文书材料核对其后续落实。

《宋史》为元末官修，本纪保存宋廷纪年与处置；战果、数字、地方执行及对方动机须以编年、会要、对方史料或地方材料互校。 因此，宋辽边政、使节记录与契丹碑刻研究 只能扩展问题，不能代替未保存的细节。

\noindent\eventanchor{SYD-1005-LIAO-102}{二年春正月庚戌朔，以契丹講和，大赦天下，非故鬬殺、放火、強盜、偽造符印、犯贓官典、十惡至死者，悉除之}\textbf{宋与辽。}
1005年（宋景德二年、辽统和23年；公元换算据两国年号对照） 的记载将 宋与辽 与〔二年春正月庚戌朔，以契丹講和，大赦天下，非故鬬殺、放火、強盜、偽造符印、犯贓官典、十惡至死者，悉除之〕相连。此处可确认的是这项被记录的行动；材料未将地点细化到城址，不能用中枢所在地替它补足。

前一层压力可由以下线索辨认：本纪年条记录政令或机构处置；实施背景须参照制度材料。 这一处置留下的后续条件是：可在同年及后续政令中追踪；条文不单独证明地方执行。

关于可说范围，须保留：《宋史》为元末官修，本纪保存宋廷纪年与处置；战果、数字、地方执行及对方动机须以编年、会要、对方史料或地方材料互校。 进一步讨论可参照 宋辽边政、使节记录与契丹碑刻研究

\subsection{1005年（宋景德二年、西夏天授礼法延祚-32年；公元换算据两国年号对照）}
\noindent\eventanchor{SYD-1005-XIA-008}{是歲，夏州、西涼府、邛部川蠻來貢}\textbf{宋与西夏。}
〔是歲，夏州、西涼府、邛部川蠻來貢〕见于 《宋史》卷007〈景德二年〉；西夏文书、碑刻及《辽史》或《金史》同年条（互校） 的 1005年（宋景德二年、西夏天授礼法延祚-32年；公元换算据两国年号对照） 记录。宋与西夏 的选择因此有了可回查的时间位置；题名保留的城、路、水道或机构名称，是目前可以落地的空间线索。

本纪年条提供日期和中枢可见行动；前因不据单条补定。 记录所见的结果则是：可回查同年及后续条目，地方影响仍须另证。 日期相邻不等于因果已经证实。

证据的边界是：《宋史》为元末官修，本纪保存宋廷纪年与处置；战果、数字、地方执行及对方动机须以编年、会要、对方史料或地方材料互校。 西夏文书、河西—宁夏考古与宋夏关系研究 可用于继续检验这一结论。

\subsection{1006 年（共享年序桥接）}
\noindent\eventanchor{SY-1006-CHRONOLOGY-BRIDGE}{【C级年序桥接】保留1006 年的多政权并行检索入口；本条不主张所列政权在当年发生直接互动，具体事件须另以单项材料入账。}\textbf{宋、辽与北汉／高丽（年序桥接）。}
1006 年（共享年序桥接），《宋史》；《辽史》；《续资治通鉴长编》相关年条（据事核） 所见的〔【C级年序桥接】保留1006 年的多政权并行检索入口；本条不主张所列政权在当年发生直接互动，具体事件须另以单项材料入账。〕把 宋、辽与北汉／高丽（年序桥接） 的处置留在可复核的年序中；题名保留的城、路、水道或机构名称，是目前可以落地的空间线索。

把本项放回事件链，能够看到的前提为：相邻已入账事件之间缺少该年度的共享年序入口，易使跨政权材料被单一政权叙事遮蔽。 而它所打开或收紧的下一步是：保留该年度的检索与补账位置；后续仅在获得可核单项材料时拆分为具体政治、财政、军事、社会或文化事件。

本条只确认该年位于可追踪的共享年序中；正史、地方文书和外部材料的保存密度不一，不能由桥接标签推断行动、统属、疆界或社会影响。 因此，北宋、辽与东北亚并行年序及边境网络研究 只能扩展问题，不能代替未保存的细节。

\subsection{1007 年（共享年序桥接）}
\noindent\eventanchor{SY-1007-CHRONOLOGY-BRIDGE}{【C级年序桥接】保留1007 年的多政权并行检索入口；本条不主张所列政权在当年发生直接互动，具体事件须另以单项材料入账。}\textbf{宋、辽与北汉／高丽（年序桥接）。}
宋、辽与北汉／高丽（年序桥接） 在 1007 年（共享年序桥接） 的材料痕迹是〔【C级年序桥接】保留1007 年的多政权并行检索入口；本条不主张所列政权在当年发生直接互动，具体事件须另以单项材料入账。〕，出处为 《宋史》；《辽史》；《续资治通鉴长编》相关年条（据事核）。题名保留的城、路、水道或机构名称，是目前可以落地的空间线索。

它承接的条件是：相邻已入账事件之间缺少该年度的共享年序入口，易使跨政权材料被单一政权叙事遮蔽。 随后可追踪的变化为：保留该年度的检索与补账位置；后续仅在获得可核单项材料时拆分为具体政治、财政、军事、社会或文化事件。

关于可说范围，须保留：本条只确认该年位于可追踪的共享年序中；正史、地方文书和外部材料的保存密度不一，不能由桥接标签推断行动、统属、疆界或社会影响。 进一步讨论可参照 北宋、辽与东北亚并行年序及边境网络研究

\subsection{1007年（宋景德四年、辽统和25年；公元换算据两国年号对照）}
\noindent\eventanchor{SYD-1007-LIAO-029}{十二月己亥，賜近臣、契丹錦綺綾縠等物}\textbf{宋与辽。}
从 《宋史》卷007〈景德四年〉；《辽史》本纪同年条（互校） 的 1007年（宋景德四年、辽统和25年；公元换算据两国年号对照） 条目看，〔十二月己亥，賜近臣、契丹錦綺綾縠等物〕使 宋与辽 的一个具体行动进入叙事；材料未将地点细化到城址，不能用中枢所在地替它补足。

前一层压力可由以下线索辨认：本纪年条提供日期和中枢可见行动；前因不据单条补定。 这一处置留下的后续条件是：可回查同年及后续条目，地方影响仍须另证。

证据的边界是：《宋史》为元末官修，本纪保存宋廷纪年与处置；战果、数字、地方执行及对方动机须以编年、会要、对方史料或地方材料互校。 宋辽边政、使节记录与契丹碑刻研究 可用于继续检验这一结论。

\subsection{1007年（宋景德四年、西夏天授礼法延祚-30年；公元换算据两国年号对照）}
\noindent\eventanchor{SYD-1007-XIA-009}{是歲，河西六谷、夏州、沙州、大食、占城、蒲端國、西南蕃溪峒蠻來貢}\textbf{宋与西夏。}
1007年（宋景德四年、西夏天授礼法延祚-30年；公元换算据两国年号对照） 的记载将 宋与西夏 与〔是歲，河西六谷、夏州、沙州、大食、占城、蒲端國、西南蕃溪峒蠻來貢〕相连。此处可确认的是这项被记录的行动；题名保留的城、路、水道或机构名称，是目前可以落地的空间线索。

本纪保留灾害或工程的报告地点与朝廷处置；灾情范围需另证。 记录所见的结果则是：可与地方文书、河工和环境材料比较，不将单点记录外推为全域因果。 日期相邻不等于因果已经证实。

《宋史》为元末官修，本纪保存宋廷纪年与处置；战果、数字、地方执行及对方动机须以编年、会要、对方史料或地方材料互校。 因此，西夏文书、河西—宁夏考古与宋夏关系研究 只能扩展问题，不能代替未保存的细节。

\subsection{1008年（宋大中祥符元年、辽统和26年；公元换算据两国年号对照）}
\noindent\eventanchor{SYD-1008-LIAO-030}{契丹使上將軍蕭智可等來賀}\textbf{宋与辽。}
〔契丹使上將軍蕭智可等來賀〕见于 《宋史》卷007〈大中祥符元年〉；《辽史》本纪同年条（互校） 的 1008年（宋大中祥符元年、辽统和26年；公元换算据两国年号对照） 记录。宋与辽 的选择因此有了可回查的时间位置；材料未将地点细化到城址，不能用中枢所在地替它补足。

把本项放回事件链，能够看到的前提为：本纪年条提供日期和中枢可见行动；前因不据单条补定。 而它所打开或收紧的下一步是：可回查同年及后续条目，地方影响仍须另证。

关于可说范围，须保留：《宋史》为元末官修，本纪保存宋廷纪年与处置；战果、数字、地方执行及对方动机须以编年、会要、对方史料或地方材料互校。 进一步讨论可参照 宋辽边政、使节记录与契丹碑刻研究

\subsection{1008年（宋大中祥符元年、西夏天授礼法延祚-29年；公元换算据两国年号对照）}
\noindent\eventanchor{SYD-1008-XIA-010}{夏州饑，請易粟，並許之}\textbf{宋与西夏。}
1008年（宋大中祥符元年、西夏天授礼法延祚-29年；公元换算据两国年号对照），《宋史》卷007〈大中祥符元年〉；西夏文书、碑刻及《辽史》或《金史》同年条（互校） 所见的〔夏州饑，請易粟，並許之〕把 宋与西夏 的处置留在可复核的年序中；题名保留的城、路、水道或机构名称，是目前可以落地的空间线索。

它承接的条件是：本纪所记财用、赈济或征发属于特定报告场景。 随后可追踪的变化为：可与食货志、文书和物质材料比较，不外推为全境总量。

证据的边界是：《宋史》为元末官修，本纪保存宋廷纪年与处置；战果、数字、地方执行及对方动机须以编年、会要、对方史料或地方材料互校。 西夏文书、河西—宁夏考古与宋夏关系研究 可用于继续检验这一结论。

\subsection{1009年（宋大中祥符二年、辽统和27年；公元换算据两国年号对照）}
\noindent\eventanchor{SYD-1009-LIAO-031}{命工部侍郎馮起為契丹國信使}\textbf{宋与辽。}
宋与辽 在 1009年（宋大中祥符二年、辽统和27年；公元换算据两国年号对照） 的材料痕迹是〔命工部侍郎馮起為契丹國信使〕，出处为 《宋史》卷007〈大中祥符二年〉；《辽史》本纪同年条（互校）。材料未将地点细化到城址，不能用中枢所在地替它补足。

前一层压力可由以下线索辨认：本纪年条记录政令或机构处置；实施背景须参照制度材料。 这一处置留下的后续条件是：可在同年及后续政令中追踪；条文不单独证明地方执行。

《宋史》为元末官修，本纪保存宋廷纪年与处置；战果、数字、地方执行及对方动机须以编年、会要、对方史料或地方材料互校。 因此，宋辽边政、使节记录与契丹碑刻研究 只能扩展问题，不能代替未保存的细节。

\noindent\eventanchor{SYD-1009-LIAO-103}{契丹國母蕭氏卒，輟視朝}\textbf{宋与辽。}
从 《宋史》卷007〈大中祥符二年〉；《辽史》本纪同年条（互校） 的 1009年（宋大中祥符二年、辽统和27年；公元换算据两国年号对照） 条目看，〔契丹國母蕭氏卒，輟視朝〕使 宋与辽 的一个具体行动进入叙事；材料未将地点细化到城址，不能用中枢所在地替它补足。

本纪年条提供日期和中枢可见行动；前因不据单条补定。 记录所见的结果则是：可回查同年及后续条目，地方影响仍须另证。 日期相邻不等于因果已经证实。

关于可说范围，须保留：《宋史》为元末官修，本纪保存宋廷纪年与处置；战果、数字、地方执行及对方动机须以编年、会要、对方史料或地方材料互校。 进一步讨论可参照 宋辽边政、使节记录与契丹碑刻研究

\subsection{1009年}
\noindent\eventanchor{XIA-1009-EVENT-052}{李德明接受宋廷名号并经营夏州、灵州一带}\textbf{党项李氏与宋。}
1009年 的记载将 党项李氏与宋 与〔李德明接受宋廷名号并经营夏州、灵州一带〕相连。此处可确认的是这项被记录的行动；题名保留的城、路、水道或机构名称，是目前可以落地的空间线索。

把本项放回事件链，能够看到的前提为：“李德明接受宋廷名号并经营夏州、灵州一带”发生在党项李氏在宋朝羁縻边地累积军政资源的具体压力与机会之中，相关行动者的选择不能脱离此前事件链解释。 而它所打开或收紧的下一步是：“李德明接受宋廷名号并经营夏州、灵州一带”为后续政治、边防或资源配置留下条件；其社会影响与实际覆盖范围仍须以异类材料限定。

证据的边界是：西夏无存世连续国史，汉文叙事多出自对手；西夏文文书的释读、定年和地域代表性须逐项说明。 西夏文文书、河西—宁夏考古与宋夏金蒙关系研究 可用于继续检验这一结论。

\subsection{1010年（宋大中祥符三年、辽统和28年；公元换算据两国年号对照）}
\noindent\eventanchor{SYD-1010-LIAO-032}{甲子，契丹國母葬，廢朝，禁邊城樂}\textbf{宋与辽。}
〔甲子，契丹國母葬，廢朝，禁邊城樂〕见于 《宋史》卷007〈大中祥符三年〉；《辽史》本纪同年条（互校） 的 1010年（宋大中祥符三年、辽统和28年；公元换算据两国年号对照） 记录。宋与辽 的选择因此有了可回查的时间位置；题名保留的城、路、水道或机构名称，是目前可以落地的空间线索。

它承接的条件是：本纪年条记录政令或机构处置；实施背景须参照制度材料。 随后可追踪的变化为：可在同年及后续政令中追踪；条文不单独证明地方执行。

《宋史》为元末官修，本纪保存宋廷纪年与处置；战果、数字、地方执行及对方动机须以编年、会要、对方史料或地方材料互校。 因此，宋辽边政、使节记录与契丹碑刻研究 只能扩展问题，不能代替未保存的细节。

\noindent\eventanchor{SYD-1010-LIAO-104}{六月庚戌，邊臣言契丹饑，來市糴，詔雄州糴粟二萬石振之}\textbf{宋与辽。}
1010年（宋大中祥符三年、辽统和28年；公元换算据两国年号对照），《宋史》卷007〈大中祥符三年〉；《辽史》本纪同年条（互校） 所见的〔六月庚戌，邊臣言契丹饑，來市糴，詔雄州糴粟二萬石振之〕把 宋与辽 的处置留在可复核的年序中；题名保留的城、路、水道或机构名称，是目前可以落地的空间线索。

前一层压力可由以下线索辨认：本纪年条记录政令或机构处置；实施背景须参照制度材料。 这一处置留下的后续条件是：可在同年及后续政令中追踪；条文不单独证明地方执行。

关于可说范围，须保留：《宋史》为元末官修，本纪保存宋廷纪年与处置；战果、数字、地方执行及对方动机须以编年、会要、对方史料或地方材料互校。 进一步讨论可参照 宋辽边政、使节记录与契丹碑刻研究

\subsection{1011 年（共享年序桥接）}
\noindent\eventanchor{SY-1011-CHRONOLOGY-BRIDGE}{【C级年序桥接】保留1011 年的多政权并行检索入口；本条不主张所列政权在当年发生直接互动，具体事件须另以单项材料入账。}\textbf{宋、辽与北汉／高丽（年序桥接）。}
宋、辽与北汉／高丽（年序桥接） 在 1011 年（共享年序桥接） 的材料痕迹是〔【C级年序桥接】保留1011 年的多政权并行检索入口；本条不主张所列政权在当年发生直接互动，具体事件须另以单项材料入账。〕，出处为 《宋史》；《辽史》；《续资治通鉴长编》相关年条（据事核）。题名保留的城、路、水道或机构名称，是目前可以落地的空间线索。

相邻已入账事件之间缺少该年度的共享年序入口，易使跨政权材料被单一政权叙事遮蔽。 记录所见的结果则是：保留该年度的检索与补账位置；后续仅在获得可核单项材料时拆分为具体政治、财政、军事、社会或文化事件。 日期相邻不等于因果已经证实。

证据的边界是：本条只确认该年位于可追踪的共享年序中；正史、地方文书和外部材料的保存密度不一，不能由桥接标签推断行动、统属、疆界或社会影响。 北宋、辽与东北亚并行年序及边境网络研究 可用于继续检验这一结论。

\subsection{1013 年（共享年序桥接）}
\noindent\eventanchor{SY-1013-CHRONOLOGY-BRIDGE}{【C级年序桥接】保留1013 年的多政权并行检索入口；本条不主张所列政权在当年发生直接互动，具体事件须另以单项材料入账。}\textbf{宋、辽与北汉／高丽（年序桥接）。}
从 《宋史》；《辽史》；《续资治通鉴长编》相关年条（据事核） 的 1013 年（共享年序桥接） 条目看，〔【C级年序桥接】保留1013 年的多政权并行检索入口；本条不主张所列政权在当年发生直接互动，具体事件须另以单项材料入账。〕使 宋、辽与北汉／高丽（年序桥接） 的一个具体行动进入叙事；题名保留的城、路、水道或机构名称，是目前可以落地的空间线索。

把本项放回事件链，能够看到的前提为：相邻已入账事件之间缺少该年度的共享年序入口，易使跨政权材料被单一政权叙事遮蔽。 而它所打开或收紧的下一步是：保留该年度的检索与补账位置；后续仅在获得可核单项材料时拆分为具体政治、财政、军事、社会或文化事件。

本条只确认该年位于可追踪的共享年序中；正史、地方文书和外部材料的保存密度不一，不能由桥接标签推断行动、统属、疆界或社会影响。 因此，北宋、辽与东北亚并行年序及边境网络研究 只能扩展问题，不能代替未保存的细节。

\subsection{1014 年（共享年序桥接）}
\noindent\eventanchor{SY-1014-CHRONOLOGY-BRIDGE}{【C级年序桥接】保留1014 年的多政权并行检索入口；本条不主张所列政权在当年发生直接互动，具体事件须另以单项材料入账。}\textbf{宋、辽与北汉／高丽（年序桥接）。}
1014 年（共享年序桥接） 的记载将 宋、辽与北汉／高丽（年序桥接） 与〔【C级年序桥接】保留1014 年的多政权并行检索入口；本条不主张所列政权在当年发生直接互动，具体事件须另以单项材料入账。〕相连。此处可确认的是这项被记录的行动；题名保留的城、路、水道或机构名称，是目前可以落地的空间线索。

它承接的条件是：相邻已入账事件之间缺少该年度的共享年序入口，易使跨政权材料被单一政权叙事遮蔽。 随后可追踪的变化为：保留该年度的检索与补账位置；后续仅在获得可核单项材料时拆分为具体政治、财政、军事、社会或文化事件。

关于可说范围，须保留：本条只确认该年位于可追踪的共享年序中；正史、地方文书和外部材料的保存密度不一，不能由桥接标签推断行动、统属、疆界或社会影响。 进一步讨论可参照 北宋、辽与东北亚并行年序及边境网络研究

\subsection{1014年（宋大中祥符七年、西夏天授礼法延祚-23年；公元换算据两国年号对照）}
\noindent\eventanchor{SYD-1014-XIA-011}{庚申，夏州趙德明遣使詣行闕朝貢}\textbf{宋与西夏。}
〔庚申，夏州趙德明遣使詣行闕朝貢〕见于 《宋史》卷008〈大中祥符七年〉；西夏文书、碑刻及《辽史》或《金史》同年条（互校） 的 1014年（宋大中祥符七年、西夏天授礼法延祚-23年；公元换算据两国年号对照） 记录。宋与西夏 的选择因此有了可回查的时间位置；题名保留的城、路、水道或机构名称，是目前可以落地的空间线索。

前一层压力可由以下线索辨认：本纪记录使节、贡赐或往来，不等于稳定统属关系。 这一处置留下的后续条件是：可与对方编年和交通、文书材料核对其后续落实。

证据的边界是：《宋史》为元末官修，本纪保存宋廷纪年与处置；战果、数字、地方执行及对方动机须以编年、会要、对方史料或地方材料互校。 西夏文书、河西—宁夏考古与宋夏关系研究 可用于继续检验这一结论。

\noindent\eventanchor{SYD-1014-XIA-054}{是歲，夏州、西涼府、高麗、女真來貢}\textbf{宋与西夏。}
1014年（宋大中祥符七年、西夏天授礼法延祚-23年；公元换算据两国年号对照），《宋史》卷008〈大中祥符七年〉；西夏文书、碑刻及《辽史》或《金史》同年条（互校） 所见的〔是歲，夏州、西涼府、高麗、女真來貢〕把 宋与西夏 的处置留在可复核的年序中；题名保留的城、路、水道或机构名称，是目前可以落地的空间线索。

本纪年条提供日期和中枢可见行动；前因不据单条补定。 记录所见的结果则是：可回查同年及后续条目，地方影响仍须另证。 日期相邻不等于因果已经证实。

《宋史》为元末官修，本纪保存宋廷纪年与处置；战果、数字、地方执行及对方动机须以编年、会要、对方史料或地方材料互校。 因此，西夏文书、河西—宁夏考古与宋夏关系研究 只能扩展问题，不能代替未保存的细节。

\subsection{1015 年（共享年序桥接）}
\noindent\eventanchor{SY-1015-CHRONOLOGY-BRIDGE}{【C级年序桥接】保留1015 年的多政权并行检索入口；本条不主张所列政权在当年发生直接互动，具体事件须另以单项材料入账。}\textbf{宋、辽与北汉／高丽（年序桥接）。}
宋、辽与北汉／高丽（年序桥接） 在 1015 年（共享年序桥接） 的材料痕迹是〔【C级年序桥接】保留1015 年的多政权并行检索入口；本条不主张所列政权在当年发生直接互动，具体事件须另以单项材料入账。〕，出处为 《宋史》；《辽史》；《续资治通鉴长编》相关年条（据事核）。题名保留的城、路、水道或机构名称，是目前可以落地的空间线索。

把本项放回事件链，能够看到的前提为：相邻已入账事件之间缺少该年度的共享年序入口，易使跨政权材料被单一政权叙事遮蔽。 而它所打开或收紧的下一步是：保留该年度的检索与补账位置；后续仅在获得可核单项材料时拆分为具体政治、财政、军事、社会或文化事件。

关于可说范围，须保留：本条只确认该年位于可追踪的共享年序中；正史、地方文书和外部材料的保存密度不一，不能由桥接标签推断行动、统属、疆界或社会影响。 进一步讨论可参照 北宋、辽与东北亚并行年序及边境网络研究

\subsection{1015年（宋大中祥符八年、西夏天授礼法延祚-22年；公元换算据两国年号对照）}
\noindent\eventanchor{SYD-1015-XIA-012}{甲寅，唃廝囉聚眾數十萬，請討平夏人以自效}\textbf{宋与西夏。}
从 《宋史》卷008〈大中祥符八年〉；西夏文书、碑刻及《辽史》或《金史》同年条（互校） 的 1015年（宋大中祥符八年、西夏天授礼法延祚-22年；公元换算据两国年号对照） 条目看，〔甲寅，唃廝囉聚眾數十萬，請討平夏人以自效〕使 宋与西夏 的一个具体行动进入叙事；材料未将地点细化到城址，不能用中枢所在地替它补足。

它承接的条件是：本纪年条提供日期和中枢可见行动；前因不据单条补定。 随后可追踪的变化为：可回查同年及后续条目，地方影响仍须另证。

证据的边界是：《宋史》为元末官修，本纪保存宋廷纪年与处置；战果、数字、地方执行及对方动机须以编年、会要、对方史料或地方材料互校。 西夏文书、河西—宁夏考古与宋夏关系研究 可用于继续检验这一结论。

\subsection{1016 年（共享年序桥接）}
\noindent\eventanchor{SY-1016-CHRONOLOGY-BRIDGE}{【C级年序桥接】保留1016 年的多政权并行检索入口；本条不主张所列政权在当年发生直接互动，具体事件须另以单项材料入账。}\textbf{宋、辽与北汉／高丽（年序桥接）。}
1016 年（共享年序桥接） 的记载将 宋、辽与北汉／高丽（年序桥接） 与〔【C级年序桥接】保留1016 年的多政权并行检索入口；本条不主张所列政权在当年发生直接互动，具体事件须另以单项材料入账。〕相连。此处可确认的是这项被记录的行动；题名保留的城、路、水道或机构名称，是目前可以落地的空间线索。

前一层压力可由以下线索辨认：相邻已入账事件之间缺少该年度的共享年序入口，易使跨政权材料被单一政权叙事遮蔽。 这一处置留下的后续条件是：保留该年度的检索与补账位置；后续仅在获得可核单项材料时拆分为具体政治、财政、军事、社会或文化事件。

本条只确认该年位于可追踪的共享年序中；正史、地方文书和外部材料的保存密度不一，不能由桥接标签推断行动、统属、疆界或社会影响。 因此，北宋、辽与东北亚并行年序及边境网络研究 只能扩展问题，不能代替未保存的细节。

\subsection{1016年（宋大中祥符九年、西夏天授礼法延祚-21年；公元换算据两国年号对照）}
\noindent\eventanchor{SYD-1016-XIA-013}{是歲，西蕃宗哥族、邛部山后蠻、夏州、甘州來貢}\textbf{宋与西夏。}
〔是歲，西蕃宗哥族、邛部山后蠻、夏州、甘州來貢〕见于 《宋史》卷008〈大中祥符九年〉；西夏文书、碑刻及《辽史》或《金史》同年条（互校） 的 1016年（宋大中祥符九年、西夏天授礼法延祚-21年；公元换算据两国年号对照） 记录。宋与西夏 的选择因此有了可回查的时间位置；题名保留的城、路、水道或机构名称，是目前可以落地的空间线索。

本纪年条提供日期和中枢可见行动；前因不据单条补定。 记录所见的结果则是：可回查同年及后续条目，地方影响仍须另证。 日期相邻不等于因果已经证实。

关于可说范围，须保留：《宋史》为元末官修，本纪保存宋廷纪年与处置；战果、数字、地方执行及对方动机须以编年、会要、对方史料或地方材料互校。 进一步讨论可参照 西夏文书、河西—宁夏考古与宋夏关系研究

\noindent\eventanchor{SYD-1016-XIA-055}{五月乙巳，邠甯環慶部署王守斌言夏州蕃騎千五百來寇慶州，內屬蕃部擊走之}\textbf{宋与西夏。}
1016年（宋大中祥符九年、西夏天授礼法延祚-21年；公元换算据两国年号对照），《宋史》卷008〈大中祥符九年〉；西夏文书、碑刻及《辽史》或《金史》同年条（互校） 所见的〔五月乙巳，邠甯環慶部署王守斌言夏州蕃騎千五百來寇慶州，內屬蕃部擊走之〕把 宋与西夏 的处置留在可复核的年序中；题名保留的城、路、水道或机构名称，是目前可以落地的空间线索。

把本项放回事件链，能够看到的前提为：本纪从朝廷视角记录军政行动；出兵与战果需分辨。 而它所打开或收紧的下一步是：后续路线、控制与损失应与对方记录和遗址材料复核。

证据的边界是：《宋史》为元末官修，本纪保存宋廷纪年与处置；战果、数字、地方执行及对方动机须以编年、会要、对方史料或地方材料互校。 西夏文书、河西—宁夏考古与宋夏关系研究 可用于继续检验这一结论。

\subsection{1017 年（共享年序桥接）}
\noindent\eventanchor{SY-1017-CHRONOLOGY-BRIDGE}{【C级年序桥接】保留1017 年的多政权并行检索入口；本条不主张所列政权在当年发生直接互动，具体事件须另以单项材料入账。}\textbf{宋、辽与北汉／高丽（年序桥接）。}
宋、辽与北汉／高丽（年序桥接） 在 1017 年（共享年序桥接） 的材料痕迹是〔【C级年序桥接】保留1017 年的多政权并行检索入口；本条不主张所列政权在当年发生直接互动，具体事件须另以单项材料入账。〕，出处为 《宋史》；《辽史》；《续资治通鉴长编》相关年条（据事核）。题名保留的城、路、水道或机构名称，是目前可以落地的空间线索。

它承接的条件是：相邻已入账事件之间缺少该年度的共享年序入口，易使跨政权材料被单一政权叙事遮蔽。 随后可追踪的变化为：保留该年度的检索与补账位置；后续仅在获得可核单项材料时拆分为具体政治、财政、军事、社会或文化事件。

本条只确认该年位于可追踪的共享年序中；正史、地方文书和外部材料的保存密度不一，不能由桥接标签推断行动、统属、疆界或社会影响。 因此，北宋、辽与东北亚并行年序及边境网络研究 只能扩展问题，不能代替未保存的细节。

\subsection{1017年（宋天禧元年、辽开泰6年；公元换算据两国年号对照）}
\noindent\eventanchor{SYD-1017-LIAO-033}{壬戌，契丹使耶律准來賀承天節}\textbf{宋与辽。}
从 《宋史》卷008〈天禧元年〉；《辽史》本纪同年条（互校） 的 1017年（宋天禧元年、辽开泰6年；公元换算据两国年号对照） 条目看，〔壬戌，契丹使耶律准來賀承天節〕使 宋与辽 的一个具体行动进入叙事；材料未将地点细化到城址，不能用中枢所在地替它补足。

前一层压力可由以下线索辨认：本纪年条提供日期和中枢可见行动；前因不据单条补定。 这一处置留下的后续条件是：可回查同年及后续条目，地方影响仍须另证。

关于可说范围，须保留：《宋史》为元末官修，本纪保存宋廷纪年与处置；战果、数字、地方执行及对方动机须以编年、会要、对方史料或地方材料互校。 进一步讨论可参照 宋辽边政、使节记录与契丹碑刻研究

\subsection{1018 年（共享年序桥接）}
\noindent\eventanchor{SY-1018-CHRONOLOGY-BRIDGE}{【C级年序桥接】保留1018 年的多政权并行检索入口；本条不主张所列政权在当年发生直接互动，具体事件须另以单项材料入账。}\textbf{宋、辽与北汉／高丽（年序桥接）。}
1018 年（共享年序桥接） 的记载将 宋、辽与北汉／高丽（年序桥接） 与〔【C级年序桥接】保留1018 年的多政权并行检索入口；本条不主张所列政权在当年发生直接互动，具体事件须另以单项材料入账。〕相连。此处可确认的是这项被记录的行动；题名保留的城、路、水道或机构名称，是目前可以落地的空间线索。

相邻已入账事件之间缺少该年度的共享年序入口，易使跨政权材料被单一政权叙事遮蔽。 记录所见的结果则是：保留该年度的检索与补账位置；后续仅在获得可核单项材料时拆分为具体政治、财政、军事、社会或文化事件。 日期相邻不等于因果已经证实。

证据的边界是：本条只确认该年位于可追踪的共享年序中；正史、地方文书和外部材料的保存密度不一，不能由桥接标签推断行动、统属、疆界或社会影响。 北宋、辽与东北亚并行年序及边境网络研究 可用于继续检验这一结论。

\subsection{1019 年（共享年序桥接）}
\noindent\eventanchor{SY-1019-CHRONOLOGY-BRIDGE}{【C级年序桥接】保留1019 年的多政权并行检索入口；本条不主张所列政权在当年发生直接互动，具体事件须另以单项材料入账。}\textbf{宋、辽与北汉／高丽（年序桥接）。}
〔【C级年序桥接】保留1019 年的多政权并行检索入口；本条不主张所列政权在当年发生直接互动，具体事件须另以单项材料入账。〕见于 《宋史》；《辽史》；《续资治通鉴长编》相关年条（据事核） 的 1019 年（共享年序桥接） 记录。宋、辽与北汉／高丽（年序桥接） 的选择因此有了可回查的时间位置；题名保留的城、路、水道或机构名称，是目前可以落地的空间线索。

把本项放回事件链，能够看到的前提为：相邻已入账事件之间缺少该年度的共享年序入口，易使跨政权材料被单一政权叙事遮蔽。 而它所打开或收紧的下一步是：保留该年度的检索与补账位置；后续仅在获得可核单项材料时拆分为具体政治、财政、军事、社会或文化事件。

本条只确认该年位于可追踪的共享年序中；正史、地方文书和外部材料的保存密度不一，不能由桥接标签推断行动、统属、疆界或社会影响。 因此，北宋、辽与东北亚并行年序及边境网络研究 只能扩展问题，不能代替未保存的细节。

\subsection{1020 年（共享年序桥接）}
\noindent\eventanchor{SY-1020-CHRONOLOGY-BRIDGE}{【C级年序桥接】保留1020 年的多政权并行检索入口；本条不主张所列政权在当年发生直接互动，具体事件须另以单项材料入账。}\textbf{宋、辽与北汉／高丽（年序桥接）。}
1020 年（共享年序桥接），《宋史》；《辽史》；《续资治通鉴长编》相关年条（据事核） 所见的〔【C级年序桥接】保留1020 年的多政权并行检索入口；本条不主张所列政权在当年发生直接互动，具体事件须另以单项材料入账。〕把 宋、辽与北汉／高丽（年序桥接） 的处置留在可复核的年序中；题名保留的城、路、水道或机构名称，是目前可以落地的空间线索。

它承接的条件是：相邻已入账事件之间缺少该年度的共享年序入口，易使跨政权材料被单一政权叙事遮蔽。 随后可追踪的变化为：保留该年度的检索与补账位置；后续仅在获得可核单项材料时拆分为具体政治、财政、军事、社会或文化事件。

关于可说范围，须保留：本条只确认该年位于可追踪的共享年序中；正史、地方文书和外部材料的保存密度不一，不能由桥接标签推断行动、统属、疆界或社会影响。 进一步讨论可参照 北宋、辽与东北亚并行年序及边境网络研究

\subsection{1021 年（共享年序桥接）}
\noindent\eventanchor{SY-1021-CHRONOLOGY-BRIDGE}{【C级年序桥接】保留1021 年的多政权并行检索入口；本条不主张所列政权在当年发生直接互动，具体事件须另以单项材料入账。}\textbf{宋、辽与北汉／高丽（年序桥接）。}
宋、辽与北汉／高丽（年序桥接） 在 1021 年（共享年序桥接） 的材料痕迹是〔【C级年序桥接】保留1021 年的多政权并行检索入口；本条不主张所列政权在当年发生直接互动，具体事件须另以单项材料入账。〕，出处为 《宋史》；《辽史》；《续资治通鉴长编》相关年条（据事核）。题名保留的城、路、水道或机构名称，是目前可以落地的空间线索。

前一层压力可由以下线索辨认：相邻已入账事件之间缺少该年度的共享年序入口，易使跨政权材料被单一政权叙事遮蔽。 这一处置留下的后续条件是：保留该年度的检索与补账位置；后续仅在获得可核单项材料时拆分为具体政治、财政、军事、社会或文化事件。

证据的边界是：本条只确认该年位于可追踪的共享年序中；正史、地方文书和外部材料的保存密度不一，不能由桥接标签推断行动、统属、疆界或社会影响。 北宋、辽与东北亚并行年序及边境网络研究 可用于继续检验这一结论。

\subsection{1022年（宋乾興元年、辽太平2年；公元换算据两国年号对照）}
\noindent\eventanchor{SYD-1022-LIAO-034}{八月壬寅，遣使賀契丹主及其妻生日、正旦}\textbf{宋与辽。}
从 《宋史》卷009〈乾興元年〉；《辽史》本纪同年条（互校） 的 1022年（宋乾興元年、辽太平2年；公元换算据两国年号对照） 条目看，〔八月壬寅，遣使賀契丹主及其妻生日、正旦〕使 宋与辽 的一个具体行动进入叙事；材料未将地点细化到城址，不能用中枢所在地替它补足。

本纪记录使节、贡赐或往来，不等于稳定统属关系。 记录所见的结果则是：可与对方编年和交通、文书材料核对其后续落实。 日期相邻不等于因果已经证实。

《宋史》为元末官修，本纪保存宋廷纪年与处置；战果、数字、地方执行及对方动机须以编年、会要、对方史料或地方材料互校。 因此，宋辽边政、使节记录与契丹碑刻研究 只能扩展问题，不能代替未保存的细节。

\noindent\eventanchor{SYD-1022-LIAO-105}{丙寅，遣使以先帝遺留物遺契丹}\textbf{宋与辽。}
1022年（宋乾興元年、辽太平2年；公元换算据两国年号对照） 的记载将 宋与辽 与〔丙寅，遣使以先帝遺留物遺契丹〕相连。此处可确认的是这项被记录的行动；材料未将地点细化到城址，不能用中枢所在地替它补足。

把本项放回事件链，能够看到的前提为：本纪记录使节、贡赐或往来，不等于稳定统属关系。 而它所打开或收紧的下一步是：可与对方编年和交通、文书材料核对其后续落实。

关于可说范围，须保留：《宋史》为元末官修，本纪保存宋廷纪年与处置；战果、数字、地方执行及对方动机须以编年、会要、对方史料或地方材料互校。 进一步讨论可参照 宋辽边政、使节记录与契丹碑刻研究

\subsection{1023年（宋天聖元年、辽太平3年；公元换算据两国年号对照）}
\noindent\eventanchor{SYD-1023-LIAO-035}{庚午，契丹使初來賀長寧節}\textbf{宋与辽。}
〔庚午，契丹使初來賀長寧節〕见于 《宋史》卷009〈天聖元年〉；《辽史》本纪同年条（互校） 的 1023年（宋天聖元年、辽太平3年；公元换算据两国年号对照） 记录。宋与辽 的选择因此有了可回查的时间位置；材料未将地点细化到城址，不能用中枢所在地替它补足。

它承接的条件是：本纪年条提供日期和中枢可见行动；前因不据单条补定。 随后可追踪的变化为：可回查同年及后续条目，地方影响仍须另证。

证据的边界是：《宋史》为元末官修，本纪保存宋廷纪年与处置；战果、数字、地方执行及对方动机须以编年、会要、对方史料或地方材料互校。 宋辽边政、使节记录与契丹碑刻研究 可用于继续检验这一结论。

\subsection{1024 年（共享年序桥接）}
\noindent\eventanchor{SY-1024-CHRONOLOGY-BRIDGE}{【C级年序桥接】保留1024 年的多政权并行检索入口；本条不主张所列政权在当年发生直接互动，具体事件须另以单项材料入账。}\textbf{宋、辽与北汉／高丽（年序桥接）。}
1024 年（共享年序桥接），《宋史》；《辽史》；《续资治通鉴长编》相关年条（据事核） 所见的〔【C级年序桥接】保留1024 年的多政权并行检索入口；本条不主张所列政权在当年发生直接互动，具体事件须另以单项材料入账。〕把 宋、辽与北汉／高丽（年序桥接） 的处置留在可复核的年序中；题名保留的城、路、水道或机构名称，是目前可以落地的空间线索。

前一层压力可由以下线索辨认：相邻已入账事件之间缺少该年度的共享年序入口，易使跨政权材料被单一政权叙事遮蔽。 这一处置留下的后续条件是：保留该年度的检索与补账位置；后续仅在获得可核单项材料时拆分为具体政治、财政、军事、社会或文化事件。

本条只确认该年位于可追踪的共享年序中；正史、地方文书和外部材料的保存密度不一，不能由桥接标签推断行动、统属、疆界或社会影响。 因此，北宋、辽与东北亚并行年序及边境网络研究 只能扩展问题，不能代替未保存的细节。

\subsection{1025 年（共享年序桥接）}
\noindent\eventanchor{SY-1025-CHRONOLOGY-BRIDGE}{【C级年序桥接】保留1025 年的多政权并行检索入口；本条不主张所列政权在当年发生直接互动，具体事件须另以单项材料入账。}\textbf{宋、辽与北汉／高丽（年序桥接）。}
宋、辽与北汉／高丽（年序桥接） 在 1025 年（共享年序桥接） 的材料痕迹是〔【C级年序桥接】保留1025 年的多政权并行检索入口；本条不主张所列政权在当年发生直接互动，具体事件须另以单项材料入账。〕，出处为 《宋史》；《辽史》；《续资治通鉴长编》相关年条（据事核）。题名保留的城、路、水道或机构名称，是目前可以落地的空间线索。

相邻已入账事件之间缺少该年度的共享年序入口，易使跨政权材料被单一政权叙事遮蔽。 记录所见的结果则是：保留该年度的检索与补账位置；后续仅在获得可核单项材料时拆分为具体政治、财政、军事、社会或文化事件。 日期相邻不等于因果已经证实。

关于可说范围，须保留：本条只确认该年位于可追踪的共享年序中；正史、地方文书和外部材料的保存密度不一，不能由桥接标签推断行动、统属、疆界或社会影响。 进一步讨论可参照 北宋、辽与东北亚并行年序及边境网络研究

\subsection{1026 年（共享年序桥接）}
\noindent\eventanchor{SY-1026-CHRONOLOGY-BRIDGE}{【C级年序桥接】保留1026 年的多政权并行检索入口；本条不主张所列政权在当年发生直接互动，具体事件须另以单项材料入账。}\textbf{宋、辽与北汉／高丽（年序桥接）。}
从 《宋史》；《辽史》；《续资治通鉴长编》相关年条（据事核） 的 1026 年（共享年序桥接） 条目看，〔【C级年序桥接】保留1026 年的多政权并行检索入口；本条不主张所列政权在当年发生直接互动，具体事件须另以单项材料入账。〕使 宋、辽与北汉／高丽（年序桥接） 的一个具体行动进入叙事；题名保留的城、路、水道或机构名称，是目前可以落地的空间线索。

把本项放回事件链，能够看到的前提为：相邻已入账事件之间缺少该年度的共享年序入口，易使跨政权材料被单一政权叙事遮蔽。 而它所打开或收紧的下一步是：保留该年度的检索与补账位置；后续仅在获得可核单项材料时拆分为具体政治、财政、军事、社会或文化事件。

证据的边界是：本条只确认该年位于可追踪的共享年序中；正史、地方文书和外部材料的保存密度不一，不能由桥接标签推断行动、统属、疆界或社会影响。 北宋、辽与东北亚并行年序及边境网络研究 可用于继续检验这一结论。

\subsection{1027 年（共享年序桥接）}
\noindent\eventanchor{SY-1027-CHRONOLOGY-BRIDGE}{【C级年序桥接】保留1027 年的多政权并行检索入口；本条不主张所列政权在当年发生直接互动，具体事件须另以单项材料入账。}\textbf{宋、辽与北汉／高丽（年序桥接）。}
1027 年（共享年序桥接） 的记载将 宋、辽与北汉／高丽（年序桥接） 与〔【C级年序桥接】保留1027 年的多政权并行检索入口；本条不主张所列政权在当年发生直接互动，具体事件须另以单项材料入账。〕相连。此处可确认的是这项被记录的行动；题名保留的城、路、水道或机构名称，是目前可以落地的空间线索。

它承接的条件是：相邻已入账事件之间缺少该年度的共享年序入口，易使跨政权材料被单一政权叙事遮蔽。 随后可追踪的变化为：保留该年度的检索与补账位置；后续仅在获得可核单项材料时拆分为具体政治、财政、军事、社会或文化事件。

本条只确认该年位于可追踪的共享年序中；正史、地方文书和外部材料的保存密度不一，不能由桥接标签推断行动、统属、疆界或社会影响。 因此，北宋、辽与东北亚并行年序及边境网络研究 只能扩展问题，不能代替未保存的细节。

\subsection{1028 年（共享年序桥接）}
\noindent\eventanchor{SY-1028-CHRONOLOGY-BRIDGE}{【C级年序桥接】保留1028 年的多政权并行检索入口；本条不主张所列政权在当年发生直接互动，具体事件须另以单项材料入账。}\textbf{宋、辽与北汉／高丽（年序桥接）。}
〔【C级年序桥接】保留1028 年的多政权并行检索入口；本条不主张所列政权在当年发生直接互动，具体事件须另以单项材料入账。〕见于 《宋史》；《辽史》；《续资治通鉴长编》相关年条（据事核） 的 1028 年（共享年序桥接） 记录。宋、辽与北汉／高丽（年序桥接） 的选择因此有了可回查的时间位置；题名保留的城、路、水道或机构名称，是目前可以落地的空间线索。

前一层压力可由以下线索辨认：相邻已入账事件之间缺少该年度的共享年序入口，易使跨政权材料被单一政权叙事遮蔽。 这一处置留下的后续条件是：保留该年度的检索与补账位置；后续仅在获得可核单项材料时拆分为具体政治、财政、军事、社会或文化事件。

关于可说范围，须保留：本条只确认该年位于可追踪的共享年序中；正史、地方文书和外部材料的保存密度不一，不能由桥接标签推断行动、统属、疆界或社会影响。 进一步讨论可参照 北宋、辽与东北亚并行年序及边境网络研究

\subsection{1029 年（共享年序桥接）}
\noindent\eventanchor{SY-1029-CHRONOLOGY-BRIDGE}{【C级年序桥接】保留1029 年的多政权并行检索入口；本条不主张所列政权在当年发生直接互动，具体事件须另以单项材料入账。}\textbf{宋、辽与北汉／高丽（年序桥接）。}
1029 年（共享年序桥接），《宋史》；《辽史》；《续资治通鉴长编》相关年条（据事核） 所见的〔【C级年序桥接】保留1029 年的多政权并行检索入口；本条不主张所列政权在当年发生直接互动，具体事件须另以单项材料入账。〕把 宋、辽与北汉／高丽（年序桥接） 的处置留在可复核的年序中；题名保留的城、路、水道或机构名称，是目前可以落地的空间线索。

相邻已入账事件之间缺少该年度的共享年序入口，易使跨政权材料被单一政权叙事遮蔽。 记录所见的结果则是：保留该年度的检索与补账位置；后续仅在获得可核单项材料时拆分为具体政治、财政、军事、社会或文化事件。 日期相邻不等于因果已经证实。

证据的边界是：本条只确认该年位于可追踪的共享年序中；正史、地方文书和外部材料的保存密度不一，不能由桥接标签推断行动、统属、疆界或社会影响。 北宋、辽与东北亚并行年序及边境网络研究 可用于继续检验这一结论。

\subsection{1029年（宋天聖七年、辽太平9年；公元换算据两国年号对照）}
\noindent\eventanchor{SYD-1029-LIAO-036}{辛巳，詔契丹饑民所過給米，分送唐、鄧等州，以閒田處之}\textbf{宋与辽。}
宋与辽 在 1029年（宋天聖七年、辽太平9年；公元换算据两国年号对照） 的材料痕迹是〔辛巳，詔契丹饑民所過給米，分送唐、鄧等州，以閒田處之〕，出处为 《宋史》卷009〈天聖七年〉；《辽史》本纪同年条（互校）。题名保留的城、路、水道或机构名称，是目前可以落地的空间线索。

把本项放回事件链，能够看到的前提为：本纪年条记录政令或机构处置；实施背景须参照制度材料。 而它所打开或收紧的下一步是：可在同年及后续政令中追踪；条文不单独证明地方执行。

《宋史》为元末官修，本纪保存宋廷纪年与处置；战果、数字、地方执行及对方动机须以编年、会要、对方史料或地方材料互校。 因此，宋辽边政、使节记录与契丹碑刻研究 只能扩展问题，不能代替未保存的细节。

\subsection{1030 年（共享年序桥接）}
\noindent\eventanchor{SY-1030-CHRONOLOGY-BRIDGE}{【C级年序桥接】保留1030 年的多政权并行检索入口；本条不主张所列政权在当年发生直接互动，具体事件须另以单项材料入账。}\textbf{宋、辽与北汉／高丽（年序桥接）。}
从 《宋史》；《辽史》；《续资治通鉴长编》相关年条（据事核） 的 1030 年（共享年序桥接） 条目看，〔【C级年序桥接】保留1030 年的多政权并行检索入口；本条不主张所列政权在当年发生直接互动，具体事件须另以单项材料入账。〕使 宋、辽与北汉／高丽（年序桥接） 的一个具体行动进入叙事；题名保留的城、路、水道或机构名称，是目前可以落地的空间线索。

它承接的条件是：相邻已入账事件之间缺少该年度的共享年序入口，易使跨政权材料被单一政权叙事遮蔽。 随后可追踪的变化为：保留该年度的检索与补账位置；后续仅在获得可核单项材料时拆分为具体政治、财政、军事、社会或文化事件。

关于可说范围，须保留：本条只确认该年位于可追踪的共享年序中；正史、地方文书和外部材料的保存密度不一，不能由桥接标签推断行动、统属、疆界或社会影响。 进一步讨论可参照 北宋、辽与东北亚并行年序及边境网络研究

\subsection{1031年}
\noindent\eventanchor{LIAO-1031-EVENT-039}{辽圣宗去世，兴宗即位}\textbf{辽。}
1031年 的记载将 辽 与〔辽圣宗去世，兴宗即位〕相连。此处可确认的是这项被记录的行动；材料未将地点细化到城址，不能用中枢所在地替它补足。

前一层压力可由以下线索辨认：辽中期继承及辽宋夏关系中前任统治的结束与宫廷、部族或官僚的权力安排相互交织，促成“辽圣宗去世，兴宗即位”这一继承节点。 这一处置留下的后续条件是：继承并不自动消除既有矛盾；“辽圣宗去世，兴宗即位”后的任官、军权和对外策略需由后续条目追踪。

证据的边界是：《辽史》为元末官修且契丹本方连续文书有限；宋使与碑刻的观察位置不同，战果、族称及制度覆盖范围不得互推。 辽代契丹国家、都城、部族与跨境关系研究 可用于继续检验这一结论。

\subsection{1031年（宋天聖九年、辽景福1年；公元换算据两国年号对照）}
\noindent\eventanchor{SYD-1031-LIAO-037}{秋七月丙午朔，契丹使來告其主隆緒殂，遣使祭奠吊慰，及賀宗真立}\textbf{宋与辽。}
〔秋七月丙午朔，契丹使來告其主隆緒殂，遣使祭奠吊慰，及賀宗真立〕见于 《宋史》卷009〈天聖九年〉；《辽史》本纪同年条（互校） 的 1031年（宋天聖九年、辽景福1年；公元换算据两国年号对照） 记录。宋与辽 的选择因此有了可回查的时间位置；材料未将地点细化到城址，不能用中枢所在地替它补足。

本纪记录使节、贡赐或往来，不等于稳定统属关系。 记录所见的结果则是：可与对方编年和交通、文书材料核对其后续落实。 日期相邻不等于因果已经证实。

《宋史》为元末官修，本纪保存宋廷纪年与处置；战果、数字、地方执行及对方动机须以编年、会要、对方史料或地方材料互校。 因此，宋辽边政、使节记录与契丹碑刻研究 只能扩展问题，不能代替未保存的细节。

\noindent\eventanchor{SYD-1031-LIAO-106}{是歲，契丹主及其國母遣使來致遺留物及謝弔祭}\textbf{宋与辽。}
1031年（宋天聖九年、辽景福1年；公元换算据两国年号对照），《宋史》卷009〈天聖九年〉；《辽史》本纪同年条（互校） 所见的〔是歲，契丹主及其國母遣使來致遺留物及謝弔祭〕把 宋与辽 的处置留在可复核的年序中；材料未将地点细化到城址，不能用中枢所在地替它补足。

把本项放回事件链，能够看到的前提为：本纪记录使节、贡赐或往来，不等于稳定统属关系。 而它所打开或收紧的下一步是：可与对方编年和交通、文书材料核对其后续落实。

关于可说范围，须保留：《宋史》为元末官修，本纪保存宋廷纪年与处置；战果、数字、地方执行及对方动机须以编年、会要、对方史料或地方材料互校。 进一步讨论可参照 宋辽边政、使节记录与契丹碑刻研究

\subsection{1032年}
\noindent\eventanchor{XIA-1032-EVENT-053}{李德明去世，李元昊继位}\textbf{党项李氏。}
党项李氏 在 1032年 的材料痕迹是〔李德明去世，李元昊继位〕，出处为 《宋史》〈夏国传〉及同年本纪；《辽史》《金史》或《元史》相关条；西夏文文书、碑刻与城址材料（据事核）。材料未将地点细化到城址，不能用中枢所在地替它补足。

它承接的条件是：党项李氏在宋朝羁縻边地累积军政资源中前任统治的结束与宫廷、部族或官僚的权力安排相互交织，促成“李德明去世，李元昊继位”这一继承节点。 随后可追踪的变化为：继承并不自动消除既有矛盾；“李德明去世，李元昊继位”后的任官、军权和对外策略需由后续条目追踪。

证据的边界是：西夏无存世连续国史，汉文叙事多出自对手；西夏文文书的释读、定年和地域代表性须逐项说明。 西夏文文书、河西—宁夏考古与宋夏金蒙关系研究 可用于继续检验这一结论。

\subsection{1033年（宋明道二年、辽重熙2年；公元换算据两国年号对照）}
\noindent\eventanchor{SYD-1033-LIAO-038}{八月甲午朔，契丹使來吊慰祭奠}\textbf{宋与辽。}
从 《宋史》卷010〈明道二年〉；《辽史》本纪同年条（互校） 的 1033年（宋明道二年、辽重熙2年；公元换算据两国年号对照） 条目看，〔八月甲午朔，契丹使來吊慰祭奠〕使 宋与辽 的一个具体行动进入叙事；材料未将地点细化到城址，不能用中枢所在地替它补足。

前一层压力可由以下线索辨认：本纪年条提供日期和中枢可见行动；前因不据单条补定。 这一处置留下的后续条件是：可回查同年及后续条目，地方影响仍须另证。

《宋史》为元末官修，本纪保存宋廷纪年与处置；战果、数字、地方执行及对方动机须以编年、会要、对方史料或地方材料互校。 因此，宋辽边政、使节记录与契丹碑刻研究 只能扩展问题，不能代替未保存的细节。

\subsection{1034年}
\noindent\eventanchor{XIA-1034-EVENT-054}{李元昊改国号、年号并改革王权象征}\textbf{党项与宋。}
1034年 的记载将 党项与宋 与〔李元昊改国号、年号并改革王权象征〕相连。此处可确认的是这项被记录的行动；材料未将地点细化到城址，不能用中枢所在地替它补足。

“李元昊改国号、年号并改革王权象征”发生在党项李氏在宋朝羁縻边地累积军政资源的具体压力与机会之中，相关行动者的选择不能脱离此前事件链解释。 记录所见的结果则是：“李元昊改国号、年号并改革王权象征”为后续政治、边防或资源配置留下条件；其社会影响与实际覆盖范围仍须以异类材料限定。 日期相邻不等于因果已经证实。

关于可说范围，须保留：西夏无存世连续国史，汉文叙事多出自对手；西夏文文书的释读、定年和地域代表性须逐项说明。 进一步讨论可参照 西夏文文书、河西—宁夏考古与宋夏金蒙关系研究

\subsection{1035 年（共享年序桥接）}
\noindent\eventanchor{SY-1035-CHRONOLOGY-BRIDGE}{【C级年序桥接】保留1035 年的多政权并行检索入口；本条不主张所列政权在当年发生直接互动，具体事件须另以单项材料入账。}\textbf{宋、辽与北汉／高丽（年序桥接）。}
〔【C级年序桥接】保留1035 年的多政权并行检索入口；本条不主张所列政权在当年发生直接互动，具体事件须另以单项材料入账。〕见于 《宋史》；《辽史》；《续资治通鉴长编》相关年条（据事核） 的 1035 年（共享年序桥接） 记录。宋、辽与北汉／高丽（年序桥接） 的选择因此有了可回查的时间位置；题名保留的城、路、水道或机构名称，是目前可以落地的空间线索。

把本项放回事件链，能够看到的前提为：相邻已入账事件之间缺少该年度的共享年序入口，易使跨政权材料被单一政权叙事遮蔽。 而它所打开或收紧的下一步是：保留该年度的检索与补账位置；后续仅在获得可核单项材料时拆分为具体政治、财政、军事、社会或文化事件。

证据的边界是：本条只确认该年位于可追踪的共享年序中；正史、地方文书和外部材料的保存密度不一，不能由桥接标签推断行动、统属、疆界或社会影响。 北宋、辽与东北亚并行年序及边境网络研究 可用于继续检验这一结论。

\subsection{1036年}
\noindent\eventanchor{XIA-1036-EVENT-055}{党项军向甘州、凉州方向扩张，控制河西要点}\textbf{党项。}
1036年，《宋史》〈夏国传〉及同年本纪；《辽史》《金史》或《元史》相关条；西夏文文书、碑刻与城址材料（据事核） 所见的〔党项军向甘州、凉州方向扩张，控制河西要点〕把 党项 的处置留在可复核的年序中；题名保留的城、路、水道或机构名称，是目前可以落地的空间线索。

它承接的条件是：“党项军向甘州、凉州方向扩张，控制河西要点”发生在党项李氏在宋朝羁縻边地累积军政资源的具体压力与机会之中，相关行动者的选择不能脱离此前事件链解释。 随后可追踪的变化为：“党项军向甘州、凉州方向扩张，控制河西要点”为后续政治、边防或资源配置留下条件；其社会影响与实际覆盖范围仍须以异类材料限定。

西夏无存世连续国史，汉文叙事多出自对手；西夏文文书的释读、定年和地域代表性须逐项说明。 因此，西夏文文书、河西—宁夏考古与宋夏金蒙关系研究 只能扩展问题，不能代替未保存的细节。

\subsection{1037 年（共享年序桥接）}
\noindent\eventanchor{SY-1037-CHRONOLOGY-BRIDGE}{【C级年序桥接】保留1037 年的多政权并行检索入口；本条不主张所列政权在当年发生直接互动，具体事件须另以单项材料入账。}\textbf{宋、辽与北汉／高丽（年序桥接）。}
宋、辽与北汉／高丽（年序桥接） 在 1037 年（共享年序桥接） 的材料痕迹是〔【C级年序桥接】保留1037 年的多政权并行检索入口；本条不主张所列政权在当年发生直接互动，具体事件须另以单项材料入账。〕，出处为 《宋史》；《辽史》；《续资治通鉴长编》相关年条（据事核）。题名保留的城、路、水道或机构名称，是目前可以落地的空间线索。

前一层压力可由以下线索辨认：相邻已入账事件之间缺少该年度的共享年序入口，易使跨政权材料被单一政权叙事遮蔽。 这一处置留下的后续条件是：保留该年度的检索与补账位置；后续仅在获得可核单项材料时拆分为具体政治、财政、军事、社会或文化事件。

关于可说范围，须保留：本条只确认该年位于可追踪的共享年序中；正史、地方文书和外部材料的保存密度不一，不能由桥接标签推断行动、统属、疆界或社会影响。 进一步讨论可参照 北宋、辽与东北亚并行年序及边境网络研究

\subsection{1038年（西夏天授礼法延祚元年、宋宝元元年、辽重熙七年；干支、月日待核；公元换算据各方本纪年号对照）}
\noindent\eventanchor{XIA-1038-EVENT-056}{李元昊称帝，西夏正式建国}\textbf{西夏。}
从 《宋史》〈夏国传〉及同年本纪；《辽史》《金史》或《元史》相关条；西夏文文书、碑刻与城址材料（据事核） 的 1038年（西夏天授礼法延祚元年、宋宝元元年、辽重熙七年；干支、月日待核；公元换算据各方本纪年号对照） 条目看，〔李元昊称帝，西夏正式建国〕使 西夏 的一个具体行动进入叙事；材料未将地点细化到城址，不能用中枢所在地替它补足。

西夏建国与宋辽边境战争中前任统治的结束与宫廷、部族或官僚的权力安排相互交织，促成“李元昊称帝，西夏正式建国”这一继承节点。 记录所见的结果则是：继承并不自动消除既有矛盾；“李元昊称帝，西夏正式建国”后的任官、军权和对外策略需由后续条目追踪。 日期相邻不等于因果已经证实。

证据的边界是：西夏无存世连续国史，汉文叙事多出自对手；西夏文文书的释读、定年和地域代表性须逐项说明。 西夏文文书、河西—宁夏考古与宋夏金蒙关系研究 可用于继续检验这一结论。

\subsection{1039 年（共享年序桥接）}
\noindent\eventanchor{SY-1039-CHRONOLOGY-BRIDGE}{【C级年序桥接】保留1039 年的多政权并行检索入口；本条不主张所列政权在当年发生直接互动，具体事件须另以单项材料入账。}\textbf{宋、辽与西夏（年序桥接）。}
1039 年（共享年序桥接） 的记载将 宋、辽与西夏（年序桥接） 与〔【C级年序桥接】保留1039 年的多政权并行检索入口；本条不主张所列政权在当年发生直接互动，具体事件须另以单项材料入账。〕相连。此处可确认的是这项被记录的行动；题名保留的城、路、水道或机构名称，是目前可以落地的空间线索。

把本项放回事件链，能够看到的前提为：相邻已入账事件之间缺少该年度的共享年序入口，易使跨政权材料被单一政权叙事遮蔽。 而它所打开或收紧的下一步是：保留该年度的检索与补账位置；后续仅在获得可核单项材料时拆分为具体政治、财政、军事、社会或文化事件。

本条只确认该年位于可追踪的共享年序中；正史、地方文书和外部材料的保存密度不一，不能由桥接标签推断行动、统属、疆界或社会影响。 因此，宋辽夏边境、财政与文书材料的并行年序研究 只能扩展问题，不能代替未保存的细节。

\subsection{1040年（宋康定元年、辽重熙9年；公元换算据两国年号对照）}
\noindent\eventanchor{SYD-1040-LIAO-039}{秋七月乙丑，遣使以討元昊告契丹}\textbf{宋与辽。}
〔秋七月乙丑，遣使以討元昊告契丹〕见于 《宋史》卷010〈康定元年〉；《辽史》本纪同年条（互校） 的 1040年（宋康定元年、辽重熙9年；公元换算据两国年号对照） 记录。宋与辽 的选择因此有了可回查的时间位置；材料未将地点细化到城址，不能用中枢所在地替它补足。

它承接的条件是：本纪记录使节、贡赐或往来，不等于稳定统属关系。 随后可追踪的变化为：可与对方编年和交通、文书材料核对其后续落实。

关于可说范围，须保留：《宋史》为元末官修，本纪保存宋廷纪年与处置；战果、数字、地方执行及对方动机须以编年、会要、对方史料或地方材料互校。 进一步讨论可参照 宋辽边政、使节记录与契丹碑刻研究

\noindent\eventanchor{SYD-1040-LIAO-107}{乙未，契丹國母復遣使來賀乾元節}\textbf{宋与辽。}
1040年（宋康定元年、辽重熙9年；公元换算据两国年号对照），《宋史》卷010〈康定元年〉；《辽史》本纪同年条（互校） 所见的〔乙未，契丹國母復遣使來賀乾元節〕把 宋与辽 的处置留在可复核的年序中；材料未将地点细化到城址，不能用中枢所在地替它补足。

前一层压力可由以下线索辨认：本纪年条记录政令或机构处置；实施背景须参照制度材料。 这一处置留下的后续条件是：可在同年及后续政令中追踪；条文不单独证明地方执行。

证据的边界是：《宋史》为元末官修，本纪保存宋廷纪年与处置；战果、数字、地方执行及对方动机须以编年、会要、对方史料或地方材料互校。 宋辽边政、使节记录与契丹碑刻研究 可用于继续检验这一结论。

\subsection{1040年（宋康定元年、西夏天授礼法延祚3年；公元换算据两国年号对照）}
\noindent\eventanchor{SYD-1040-XIA-014}{賜京師、河北、陝西、河東諸軍緡錢，蠲陝西夏稅十之二，減河東所科粟}\textbf{宋与西夏。}
宋与西夏 在 1040年（宋康定元年、西夏天授礼法延祚3年；公元换算据两国年号对照） 的材料痕迹是〔賜京師、河北、陝西、河東諸軍緡錢，蠲陝西夏稅十之二，減河東所科粟〕，出处为 《宋史》卷010〈康定元年〉；西夏文书、碑刻及《辽史》或《金史》同年条（互校）。题名保留的城、路、水道或机构名称，是目前可以落地的空间线索。

本纪所记财用、赈济或征发属于特定报告场景。 记录所见的结果则是：可与食货志、文书和物质材料比较，不外推为全境总量。 日期相邻不等于因果已经证实。

《宋史》为元末官修，本纪保存宋廷纪年与处置；战果、数字、地方执行及对方动机须以编年、会要、对方史料或地方材料互校。 因此，西夏文书、河西—宁夏考古与宋夏关系研究 只能扩展问题，不能代替未保存的细节。

\subsection{1040年}
\noindent\eventanchor{XIA-1040-EVENT-057}{三川口之战中宋军失利}\textbf{西夏与宋。}
从 《宋史》〈夏国传〉及同年本纪；《辽史》《金史》或《元史》相关条；西夏文文书、碑刻与城址材料（据事核） 的 1040年 条目看，〔三川口之战中宋军失利〕使 西夏与宋 的一个具体行动进入叙事；题名保留的城、路、水道或机构名称，是目前可以落地的空间线索。

把本项放回事件链，能够看到的前提为：西夏建国与宋辽边境战争中既有的联盟、交通与补给压力，为“三川口之战中宋军失利”提供了直接的军事或动员背景。 而它所打开或收紧的下一步是：“三川口之战中宋军失利”改变了相关城镇、道路或政权间的军事选择；伤亡、俘获和控制范围仅可按各方材料分别确认。

关于可说范围，须保留：西夏无存世连续国史，汉文叙事多出自对手；西夏文文书的释读、定年和地域代表性须逐项说明。 进一步讨论可参照 西夏文文书、河西—宁夏考古与宋夏金蒙关系研究

\subsection{1041年}
\noindent\eventanchor{XIA-1041-EVENT-058}{好水川之战使宋军再遭挫败}\textbf{西夏与宋。}
1041年 的记载将 西夏与宋 与〔好水川之战使宋军再遭挫败〕相连。此处可确认的是这项被记录的行动；题名保留的城、路、水道或机构名称，是目前可以落地的空间线索。

它承接的条件是：西夏建国与宋辽边境战争中既有的联盟、交通与补给压力，为“好水川之战使宋军再遭挫败”提供了直接的军事或动员背景。 随后可追踪的变化为：“好水川之战使宋军再遭挫败”改变了相关城镇、道路或政权间的军事选择；伤亡、俘获和控制范围仅可按各方材料分别确认。

证据的边界是：西夏无存世连续国史，汉文叙事多出自对手；西夏文文书的释读、定年和地域代表性须逐项说明。 西夏文文书、河西—宁夏考古与宋夏金蒙关系研究 可用于继续检验这一结论。

\subsection{1042年（宋慶曆二年、辽重熙11年；公元换算据两国年号对照）}
\noindent\eventanchor{SYD-1042-LIAO-040}{丙寅，契丹遣使來再致誓書，報徹兵}\textbf{宋与辽。}
〔丙寅，契丹遣使來再致誓書，報徹兵〕见于 《宋史》卷011〈慶曆二年〉；《辽史》本纪同年条（互校） 的 1042年（宋慶曆二年、辽重熙11年；公元换算据两国年号对照） 记录。宋与辽 的选择因此有了可回查的时间位置；材料未将地点细化到城址，不能用中枢所在地替它补足。

前一层压力可由以下线索辨认：本纪记录使节、贡赐或往来，不等于稳定统属关系。 这一处置留下的后续条件是：可与对方编年和交通、文书材料核对其后续落实。

《宋史》为元末官修，本纪保存宋廷纪年与处置；战果、数字、地方执行及对方动机须以编年、会要、对方史料或地方材料互校。 因此，宋辽边政、使节记录与契丹碑刻研究 只能扩展问题，不能代替未保存的细节。

\noindent\eventanchor{SYD-1042-LIAO-108}{己巳，契丹遣蕭英、劉六符來致書求割地}\textbf{宋与辽。}
1042年（宋慶曆二年、辽重熙11年；公元换算据两国年号对照），《宋史》卷011〈慶曆二年〉；《辽史》本纪同年条（互校） 所见的〔己巳，契丹遣蕭英、劉六符來致書求割地〕把 宋与辽 的处置留在可复核的年序中；材料未将地点细化到城址，不能用中枢所在地替它补足。

本纪年条提供日期和中枢可见行动；前因不据单条补定。 记录所见的结果则是：可回查同年及后续条目，地方影响仍须另证。 日期相邻不等于因果已经证实。

关于可说范围，须保留：《宋史》为元末官修，本纪保存宋廷纪年与处置；战果、数字、地方执行及对方动机须以编年、会要、对方史料或地方材料互校。 进一步讨论可参照 宋辽边政、使节记录与契丹碑刻研究

\subsection{1042年}
\noindent\eventanchor{XIA-1042-EVENT-059}{定川寨之战加剧宋西北防务危机}\textbf{西夏与宋。}
西夏与宋 在 1042年 的材料痕迹是〔定川寨之战加剧宋西北防务危机〕，出处为 《宋史》〈夏国传〉及同年本纪；《辽史》《金史》或《元史》相关条；西夏文文书、碑刻与城址材料（据事核）。题名保留的城、路、水道或机构名称，是目前可以落地的空间线索。

把本项放回事件链，能够看到的前提为：西夏建国与宋辽边境战争中既有的联盟、交通与补给压力，为“定川寨之战加剧宋西北防务危机”提供了直接的军事或动员背景。 而它所打开或收紧的下一步是：“定川寨之战加剧宋西北防务危机”改变了相关城镇、道路或政权间的军事选择；伤亡、俘获和控制范围仅可按各方材料分别确认。

证据的边界是：西夏无存世连续国史，汉文叙事多出自对手；西夏文文书的释读、定年和地域代表性须逐项说明。 西夏文文书、河西—宁夏考古与宋夏金蒙关系研究 可用于继续检验这一结论。

\subsection{1043年（宋慶曆三年、西夏天授礼法延祚6年；公元换算据两国年号对照）}
\noindent\eventanchor{SYD-1043-XIA-015}{癸卯，遣保安軍判官邵良佐使元昊，許封冊為夏國主，歲賜絹十萬匹、茶三萬斤}\textbf{宋与西夏。}
从 《宋史》卷011〈慶曆三年〉；西夏文书、碑刻及《辽史》或《金史》同年条（互校） 的 1043年（宋慶曆三年、西夏天授礼法延祚6年；公元换算据两国年号对照） 条目看，〔癸卯，遣保安軍判官邵良佐使元昊，許封冊為夏國主，歲賜絹十萬匹、茶三萬斤〕使 宋与西夏 的一个具体行动进入叙事；材料未将地点细化到城址，不能用中枢所在地替它补足。

它承接的条件是：本纪记录使节、贡赐或往来，不等于稳定统属关系。 随后可追踪的变化为：可与对方编年和交通、文书材料核对其后续落实。

《宋史》为元末官修，本纪保存宋廷纪年与处置；战果、数字、地方执行及对方动机须以编年、会要、对方史料或地方材料互校。 因此，西夏文书、河西—宁夏考古与宋夏关系研究 只能扩展问题，不能代替未保存的细节。

\noindent\eventanchor{SYD-1043-XIA-056}{癸巳，元昊自名曩霄，遣人來納款，稱夏國}\textbf{宋与西夏。}
1043年（宋慶曆三年、西夏天授礼法延祚6年；公元换算据两国年号对照） 的记载将 宋与西夏 与〔癸巳，元昊自名曩霄，遣人來納款，稱夏國〕相连。此处可确认的是这项被记录的行动；材料未将地点细化到城址，不能用中枢所在地替它补足。

前一层压力可由以下线索辨认：本纪年条提供日期和中枢可见行动；前因不据单条补定。 这一处置留下的后续条件是：可回查同年及后续条目，地方影响仍须另证。

关于可说范围，须保留：《宋史》为元末官修，本纪保存宋廷纪年与处置；战果、数字、地方执行及对方动机须以编年、会要、对方史料或地方材料互校。 进一步讨论可参照 西夏文书、河西—宁夏考古与宋夏关系研究

\subsection{1044年（辽重熙十三年、宋庆历四年、西夏天授礼法延祚七年；干支、月日待核；公元换算据各方本纪年号对照）}
\noindent\eventanchor{LIAO-1044-EVENT-040}{宋夏议和期间，辽的压力与外交姿态影响三方关系}\textbf{辽宋西夏。}
〔宋夏议和期间，辽的压力与外交姿态影响三方关系〕见于 《辽史》相关本纪及志；《宋史》辽国传、同年本纪；宋使行录、契丹／汉文碑刻与考古材料（据事核） 的 1044年（辽重熙十三年、宋庆历四年、西夏天授礼法延祚七年；干支、月日待核；公元换算据各方本纪年号对照） 记录。辽宋西夏 的选择因此有了可回查的时间位置；题名保留的城、路、水道或机构名称，是目前可以落地的空间线索。

辽中期继承及辽宋夏关系里的军事消耗、边境供给与使节往来构成谈判条件，促成“宋夏议和期间，辽的压力与外交姿态影响三方关系”所涉的协议或名义关系。 记录所见的结果则是：“宋夏议和期间，辽的压力与外交姿态影响三方关系”重排了贡赐、使节或停战的制度接口；条款是否执行及边地实际控制不能由盟约文字直接推定。 日期相邻不等于因果已经证实。

证据的边界是：《辽史》为元末官修且契丹本方连续文书有限；宋使与碑刻的观察位置不同，战果、族称及制度覆盖范围不得互推。 辽代契丹国家、都城、部族与跨境关系研究 可用于继续检验这一结论。

\subsection{1044年（宋慶曆四年、辽重熙13年；公元换算据两国年号对照）}
\noindent\eventanchor{SYD-1044-LIAO-041}{癸未，契丹遣使來告伐夏國}\textbf{宋与辽。}
1044年（宋慶曆四年、辽重熙13年；公元换算据两国年号对照），《宋史》卷011〈慶曆四年〉；《辽史》本纪同年条（互校） 所见的〔癸未，契丹遣使來告伐夏國〕把 宋与辽 的处置留在可复核的年序中；材料未将地点细化到城址，不能用中枢所在地替它补足。

把本项放回事件链，能够看到的前提为：本纪记录使节、贡赐或往来，不等于稳定统属关系。 而它所打开或收紧的下一步是：可与对方编年和交通、文书材料核对其后续落实。

《宋史》为元末官修，本纪保存宋廷纪年与处置；战果、数字、地方执行及对方动机须以编年、会要、对方史料或地方材料互校。 因此，宋辽边政、使节记录与契丹碑刻研究 只能扩展问题，不能代替未保存的细节。

\noindent\eventanchor{SYD-1044-LIAO-109}{戊戌，命右正言余靖報使契丹}\textbf{宋与辽。}
宋与辽 在 1044年（宋慶曆四年、辽重熙13年；公元换算据两国年号对照） 的材料痕迹是〔戊戌，命右正言余靖報使契丹〕，出处为 《宋史》卷011〈慶曆四年〉；《辽史》本纪同年条（互校）。材料未将地点细化到城址，不能用中枢所在地替它补足。

它承接的条件是：本纪年条记录政令或机构处置；实施背景须参照制度材料。 随后可追踪的变化为：可在同年及后续政令中追踪；条文不单独证明地方执行。

关于可说范围，须保留：《宋史》为元末官修，本纪保存宋廷纪年与处置；战果、数字、地方执行及对方动机须以编年、会要、对方史料或地方材料互校。 进一步讨论可参照 宋辽边政、使节记录与契丹碑刻研究

\subsection{1044年（宋慶曆四年、西夏天授礼法延祚7年；公元换算据两国年号对照）}
\noindent\eventanchor{SYD-1044-XIA-016}{乙未，封曩霄為夏國主}\textbf{宋与西夏。}
从 《宋史》卷011〈慶曆四年〉；西夏文书、碑刻及《辽史》或《金史》同年条（互校） 的 1044年（宋慶曆四年、西夏天授礼法延祚7年；公元换算据两国年号对照） 条目看，〔乙未，封曩霄為夏國主〕使 宋与西夏 的一个具体行动进入叙事；材料未将地点细化到城址，不能用中枢所在地替它补足。

前一层压力可由以下线索辨认：本纪年条提供日期和中枢可见行动；前因不据单条补定。 这一处置留下的后续条件是：可回查同年及后续条目，地方影响仍须另证。

证据的边界是：《宋史》为元末官修，本纪保存宋廷纪年与处置；战果、数字、地方执行及对方动机须以编年、会要、对方史料或地方材料互校。 西夏文书、河西—宁夏考古与宋夏关系研究 可用于继续检验这一结论。

\subsection{1044年（西夏天授礼法延祚七年、宋庆历四年、辽重熙十三年；干支、月日待核；公元换算据各方本纪年号对照）}
\noindent\eventanchor{XIA-1044-EVENT-060}{宋夏达成庆历和议}\textbf{西夏与宋。}
1044年（西夏天授礼法延祚七年、宋庆历四年、辽重熙十三年；干支、月日待核；公元换算据各方本纪年号对照） 的记载将 西夏与宋 与〔宋夏达成庆历和议〕相连。此处可确认的是这项被记录的行动；材料未将地点细化到城址，不能用中枢所在地替它补足。

西夏建国与宋辽边境战争里的军事消耗、边境供给与使节往来构成谈判条件，促成“宋夏达成庆历和议”所涉的协议或名义关系。 记录所见的结果则是：“宋夏达成庆历和议”重排了贡赐、使节或停战的制度接口；条款是否执行及边地实际控制不能由盟约文字直接推定。 日期相邻不等于因果已经证实。

西夏无存世连续国史，汉文叙事多出自对手；西夏文文书的释读、定年和地域代表性须逐项说明。 因此，西夏文文书、河西—宁夏考古与宋夏金蒙关系研究 只能扩展问题，不能代替未保存的细节。

\subsection{1045年（宋慶曆五年、辽重熙14年；公元换算据两国年号对照）}
\noindent\eventanchor{SYD-1045-LIAO-042}{丙子，契丹遣使來告伐夏國還}\textbf{宋与辽。}
〔丙子，契丹遣使來告伐夏國還〕见于 《宋史》卷011〈慶曆五年〉；《辽史》本纪同年条（互校） 的 1045年（宋慶曆五年、辽重熙14年；公元换算据两国年号对照） 记录。宋与辽 的选择因此有了可回查的时间位置；材料未将地点细化到城址，不能用中枢所在地替它补足。

把本项放回事件链，能够看到的前提为：本纪记录使节、贡赐或往来，不等于稳定统属关系。 而它所打开或收紧的下一步是：可与对方编年和交通、文书材料核对其后续落实。

关于可说范围，须保留：《宋史》为元末官修，本纪保存宋廷纪年与处置；战果、数字、地方执行及对方动机须以编年、会要、对方史料或地方材料互校。 进一步讨论可参照 宋辽边政、使节记录与契丹碑刻研究

\noindent\eventanchor{SYD-1045-LIAO-110}{冬十月乙卯，契丹遣使來獻九龍車及所獲夏國羊馬}\textbf{宋与辽。}
1045年（宋慶曆五年、辽重熙14年；公元换算据两国年号对照），《宋史》卷011〈慶曆五年〉；《辽史》本纪同年条（互校） 所见的〔冬十月乙卯，契丹遣使來獻九龍車及所獲夏國羊馬〕把 宋与辽 的处置留在可复核的年序中；材料未将地点细化到城址，不能用中枢所在地替它补足。

它承接的条件是：本纪记录使节、贡赐或往来，不等于稳定统属关系。 随后可追踪的变化为：可与对方编年和交通、文书材料核对其后续落实。

证据的边界是：《宋史》为元末官修，本纪保存宋廷纪年与处置；战果、数字、地方执行及对方动机须以编年、会要、对方史料或地方材料互校。 宋辽边政、使节记录与契丹碑刻研究 可用于继续检验这一结论。

\subsection{1046 年（共享年序桥接）}
\noindent\eventanchor{SY-1046-CHRONOLOGY-BRIDGE}{【C级年序桥接】保留1046 年的多政权并行检索入口；本条不主张所列政权在当年发生直接互动，具体事件须另以单项材料入账。}\textbf{宋、辽与西夏（年序桥接）。}
宋、辽与西夏（年序桥接） 在 1046 年（共享年序桥接） 的材料痕迹是〔【C级年序桥接】保留1046 年的多政权并行检索入口；本条不主张所列政权在当年发生直接互动，具体事件须另以单项材料入账。〕，出处为 《宋史》；《辽史》；《续资治通鉴长编》；西夏文书与碑刻（据事核）。题名保留的城、路、水道或机构名称，是目前可以落地的空间线索。

前一层压力可由以下线索辨认：相邻已入账事件之间缺少该年度的共享年序入口，易使跨政权材料被单一政权叙事遮蔽。 这一处置留下的后续条件是：保留该年度的检索与补账位置；后续仅在获得可核单项材料时拆分为具体政治、财政、军事、社会或文化事件。

本条只确认该年位于可追踪的共享年序中；正史、地方文书和外部材料的保存密度不一，不能由桥接标签推断行动、统属、疆界或社会影响。 因此，宋辽夏边境、财政与文书材料的并行年序研究 只能扩展问题，不能代替未保存的细节。

\subsection{1046年（宋慶曆六年、西夏天授礼法延祚9年；公元换算据两国年号对照）}
\noindent\eventanchor{SYD-1046-XIA-017}{十一月己卯，遣官議夏國公封界}\textbf{宋与西夏。}
从 《宋史》卷011〈慶曆六年〉；西夏文书、碑刻及《辽史》或《金史》同年条（互校） 的 1046年（宋慶曆六年、西夏天授礼法延祚9年；公元换算据两国年号对照） 条目看，〔十一月己卯，遣官議夏國公封界〕使 宋与西夏 的一个具体行动进入叙事；材料未将地点细化到城址，不能用中枢所在地替它补足。

本纪年条提供日期和中枢可见行动；前因不据单条补定。 记录所见的结果则是：可回查同年及后续条目，地方影响仍须另证。 日期相邻不等于因果已经证实。

关于可说范围，须保留：《宋史》为元末官修，本纪保存宋廷纪年与处置；战果、数字、地方执行及对方动机须以编年、会要、对方史料或地方材料互校。 进一步讨论可参照 西夏文书、河西—宁夏考古与宋夏关系研究

\subsection{1048年（宋慶曆八年、西夏天授礼法延祚11年；公元换算据两国年号对照）}
\noindent\eventanchor{SYD-1048-XIA-018}{夏四月己巳朔，封曩霄子諒祚為夏國主}\textbf{宋与西夏。}
1048年（宋慶曆八年、西夏天授礼法延祚11年；公元换算据两国年号对照） 的记载将 宋与西夏 与〔夏四月己巳朔，封曩霄子諒祚為夏國主〕相连。此处可确认的是这项被记录的行动；材料未将地点细化到城址，不能用中枢所在地替它补足。

把本项放回事件链，能够看到的前提为：本纪年条提供日期和中枢可见行动；前因不据单条补定。 而它所打开或收紧的下一步是：可回查同年及后续条目，地方影响仍须另证。

证据的边界是：《宋史》为元末官修，本纪保存宋廷纪年与处置；战果、数字、地方执行及对方动机须以编年、会要、对方史料或地方材料互校。 西夏文书、河西—宁夏考古与宋夏关系研究 可用于继续检验这一结论。

\subsection{1048年}
\noindent\eventanchor{XIA-1048-EVENT-061}{李元昊遇害，幼主继位并出现后族、权臣主政}\textbf{西夏。}
〔李元昊遇害，幼主继位并出现后族、权臣主政〕见于 《宋史》〈夏国传〉及同年本纪；《辽史》《金史》或《元史》相关条；西夏文文书、碑刻与城址材料（据事核） 的 1048年 记录。西夏 的选择因此有了可回查的时间位置；材料未将地点细化到城址，不能用中枢所在地替它补足。

它承接的条件是：西夏建国与宋辽边境战争中前任统治的结束与宫廷、部族或官僚的权力安排相互交织，促成“李元昊遇害，幼主继位并出现后族、权臣主政”这一继承节点。 随后可追踪的变化为：继承并不自动消除既有矛盾；“李元昊遇害，幼主继位并出现后族、权臣主政”后的任官、军权和对外策略需由后续条目追踪。

西夏无存世连续国史，汉文叙事多出自对手；西夏文文书的释读、定年和地域代表性须逐项说明。 因此，西夏文文书、河西—宁夏考古与宋夏金蒙关系研究 只能扩展问题，不能代替未保存的细节。

\subsection{1049年（宋皇祐元年、辽重熙18年；公元换算据两国年号对照）}
\noindent\eventanchor{SYD-1049-LIAO-043}{己未，契丹遣使來告伐夏國}\textbf{宋与辽。}
1049年（宋皇祐元年、辽重熙18年；公元换算据两国年号对照），《宋史》卷011〈皇祐元年〉；《辽史》本纪同年条（互校） 所见的〔己未，契丹遣使來告伐夏國〕把 宋与辽 的处置留在可复核的年序中；材料未将地点细化到城址，不能用中枢所在地替它补足。

前一层压力可由以下线索辨认：本纪记录使节、贡赐或往来，不等于稳定统属关系。 这一处置留下的后续条件是：可与对方编年和交通、文书材料核对其后续落实。

关于可说范围，须保留：《宋史》为元末官修，本纪保存宋廷纪年与处置；战果、数字、地方执行及对方动机须以编年、会要、对方史料或地方材料互校。 进一步讨论可参照 宋辽边政、使节记录与契丹碑刻研究

\subsection{1049年}
\noindent\eventanchor{XIA-1049-EVENT-062}{李谅祚即位}\textbf{西夏。}
西夏 在 1049年 的材料痕迹是〔李谅祚即位〕，出处为 《宋史》〈夏国传〉及同年本纪；《辽史》《金史》或《元史》相关条；西夏文文书、碑刻与城址材料（据事核）。材料未将地点细化到城址，不能用中枢所在地替它补足。

西夏建国与宋辽边境战争中前任统治的结束与宫廷、部族或官僚的权力安排相互交织，促成“李谅祚即位”这一继承节点。 记录所见的结果则是：继承并不自动消除既有矛盾；“李谅祚即位”后的任官、军权和对外策略需由后续条目追踪。 日期相邻不等于因果已经证实。

证据的边界是：西夏无存世连续国史，汉文叙事多出自对手；西夏文文书的释读、定年和地域代表性须逐项说明。 西夏文文书、河西—宁夏考古与宋夏金蒙关系研究 可用于继续检验这一结论。

\subsection{1050 年（共享年序桥接）}
\noindent\eventanchor{SY-1050-CHRONOLOGY-BRIDGE}{【C级年序桥接】保留1050 年的多政权并行检索入口；本条不主张所列政权在当年发生直接互动，具体事件须另以单项材料入账。}\textbf{宋、辽与西夏（年序桥接）。}
从 《宋史》；《辽史》；《续资治通鉴长编》；西夏文书与碑刻（据事核） 的 1050 年（共享年序桥接） 条目看，〔【C级年序桥接】保留1050 年的多政权并行检索入口；本条不主张所列政权在当年发生直接互动，具体事件须另以单项材料入账。〕使 宋、辽与西夏（年序桥接） 的一个具体行动进入叙事；题名保留的城、路、水道或机构名称，是目前可以落地的空间线索。

把本项放回事件链，能够看到的前提为：相邻已入账事件之间缺少该年度的共享年序入口，易使跨政权材料被单一政权叙事遮蔽。 而它所打开或收紧的下一步是：保留该年度的检索与补账位置；后续仅在获得可核单项材料时拆分为具体政治、财政、军事、社会或文化事件。

本条只确认该年位于可追踪的共享年序中；正史、地方文书和外部材料的保存密度不一，不能由桥接标签推断行动、统属、疆界或社会影响。 因此，宋辽夏边境、财政与文书材料的并行年序研究 只能扩展问题，不能代替未保存的细节。

\subsection{1050年（宋皇祐二年、辽重熙19年；公元换算据两国年号对照）}
\noindent\eventanchor{SYD-1050-LIAO-044}{庚子，契丹遣使以伐夏師還來告}\textbf{宋与辽。}
1050年（宋皇祐二年、辽重熙19年；公元换算据两国年号对照） 的记载将 宋与辽 与〔庚子，契丹遣使以伐夏師還來告〕相连。此处可确认的是这项被记录的行动；材料未将地点细化到城址，不能用中枢所在地替它补足。

它承接的条件是：本纪记录使节、贡赐或往来，不等于稳定统属关系。 随后可追踪的变化为：可与对方编年和交通、文书材料核对其后续落实。

关于可说范围，须保留：《宋史》为元末官修，本纪保存宋廷纪年与处置；战果、数字、地方执行及对方动机须以编年、会要、对方史料或地方材料互校。 进一步讨论可参照 宋辽边政、使节记录与契丹碑刻研究

\subsection{1051 年（共享年序桥接）}
\noindent\eventanchor{SY-1051-CHRONOLOGY-BRIDGE}{【C级年序桥接】保留1051 年的多政权并行检索入口；本条不主张所列政权在当年发生直接互动，具体事件须另以单项材料入账。}\textbf{宋、辽与西夏（年序桥接）。}
〔【C级年序桥接】保留1051 年的多政权并行检索入口；本条不主张所列政权在当年发生直接互动，具体事件须另以单项材料入账。〕见于 《宋史》；《辽史》；《续资治通鉴长编》；西夏文书与碑刻（据事核） 的 1051 年（共享年序桥接） 记录。宋、辽与西夏（年序桥接） 的选择因此有了可回查的时间位置；题名保留的城、路、水道或机构名称，是目前可以落地的空间线索。

前一层压力可由以下线索辨认：相邻已入账事件之间缺少该年度的共享年序入口，易使跨政权材料被单一政权叙事遮蔽。 这一处置留下的后续条件是：保留该年度的检索与补账位置；后续仅在获得可核单项材料时拆分为具体政治、财政、军事、社会或文化事件。

证据的边界是：本条只确认该年位于可追踪的共享年序中；正史、地方文书和外部材料的保存密度不一，不能由桥接标签推断行动、统属、疆界或社会影响。 宋辽夏边境、财政与文书材料的并行年序研究 可用于继续检验这一结论。

\subsection{1054 年（共享年序桥接）}
\noindent\eventanchor{SY-1054-CHRONOLOGY-BRIDGE}{【C级年序桥接】保留1054 年的多政权并行检索入口；本条不主张所列政权在当年发生直接互动，具体事件须另以单项材料入账。}\textbf{宋、辽与西夏（年序桥接）。}
1054 年（共享年序桥接），《宋史》；《辽史》；《续资治通鉴长编》；西夏文书与碑刻（据事核） 所见的〔【C级年序桥接】保留1054 年的多政权并行检索入口；本条不主张所列政权在当年发生直接互动，具体事件须另以单项材料入账。〕把 宋、辽与西夏（年序桥接） 的处置留在可复核的年序中；题名保留的城、路、水道或机构名称，是目前可以落地的空间线索。

相邻已入账事件之间缺少该年度的共享年序入口，易使跨政权材料被单一政权叙事遮蔽。 记录所见的结果则是：保留该年度的检索与补账位置；后续仅在获得可核单项材料时拆分为具体政治、财政、军事、社会或文化事件。 日期相邻不等于因果已经证实。

本条只确认该年位于可追踪的共享年序中；正史、地方文书和外部材料的保存密度不一，不能由桥接标签推断行动、统属、疆界或社会影响。 因此，宋辽夏边境、财政与文书材料的并行年序研究 只能扩展问题，不能代替未保存的细节。

\subsection{1054年（宋至和元年、辽重熙23年；公元换算据两国年号对照）}
\noindent\eventanchor{SYD-1054-LIAO-045}{九月乙亥，契丹遣使來告夏國平}\textbf{宋与辽。}
宋与辽 在 1054年（宋至和元年、辽重熙23年；公元换算据两国年号对照） 的材料痕迹是〔九月乙亥，契丹遣使來告夏國平〕，出处为 《宋史》卷012〈至和元年〉；《辽史》本纪同年条（互校）。材料未将地点细化到城址，不能用中枢所在地替它补足。

把本项放回事件链，能够看到的前提为：本纪记录使节、贡赐或往来，不等于稳定统属关系。 而它所打开或收紧的下一步是：可与对方编年和交通、文书材料核对其后续落实。

关于可说范围，须保留：《宋史》为元末官修，本纪保存宋廷纪年与处置；战果、数字、地方执行及对方动机须以编年、会要、对方史料或地方材料互校。 进一步讨论可参照 宋辽边政、使节记录与契丹碑刻研究

\noindent\eventanchor{SYD-1054-LIAO-111}{辛巳，遣三司使王拱辰報使契丹}\textbf{宋与辽。}
从 《宋史》卷012〈至和元年〉；《辽史》本纪同年条（互校） 的 1054年（宋至和元年、辽重熙23年；公元换算据两国年号对照） 条目看，〔辛巳，遣三司使王拱辰報使契丹〕使 宋与辽 的一个具体行动进入叙事；材料未将地点细化到城址，不能用中枢所在地替它补足。

它承接的条件是：本纪记录使节、贡赐或往来，不等于稳定统属关系。 随后可追踪的变化为：可与对方编年和交通、文书材料核对其后续落实。

证据的边界是：《宋史》为元末官修，本纪保存宋廷纪年与处置；战果、数字、地方执行及对方动机须以编年、会要、对方史料或地方材料互校。 宋辽边政、使节记录与契丹碑刻研究 可用于继续检验这一结论。

\subsection{1055年}
\noindent\eventanchor{LIAO-1055-EVENT-041}{辽兴宗去世，道宗即位}\textbf{辽。}
1055年 的记载将 辽 与〔辽兴宗去世，道宗即位〕相连。此处可确认的是这项被记录的行动；材料未将地点细化到城址，不能用中枢所在地替它补足。

前一层压力可由以下线索辨认：辽中期继承及辽宋夏关系中前任统治的结束与宫廷、部族或官僚的权力安排相互交织，促成“辽兴宗去世，道宗即位”这一继承节点。 这一处置留下的后续条件是：继承并不自动消除既有矛盾；“辽兴宗去世，道宗即位”后的任官、军权和对外策略需由后续条目追踪。

《辽史》为元末官修且契丹本方连续文书有限；宋使与碑刻的观察位置不同，战果、族称及制度覆盖范围不得互推。 因此，辽代契丹国家、都城、部族与跨境关系研究 只能扩展问题，不能代替未保存的细节。

\subsection{1055年（宋至和二年、辽清宁1年；公元换算据两国年号对照）}
\noindent\eventanchor{SYD-1055-LIAO-046}{庚子，契丹遣使致其主宗真遺留物及謝弔祭}\textbf{宋与辽。}
〔庚子，契丹遣使致其主宗真遺留物及謝弔祭〕见于 《宋史》卷012〈至和二年〉；《辽史》本纪同年条（互校） 的 1055年（宋至和二年、辽清宁1年；公元换算据两国年号对照） 记录。宋与辽 的选择因此有了可回查的时间位置；材料未将地点细化到城址，不能用中枢所在地替它补足。

本纪记录使节、贡赐或往来，不等于稳定统属关系。 记录所见的结果则是：可与对方编年和交通、文书材料核对其后续落实。 日期相邻不等于因果已经证实。

关于可说范围，须保留：《宋史》为元末官修，本纪保存宋廷纪年与处置；战果、数字、地方执行及对方动机须以编年、会要、对方史料或地方材料互校。 进一步讨论可参照 宋辽边政、使节记录与契丹碑刻研究

\noindent\eventanchor{SYD-1055-LIAO-112}{九月戊午，契丹使來告其國主宗真殂，帝為發哀，成服於內東門幕次，遣使祭奠、吊慰及賀其子洪基立}\textbf{宋与辽。}
1055年（宋至和二年、辽清宁1年；公元换算据两国年号对照），《宋史》卷012〈至和二年〉；《辽史》本纪同年条（互校） 所见的〔九月戊午，契丹使來告其國主宗真殂，帝為發哀，成服於內東門幕次，遣使祭奠、吊慰及賀其子洪基立〕把 宋与辽 的处置留在可复核的年序中；材料未将地点细化到城址，不能用中枢所在地替它补足。

把本项放回事件链，能够看到的前提为：本纪记录使节、贡赐或往来，不等于稳定统属关系。 而它所打开或收紧的下一步是：可与对方编年和交通、文书材料核对其后续落实。

证据的边界是：《宋史》为元末官修，本纪保存宋廷纪年与处置；战果、数字、地方执行及对方动机须以编年、会要、对方史料或地方材料互校。 宋辽边政、使节记录与契丹碑刻研究 可用于继续检验这一结论。

\subsection{1056 年（共享年序桥接）}
\noindent\eventanchor{SY-1056-CHRONOLOGY-BRIDGE}{【C级年序桥接】保留1056 年的多政权并行检索入口；本条不主张所列政权在当年发生直接互动，具体事件须另以单项材料入账。}\textbf{宋、辽与西夏（年序桥接）。}
宋、辽与西夏（年序桥接） 在 1056 年（共享年序桥接） 的材料痕迹是〔【C级年序桥接】保留1056 年的多政权并行检索入口；本条不主张所列政权在当年发生直接互动，具体事件须另以单项材料入账。〕，出处为 《宋史》；《辽史》；《续资治通鉴长编》；西夏文书与碑刻（据事核）。题名保留的城、路、水道或机构名称，是目前可以落地的空间线索。

它承接的条件是：相邻已入账事件之间缺少该年度的共享年序入口，易使跨政权材料被单一政权叙事遮蔽。 随后可追踪的变化为：保留该年度的检索与补账位置；后续仅在获得可核单项材料时拆分为具体政治、财政、军事、社会或文化事件。

本条只确认该年位于可追踪的共享年序中；正史、地方文书和外部材料的保存密度不一，不能由桥接标签推断行动、统属、疆界或社会影响。 因此，宋辽夏边境、财政与文书材料的并行年序研究 只能扩展问题，不能代替未保存的细节。

\subsection{1057年（宋嘉祐二年、辽清宁3年；公元换算据两国年号对照）}
\noindent\eventanchor{SYD-1057-LIAO-047}{冬十月乙巳，遣胡宿報使契丹}\textbf{宋与辽。}
从 《宋史》卷012〈嘉祐二年〉；《辽史》本纪同年条（互校） 的 1057年（宋嘉祐二年、辽清宁3年；公元换算据两国年号对照） 条目看，〔冬十月乙巳，遣胡宿報使契丹〕使 宋与辽 的一个具体行动进入叙事；材料未将地点细化到城址，不能用中枢所在地替它补足。

前一层压力可由以下线索辨认：本纪记录使节、贡赐或往来，不等于稳定统属关系。 这一处置留下的后续条件是：可与对方编年和交通、文书材料核对其后续落实。

关于可说范围，须保留：《宋史》为元末官修，本纪保存宋廷纪年与处置；战果、数字、地方执行及对方动机须以编年、会要、对方史料或地方材料互校。 进一步讨论可参照 宋辽边政、使节记录与契丹碑刻研究

\noindent\eventanchor{SYD-1057-LIAO-113}{九月庚子，契丹再使蕭扈、吳湛來求御容}\textbf{宋与辽。}
1057年（宋嘉祐二年、辽清宁3年；公元换算据两国年号对照） 的记载将 宋与辽 与〔九月庚子，契丹再使蕭扈、吳湛來求御容〕相连。此处可确认的是这项被记录的行动；材料未将地点细化到城址，不能用中枢所在地替它补足。

本纪年条提供日期和中枢可见行动；前因不据单条补定。 记录所见的结果则是：可回查同年及后续条目，地方影响仍须另证。 日期相邻不等于因果已经证实。

证据的边界是：《宋史》为元末官修，本纪保存宋廷纪年与处置；战果、数字、地方执行及对方动机须以编年、会要、对方史料或地方材料互校。 宋辽边政、使节记录与契丹碑刻研究 可用于继续检验这一结论。

\subsection{1057年（宋嘉祐二年、西夏奲都2年；公元换算据两国年号对照）}
\noindent\eventanchor{SYD-1057-XIA-019}{六月戊午，夏國主諒祚遣人來謝使弔祭}\textbf{宋与西夏。}
〔六月戊午，夏國主諒祚遣人來謝使弔祭〕见于 《宋史》卷012〈嘉祐二年〉；西夏文书、碑刻及《辽史》或《金史》同年条（互校） 的 1057年（宋嘉祐二年、西夏奲都2年；公元换算据两国年号对照） 记录。宋与西夏 的选择因此有了可回查的时间位置；材料未将地点细化到城址，不能用中枢所在地替它补足。

把本项放回事件链，能够看到的前提为：本纪记录使节、贡赐或往来，不等于稳定统属关系。 而它所打开或收紧的下一步是：可与对方编年和交通、文书材料核对其后续落实。

《宋史》为元末官修，本纪保存宋廷纪年与处置；战果、数字、地方执行及对方动机须以编年、会要、对方史料或地方材料互校。 因此，西夏文书、河西—宁夏考古与宋夏关系研究 只能扩展问题，不能代替未保存的细节。

\subsection{1058 年（共享年序桥接）}
\noindent\eventanchor{SY-1058-CHRONOLOGY-BRIDGE}{【C级年序桥接】保留1058 年的多政权并行检索入口；本条不主张所列政权在当年发生直接互动，具体事件须另以单项材料入账。}\textbf{宋、辽与西夏（年序桥接）。}
1058 年（共享年序桥接），《宋史》；《辽史》；《续资治通鉴长编》；西夏文书与碑刻（据事核） 所见的〔【C级年序桥接】保留1058 年的多政权并行检索入口；本条不主张所列政权在当年发生直接互动，具体事件须另以单项材料入账。〕把 宋、辽与西夏（年序桥接） 的处置留在可复核的年序中；题名保留的城、路、水道或机构名称，是目前可以落地的空间线索。

它承接的条件是：相邻已入账事件之间缺少该年度的共享年序入口，易使跨政权材料被单一政权叙事遮蔽。 随后可追踪的变化为：保留该年度的检索与补账位置；后续仅在获得可核单项材料时拆分为具体政治、财政、军事、社会或文化事件。

关于可说范围，须保留：本条只确认该年位于可追踪的共享年序中；正史、地方文书和外部材料的保存密度不一，不能由桥接标签推断行动、统属、疆界或社会影响。 进一步讨论可参照 宋辽夏边境、财政与文书材料的并行年序研究

\subsection{1058年（宋嘉祐三年、辽清宁4年；公元换算据两国年号对照）}
\noindent\eventanchor{SYD-1058-LIAO-048}{二月癸卯，契丹使來告其祖母哀，輟視朝七日，遣使祭奠吊慰}\textbf{宋与辽。}
宋与辽 在 1058年（宋嘉祐三年、辽清宁4年；公元换算据两国年号对照） 的材料痕迹是〔二月癸卯，契丹使來告其祖母哀，輟視朝七日，遣使祭奠吊慰〕，出处为 《宋史》卷012〈嘉祐三年〉；《辽史》本纪同年条（互校）。材料未将地点细化到城址，不能用中枢所在地替它补足。

前一层压力可由以下线索辨认：本纪记录使节、贡赐或往来，不等于稳定统属关系。 这一处置留下的后续条件是：可与对方编年和交通、文书材料核对其后续落实。

证据的边界是：《宋史》为元末官修，本纪保存宋廷纪年与处置；战果、数字、地方执行及对方动机须以编年、会要、对方史料或地方材料互校。 宋辽边政、使节记录与契丹碑刻研究 可用于继续检验这一结论。

\noindent\eventanchor{SYD-1058-LIAO-114}{己-{丑}-，契丹遣使來謝}\textbf{宋与辽。}
从 《宋史》卷012〈嘉祐三年〉；《辽史》本纪同年条（互校） 的 1058年（宋嘉祐三年、辽清宁4年；公元换算据两国年号对照） 条目看，〔己-{丑}-，契丹遣使來謝〕使 宋与辽 的一个具体行动进入叙事；材料未将地点细化到城址，不能用中枢所在地替它补足。

本纪记录使节、贡赐或往来，不等于稳定统属关系。 记录所见的结果则是：可与对方编年和交通、文书材料核对其后续落实。 日期相邻不等于因果已经证实。

《宋史》为元末官修，本纪保存宋廷纪年与处置；战果、数字、地方执行及对方动机须以编年、会要、对方史料或地方材料互校。 因此，宋辽边政、使节记录与契丹碑刻研究 只能扩展问题，不能代替未保存的细节。

\subsection{1059 年（共享年序桥接）}
\noindent\eventanchor{SY-1059-CHRONOLOGY-BRIDGE}{【C级年序桥接】保留1059 年的多政权并行检索入口；本条不主张所列政权在当年发生直接互动，具体事件须另以单项材料入账。}\textbf{宋、辽与西夏（年序桥接）。}
1059 年（共享年序桥接） 的记载将 宋、辽与西夏（年序桥接） 与〔【C级年序桥接】保留1059 年的多政权并行检索入口；本条不主张所列政权在当年发生直接互动，具体事件须另以单项材料入账。〕相连。此处可确认的是这项被记录的行动；题名保留的城、路、水道或机构名称，是目前可以落地的空间线索。

把本项放回事件链，能够看到的前提为：相邻已入账事件之间缺少该年度的共享年序入口，易使跨政权材料被单一政权叙事遮蔽。 而它所打开或收紧的下一步是：保留该年度的检索与补账位置；后续仅在获得可核单项材料时拆分为具体政治、财政、军事、社会或文化事件。

关于可说范围，须保留：本条只确认该年位于可追踪的共享年序中；正史、地方文书和外部材料的保存密度不一，不能由桥接标签推断行动、统属、疆界或社会影响。 进一步讨论可参照 宋辽夏边境、财政与文书材料的并行年序研究

\subsection{1060 年（共享年序桥接）}
\noindent\eventanchor{SY-1060-CHRONOLOGY-BRIDGE}{【C级年序桥接】保留1060 年的多政权并行检索入口；本条不主张所列政权在当年发生直接互动，具体事件须另以单项材料入账。}\textbf{宋、辽与西夏（年序桥接）。}
〔【C级年序桥接】保留1060 年的多政权并行检索入口；本条不主张所列政权在当年发生直接互动，具体事件须另以单项材料入账。〕见于 《宋史》；《辽史》；《续资治通鉴长编》；西夏文书与碑刻（据事核） 的 1060 年（共享年序桥接） 记录。宋、辽与西夏（年序桥接） 的选择因此有了可回查的时间位置；题名保留的城、路、水道或机构名称，是目前可以落地的空间线索。

它承接的条件是：相邻已入账事件之间缺少该年度的共享年序入口，易使跨政权材料被单一政权叙事遮蔽。 随后可追踪的变化为：保留该年度的检索与补账位置；后续仅在获得可核单项材料时拆分为具体政治、财政、军事、社会或文化事件。

证据的边界是：本条只确认该年位于可追踪的共享年序中；正史、地方文书和外部材料的保存密度不一，不能由桥接标签推断行动、统属、疆界或社会影响。 宋辽夏边境、财政与文书材料的并行年序研究 可用于继续检验这一结论。

\subsection{1061 年（共享年序桥接）}
\noindent\eventanchor{SY-1061-CHRONOLOGY-BRIDGE}{【C级年序桥接】保留1061 年的多政权并行检索入口；本条不主张所列政权在当年发生直接互动，具体事件须另以单项材料入账。}\textbf{宋、辽与西夏（年序桥接）。}
1061 年（共享年序桥接），《宋史》；《辽史》；《续资治通鉴长编》；西夏文书与碑刻（据事核） 所见的〔【C级年序桥接】保留1061 年的多政权并行检索入口；本条不主张所列政权在当年发生直接互动，具体事件须另以单项材料入账。〕把 宋、辽与西夏（年序桥接） 的处置留在可复核的年序中；题名保留的城、路、水道或机构名称，是目前可以落地的空间线索。

前一层压力可由以下线索辨认：相邻已入账事件之间缺少该年度的共享年序入口，易使跨政权材料被单一政权叙事遮蔽。 这一处置留下的后续条件是：保留该年度的检索与补账位置；后续仅在获得可核单项材料时拆分为具体政治、财政、军事、社会或文化事件。

本条只确认该年位于可追踪的共享年序中；正史、地方文书和外部材料的保存密度不一，不能由桥接标签推断行动、统属、疆界或社会影响。 因此，宋辽夏边境、财政与文书材料的并行年序研究 只能扩展问题，不能代替未保存的细节。

\subsection{1062 年（共享年序桥接）}
\noindent\eventanchor{SY-1062-CHRONOLOGY-BRIDGE}{【C级年序桥接】保留1062 年的多政权并行检索入口；本条不主张所列政权在当年发生直接互动，具体事件须另以单项材料入账。}\textbf{宋、辽与西夏（年序桥接）。}
宋、辽与西夏（年序桥接） 在 1062 年（共享年序桥接） 的材料痕迹是〔【C级年序桥接】保留1062 年的多政权并行检索入口；本条不主张所列政权在当年发生直接互动，具体事件须另以单项材料入账。〕，出处为 《宋史》；《辽史》；《续资治通鉴长编》；西夏文书与碑刻（据事核）。题名保留的城、路、水道或机构名称，是目前可以落地的空间线索。

相邻已入账事件之间缺少该年度的共享年序入口，易使跨政权材料被单一政权叙事遮蔽。 记录所见的结果则是：保留该年度的检索与补账位置；后续仅在获得可核单项材料时拆分为具体政治、财政、军事、社会或文化事件。 日期相邻不等于因果已经证实。

关于可说范围，须保留：本条只确认该年位于可追踪的共享年序中；正史、地方文书和外部材料的保存密度不一，不能由桥接标签推断行动、统属、疆界或社会影响。 进一步讨论可参照 宋辽夏边境、财政与文书材料的并行年序研究

\subsection{1062年（宋嘉祐七年、西夏奲都7年；公元换算据两国年号对照）}
\noindent\eventanchor{SYD-1062-XIA-020}{己-{丑}-，夏國主諒祚進馬，求賜書，詔賜《九經》，還其馬}\textbf{宋与西夏。}
从 《宋史》卷012〈嘉祐七年〉；西夏文书、碑刻及《辽史》或《金史》同年条（互校） 的 1062年（宋嘉祐七年、西夏奲都7年；公元换算据两国年号对照） 条目看，〔己-{丑}-，夏國主諒祚進馬，求賜書，詔賜《九經》，還其馬〕使 宋与西夏 的一个具体行动进入叙事；材料未将地点细化到城址，不能用中枢所在地替它补足。

把本项放回事件链，能够看到的前提为：本纪年条记录政令或机构处置；实施背景须参照制度材料。 而它所打开或收紧的下一步是：可在同年及后续政令中追踪；条文不单独证明地方执行。

证据的边界是：《宋史》为元末官修，本纪保存宋廷纪年与处置；战果、数字、地方执行及对方动机须以编年、会要、对方史料或地方材料互校。 西夏文书、河西—宁夏考古与宋夏关系研究 可用于继续检验这一结论。

\subsection{1063年（宋嘉祐八年、辽清宁9年；公元换算据两国年号对照）}
\noindent\eventanchor{SYD-1063-LIAO-049}{六月辛卯，契丹遣蕭福延等來祭吊}\textbf{宋与辽。}
1063年（宋嘉祐八年、辽清宁9年；公元换算据两国年号对照） 的记载将 宋与辽 与〔六月辛卯，契丹遣蕭福延等來祭吊〕相连。此处可确认的是这项被记录的行动；材料未将地点细化到城址，不能用中枢所在地替它补足。

它承接的条件是：本纪年条提供日期和中枢可见行动；前因不据单条补定。 随后可追踪的变化为：可回查同年及后续条目，地方影响仍须另证。

《宋史》为元末官修，本纪保存宋廷纪年与处置；战果、数字、地方执行及对方动机须以编年、会要、对方史料或地方材料互校。 因此，宋辽边政、使节记录与契丹碑刻研究 只能扩展问题，不能代替未保存的细节。

\noindent\eventanchor{SYD-1063-LIAO-115}{遣韓贄等告即位於契丹}\textbf{宋与辽。}
〔遣韓贄等告即位於契丹〕见于 《宋史》卷013〈嘉祐八年〉；《辽史》本纪同年条（互校） 的 1063年（宋嘉祐八年、辽清宁9年；公元换算据两国年号对照） 记录。宋与辽 的选择因此有了可回查的时间位置；材料未将地点细化到城址，不能用中枢所在地替它补足。

前一层压力可由以下线索辨认：本纪年条提供日期和中枢可见行动；前因不据单条补定。 这一处置留下的后续条件是：可回查同年及后续条目，地方影响仍须另证。

关于可说范围，须保留：《宋史》为元末官修，本纪保存宋廷纪年与处置；战果、数字、地方执行及对方动机须以编年、会要、对方史料或地方材料互校。 进一步讨论可参照 宋辽边政、使节记录与契丹碑刻研究

\subsection{1064年（宋治平元年、辽清宁10年；公元换算据两国年号对照）}
\noindent\eventanchor{SYD-1064-LIAO-050}{丙辰，契丹遣耶律烈等來賀壽聖節，蕭禧等來賀明年正旦}\textbf{宋与辽。}
1064年（宋治平元年、辽清宁10年；公元换算据两国年号对照），《宋史》卷013〈治平元年〉；《辽史》本纪同年条（互校） 所见的〔丙辰，契丹遣耶律烈等來賀壽聖節，蕭禧等來賀明年正旦〕把 宋与辽 的处置留在可复核的年序中；材料未将地点细化到城址，不能用中枢所在地替它补足。

本纪年条提供日期和中枢可见行动；前因不据单条补定。 记录所见的结果则是：可回查同年及后续条目，地方影响仍须另证。 日期相邻不等于因果已经证实。

证据的边界是：《宋史》为元末官修，本纪保存宋廷纪年与处置；战果、数字、地方执行及对方动机须以编年、会要、对方史料或地方材料互校。 宋辽边政、使节记录与契丹碑刻研究 可用于继续检验这一结论。

\noindent\eventanchor{SYD-1064-LIAO-116}{乙卯，遣兵部員外郎呂誨等四人充賀契丹太后生辰、正旦使，刑部郎中章岷等四人充賀契丹主生辰、正旦使}\textbf{宋与辽。}
宋与辽 在 1064年（宋治平元年、辽清宁10年；公元换算据两国年号对照） 的材料痕迹是〔乙卯，遣兵部員外郎呂誨等四人充賀契丹太后生辰、正旦使，刑部郎中章岷等四人充賀契丹主生辰、正旦使〕，出处为 《宋史》卷013〈治平元年〉；《辽史》本纪同年条（互校）。材料未将地点细化到城址，不能用中枢所在地替它补足。

把本项放回事件链，能够看到的前提为：本纪记录使节、贡赐或往来，不等于稳定统属关系。 而它所打开或收紧的下一步是：可与对方编年和交通、文书材料核对其后续落实。

《宋史》为元末官修，本纪保存宋廷纪年与处置；战果、数字、地方执行及对方动机须以编年、会要、对方史料或地方材料互校。 因此，宋辽边政、使节记录与契丹碑刻研究 只能扩展问题，不能代替未保存的细节。

\subsection{1064年（宋治平元年、西夏奲都9年；公元换算据两国年号对照）}
\noindent\eventanchor{SYD-1064-XIA-021}{庚午，詔夏國精擇使人，戒勵毋紊彝章}\textbf{宋与西夏。}
从 《宋史》卷013〈治平元年〉；西夏文书、碑刻及《辽史》或《金史》同年条（互校） 的 1064年（宋治平元年、西夏奲都9年；公元换算据两国年号对照） 条目看，〔庚午，詔夏國精擇使人，戒勵毋紊彝章〕使 宋与西夏 的一个具体行动进入叙事；材料未将地点细化到城址，不能用中枢所在地替它补足。

它承接的条件是：本纪年条记录政令或机构处置；实施背景须参照制度材料。 随后可追踪的变化为：可在同年及后续政令中追踪；条文不单独证明地方执行。

关于可说范围，须保留：《宋史》为元末官修，本纪保存宋廷纪年与处置；战果、数字、地方执行及对方动机须以编年、会要、对方史料或地方材料互校。 进一步讨论可参照 西夏文书、河西—宁夏考古与宋夏关系研究

\subsection{1066 年（共享年序桥接）}
\noindent\eventanchor{SY-1066-CHRONOLOGY-BRIDGE}{【C级年序桥接】保留1066 年的多政权并行检索入口；本条不主张所列政权在当年发生直接互动，具体事件须另以单项材料入账。}\textbf{宋、辽与西夏（年序桥接）。}
1066 年（共享年序桥接） 的记载将 宋、辽与西夏（年序桥接） 与〔【C级年序桥接】保留1066 年的多政权并行检索入口；本条不主张所列政权在当年发生直接互动，具体事件须另以单项材料入账。〕相连。此处可确认的是这项被记录的行动；题名保留的城、路、水道或机构名称，是目前可以落地的空间线索。

前一层压力可由以下线索辨认：相邻已入账事件之间缺少该年度的共享年序入口，易使跨政权材料被单一政权叙事遮蔽。 这一处置留下的后续条件是：保留该年度的检索与补账位置；后续仅在获得可核单项材料时拆分为具体政治、财政、军事、社会或文化事件。

证据的边界是：本条只确认该年位于可追踪的共享年序中；正史、地方文书和外部材料的保存密度不一，不能由桥接标签推断行动、统属、疆界或社会影响。 宋辽夏边境、财政与文书材料的并行年序研究 可用于继续检验这一结论。

\subsection{1066年（宋治平三年、辽咸雍2年；公元换算据两国年号对照）}
\noindent\eventanchor{SYD-1066-LIAO-051}{八月庚子，遣傅卞等賀遼主生辰，張師顏等賀正旦}\textbf{宋与辽。}
〔八月庚子，遣傅卞等賀遼主生辰，張師顏等賀正旦〕见于 《宋史》卷013〈治平三年〉；《辽史》本纪同年条（互校） 的 1066年（宋治平三年、辽咸雍2年；公元换算据两国年号对照） 记录。宋与辽 的选择因此有了可回查的时间位置；材料未将地点细化到城址，不能用中枢所在地替它补足。

本纪年条提供日期和中枢可见行动；前因不据单条补定。 记录所见的结果则是：可回查同年及后续条目，地方影响仍须另证。 日期相邻不等于因果已经证实。

《宋史》为元末官修，本纪保存宋廷纪年与处置；战果、数字、地方执行及对方动机须以编年、会要、对方史料或地方材料互校。 因此，宋辽边政、使节记录与契丹碑刻研究 只能扩展问题，不能代替未保存的细节。

\noindent\eventanchor{SYD-1066-LIAO-117}{癸酉，契丹改國號爲遼}\textbf{宋与辽。}
1066年（宋治平三年、辽咸雍2年；公元换算据两国年号对照），《宋史》卷013〈治平三年〉；《辽史》本纪同年条（互校） 所见的〔癸酉，契丹改國號爲遼〕把 宋与辽 的处置留在可复核的年序中；材料未将地点细化到城址，不能用中枢所在地替它补足。

把本项放回事件链，能够看到的前提为：本纪年条记录政令或机构处置；实施背景须参照制度材料。 而它所打开或收紧的下一步是：可在同年及后续政令中追踪；条文不单独证明地方执行。

关于可说范围，须保留：《宋史》为元末官修，本纪保存宋廷纪年与处置；战果、数字、地方执行及对方动机须以编年、会要、对方史料或地方材料互校。 进一步讨论可参照 宋辽边政、使节记录与契丹碑刻研究

\subsection{1066年（宋治平三年、西夏奲都11年；公元换算据两国年号对照）}
\noindent\eventanchor{SYD-1066-XIA-022}{是歲，遣使以違約數寇責夏國，諒詐獻方物謝罪}\textbf{宋与西夏。}
宋与西夏 在 1066年（宋治平三年、西夏奲都11年；公元换算据两国年号对照） 的材料痕迹是〔是歲，遣使以違約數寇責夏國，諒詐獻方物謝罪〕，出处为 《宋史》卷013〈治平三年〉；西夏文书、碑刻及《辽史》或《金史》同年条（互校）。材料未将地点细化到城址，不能用中枢所在地替它补足。

它承接的条件是：本纪从朝廷视角记录军政行动；出兵与战果需分辨。 随后可追踪的变化为：后续路线、控制与损失应与对方记录和遗址材料复核。

证据的边界是：《宋史》为元末官修，本纪保存宋廷纪年与处置；战果、数字、地方执行及对方动机须以编年、会要、对方史料或地方材料互校。 西夏文书、河西—宁夏考古与宋夏关系研究 可用于继续检验这一结论。

\subsection{1067年（宋治平四年、辽咸雍3年；公元换算据两国年号对照）}
\noindent\eventanchor{SYD-1067-LIAO-052}{己巳，遼遣蕭傑等來賀正旦}\textbf{宋与辽。}
从 《宋史》卷014〈治平四年〉；《辽史》本纪同年条（互校） 的 1067年（宋治平四年、辽咸雍3年；公元换算据两国年号对照） 条目看，〔己巳，遼遣蕭傑等來賀正旦〕使 宋与辽 的一个具体行动进入叙事；材料未将地点细化到城址，不能用中枢所在地替它补足。

前一层压力可由以下线索辨认：本纪年条提供日期和中枢可见行动；前因不据单条补定。 这一处置留下的后续条件是：可回查同年及后续条目，地方影响仍须另证。

《宋史》为元末官修，本纪保存宋廷纪年与处置；战果、数字、地方执行及对方动机须以编年、会要、对方史料或地方材料互校。 因此，宋辽边政、使节记录与契丹碑刻研究 只能扩展问题，不能代替未保存的细节。

\noindent\eventanchor{SYD-1067-LIAO-118}{甲午，遼遣耶律好謀等來賀即位}\textbf{宋与辽。}
1067年（宋治平四年、辽咸雍3年；公元换算据两国年号对照） 的记载将 宋与辽 与〔甲午，遼遣耶律好謀等來賀即位〕相连。此处可确认的是这项被记录的行动；材料未将地点细化到城址，不能用中枢所在地替它补足。

本纪年条提供日期和中枢可见行动；前因不据单条补定。 记录所见的结果则是：可回查同年及后续条目，地方影响仍须另证。 日期相邻不等于因果已经证实。

关于可说范围，须保留：《宋史》为元末官修，本纪保存宋廷纪年与处置；战果、数字、地方执行及对方动机须以编年、会要、对方史料或地方材料互校。 进一步讨论可参照 宋辽边政、使节记录与契丹碑刻研究

\subsection{1067年（宋治平四年、西夏拱化1年；公元换算据两国年号对照）}
\noindent\eventanchor{SYD-1067-XIA-023}{甲申，夏國主諒祚遣使謝罪}\textbf{宋与西夏。}
〔甲申，夏國主諒祚遣使謝罪〕见于 《宋史》卷014〈治平四年〉；西夏文书、碑刻及《辽史》或《金史》同年条（互校） 的 1067年（宋治平四年、西夏拱化1年；公元换算据两国年号对照） 记录。宋与西夏 的选择因此有了可回查的时间位置；材料未将地点细化到城址，不能用中枢所在地替它补足。

把本项放回事件链，能够看到的前提为：本纪记录使节、贡赐或往来，不等于稳定统属关系。 而它所打开或收紧的下一步是：可与对方编年和交通、文书材料核对其后续落实。

证据的边界是：《宋史》为元末官修，本纪保存宋廷纪年与处置；战果、数字、地方执行及对方动机须以编年、会要、对方史料或地方材料互校。 西夏文书、河西—宁夏考古与宋夏关系研究 可用于继续检验这一结论。

\subsection{1067年}
\noindent\eventanchor{XIA-1067-EVENT-063}{李秉常即位，西夏再度出现幼主政治}\textbf{西夏。}
1067年，《宋史》〈夏国传〉及同年本纪；《辽史》《金史》或《元史》相关条；西夏文文书、碑刻与城址材料（据事核） 所见的〔李秉常即位，西夏再度出现幼主政治〕把 西夏 的处置留在可复核的年序中；材料未将地点细化到城址，不能用中枢所在地替它补足。

它承接的条件是：西夏建国与宋辽边境战争中前任统治的结束与宫廷、部族或官僚的权力安排相互交织，促成“李秉常即位，西夏再度出现幼主政治”这一继承节点。 随后可追踪的变化为：继承并不自动消除既有矛盾；“李秉常即位，西夏再度出现幼主政治”后的任官、军权和对外策略需由后续条目追踪。

西夏无存世连续国史，汉文叙事多出自对手；西夏文文书的释读、定年和地域代表性须逐项说明。 因此，西夏文文书、河西—宁夏考古与宋夏金蒙关系研究 只能扩展问题，不能代替未保存的细节。

\subsection{1068年（宋熙寧元年、辽咸雍4年；公元换算据两国年号对照）}
\noindent\eventanchor{SYD-1068-LIAO-053}{丁卯，遣張宗益等賀遼主生辰、正旦}\textbf{宋与辽。}
宋与辽 在 1068年（宋熙寧元年、辽咸雍4年；公元换算据两国年号对照） 的材料痕迹是〔丁卯，遣張宗益等賀遼主生辰、正旦〕，出处为 《宋史》卷014〈熙寧元年〉；《辽史》本纪同年条（互校）。材料未将地点细化到城址，不能用中枢所在地替它补足。

前一层压力可由以下线索辨认：本纪年条提供日期和中枢可见行动；前因不据单条补定。 这一处置留下的后续条件是：可回查同年及后续条目，地方影响仍须另证。

关于可说范围，须保留：《宋史》为元末官修，本纪保存宋廷纪年与处置；战果、数字、地方执行及对方动机须以编年、会要、对方史料或地方材料互校。 进一步讨论可参照 宋辽边政、使节记录与契丹碑刻研究

\noindent\eventanchor{SYD-1068-LIAO-119}{甲子，遼遣耶律公質等來賀正旦}\textbf{宋与辽。}
从 《宋史》卷014〈熙寧元年〉；《辽史》本纪同年条（互校） 的 1068年（宋熙寧元年、辽咸雍4年；公元换算据两国年号对照） 条目看，〔甲子，遼遣耶律公質等來賀正旦〕使 宋与辽 的一个具体行动进入叙事；材料未将地点细化到城址，不能用中枢所在地替它补足。

本纪年条提供日期和中枢可见行动；前因不据单条补定。 记录所见的结果则是：可回查同年及后续条目，地方影响仍须另证。 日期相邻不等于因果已经证实。

证据的边界是：《宋史》为元末官修，本纪保存宋廷纪年与处置；战果、数字、地方执行及对方动机须以编年、会要、对方史料或地方材料互校。 宋辽边政、使节记录与契丹碑刻研究 可用于继续检验这一结论。

\subsection{1068年（宋熙寧元年、西夏拱化2年；公元换算据两国年号对照）}
\noindent\eventanchor{SYD-1068-XIA-024}{庚戌，賜夏國主秉常詔，許納塞門、安遠二砦歸其綏州}\textbf{宋与西夏。}
1068年（宋熙寧元年、西夏拱化2年；公元换算据两国年号对照） 的记载将 宋与西夏 与〔庚戌，賜夏國主秉常詔，許納塞門、安遠二砦歸其綏州〕相连。此处可确认的是这项被记录的行动；题名保留的城、路、水道或机构名称，是目前可以落地的空间线索。

把本项放回事件链，能够看到的前提为：本纪年条记录政令或机构处置；实施背景须参照制度材料。 而它所打开或收紧的下一步是：可在同年及后续政令中追踪；条文不单独证明地方执行。

《宋史》为元末官修，本纪保存宋廷纪年与处置；战果、数字、地方执行及对方动机须以编年、会要、对方史料或地方材料互校。 因此，西夏文书、河西—宁夏考古与宋夏关系研究 只能扩展问题，不能代替未保存的细节。

\noindent\eventanchor{SYD-1068-XIA-057}{三月庚辰，夏主諒詐卒，遣使來告哀}\textbf{宋与西夏。}
〔三月庚辰，夏主諒詐卒，遣使來告哀〕见于 《宋史》卷014〈熙寧元年〉；西夏文书、碑刻及《辽史》或《金史》同年条（互校） 的 1068年（宋熙寧元年、西夏拱化2年；公元换算据两国年号对照） 记录。宋与西夏 的选择因此有了可回查的时间位置；材料未将地点细化到城址，不能用中枢所在地替它补足。

它承接的条件是：本纪记录使节、贡赐或往来，不等于稳定统属关系。 随后可追踪的变化为：可与对方编年和交通、文书材料核对其后续落实。

关于可说范围，须保留：《宋史》为元末官修，本纪保存宋廷纪年与处置；战果、数字、地方执行及对方动机须以编年、会要、对方史料或地方材料互校。 进一步讨论可参照 西夏文书、河西—宁夏考古与宋夏关系研究

\subsection{1069年（宋熙寧二年、辽咸雍5年；公元换算据两国年号对照）}
\noindent\eventanchor{SYD-1069-LIAO-054}{丁-{丑}-，遣孫固等賀遼主生辰、正旦}\textbf{宋与辽。}
1069年（宋熙寧二年、辽咸雍5年；公元换算据两国年号对照），《宋史》卷014〈熙寧二年〉；《辽史》本纪同年条（互校） 所见的〔丁-{丑}-，遣孫固等賀遼主生辰、正旦〕把 宋与辽 的处置留在可复核的年序中；材料未将地点细化到城址，不能用中枢所在地替它补足。

前一层压力可由以下线索辨认：本纪年条提供日期和中枢可见行动；前因不据单条补定。 这一处置留下的后续条件是：可回查同年及后续条目，地方影响仍须另证。

证据的边界是：《宋史》为元末官修，本纪保存宋廷纪年与处置；战果、数字、地方执行及对方动机须以编年、会要、对方史料或地方材料互校。 宋辽边政、使节记录与契丹碑刻研究 可用于继续检验这一结论。

\noindent\eventanchor{SYD-1069-LIAO-120}{壬寅，遼遣耶律昌等來賀同天節}\textbf{宋与辽。}
宋与辽 在 1069年（宋熙寧二年、辽咸雍5年；公元换算据两国年号对照） 的材料痕迹是〔壬寅，遼遣耶律昌等來賀同天節〕，出处为 《宋史》卷014〈熙寧二年〉；《辽史》本纪同年条（互校）。材料未将地点细化到城址，不能用中枢所在地替它补足。

本纪年条提供日期和中枢可见行动；前因不据单条补定。 记录所见的结果则是：可回查同年及后续条目，地方影响仍须另证。 日期相邻不等于因果已经证实。

《宋史》为元末官修，本纪保存宋廷纪年与处置；战果、数字、地方执行及对方动机须以编年、会要、对方史料或地方材料互校。 因此，宋辽边政、使节记录与契丹碑刻研究 只能扩展问题，不能代替未保存的细节。

\subsection{1069年（宋熙寧二年、西夏乾道1年；公元换算据两国年号对照）}
\noindent\eventanchor{SYD-1069-XIA-025}{逵遣趙禼交夏人所納安遠、塞門二砦，就定地界}\textbf{宋与西夏。}
从 《宋史》卷014〈熙寧二年〉；西夏文书、碑刻及《辽史》或《金史》同年条（互校） 的 1069年（宋熙寧二年、西夏乾道1年；公元换算据两国年号对照） 条目看，〔逵遣趙禼交夏人所納安遠、塞門二砦，就定地界〕使 宋与西夏 的一个具体行动进入叙事；材料未将地点细化到城址，不能用中枢所在地替它补足。

把本项放回事件链，能够看到的前提为：本纪年条提供日期和中枢可见行动；前因不据单条补定。 而它所打开或收紧的下一步是：可回查同年及后续条目，地方影响仍须另证。

关于可说范围，须保留：《宋史》为元末官修，本纪保存宋廷纪年与处置；战果、数字、地方执行及对方动机须以编年、会要、对方史料或地方材料互校。 进一步讨论可参照 西夏文书、河西—宁夏考古与宋夏关系研究

\noindent\eventanchor{SYD-1069-XIA-058}{夏國請從舊蕃儀，詔許之}\textbf{宋与西夏。}
1069年（宋熙寧二年、西夏乾道1年；公元换算据两国年号对照） 的记载将 宋与西夏 与〔夏國請從舊蕃儀，詔許之〕相连。此处可确认的是这项被记录的行动；材料未将地点细化到城址，不能用中枢所在地替它补足。

它承接的条件是：本纪年条记录政令或机构处置；实施背景须参照制度材料。 随后可追踪的变化为：可在同年及后续政令中追踪；条文不单独证明地方执行。

证据的边界是：《宋史》为元末官修，本纪保存宋廷纪年与处置；战果、数字、地方执行及对方动机须以编年、会要、对方史料或地方材料互校。 西夏文书、河西—宁夏考古与宋夏关系研究 可用于继续检验这一结论。

\noindent\eventanchor{SYD-1069-XIA-076}{夏人渝初盟，禼請城綏州，不以易二砦，因改名綏德城}\textbf{宋与西夏。}
〔夏人渝初盟，禼請城綏州，不以易二砦，因改名綏德城〕见于 《宋史》卷014〈熙寧二年〉；西夏文书、碑刻及《辽史》或《金史》同年条（互校） 的 1069年（宋熙寧二年、西夏乾道1年；公元换算据两国年号对照） 记录。宋与西夏 的选择因此有了可回查的时间位置；题名保留的城、路、水道或机构名称，是目前可以落地的空间线索。

前一层压力可由以下线索辨认：本纪年条记录政令或机构处置；实施背景须参照制度材料。 这一处置留下的后续条件是：可在同年及后续政令中追踪；条文不单独证明地方执行。

《宋史》为元末官修，本纪保存宋廷纪年与处置；战果、数字、地方执行及对方动机须以编年、会要、对方史料或地方材料互校。 因此，西夏文书、河西—宁夏考古与宋夏关系研究 只能扩展问题，不能代替未保存的细节。

\subsection{1070年（宋熙寧三年、辽咸雍6年；公元换算据两国年号对照）}
\noindent\eventanchor{SYD-1070-LIAO-055}{丙寅，遼遣耶律寬來賀同天節}\textbf{宋与辽。}
1070年（宋熙寧三年、辽咸雍6年；公元换算据两国年号对照），《宋史》卷015〈熙寧三年〉；《辽史》本纪同年条（互校） 所见的〔丙寅，遼遣耶律寬來賀同天節〕把 宋与辽 的处置留在可复核的年序中；材料未将地点细化到城址，不能用中枢所在地替它补足。

本纪年条提供日期和中枢可见行动；前因不据单条补定。 记录所见的结果则是：可回查同年及后续条目，地方影响仍须另证。 日期相邻不等于因果已经证实。

关于可说范围，须保留：《宋史》为元末官修，本纪保存宋廷纪年与处置；战果、数字、地方执行及对方动机须以编年、会要、对方史料或地方材料互校。 进一步讨论可参照 宋辽边政、使节记录与契丹碑刻研究

\noindent\eventanchor{SYD-1070-LIAO-121}{遣張景憲等賀遼主生辰、正旦}\textbf{宋与辽。}
宋与辽 在 1070年（宋熙寧三年、辽咸雍6年；公元换算据两国年号对照） 的材料痕迹是〔遣張景憲等賀遼主生辰、正旦〕，出处为 《宋史》卷015〈熙寧三年〉；《辽史》本纪同年条（互校）。材料未将地点细化到城址，不能用中枢所在地替它补足。

把本项放回事件链，能够看到的前提为：本纪年条提供日期和中枢可见行动；前因不据单条补定。 而它所打开或收紧的下一步是：可回查同年及后续条目，地方影响仍须另证。

证据的边界是：《宋史》为元末官修，本纪保存宋廷纪年与处置；战果、数字、地方执行及对方动机须以编年、会要、对方史料或地方材料互校。 宋辽边政、使节记录与契丹碑刻研究 可用于继续检验这一结论。

\subsection{1070年（宋熙寧三年、西夏天赐礼盛国庆1年；公元换算据两国年号对照）}
\noindent\eventanchor{SYD-1070-XIA-026}{冬十月辛酉，詔延州毋納夏使}\textbf{宋与西夏。}
从 《宋史》卷015〈熙寧三年〉；西夏文书、碑刻及《辽史》或《金史》同年条（互校） 的 1070年（宋熙寧三年、西夏天赐礼盛国庆1年；公元换算据两国年号对照） 条目看，〔冬十月辛酉，詔延州毋納夏使〕使 宋与西夏 的一个具体行动进入叙事；题名保留的城、路、水道或机构名称，是目前可以落地的空间线索。

它承接的条件是：本纪年条记录政令或机构处置；实施背景须参照制度材料。 随后可追踪的变化为：可在同年及后续政令中追踪；条文不单独证明地方执行。

《宋史》为元末官修，本纪保存宋廷纪年与处置；战果、数字、地方执行及对方动机须以编年、会要、对方史料或地方材料互校。 因此，西夏文书、河西—宁夏考古与宋夏关系研究 只能扩展问题，不能代替未保存的细节。

\noindent\eventanchor{SYD-1070-XIA-059}{庚午，夏人寇鎮戎軍三川砦，巡檢趙普伏兵邀擊，敗之}\textbf{宋与西夏。}
1070年（宋熙寧三年、西夏天赐礼盛国庆1年；公元换算据两国年号对照） 的记载将 宋与西夏 与〔庚午，夏人寇鎮戎軍三川砦，巡檢趙普伏兵邀擊，敗之〕相连。此处可确认的是这项被记录的行动；题名保留的城、路、水道或机构名称，是目前可以落地的空间线索。

前一层压力可由以下线索辨认：本纪从朝廷视角记录军政行动；出兵与战果需分辨。 这一处置留下的后续条件是：后续路线、控制与损失应与对方记录和遗址材料复核。

关于可说范围，须保留：《宋史》为元末官修，本纪保存宋廷纪年与处置；战果、数字、地方执行及对方动机须以编年、会要、对方史料或地方材料互校。 进一步讨论可参照 西夏文书、河西—宁夏考古与宋夏关系研究

\noindent\eventanchor{SYD-1070-XIA-077}{己卯，夏人犯大順城，知慶州李復圭以方略授環慶路鈐轄李信、慶州東路都巡檢劉甫、監押種詠出戰，兵少取敗}\textbf{宋与西夏。}
〔己卯，夏人犯大順城，知慶州李復圭以方略授環慶路鈐轄李信、慶州東路都巡檢劉甫、監押種詠出戰，兵少取敗〕见于 《宋史》卷015〈熙寧三年〉；西夏文书、碑刻及《辽史》或《金史》同年条（互校） 的 1070年（宋熙寧三年、西夏天赐礼盛国庆1年；公元换算据两国年号对照） 记录。宋与西夏 的选择因此有了可回查的时间位置；题名保留的城、路、水道或机构名称，是目前可以落地的空间线索。

本纪年条记录政令或机构处置；实施背景须参照制度材料。 记录所见的结果则是：可在同年及后续政令中追踪；条文不单独证明地方执行。 日期相邻不等于因果已经证实。

证据的边界是：《宋史》为元末官修，本纪保存宋廷纪年与处置；战果、数字、地方执行及对方动机须以编年、会要、对方史料或地方材料互校。 西夏文书、河西—宁夏考古与宋夏关系研究 可用于继续检验这一结论。

\noindent\eventanchor{SYD-1070-XIA-089}{甲辰，夏人寇大順城，都監燕達等擊走之}\textbf{宋与西夏。}
1070年（宋熙寧三年、西夏天赐礼盛国庆1年；公元换算据两国年号对照），《宋史》卷015〈熙寧三年〉；西夏文书、碑刻及《辽史》或《金史》同年条（互校） 所见的〔甲辰，夏人寇大順城，都監燕達等擊走之〕把 宋与西夏 的处置留在可复核的年序中；题名保留的城、路、水道或机构名称，是目前可以落地的空间线索。

把本项放回事件链，能够看到的前提为：本纪从朝廷视角记录军政行动；出兵与战果需分辨。 而它所打开或收紧的下一步是：后续路线、控制与损失应与对方记录和遗址材料复核。

《宋史》为元末官修，本纪保存宋廷纪年与处置；战果、数字、地方执行及对方动机须以编年、会要、对方史料或地方材料互校。 因此，西夏文书、河西—宁夏考古与宋夏关系研究 只能扩展问题，不能代替未保存的细节。

\noindent\eventanchor{SYD-1070-XIA-095}{是月，慶州巡檢姚兕敗夏人于荔原堡}\textbf{宋与西夏。}
宋与西夏 在 1070年（宋熙寧三年、西夏天赐礼盛国庆1年；公元换算据两国年号对照） 的材料痕迹是〔是月，慶州巡檢姚兕敗夏人于荔原堡〕，出处为 《宋史》卷015〈熙寧三年〉；西夏文书、碑刻及《辽史》或《金史》同年条（互校）。题名保留的城、路、水道或机构名称，是目前可以落地的空间线索。

它承接的条件是：本纪从朝廷视角记录军政行动；出兵与战果需分辨。 随后可追踪的变化为：后续路线、控制与损失应与对方记录和遗址材料复核。

关于可说范围，须保留：《宋史》为元末官修，本纪保存宋廷纪年与处置；战果、数字、地方执行及对方动机须以编年、会要、对方史料或地方材料互校。 进一步讨论可参照 西夏文书、河西—宁夏考古与宋夏关系研究

\subsection{1071年（宋熙寧四年、辽咸雍7年；公元换算据两国年号对照）}
\noindent\eventanchor{SYD-1071-LIAO-056}{丙子，遼遣耶律紀等來賀正旦}\textbf{宋与辽。}
从 《宋史》卷015〈熙寧四年〉；《辽史》本纪同年条（互校） 的 1071年（宋熙寧四年、辽咸雍7年；公元换算据两国年号对照） 条目看，〔丙子，遼遣耶律紀等來賀正旦〕使 宋与辽 的一个具体行动进入叙事；材料未将地点细化到城址，不能用中枢所在地替它补足。

前一层压力可由以下线索辨认：本纪年条提供日期和中枢可见行动；前因不据单条补定。 这一处置留下的后续条件是：可回查同年及后续条目，地方影响仍须另证。

证据的边界是：《宋史》为元末官修，本纪保存宋廷纪年与处置；战果、数字、地方执行及对方动机须以编年、会要、对方史料或地方材料互校。 宋辽边政、使节记录与契丹碑刻研究 可用于继续检验这一结论。

\noindent\eventanchor{SYD-1071-LIAO-122}{癸酉，遣楚建中等賀遼主生辰、正旦}\textbf{宋与辽。}
1071年（宋熙寧四年、辽咸雍7年；公元换算据两国年号对照） 的记载将 宋与辽 与〔癸酉，遣楚建中等賀遼主生辰、正旦〕相连。此处可确认的是这项被记录的行动；材料未将地点细化到城址，不能用中枢所在地替它补足。

本纪年条提供日期和中枢可见行动；前因不据单条补定。 记录所见的结果则是：可回查同年及后续条目，地方影响仍须另证。 日期相邻不等于因果已经证实。

《宋史》为元末官修，本纪保存宋廷纪年与处置；战果、数字、地方执行及对方动机须以编年、会要、对方史料或地方材料互校。 因此，宋辽边政、使节记录与契丹碑刻研究 只能扩展问题，不能代替未保存的细节。

\subsection{1071年（宋熙寧四年、西夏天赐礼盛国庆2年；公元换算据两国年号对照）}
\noindent\eventanchor{SYD-1071-XIA-027}{己-{丑}-，種諤襲夏兵於囉兀北，大敗之，遂城囉兀}\textbf{宋与西夏。}
〔己-{丑}-，種諤襲夏兵於囉兀北，大敗之，遂城囉兀〕见于 《宋史》卷015〈熙寧四年〉；西夏文书、碑刻及《辽史》或《金史》同年条（互校） 的 1071年（宋熙寧四年、西夏天赐礼盛国庆2年；公元换算据两国年号对照） 记录。宋与西夏 的选择因此有了可回查的时间位置；题名保留的城、路、水道或机构名称，是目前可以落地的空间线索。

把本项放回事件链，能够看到的前提为：本纪从朝廷视角记录军政行动；出兵与战果需分辨。 而它所打开或收紧的下一步是：后续路线、控制与损失应与对方记录和遗址材料复核。

关于可说范围，须保留：《宋史》为元末官修，本纪保存宋廷纪年与处置；战果、数字、地方执行及对方动机须以编年、会要、对方史料或地方材料互校。 进一步讨论可参照 西夏文书、河西—宁夏考古与宋夏关系研究

\noindent\eventanchor{SYD-1071-XIA-060}{壬戌，遣環慶都鈐轄幵贇以兵屯邠、涇、河中，以備西夏}\textbf{宋与西夏。}
1071年（宋熙寧四年、西夏天赐礼盛国庆2年；公元换算据两国年号对照），《宋史》卷015〈熙寧四年〉；西夏文书、碑刻及《辽史》或《金史》同年条（互校） 所见的〔壬戌，遣環慶都鈐轄幵贇以兵屯邠、涇、河中，以備西夏〕把 宋与西夏 的处置留在可复核的年序中；题名保留的城、路、水道或机构名称，是目前可以落地的空间线索。

它承接的条件是：本纪保留灾害或工程的报告地点与朝廷处置；灾情范围需另证。 随后可追踪的变化为：可与地方文书、河工和环境材料比较，不将单点记录外推为全域因果。

证据的边界是：《宋史》为元末官修，本纪保存宋廷纪年与处置；战果、数字、地方执行及对方动机须以编年、会要、对方史料或地方材料互校。 西夏文书、河西—宁夏考古与宋夏关系研究 可用于继续检验这一结论。

\noindent\eventanchor{SYD-1071-XIA-078}{自是夏人日聚兵爲報復計，言者以諤爲稔邊患不便}\textbf{宋与西夏。}
宋与西夏 在 1071年（宋熙寧四年、西夏天赐礼盛国庆2年；公元换算据两国年号对照） 的材料痕迹是〔自是夏人日聚兵爲報復計，言者以諤爲稔邊患不便〕，出处为 《宋史》卷015〈熙寧四年〉；西夏文书、碑刻及《辽史》或《金史》同年条（互校）。材料未将地点细化到城址，不能用中枢所在地替它补足。

前一层压力可由以下线索辨认：本纪年条记录政令或机构处置；实施背景须参照制度材料。 这一处置留下的后续条件是：可在同年及后续政令中追踪；条文不单独证明地方执行。

《宋史》为元末官修，本纪保存宋廷纪年与处置；战果、数字、地方执行及对方动机须以编年、会要、对方史料或地方材料互校。 因此，西夏文书、河西—宁夏考古与宋夏关系研究 只能扩展问题，不能代替未保存的细节。

\subsection{1072年（宋熙寧五年、辽咸雍8年；公元换算据两国年号对照）}
\noindent\eventanchor{SYD-1072-LIAO-057}{癸巳，遣崔台符等賀遼主生辰、正旦}\textbf{宋与辽。}
从 《宋史》卷015〈熙寧五年〉；《辽史》本纪同年条（互校） 的 1072年（宋熙寧五年、辽咸雍8年；公元换算据两国年号对照） 条目看，〔癸巳，遣崔台符等賀遼主生辰、正旦〕使 宋与辽 的一个具体行动进入叙事；材料未将地点细化到城址，不能用中枢所在地替它补足。

本纪年条提供日期和中枢可见行动；前因不据单条补定。 记录所见的结果则是：可回查同年及后续条目，地方影响仍须另证。 日期相邻不等于因果已经证实。

关于可说范围，须保留：《宋史》为元末官修，本纪保存宋廷纪年与处置；战果、数字、地方执行及对方动机须以编年、会要、对方史料或地方材料互校。 进一步讨论可参照 宋辽边政、使节记录与契丹碑刻研究

\noindent\eventanchor{SYD-1072-LIAO-123}{己亥，遼遣蕭瑜等來賀正旦}\textbf{宋与辽。}
1072年（宋熙寧五年、辽咸雍8年；公元换算据两国年号对照） 的记载将 宋与辽 与〔己亥，遼遣蕭瑜等來賀正旦〕相连。此处可确认的是这项被记录的行动；材料未将地点细化到城址，不能用中枢所在地替它补足。

把本项放回事件链，能够看到的前提为：本纪年条提供日期和中枢可见行动；前因不据单条补定。 而它所打开或收紧的下一步是：可回查同年及后续条目，地方影响仍须另证。

证据的边界是：《宋史》为元末官修，本纪保存宋廷纪年与处置；战果、数字、地方执行及对方动机须以编年、会要、对方史料或地方材料互校。 宋辽边政、使节记录与契丹碑刻研究 可用于继续检验这一结论。

\subsection{1073 年（共享年序桥接）}
\noindent\eventanchor{SY-1073-CHRONOLOGY-BRIDGE}{【C级年序桥接】保留1073 年的多政权并行检索入口；本条不主张所列政权在当年发生直接互动，具体事件须另以单项材料入账。}\textbf{宋、辽与西夏（年序桥接）。}
〔【C级年序桥接】保留1073 年的多政权并行检索入口；本条不主张所列政权在当年发生直接互动，具体事件须另以单项材料入账。〕见于 《宋史》；《辽史》；《续资治通鉴长编》；西夏文书与碑刻（据事核） 的 1073 年（共享年序桥接） 记录。宋、辽与西夏（年序桥接） 的选择因此有了可回查的时间位置；题名保留的城、路、水道或机构名称，是目前可以落地的空间线索。

它承接的条件是：相邻已入账事件之间缺少该年度的共享年序入口，易使跨政权材料被单一政权叙事遮蔽。 随后可追踪的变化为：保留该年度的检索与补账位置；后续仅在获得可核单项材料时拆分为具体政治、财政、军事、社会或文化事件。

本条只确认该年位于可追踪的共享年序中；正史、地方文书和外部材料的保存密度不一，不能由桥接标签推断行动、统属、疆界或社会影响。 因此，宋辽夏边境、财政与文书材料的并行年序研究 只能扩展问题，不能代替未保存的细节。

\subsection{1073年（宋熙寧六年、辽咸雍9年；公元换算据两国年号对照）}
\noindent\eventanchor{SYD-1073-LIAO-058}{八月壬申朔，遣賈昌衡等賀遼主生辰、正旦}\textbf{宋与辽。}
1073年（宋熙寧六年、辽咸雍9年；公元换算据两国年号对照），《宋史》卷015〈熙寧六年〉；《辽史》本纪同年条（互校） 所见的〔八月壬申朔，遣賈昌衡等賀遼主生辰、正旦〕把 宋与辽 的处置留在可复核的年序中；材料未将地点细化到城址，不能用中枢所在地替它补足。

前一层压力可由以下线索辨认：本纪年条提供日期和中枢可见行动；前因不据单条补定。 这一处置留下的后续条件是：可回查同年及后续条目，地方影响仍须另证。

关于可说范围，须保留：《宋史》为元末官修，本纪保存宋廷纪年与处置；战果、数字、地方执行及对方动机须以编年、会要、对方史料或地方材料互校。 进一步讨论可参照 宋辽边政、使节记录与契丹碑刻研究

\noindent\eventanchor{SYD-1073-LIAO-124}{丙申，遼遣耶律洞等來賀正旦}\textbf{宋与辽。}
宋与辽 在 1073年（宋熙寧六年、辽咸雍9年；公元换算据两国年号对照） 的材料痕迹是〔丙申，遼遣耶律洞等來賀正旦〕，出处为 《宋史》卷015〈熙寧六年〉；《辽史》本纪同年条（互校）。材料未将地点细化到城址，不能用中枢所在地替它补足。

本纪年条提供日期和中枢可见行动；前因不据单条补定。 记录所见的结果则是：可回查同年及后续条目，地方影响仍须另证。 日期相邻不等于因果已经证实。

证据的边界是：《宋史》为元末官修，本纪保存宋廷纪年与处置；战果、数字、地方执行及对方动机须以编年、会要、对方史料或地方材料互校。 宋辽边政、使节记录与契丹碑刻研究 可用于继续检验这一结论。

\subsection{1073年（宋熙寧六年、西夏天赐礼盛国庆4年；公元换算据两国年号对照）}
\noindent\eventanchor{SYD-1073-XIA-028}{二月辛卯，夏人寇秦州，都巡檢使劉惟吉敗之}\textbf{宋与西夏。}
从 《宋史》卷015〈熙寧六年〉；西夏文书、碑刻及《辽史》或《金史》同年条（互校） 的 1073年（宋熙寧六年、西夏天赐礼盛国庆4年；公元换算据两国年号对照） 条目看，〔二月辛卯，夏人寇秦州，都巡檢使劉惟吉敗之〕使 宋与西夏 的一个具体行动进入叙事；题名保留的城、路、水道或机构名称，是目前可以落地的空间线索。

把本项放回事件链，能够看到的前提为：本纪从朝廷视角记录军政行动；出兵与战果需分辨。 而它所打开或收紧的下一步是：后续路线、控制与损失应与对方记录和遗址材料复核。

《宋史》为元末官修，本纪保存宋廷纪年与处置；战果、数字、地方执行及对方动机须以编年、会要、对方史料或地方材料互校。 因此，西夏文书、河西—宁夏考古与宋夏关系研究 只能扩展问题，不能代替未保存的细节。

\subsection{1074年（宋熙寧七年、辽咸雍10年；公元换算据两国年号对照）}
\noindent\eventanchor{SYD-1074-LIAO-059}{丙辰，遼遣林牙蕭禧來言河東疆界，命太常少卿劉忱議之}\textbf{宋与辽。}
1074年（宋熙寧七年、辽咸雍10年；公元换算据两国年号对照） 的记载将 宋与辽 与〔丙辰，遼遣林牙蕭禧來言河東疆界，命太常少卿劉忱議之〕相连。此处可确认的是这项被记录的行动；题名保留的城、路、水道或机构名称，是目前可以落地的空间线索。

它承接的条件是：本纪年条记录政令或机构处置；实施背景须参照制度材料。 随后可追踪的变化为：可在同年及后续政令中追踪；条文不单独证明地方执行。

关于可说范围，须保留：《宋史》为元末官修，本纪保存宋廷纪年与处置；战果、数字、地方执行及对方动机须以编年、会要、对方史料或地方材料互校。 进一步讨论可参照 宋辽边政、使节记录与契丹碑刻研究

\noindent\eventanchor{SYD-1074-LIAO-125}{己-{丑}-，遼遣耶律甯等來賀正旦}\textbf{宋与辽。}
〔己-{丑}-，遼遣耶律甯等來賀正旦〕见于 《宋史》卷015〈熙寧七年〉；《辽史》本纪同年条（互校） 的 1074年（宋熙寧七年、辽咸雍10年；公元换算据两国年号对照） 记录。宋与辽 的选择因此有了可回查的时间位置；材料未将地点细化到城址，不能用中枢所在地替它补足。

前一层压力可由以下线索辨认：本纪年条提供日期和中枢可见行动；前因不据单条补定。 这一处置留下的后续条件是：可回查同年及后续条目，地方影响仍须另证。

证据的边界是：《宋史》为元末官修，本纪保存宋廷纪年与处置；战果、数字、地方执行及对方动机须以编年、会要、对方史料或地方材料互校。 宋辽边政、使节记录与契丹碑刻研究 可用于继续检验这一结论。

\subsection{1075年（宋熙寧八年、辽太康1年；公元换算据两国年号对照）}
\noindent\eventanchor{SYD-1075-LIAO-060}{丙申，遣謝景溫等賀遼主生辰、正旦}\textbf{宋与辽。}
1075年（宋熙寧八年、辽太康1年；公元换算据两国年号对照），《宋史》卷015〈熙寧八年〉；《辽史》本纪同年条（互校） 所见的〔丙申，遣謝景溫等賀遼主生辰、正旦〕把 宋与辽 的处置留在可复核的年序中；材料未将地点细化到城址，不能用中枢所在地替它补足。

本纪年条提供日期和中枢可见行动；前因不据单条补定。 记录所见的结果则是：可回查同年及后续条目，地方影响仍须另证。 日期相邻不等于因果已经证实。

《宋史》为元末官修，本纪保存宋廷纪年与处置；战果、数字、地方执行及对方动机须以编年、会要、对方史料或地方材料互校。 因此，宋辽边政、使节记录与契丹碑刻研究 只能扩展问题，不能代替未保存的细节。

\noindent\eventanchor{SYD-1075-LIAO-126}{丁卯，遼遣耶律景熙等來賀同天節}\textbf{宋与辽。}
宋与辽 在 1075年（宋熙寧八年、辽太康1年；公元换算据两国年号对照） 的材料痕迹是〔丁卯，遼遣耶律景熙等來賀同天節〕，出处为 《宋史》卷015〈熙寧八年〉；《辽史》本纪同年条（互校）。材料未将地点细化到城址，不能用中枢所在地替它补足。

把本项放回事件链，能够看到的前提为：本纪年条提供日期和中枢可见行动；前因不据单条补定。 而它所打开或收紧的下一步是：可回查同年及后续条目，地方影响仍须另证。

关于可说范围，须保留：《宋史》为元末官修，本纪保存宋廷纪年与处置；战果、数字、地方执行及对方动机须以编年、会要、对方史料或地方材料互校。 进一步讨论可参照 宋辽边政、使节记录与契丹碑刻研究

\subsection{1076年（宋熙寧九年、辽太康2年；公元换算据两国年号对照）}
\noindent\eventanchor{SYD-1076-LIAO-061}{丙午，遣王克臣等吊慰於遼}\textbf{宋与辽。}
从 《宋史》卷015〈熙寧九年〉；《辽史》本纪同年条（互校） 的 1076年（宋熙寧九年、辽太康2年；公元换算据两国年号对照） 条目看，〔丙午，遣王克臣等吊慰於遼〕使 宋与辽 的一个具体行动进入叙事；材料未将地点细化到城址，不能用中枢所在地替它补足。

它承接的条件是：本纪年条提供日期和中枢可见行动；前因不据单条补定。 随后可追踪的变化为：可回查同年及后续条目，地方影响仍须另证。

证据的边界是：《宋史》为元末官修，本纪保存宋廷纪年与处置；战果、数字、地方执行及对方动机须以编年、会要、对方史料或地方材料互校。 宋辽边政、使节记录与契丹碑刻研究 可用于继续检验这一结论。

\noindent\eventanchor{SYD-1076-LIAO-127}{丁未，遼遣耶律運等來賀正旦}\textbf{宋与辽。}
1076年（宋熙寧九年、辽太康2年；公元换算据两国年号对照） 的记载将 宋与辽 与〔丁未，遼遣耶律運等來賀正旦〕相连。此处可确认的是这项被记录的行动；材料未将地点细化到城址，不能用中枢所在地替它补足。

前一层压力可由以下线索辨认：本纪年条提供日期和中枢可见行动；前因不据单条补定。 这一处置留下的后续条件是：可回查同年及后续条目，地方影响仍须另证。

《宋史》为元末官修，本纪保存宋廷纪年与处置；战果、数字、地方执行及对方动机须以编年、会要、对方史料或地方材料互校。 因此，宋辽边政、使节记录与契丹碑刻研究 只能扩展问题，不能代替未保存的细节。

\subsection{1077 年（共享年序桥接）}
\noindent\eventanchor{SY-1077-CHRONOLOGY-BRIDGE}{【C级年序桥接】保留1077 年的多政权并行检索入口；本条不主张所列政权在当年发生直接互动，具体事件须另以单项材料入账。}\textbf{宋、辽与西夏（年序桥接）。}
〔【C级年序桥接】保留1077 年的多政权并行检索入口；本条不主张所列政权在当年发生直接互动，具体事件须另以单项材料入账。〕见于 《宋史》；《辽史》；《续资治通鉴长编》；西夏文书与碑刻（据事核） 的 1077 年（共享年序桥接） 记录。宋、辽与西夏（年序桥接） 的选择因此有了可回查的时间位置；题名保留的城、路、水道或机构名称，是目前可以落地的空间线索。

相邻已入账事件之间缺少该年度的共享年序入口，易使跨政权材料被单一政权叙事遮蔽。 记录所见的结果则是：保留该年度的检索与补账位置；后续仅在获得可核单项材料时拆分为具体政治、财政、军事、社会或文化事件。 日期相邻不等于因果已经证实。

关于可说范围，须保留：本条只确认该年位于可追踪的共享年序中；正史、地方文书和外部材料的保存密度不一，不能由桥接标签推断行动、统属、疆界或社会影响。 进一步讨论可参照 宋辽夏边境、财政与文书材料的并行年序研究

\subsection{1077年（宋熙寧十年、辽太康3年；公元换算据两国年号对照）}
\noindent\eventanchor{SYD-1077-LIAO-062}{遣蘇頌等賀遼主生辰、正旦}\textbf{宋与辽。}
1077年（宋熙寧十年、辽太康3年；公元换算据两国年号对照），《宋史》卷015〈熙寧十年〉；《辽史》本纪同年条（互校） 所见的〔遣蘇頌等賀遼主生辰、正旦〕把 宋与辽 的处置留在可复核的年序中；材料未将地点细化到城址，不能用中枢所在地替它补足。

把本项放回事件链，能够看到的前提为：本纪年条提供日期和中枢可见行动；前因不据单条补定。 而它所打开或收紧的下一步是：可回查同年及后续条目，地方影响仍须另证。

证据的边界是：《宋史》为元末官修，本纪保存宋廷纪年与处置；战果、数字、地方执行及对方动机须以编年、会要、对方史料或地方材料互校。 宋辽边政、使节记录与契丹碑刻研究 可用于继续检验这一结论。

\noindent\eventanchor{SYD-1077-LIAO-128}{辛-{丑}-，遼遣耶律孝淳等來賀正旦}\textbf{宋与辽。}
宋与辽 在 1077年（宋熙寧十年、辽太康3年；公元换算据两国年号对照） 的材料痕迹是〔辛-{丑}-，遼遣耶律孝淳等來賀正旦〕，出处为 《宋史》卷015〈熙寧十年〉；《辽史》本纪同年条（互校）。材料未将地点细化到城址，不能用中枢所在地替它补足。

它承接的条件是：本纪年条提供日期和中枢可见行动；前因不据单条补定。 随后可追踪的变化为：可回查同年及后续条目，地方影响仍须另证。

《宋史》为元末官修，本纪保存宋廷纪年与处置；战果、数字、地方执行及对方动机须以编年、会要、对方史料或地方材料互校。 因此，宋辽边政、使节记录与契丹碑刻研究 只能扩展问题，不能代替未保存的细节。

\subsection{1078 年（共享年序桥接）}
\noindent\eventanchor{SY-1078-CHRONOLOGY-BRIDGE}{【C级年序桥接】保留1078 年的多政权并行检索入口；本条不主张所列政权在当年发生直接互动，具体事件须另以单项材料入账。}\textbf{宋、辽与西夏（年序桥接）。}
从 《宋史》；《辽史》；《续资治通鉴长编》；西夏文书与碑刻（据事核） 的 1078 年（共享年序桥接） 条目看，〔【C级年序桥接】保留1078 年的多政权并行检索入口；本条不主张所列政权在当年发生直接互动，具体事件须另以单项材料入账。〕使 宋、辽与西夏（年序桥接） 的一个具体行动进入叙事；题名保留的城、路、水道或机构名称，是目前可以落地的空间线索。

前一层压力可由以下线索辨认：相邻已入账事件之间缺少该年度的共享年序入口，易使跨政权材料被单一政权叙事遮蔽。 这一处置留下的后续条件是：保留该年度的检索与补账位置；后续仅在获得可核单项材料时拆分为具体政治、财政、军事、社会或文化事件。

关于可说范围，须保留：本条只确认该年位于可追踪的共享年序中；正史、地方文书和外部材料的保存密度不一，不能由桥接标签推断行动、统属、疆界或社会影响。 进一步讨论可参照 宋辽夏边境、财政与文书材料的并行年序研究

\subsection{1078年（宋元豐元年、辽太康4年；公元换算据两国年号对照）}
\noindent\eventanchor{SYD-1078-LIAO-063}{丙寅，遼遣耶律隆等來賀正旦}\textbf{宋与辽。}
1078年（宋元豐元年、辽太康4年；公元换算据两国年号对照） 的记载将 宋与辽 与〔丙寅，遼遣耶律隆等來賀正旦〕相连。此处可确认的是这项被记录的行动；材料未将地点细化到城址，不能用中枢所在地替它补足。

本纪年条提供日期和中枢可见行动；前因不据单条补定。 记录所见的结果则是：可回查同年及后续条目，地方影响仍须另证。 日期相邻不等于因果已经证实。

证据的边界是：《宋史》为元末官修，本纪保存宋廷纪年与处置；战果、数字、地方执行及对方动机须以编年、会要、对方史料或地方材料互校。 宋辽边政、使节记录与契丹碑刻研究 可用于继续检验这一结论。

\noindent\eventanchor{SYD-1078-LIAO-129}{甲寅，遣黃履等賀遼主生辰、正旦}\textbf{宋与辽。}
〔甲寅，遣黃履等賀遼主生辰、正旦〕见于 《宋史》卷015〈元豐元年〉；《辽史》本纪同年条（互校） 的 1078年（宋元豐元年、辽太康4年；公元换算据两国年号对照） 记录。宋与辽 的选择因此有了可回查的时间位置；材料未将地点细化到城址，不能用中枢所在地替它补足。

把本项放回事件链，能够看到的前提为：本纪年条提供日期和中枢可见行动；前因不据单条补定。 而它所打开或收紧的下一步是：可回查同年及后续条目，地方影响仍须另证。

《宋史》为元末官修，本纪保存宋廷纪年与处置；战果、数字、地方执行及对方动机须以编年、会要、对方史料或地方材料互校。 因此，宋辽边政、使节记录与契丹碑刻研究 只能扩展问题，不能代替未保存的细节。

\subsection{1079年（宋元豐二年、辽太康5年；公元换算据两国年号对照）}
\noindent\eventanchor{SYD-1079-LIAO-064}{庚申，遼遣蕭甯等來賀正旦}\textbf{宋与辽。}
1079年（宋元豐二年、辽太康5年；公元换算据两国年号对照），《宋史》卷015〈元豐二年〉；《辽史》本纪同年条（互校） 所见的〔庚申，遼遣蕭甯等來賀正旦〕把 宋与辽 的处置留在可复核的年序中；材料未将地点细化到城址，不能用中枢所在地替它补足。

它承接的条件是：本纪年条提供日期和中枢可见行动；前因不据单条补定。 随后可追踪的变化为：可回查同年及后续条目，地方影响仍须另证。

关于可说范围，须保留：《宋史》为元末官修，本纪保存宋廷纪年与处置；战果、数字、地方执行及对方动机须以编年、会要、对方史料或地方材料互校。 进一步讨论可参照 宋辽边政、使节记录与契丹碑刻研究

\noindent\eventanchor{SYD-1079-LIAO-130}{甲辰，遼遣蕭晟等來賀同天節}\textbf{宋与辽。}
宋与辽 在 1079年（宋元豐二年、辽太康5年；公元换算据两国年号对照） 的材料痕迹是〔甲辰，遼遣蕭晟等來賀同天節〕，出处为 《宋史》卷015〈元豐二年〉；《辽史》本纪同年条（互校）。材料未将地点细化到城址，不能用中枢所在地替它补足。

前一层压力可由以下线索辨认：本纪年条提供日期和中枢可见行动；前因不据单条补定。 这一处置留下的后续条件是：可回查同年及后续条目，地方影响仍须另证。

证据的边界是：《宋史》为元末官修，本纪保存宋廷纪年与处置；战果、数字、地方执行及对方动机须以编年、会要、对方史料或地方材料互校。 宋辽边政、使节记录与契丹碑刻研究 可用于继续检验这一结论。

\subsection{1079年（宋元豐二年、西夏大安5年；公元换算据两国年号对照）}
\noindent\eventanchor{SYD-1079-XIA-029}{八月丙申朔，夏人寇綏德城，都監李浦敗之}\textbf{宋与西夏。}
从 《宋史》卷015〈元豐二年〉；西夏文书、碑刻及《辽史》或《金史》同年条（互校） 的 1079年（宋元豐二年、西夏大安5年；公元换算据两国年号对照） 条目看，〔八月丙申朔，夏人寇綏德城，都監李浦敗之〕使 宋与西夏 的一个具体行动进入叙事；题名保留的城、路、水道或机构名称，是目前可以落地的空间线索。

本纪从朝廷视角记录军政行动；出兵与战果需分辨。 记录所见的结果则是：后续路线、控制与损失应与对方记录和遗址材料复核。 日期相邻不等于因果已经证实。

《宋史》为元末官修，本纪保存宋廷纪年与处置；战果、数字、地方执行及对方动机须以编年、会要、对方史料或地方材料互校。 因此，西夏文书、河西—宁夏考古与宋夏关系研究 只能扩展问题，不能代替未保存的细节。

\subsection{1080年（宋元豐三年、辽太康6年；公元换算据两国年号对照）}
\noindent\eventanchor{SYD-1080-LIAO-065}{癸-{丑}-，遣王存等賀遼主生辰、正旦}\textbf{宋与辽。}
1080年（宋元豐三年、辽太康6年；公元换算据两国年号对照） 的记载将 宋与辽 与〔癸-{丑}-，遣王存等賀遼主生辰、正旦〕相连。此处可确认的是这项被记录的行动；材料未将地点细化到城址，不能用中枢所在地替它补足。

把本项放回事件链，能够看到的前提为：本纪年条提供日期和中枢可见行动；前因不据单条补定。 而它所打开或收紧的下一步是：可回查同年及后续条目，地方影响仍须另证。

关于可说范围，须保留：《宋史》为元末官修，本纪保存宋廷纪年与处置；战果、数字、地方执行及对方动机须以编年、会要、对方史料或地方材料互校。 进一步讨论可参照 宋辽边政、使节记录与契丹碑刻研究

\noindent\eventanchor{SYD-1080-LIAO-131}{己亥，遼遣耶律永芳等來賀同天節}\textbf{宋与辽。}
〔己亥，遼遣耶律永芳等來賀同天節〕见于 《宋史》卷016〈元豐三年〉；《辽史》本纪同年条（互校） 的 1080年（宋元豐三年、辽太康6年；公元换算据两国年号对照） 记录。宋与辽 的选择因此有了可回查的时间位置；材料未将地点细化到城址，不能用中枢所在地替它补足。

它承接的条件是：本纪年条提供日期和中枢可见行动；前因不据单条补定。 随后可追踪的变化为：可回查同年及后续条目，地方影响仍须另证。

证据的边界是：《宋史》为元末官修，本纪保存宋廷纪年与处置；战果、数字、地方执行及对方动机须以编年、会要、对方史料或地方材料互校。 宋辽边政、使节记录与契丹碑刻研究 可用于继续检验这一结论。

\subsection{1080年}
\noindent\eventanchor{XIA-1080-EVENT-064}{宋神宗西北进取使宋夏边境冲突升级}\textbf{西夏与宋。}
1080年，《宋史》〈夏国传〉及同年本纪；《辽史》《金史》或《元史》相关条；西夏文文书、碑刻与城址材料（据事核） 所见的〔宋神宗西北进取使宋夏边境冲突升级〕把 西夏与宋 的处置留在可复核的年序中；材料未将地点细化到城址，不能用中枢所在地替它补足。

前一层压力可由以下线索辨认：“宋神宗西北进取使宋夏边境冲突升级”发生在西夏建国与宋辽边境战争的具体压力与机会之中，相关行动者的选择不能脱离此前事件链解释。 这一处置留下的后续条件是：“宋神宗西北进取使宋夏边境冲突升级”为后续政治、边防或资源配置留下条件；其社会影响与实际覆盖范围仍须以异类材料限定。

西夏无存世连续国史，汉文叙事多出自对手；西夏文文书的释读、定年和地域代表性须逐项说明。 因此，西夏文文书、河西—宁夏考古与宋夏金蒙关系研究 只能扩展问题，不能代替未保存的细节。

\subsection{1081年（宋元豐四年、辽太康7年；公元换算据两国年号对照）}
\noindent\eventanchor{SYD-1081-LIAO-066}{戊寅，遼遣蕭福全等來賀正旦}\textbf{宋与辽。}
宋与辽 在 1081年（宋元豐四年、辽太康7年；公元换算据两国年号对照） 的材料痕迹是〔戊寅，遼遣蕭福全等來賀正旦〕，出处为 《宋史》卷016〈元豐四年〉；《辽史》本纪同年条（互校）。材料未将地点细化到城址，不能用中枢所在地替它补足。

本纪年条提供日期和中枢可见行动；前因不据单条补定。 记录所见的结果则是：可回查同年及后续条目，地方影响仍须另证。 日期相邻不等于因果已经证实。

关于可说范围，须保留：《宋史》为元末官修，本纪保存宋廷纪年与处置；战果、数字、地方执行及对方动机须以编年、会要、对方史料或地方材料互校。 进一步讨论可参照 宋辽边政、使节记录与契丹碑刻研究

\noindent\eventanchor{SYD-1081-LIAO-132}{夏四月癸亥，遼遣耶律祐等來賀同天節}\textbf{宋与辽。}
从 《宋史》卷016〈元豐四年〉；《辽史》本纪同年条（互校） 的 1081年（宋元豐四年、辽太康7年；公元换算据两国年号对照） 条目看，〔夏四月癸亥，遼遣耶律祐等來賀同天節〕使 宋与辽 的一个具体行动进入叙事；材料未将地点细化到城址，不能用中枢所在地替它补足。

把本项放回事件链，能够看到的前提为：本纪年条提供日期和中枢可见行动；前因不据单条补定。 而它所打开或收紧的下一步是：可回查同年及后续条目，地方影响仍须另证。

证据的边界是：《宋史》为元末官修，本纪保存宋廷纪年与处置；战果、数字、地方执行及对方动机须以编年、会要、对方史料或地方材料互校。 宋辽边政、使节记录与契丹碑刻研究 可用于继续检验这一结论。

\subsection{1081年（宋元豐四年、西夏大安7年；公元换算据两国年号对照）}
\noindent\eventanchor{SYD-1081-XIA-030}{-{种}-諤遣曲珍等領兵通黑水安定堡，路遇夏人，與戰，破之，斬獲甚衆}\textbf{宋与西夏。}
1081年（宋元豐四年、西夏大安7年；公元换算据两国年号对照） 的记载将 宋与西夏 与〔-{种}-諤遣曲珍等領兵通黑水安定堡，路遇夏人，與戰，破之，斬獲甚衆〕相连。此处可确认的是这项被记录的行动；题名保留的城、路、水道或机构名称，是目前可以落地的空间线索。

它承接的条件是：本纪保留灾害或工程的报告地点与朝廷处置；灾情范围需另证。 随后可追踪的变化为：可与地方文书、河工和环境材料比较，不将单点记录外推为全域因果。

《宋史》为元末官修，本纪保存宋廷纪年与处置；战果、数字、地方执行及对方动机须以编年、会要、对方史料或地方材料互校。 因此，西夏文书、河西—宁夏考古与宋夏关系研究 只能扩展问题，不能代替未保存的细节。

\noindent\eventanchor{SYD-1081-XIA-061}{丙午，高遵裕以師還，夏人來追，遂潰}\textbf{宋与西夏。}
〔丙午，高遵裕以師還，夏人來追，遂潰〕见于 《宋史》卷016〈元豐四年〉；西夏文书、碑刻及《辽史》或《金史》同年条（互校） 的 1081年（宋元豐四年、西夏大安7年；公元换算据两国年号对照） 记录。宋与西夏 的选择因此有了可回查的时间位置；材料未将地点细化到城址，不能用中枢所在地替它补足。

前一层压力可由以下线索辨认：本纪年条提供日期和中枢可见行动；前因不据单条补定。 这一处置留下的后续条件是：可回查同年及后续条目，地方影响仍须另证。

关于可说范围，须保留：《宋史》为元末官修，本纪保存宋廷纪年与处置；战果、数字、地方执行及对方动机须以编年、会要、对方史料或地方材料互校。 进一步讨论可参照 西夏文书、河西—宁夏考古与宋夏关系研究

\noindent\eventanchor{SYD-1081-XIA-079}{丙午，詔諭夏主左右并嵬名部族諸部首領，并許自歸}\textbf{宋与西夏。}
1081年（宋元豐四年、西夏大安7年；公元换算据两国年号对照），《宋史》卷016〈元豐四年〉；西夏文书、碑刻及《辽史》或《金史》同年条（互校） 所见的〔丙午，詔諭夏主左右并嵬名部族諸部首領，并許自歸〕把 宋与西夏 的处置留在可复核的年序中；材料未将地点细化到城址，不能用中枢所在地替它补足。

本纪年条记录政令或机构处置；实施背景须参照制度材料。 记录所见的结果则是：可在同年及后续政令中追踪；条文不单独证明地方执行。 日期相邻不等于因果已经证实。

证据的边界是：《宋史》为元末官修，本纪保存宋廷纪年与处置；战果、数字、地方执行及对方动机须以编年、会要、对方史料或地方材料互校。 西夏文书、河西—宁夏考古与宋夏关系研究 可用于继续检验这一结论。

\noindent\eventanchor{SYD-1081-XIA-090}{丁-{丑}-，曲珍與夏人戰於蒲桃山，敗之}\textbf{宋与西夏。}
宋与西夏 在 1081年（宋元豐四年、西夏大安7年；公元换算据两国年号对照） 的材料痕迹是〔丁-{丑}-，曲珍與夏人戰於蒲桃山，敗之〕，出处为 《宋史》卷016〈元豐四年〉；西夏文书、碑刻及《辽史》或《金史》同年条（互校）。题名保留的城、路、水道或机构名称，是目前可以落地的空间线索。

把本项放回事件链，能够看到的前提为：本纪从朝廷视角记录军政行动；出兵与战果需分辨。 而它所打开或收紧的下一步是：后续路线、控制与损失应与对方记录和遗址材料复核。

《宋史》为元末官修，本纪保存宋廷纪年与处置；战果、数字、地方执行及对方动机须以编年、会要、对方史料或地方材料互校。 因此，西夏文书、河西—宁夏考古与宋夏关系研究 只能扩展问题，不能代替未保存的细节。

\noindent\eventanchor{SYD-1081-XIA-096}{丁-{丑}-，熙河經制李憲敗夏人於西市新城，獲酋首三人、首領二十余人}\textbf{宋与西夏。}
从 《宋史》卷016〈元豐四年〉；西夏文书、碑刻及《辽史》或《金史》同年条（互校） 的 1081年（宋元豐四年、西夏大安7年；公元换算据两国年号对照） 条目看，〔丁-{丑}-，熙河經制李憲敗夏人於西市新城，獲酋首三人、首領二十余人〕使 宋与西夏 的一个具体行动进入叙事；题名保留的城、路、水道或机构名称，是目前可以落地的空间线索。

它承接的条件是：本纪保留灾害或工程的报告地点与朝廷处置；灾情范围需另证。 随后可追踪的变化为：可与地方文书、河工和环境材料比较，不将单点记录外推为全域因果。

关于可说范围，须保留：《宋史》为元末官修，本纪保存宋廷纪年与处置；战果、数字、地方执行及对方动机须以编年、会要、对方史料或地方材料互校。 进一步讨论可参照 西夏文书、河西—宁夏考古与宋夏关系研究

\noindent\eventanchor{SYD-1081-XIA-100}{丁亥，諸軍合攻靈州，-{种}-諤敗夏人于黑水}\textbf{宋与西夏。}
1081年（宋元豐四年、西夏大安7年；公元换算据两国年号对照） 的记载将 宋与西夏 与〔丁亥，諸軍合攻靈州，-{种}-諤敗夏人于黑水〕相连。此处可确认的是这项被记录的行动；题名保留的城、路、水道或机构名称，是目前可以落地的空间线索。

前一层压力可由以下线索辨认：本纪保留灾害或工程的报告地点与朝廷处置；灾情范围需另证。 这一处置留下的后续条件是：可与地方文书、河工和环境材料比较，不将单点记录外推为全域因果。

证据的边界是：《宋史》为元末官修，本纪保存宋廷纪年与处置；战果、数字、地方执行及对方动机须以编年、会要、对方史料或地方材料互校。 西夏文书、河西—宁夏考古与宋夏关系研究 可用于继续检验这一结论。

\noindent\eventanchor{SYD-1081-XIA-104}{庚申，熙河兵至女遮谷，與夏人遇，戰敗之}\textbf{宋与西夏。}
〔庚申，熙河兵至女遮谷，與夏人遇，戰敗之〕见于 《宋史》卷016〈元豐四年〉；西夏文书、碑刻及《辽史》或《金史》同年条（互校） 的 1081年（宋元豐四年、西夏大安7年；公元换算据两国年号对照） 记录。宋与西夏 的选择因此有了可回查的时间位置；题名保留的城、路、水道或机构名称，是目前可以落地的空间线索。

本纪保留灾害或工程的报告地点与朝廷处置；灾情范围需另证。 记录所见的结果则是：可与地方文书、河工和环境材料比较，不将单点记录外推为全域因果。 日期相邻不等于因果已经证实。

《宋史》为元末官修，本纪保存宋廷纪年与处置；战果、数字、地方执行及对方动机须以编年、会要、对方史料或地方材料互校。 因此，西夏文书、河西—宁夏考古与宋夏关系研究 只能扩展问题，不能代替未保存的细节。

\noindent\eventanchor{SYD-1081-XIA-108}{庚戌，夏兵救米脂砦，鄜延經略副使-{种}-諤率衆擊破之}\textbf{宋与西夏。}
1081年（宋元豐四年、西夏大安7年；公元换算据两国年号对照），《宋史》卷016〈元豐四年〉；西夏文书、碑刻及《辽史》或《金史》同年条（互校） 所见的〔庚戌，夏兵救米脂砦，鄜延經略副使-{种}-諤率衆擊破之〕把 宋与西夏 的处置留在可复核的年序中；材料未将地点细化到城址，不能用中枢所在地替它补足。

把本项放回事件链，能够看到的前提为：本纪所记财用、赈济或征发属于特定报告场景。 而它所打开或收紧的下一步是：可与食货志、文书和物质材料比较，不外推为全境总量。

关于可说范围，须保留：《宋史》为元末官修，本纪保存宋廷纪年与处置；战果、数字、地方执行及对方动机须以编年、会要、对方史料或地方材料互校。 进一步讨论可参照 西夏文书、河西—宁夏考古与宋夏关系研究

\noindent\eventanchor{SYD-1081-XIA-110}{庚寅，西邊守臣言夏人囚其主秉常，詔陝西、河東路討之}\textbf{宋与西夏。}
宋与西夏 在 1081年（宋元豐四年、西夏大安7年；公元换算据两国年号对照） 的材料痕迹是〔庚寅，西邊守臣言夏人囚其主秉常，詔陝西、河東路討之〕，出处为 《宋史》卷016〈元豐四年〉；西夏文书、碑刻及《辽史》或《金史》同年条（互校）。题名保留的城、路、水道或机构名称，是目前可以落地的空间线索。

它承接的条件是：本纪年条记录政令或机构处置；实施背景须参照制度材料。 随后可追踪的变化为：可在同年及后续政令中追踪；条文不单独证明地方执行。

证据的边界是：《宋史》为元末官修，本纪保存宋廷纪年与处置；战果、数字、地方执行及对方动机须以编年、会要、对方史料或地方材料互校。 西夏文书、河西—宁夏考古与宋夏关系研究 可用于继续检验这一结论。

\subsection{1082年（宋元豐五年、辽太康8年；公元换算据两国年号对照）}
\noindent\eventanchor{SYD-1082-LIAO-067}{丁巳，遼遣耶律永端等來賀同天節}\textbf{宋与辽。}
从 《宋史》卷016〈元豐五年〉；《辽史》本纪同年条（互校） 的 1082年（宋元豐五年、辽太康8年；公元换算据两国年号对照） 条目看，〔丁巳，遼遣耶律永端等來賀同天節〕使 宋与辽 的一个具体行动进入叙事；材料未将地点细化到城址，不能用中枢所在地替它补足。

前一层压力可由以下线索辨认：本纪年条提供日期和中枢可见行动；前因不据单条补定。 这一处置留下的后续条件是：可回查同年及后续条目，地方影响仍须另证。

《宋史》为元末官修，本纪保存宋廷纪年与处置；战果、数字、地方执行及对方动机须以编年、会要、对方史料或地方材料互校。 因此，宋辽边政、使节记录与契丹碑刻研究 只能扩展问题，不能代替未保存的细节。

\noindent\eventanchor{SYD-1082-LIAO-133}{壬申，遼遣耶律儀等來賀正旦}\textbf{宋与辽。}
1082年（宋元豐五年、辽太康8年；公元换算据两国年号对照） 的记载将 宋与辽 与〔壬申，遼遣耶律儀等來賀正旦〕相连。此处可确认的是这项被记录的行动；材料未将地点细化到城址，不能用中枢所在地替它补足。

本纪年条提供日期和中枢可见行动；前因不据单条补定。 记录所见的结果则是：可回查同年及后续条目，地方影响仍须另证。 日期相邻不等于因果已经证实。

关于可说范围，须保留：《宋史》为元末官修，本纪保存宋廷纪年与处置；战果、数字、地方执行及对方动机须以编年、会要、对方史料或地方材料互校。 进一步讨论可参照 宋辽边政、使节记录与契丹碑刻研究

\subsection{1082年（宋元豐五年、西夏大安8年；公元换算据两国年号对照）}
\noindent\eventanchor{SYD-1082-XIA-031}{九月丁亥，夏人三十萬衆寇永樂，曲珍戰不利，裨將寇偉等死之，夏人遂圍城}\textbf{宋与西夏。}
〔九月丁亥，夏人三十萬衆寇永樂，曲珍戰不利，裨將寇偉等死之，夏人遂圍城〕见于 《宋史》卷016〈元豐五年〉；西夏文书、碑刻及《辽史》或《金史》同年条（互校） 的 1082年（宋元豐五年、西夏大安8年；公元换算据两国年号对照） 记录。宋与西夏 的选择因此有了可回查的时间位置；题名保留的城、路、水道或机构名称，是目前可以落地的空间线索。

把本项放回事件链，能够看到的前提为：本纪从朝廷视角记录军政行动；出兵与战果需分辨。 而它所打开或收紧的下一步是：后续路线、控制与损失应与对方记录和遗址材料复核。

证据的边界是：《宋史》为元末官修，本纪保存宋廷纪年与处置；战果、数字、地方执行及对方动机须以编年、会要、对方史料或地方材料互校。 西夏文书、河西—宁夏考古与宋夏关系研究 可用于继续检验这一结论。

\noindent\eventanchor{SYD-1082-XIA-062}{六月辛亥朔，環慶經略司遣將與夏人戰，破之，斬其統軍嵬名妹精嵬、副統軍訛勃遇}\textbf{宋与西夏。}
1082年（宋元豐五年、西夏大安8年；公元换算据两国年号对照），《宋史》卷016〈元豐五年〉；西夏文书、碑刻及《辽史》或《金史》同年条（互校） 所见的〔六月辛亥朔，環慶經略司遣將與夏人戰，破之，斬其統軍嵬名妹精嵬、副統軍訛勃遇〕把 宋与西夏 的处置留在可复核的年序中；材料未将地点细化到城址，不能用中枢所在地替它补足。

它承接的条件是：本纪从朝廷视角记录军政行动；出兵与战果需分辨。 随后可追踪的变化为：后续路线、控制与损失应与对方记录和遗址材料复核。

《宋史》为元末官修，本纪保存宋廷纪年与处置；战果、数字、地方执行及对方动机须以编年、会要、对方史料或地方材料互校。 因此，西夏文书、河西—宁夏考古与宋夏关系研究 只能扩展问题，不能代替未保存的细节。

\noindent\eventanchor{SYD-1082-XIA-080}{壬寅，鄜延路副總管曲珍敗夏人于金湯}\textbf{宋与西夏。}
宋与西夏 在 1082年（宋元豐五年、西夏大安8年；公元换算据两国年号对照） 的材料痕迹是〔壬寅，鄜延路副總管曲珍敗夏人于金湯〕，出处为 《宋史》卷016〈元豐五年〉；西夏文书、碑刻及《辽史》或《金史》同年条（互校）。题名保留的城、路、水道或机构名称，是目前可以落地的空间线索。

前一层压力可由以下线索辨认：本纪从朝廷视角记录军政行动；出兵与战果需分辨。 这一处置留下的后续条件是：后续路线、控制与损失应与对方记录和遗址材料复核。

关于可说范围，须保留：《宋史》为元末官修，本纪保存宋廷纪年与处置；战果、数字、地方执行及对方动机须以编年、会要、对方史料或地方材料互校。 进一步讨论可参照 西夏文书、河西—宁夏考古与宋夏关系研究

\noindent\eventanchor{SYD-1082-XIA-091}{戊寅，曲珍等敗夏人於明堂川}\textbf{宋与西夏。}
从 《宋史》卷016〈元豐五年〉；西夏文书、碑刻及《辽史》或《金史》同年条（互校） 的 1082年（宋元豐五年、西夏大安8年；公元换算据两国年号对照） 条目看，〔戊寅，曲珍等敗夏人於明堂川〕使 宋与西夏 的一个具体行动进入叙事；题名保留的城、路、水道或机构名称，是目前可以落地的空间线索。

本纪从朝廷视角记录军政行动；出兵与战果需分辨。 记录所见的结果则是：后续路线、控制与损失应与对方记录和遗址材料复核。 日期相邻不等于因果已经证实。

证据的边界是：《宋史》为元末官修，本纪保存宋廷纪年与处置；战果、数字、地方执行及对方动机须以编年、会要、对方史料或地方材料互校。 西夏文书、河西—宁夏考古与宋夏关系研究 可用于继续检验这一结论。

\subsection{1082年}
\noindent\eventanchor{XIA-1082-EVENT-065}{西夏攻破永乐城}\textbf{西夏与宋。}
1082年 的记载将 西夏与宋 与〔西夏攻破永乐城〕相连。此处可确认的是这项被记录的行动；题名保留的城、路、水道或机构名称，是目前可以落地的空间线索。

把本项放回事件链，能够看到的前提为：西夏建国与宋辽边境战争中既有的联盟、交通与补给压力，为“西夏攻破永乐城”提供了直接的军事或动员背景。 而它所打开或收紧的下一步是：“西夏攻破永乐城”改变了相关城镇、道路或政权间的军事选择；伤亡、俘获和控制范围仅可按各方材料分别确认。

西夏无存世连续国史，汉文叙事多出自对手；西夏文文书的释读、定年和地域代表性须逐项说明。 因此，西夏文文书、河西—宁夏考古与宋夏金蒙关系研究 只能扩展问题，不能代替未保存的细节。

\subsection{1083 年（共享年序桥接）}
\noindent\eventanchor{SY-1083-CHRONOLOGY-BRIDGE}{【C级年序桥接】保留1083 年的多政权并行检索入口；本条不主张所列政权在当年发生直接互动，具体事件须另以单项材料入账。}\textbf{宋、辽与西夏（年序桥接）。}
〔【C级年序桥接】保留1083 年的多政权并行检索入口；本条不主张所列政权在当年发生直接互动，具体事件须另以单项材料入账。〕见于 《宋史》；《辽史》；《续资治通鉴长编》；西夏文书与碑刻（据事核） 的 1083 年（共享年序桥接） 记录。宋、辽与西夏（年序桥接） 的选择因此有了可回查的时间位置；题名保留的城、路、水道或机构名称，是目前可以落地的空间线索。

它承接的条件是：相邻已入账事件之间缺少该年度的共享年序入口，易使跨政权材料被单一政权叙事遮蔽。 随后可追踪的变化为：保留该年度的检索与补账位置；后续仅在获得可核单项材料时拆分为具体政治、财政、军事、社会或文化事件。

关于可说范围，须保留：本条只确认该年位于可追踪的共享年序中；正史、地方文书和外部材料的保存密度不一，不能由桥接标签推断行动、统属、疆界或社会影响。 进一步讨论可参照 宋辽夏边境、财政与文书材料的并行年序研究

\subsection{1083年（宋元豐六年、辽太康9年；公元换算据两国年号对照）}
\noindent\eventanchor{SYD-1083-LIAO-068}{辛亥，遼遣蕭固等來賀同天節}\textbf{宋与辽。}
1083年（宋元豐六年、辽太康9年；公元换算据两国年号对照），《宋史》卷016〈元豐六年〉；《辽史》本纪同年条（互校） 所见的〔辛亥，遼遣蕭固等來賀同天節〕把 宋与辽 的处置留在可复核的年序中；材料未将地点细化到城址，不能用中枢所在地替它补足。

前一层压力可由以下线索辨认：本纪年条提供日期和中枢可见行动；前因不据单条补定。 这一处置留下的后续条件是：可回查同年及后续条目，地方影响仍须另证。

证据的边界是：《宋史》为元末官修，本纪保存宋廷纪年与处置；战果、数字、地方执行及对方动机须以编年、会要、对方史料或地方材料互校。 宋辽边政、使节记录与契丹碑刻研究 可用于继续检验这一结论。

\noindent\eventanchor{SYD-1083-LIAO-134}{乙酉，遣蔡京等賀遼主生辰、正旦}\textbf{宋与辽。}
宋与辽 在 1083年（宋元豐六年、辽太康9年；公元换算据两国年号对照） 的材料痕迹是〔乙酉，遣蔡京等賀遼主生辰、正旦〕，出处为 《宋史》卷016〈元豐六年〉；《辽史》本纪同年条（互校）。题名保留的城、路、水道或机构名称，是目前可以落地的空间线索。

本纪年条提供日期和中枢可见行动；前因不据单条补定。 记录所见的结果则是：可回查同年及后续条目，地方影响仍须另证。 日期相邻不等于因果已经证实。

《宋史》为元末官修，本纪保存宋廷纪年与处置；战果、数字、地方执行及对方动机须以编年、会要、对方史料或地方材料互校。 因此，宋辽边政、使节记录与契丹碑刻研究 只能扩展问题，不能代替未保存的细节。

\subsection{1083年（宋元豐六年、西夏大安9年；公元换算据两国年号对照）}
\noindent\eventanchor{SYD-1083-XIA-032}{丙申，河東將薛義敗夏人于葭蘆西嶺}\textbf{宋与西夏。}
从 《宋史》卷016〈元豐六年〉；西夏文书、碑刻及《辽史》或《金史》同年条（互校） 的 1083年（宋元豐六年、西夏大安9年；公元换算据两国年号对照） 条目看，〔丙申，河東將薛義敗夏人于葭蘆西嶺〕使 宋与西夏 的一个具体行动进入叙事；题名保留的城、路、水道或机构名称，是目前可以落地的空间线索。

把本项放回事件链，能够看到的前提为：本纪保留灾害或工程的报告地点与朝廷处置；灾情范围需另证。 而它所打开或收紧的下一步是：可与地方文书、河工和环境材料比较，不将单点记录外推为全域因果。

关于可说范围，须保留：《宋史》为元末官修，本纪保存宋廷纪年与处置；战果、数字、地方执行及对方动机须以编年、会要、对方史料或地方材料互校。 进一步讨论可参照 西夏文书、河西—宁夏考古与宋夏关系研究

\noindent\eventanchor{SYD-1083-XIA-063}{二月丁未，夏人數十萬衆攻蘭州，鈐轄王文鬱率死士七百餘人擊走之}\textbf{宋与西夏。}
1083年（宋元豐六年、西夏大安9年；公元换算据两国年号对照） 的记载将 宋与西夏 与〔二月丁未，夏人數十萬衆攻蘭州，鈐轄王文鬱率死士七百餘人擊走之〕相连。此处可确认的是这项被记录的行动；题名保留的城、路、水道或机构名称，是目前可以落地的空间线索。

它承接的条件是：本纪从朝廷视角记录军政行动；出兵与战果需分辨。 随后可追踪的变化为：后续路线、控制与损失应与对方记录和遗址材料复核。

证据的边界是：《宋史》为元末官修，本纪保存宋廷纪年与处置；战果、数字、地方执行及对方动机须以编年、会要、对方史料或地方材料互校。 西夏文书、河西—宁夏考古与宋夏关系研究 可用于继续检验这一结论。

\noindent\eventanchor{SYD-1083-XIA-081}{己亥，河東將高永翼敗夏人於真卿流部}\textbf{宋与西夏。}
〔己亥，河東將高永翼敗夏人於真卿流部〕见于 《宋史》卷016〈元豐六年〉；西夏文书、碑刻及《辽史》或《金史》同年条（互校） 的 1083年（宋元豐六年、西夏大安9年；公元换算据两国年号对照） 记录。宋与西夏 的选择因此有了可回查的时间位置；题名保留的城、路、水道或机构名称，是目前可以落地的空间线索。

前一层压力可由以下线索辨认：本纪保留灾害或工程的报告地点与朝廷处置；灾情范围需另证。 这一处置留下的后续条件是：可与地方文书、河工和环境材料比较，不将单点记录外推为全域因果。

《宋史》为元末官修，本纪保存宋廷纪年与处置；战果、数字、地方执行及对方动机须以编年、会要、对方史料或地方材料互校。 因此，西夏文书、河西—宁夏考古与宋夏关系研究 只能扩展问题，不能代替未保存的细节。

\noindent\eventanchor{SYD-1083-XIA-092}{甲午，夏人寇蘭州，右侍禁韋定死之}\textbf{宋与西夏。}
1083年（宋元豐六年、西夏大安9年；公元换算据两国年号对照），《宋史》卷016〈元豐六年〉；西夏文书、碑刻及《辽史》或《金史》同年条（互校） 所见的〔甲午，夏人寇蘭州，右侍禁韋定死之〕把 宋与西夏 的处置留在可复核的年序中；题名保留的城、路、水道或机构名称，是目前可以落地的空间线索。

本纪年条记录政令或机构处置；实施背景须参照制度材料。 记录所见的结果则是：可在同年及后续政令中追踪；条文不单独证明地方执行。 日期相邻不等于因果已经证实。

关于可说范围，须保留：《宋史》为元末官修，本纪保存宋廷纪年与处置；战果、数字、地方执行及对方动机须以编年、会要、对方史料或地方材料互校。 进一步讨论可参照 西夏文书、河西—宁夏考古与宋夏关系研究

\noindent\eventanchor{SYD-1083-XIA-097}{麟、府州將郭忠詔等敗夏人於乜離抑部，詔行賞有差}\textbf{宋与西夏。}
宋与西夏 在 1083年（宋元豐六年、西夏大安9年；公元换算据两国年号对照） 的材料痕迹是〔麟、府州將郭忠詔等敗夏人於乜離抑部，詔行賞有差〕，出处为 《宋史》卷016〈元豐六年〉；西夏文书、碑刻及《辽史》或《金史》同年条（互校）。题名保留的城、路、水道或机构名称，是目前可以落地的空间线索。

把本项放回事件链，能够看到的前提为：本纪年条记录政令或机构处置；实施背景须参照制度材料。 而它所打开或收紧的下一步是：可在同年及后续政令中追踪；条文不单独证明地方执行。

证据的边界是：《宋史》为元末官修，本纪保存宋廷纪年与处置；战果、数字、地方执行及对方动机须以编年、会要、对方史料或地方材料互校。 西夏文书、河西—宁夏考古与宋夏关系研究 可用于继续检验这一结论。

\noindent\eventanchor{SYD-1083-XIA-101}{閏月乙亥朔，夏主秉常請修貢，許之}\textbf{宋与西夏。}
从 《宋史》卷016〈元豐六年〉；西夏文书、碑刻及《辽史》或《金史》同年条（互校） 的 1083年（宋元豐六年、西夏大安9年；公元换算据两国年号对照） 条目看，〔閏月乙亥朔，夏主秉常請修貢，許之〕使 宋与西夏 的一个具体行动进入叙事；材料未将地点细化到城址，不能用中枢所在地替它补足。

它承接的条件是：本纪年条提供日期和中枢可见行动；前因不据单条补定。 随后可追踪的变化为：可回查同年及后续条目，地方影响仍须另证。

《宋史》为元末官修，本纪保存宋廷纪年与处置；战果、数字、地方执行及对方动机须以编年、会要、对方史料或地方材料互校。 因此，西夏文书、河西—宁夏考古与宋夏关系研究 只能扩展问题，不能代替未保存的细节。

\noindent\eventanchor{SYD-1083-XIA-105}{三月辛卯，夏人寇蘭州，副總管李浩以衛城有功，復隴州團練使}\textbf{宋与西夏。}
1083年（宋元豐六年、西夏大安9年；公元换算据两国年号对照） 的记载将 宋与西夏 与〔三月辛卯，夏人寇蘭州，副總管李浩以衛城有功，復隴州團練使〕相连。此处可确认的是这项被记录的行动；题名保留的城、路、水道或机构名称，是目前可以落地的空间线索。

前一层压力可由以下线索辨认：本纪年条记录政令或机构处置；实施背景须参照制度材料。 这一处置留下的后续条件是：可在同年及后续政令中追踪；条文不单独证明地方执行。

关于可说范围，须保留：《宋史》为元末官修，本纪保存宋廷纪年与处置；战果、数字、地方执行及对方动机须以编年、会要、对方史料或地方材料互校。 进一步讨论可参照 西夏文书、河西—宁夏考古与宋夏关系研究

\noindent\eventanchor{SYD-1083-XIA-109}{是月，夏人寇麟州，知州訾虎敗之}\textbf{宋与西夏。}
〔是月，夏人寇麟州，知州訾虎敗之〕见于 《宋史》卷016〈元豐六年〉；西夏文书、碑刻及《辽史》或《金史》同年条（互校） 的 1083年（宋元豐六年、西夏大安9年；公元换算据两国年号对照） 记录。宋与西夏 的选择因此有了可回查的时间位置；题名保留的城、路、水道或机构名称，是目前可以落地的空间线索。

本纪从朝廷视角记录军政行动；出兵与战果需分辨。 记录所见的结果则是：后续路线、控制与损失应与对方记录和遗址材料复核。 日期相邻不等于因果已经证实。

证据的边界是：《宋史》为元末官修，本纪保存宋廷纪年与处置；战果、数字、地方执行及对方动机须以编年、会要、对方史料或地方材料互校。 西夏文书、河西—宁夏考古与宋夏关系研究 可用于继续检验这一结论。

\subsection{1084年（宋元豐七年、辽太康10年；公元换算据两国年号对照）}
\noindent\eventanchor{SYD-1084-LIAO-069}{辛卯，遼遣耶律襄等來賀正旦}\textbf{宋与辽。}
1084年（宋元豐七年、辽太康10年；公元换算据两国年号对照），《宋史》卷016〈元豐七年〉；《辽史》本纪同年条（互校） 所见的〔辛卯，遼遣耶律襄等來賀正旦〕把 宋与辽 的处置留在可复核的年序中；材料未将地点细化到城址，不能用中枢所在地替它补足。

把本项放回事件链，能够看到的前提为：本纪年条提供日期和中枢可见行动；前因不据单条补定。 而它所打开或收紧的下一步是：可回查同年及后续条目，地方影响仍须另证。

《宋史》为元末官修，本纪保存宋廷纪年与处置；战果、数字、地方执行及对方动机须以编年、会要、对方史料或地方材料互校。 因此，宋辽边政、使节记录与契丹碑刻研究 只能扩展问题，不能代替未保存的细节。

\noindent\eventanchor{SYD-1084-LIAO-135}{辛巳，遣陳睦等賀遼主生辰、正旦}\textbf{宋与辽。}
宋与辽 在 1084年（宋元豐七年、辽太康10年；公元换算据两国年号对照） 的材料痕迹是〔辛巳，遣陳睦等賀遼主生辰、正旦〕，出处为 《宋史》卷016〈元豐七年〉；《辽史》本纪同年条（互校）。材料未将地点细化到城址，不能用中枢所在地替它补足。

它承接的条件是：本纪年条提供日期和中枢可见行动；前因不据单条补定。 随后可追踪的变化为：可回查同年及后续条目，地方影响仍须另证。

关于可说范围，须保留：《宋史》为元末官修，本纪保存宋廷纪年与处置；战果、数字、地方执行及对方动机须以编年、会要、对方史料或地方材料互校。 进一步讨论可参照 宋辽边政、使节记录与契丹碑刻研究

\subsection{1084年（宋元豐七年、西夏大安10年；公元换算据两国年号对照）}
\noindent\eventanchor{SYD-1084-XIA-033}{冬十月乙亥，夏人寇熙河}\textbf{宋与西夏。}
从 《宋史》卷016〈元豐七年〉；西夏文书、碑刻及《辽史》或《金史》同年条（互校） 的 1084年（宋元豐七年、西夏大安10年；公元换算据两国年号对照） 条目看，〔冬十月乙亥，夏人寇熙河〕使 宋与西夏 的一个具体行动进入叙事；题名保留的城、路、水道或机构名称，是目前可以落地的空间线索。

前一层压力可由以下线索辨认：本纪保留灾害或工程的报告地点与朝廷处置；灾情范围需另证。 这一处置留下的后续条件是：可与地方文书、河工和环境材料比较，不将单点记录外推为全域因果。

证据的边界是：《宋史》为元末官修，本纪保存宋廷纪年与处置；战果、数字、地方执行及对方动机须以编年、会要、对方史料或地方材料互校。 西夏文书、河西—宁夏考古与宋夏关系研究 可用于继续检验这一结论。

\noindent\eventanchor{SYD-1084-XIA-064}{癸-{丑}-，夏人寇蘭州，李憲等擊走之}\textbf{宋与西夏。}
1084年（宋元豐七年、西夏大安10年；公元换算据两国年号对照） 的记载将 宋与西夏 与〔癸-{丑}-，夏人寇蘭州，李憲等擊走之〕相连。此处可确认的是这项被记录的行动；题名保留的城、路、水道或机构名称，是目前可以落地的空间线索。

本纪从朝廷视角记录军政行动；出兵与战果需分辨。 记录所见的结果则是：后续路线、控制与损失应与对方记录和遗址材料复核。 日期相邻不等于因果已经证实。

《宋史》为元末官修，本纪保存宋廷纪年与处置；战果、数字、地方执行及对方动机须以编年、会要、对方史料或地方材料互校。 因此，西夏文书、河西—宁夏考古与宋夏关系研究 只能扩展问题，不能代替未保存的细节。

\noindent\eventanchor{SYD-1084-XIA-082}{癸巳，夏人寇延州安塞堡，將官呂真敗之}\textbf{宋与西夏。}
〔癸巳，夏人寇延州安塞堡，將官呂真敗之〕见于 《宋史》卷016〈元豐七年〉；西夏文书、碑刻及《辽史》或《金史》同年条（互校） 的 1084年（宋元豐七年、西夏大安10年；公元换算据两国年号对照） 记录。宋与西夏 的选择因此有了可回查的时间位置；题名保留的城、路、水道或机构名称，是目前可以落地的空间线索。

把本项放回事件链，能够看到的前提为：本纪从朝廷视角记录军政行动；出兵与战果需分辨。 而它所打开或收紧的下一步是：后续路线、控制与损失应与对方记录和遗址材料复核。

关于可说范围，须保留：《宋史》为元末官修，本纪保存宋廷纪年与处置；战果、数字、地方执行及对方动机须以编年、会要、对方史料或地方材料互校。 进一步讨论可参照 西夏文书、河西—宁夏考古与宋夏关系研究

\noindent\eventanchor{SYD-1084-XIA-093}{甲辰，夏國主秉常遣使來貢}\textbf{宋与西夏。}
1084年（宋元豐七年、西夏大安10年；公元换算据两国年号对照），《宋史》卷016〈元豐七年〉；西夏文书、碑刻及《辽史》或《金史》同年条（互校） 所见的〔甲辰，夏國主秉常遣使來貢〕把 宋与西夏 的处置留在可复核的年序中；材料未将地点细化到城址，不能用中枢所在地替它补足。

它承接的条件是：本纪记录使节、贡赐或往来，不等于稳定统属关系。 随后可追踪的变化为：可与对方编年和交通、文书材料核对其后续落实。

证据的边界是：《宋史》为元末官修，本纪保存宋廷纪年与处置；战果、数字、地方执行及对方动机须以编年、会要、对方史料或地方材料互校。 西夏文书、河西—宁夏考古与宋夏关系研究 可用于继续检验这一结论。

\noindent\eventanchor{SYD-1084-XIA-098}{六月丙子，夏人寇德順軍，巡檢王友死之}\textbf{宋与西夏。}
宋与西夏 在 1084年（宋元豐七年、西夏大安10年；公元换算据两国年号对照） 的材料痕迹是〔六月丙子，夏人寇德順軍，巡檢王友死之〕，出处为 《宋史》卷016〈元豐七年〉；西夏文书、碑刻及《辽史》或《金史》同年条（互校）。材料未将地点细化到城址，不能用中枢所在地替它补足。

前一层压力可由以下线索辨认：本纪从朝廷视角记录军政行动；出兵与战果需分辨。 这一处置留下的后续条件是：后续路线、控制与损失应与对方记录和遗址材料复核。

《宋史》为元末官修，本纪保存宋廷纪年与处置；战果、数字、地方执行及对方动机须以编年、会要、对方史料或地方材料互校。 因此，西夏文书、河西—宁夏考古与宋夏关系研究 只能扩展问题，不能代替未保存的细节。

\noindent\eventanchor{SYD-1084-XIA-102}{乙丑，夏人圍定西城，熙河將秦貴敗之}\textbf{宋与西夏。}
从 《宋史》卷016〈元豐七年〉；西夏文书、碑刻及《辽史》或《金史》同年条（互校） 的 1084年（宋元豐七年、西夏大安10年；公元换算据两国年号对照） 条目看，〔乙丑，夏人圍定西城，熙河將秦貴敗之〕使 宋与西夏 的一个具体行动进入叙事；题名保留的城、路、水道或机构名称，是目前可以落地的空间线索。

本纪保留灾害或工程的报告地点与朝廷处置；灾情范围需另证。 记录所见的结果则是：可与地方文书、河工和环境材料比较，不将单点记录外推为全域因果。 日期相邻不等于因果已经证实。

关于可说范围，须保留：《宋史》为元末官修，本纪保存宋廷纪年与处置；战果、数字、地方执行及对方动机须以编年、会要、对方史料或地方材料互校。 进一步讨论可参照 西夏文书、河西—宁夏考古与宋夏关系研究

\noindent\eventanchor{SYD-1084-XIA-106}{乙未，夏人寇靜邊砦，涇原將彭孫敗之}\textbf{宋与西夏。}
1084年（宋元豐七年、西夏大安10年；公元换算据两国年号对照） 的记载将 宋与西夏 与〔乙未，夏人寇靜邊砦，涇原將彭孫敗之〕相连。此处可确认的是这项被记录的行动；材料未将地点细化到城址，不能用中枢所在地替它补足。

把本项放回事件链，能够看到的前提为：本纪从朝廷视角记录军政行动；出兵与战果需分辨。 而它所打开或收紧的下一步是：后续路线、控制与损失应与对方记录和遗址材料复核。

证据的边界是：《宋史》为元末官修，本纪保存宋廷纪年与处置；战果、数字、地方执行及对方动机须以编年、会要、对方史料或地方材料互校。 西夏文书、河西—宁夏考古与宋夏关系研究 可用于继续检验这一结论。

\subsection{1085年（宋元豐八年、辽大安1年；公元换算据两国年号对照）}
\noindent\eventanchor{SYD-1085-LIAO-070}{癸酉，遣使賀遼主生辰、正旦}\textbf{宋与辽。}
〔癸酉，遣使賀遼主生辰、正旦〕见于 《宋史》卷017〈元豐八年〉；《辽史》本纪同年条（互校） 的 1085年（宋元豐八年、辽大安1年；公元换算据两国年号对照） 记录。宋与辽 的选择因此有了可回查的时间位置；材料未将地点细化到城址，不能用中枢所在地替它补足。

它承接的条件是：本纪记录使节、贡赐或往来，不等于稳定统属关系。 随后可追踪的变化为：可与对方编年和交通、文书材料核对其后续落实。

《宋史》为元末官修，本纪保存宋廷纪年与处置；战果、数字、地方执行及对方动机须以编年、会要、对方史料或地方材料互校。 因此，宋辽边政、使节记录与契丹碑刻研究 只能扩展问题，不能代替未保存的细节。

\noindent\eventanchor{SYD-1085-LIAO-136}{辛巳，遣使以先帝遺留物遺遼國及告即位，甲申，水部員外郎王諤非職言事，坐罰金}\textbf{宋与辽。}
1085年（宋元豐八年、辽大安1年；公元换算据两国年号对照），《宋史》卷017〈元豐八年〉；《辽史》本纪同年条（互校） 所见的〔辛巳，遣使以先帝遺留物遺遼國及告即位，甲申，水部員外郎王諤非職言事，坐罰金〕把 宋与辽 的处置留在可复核的年序中；材料未将地点细化到城址，不能用中枢所在地替它补足。

前一层压力可由以下线索辨认：本纪保留灾害或工程的报告地点与朝廷处置；灾情范围需另证。 这一处置留下的后续条件是：可与地方文书、河工和环境材料比较，不将单点记录外推为全域因果。

关于可说范围，须保留：《宋史》为元末官修，本纪保存宋廷纪年与处置；战果、数字、地方执行及对方动机须以编年、会要、对方史料或地方材料互校。 进一步讨论可参照 宋辽边政、使节记录与契丹碑刻研究

\subsection{1085年（宋元豐八年、西夏大安11年；公元换算据两国年号对照）}
\noindent\eventanchor{SYD-1085-XIA-034}{丙寅，夏人以其母遺留物、馬、白駝來獻}\textbf{宋与西夏。}
宋与西夏 在 1085年（宋元豐八年、西夏大安11年；公元换算据两国年号对照） 的材料痕迹是〔丙寅，夏人以其母遺留物、馬、白駝來獻〕，出处为 《宋史》卷017〈元豐八年〉；西夏文书、碑刻及《辽史》或《金史》同年条（互校）。材料未将地点细化到城址，不能用中枢所在地替它补足。

本纪年条提供日期和中枢可见行动；前因不据单条补定。 记录所见的结果则是：可回查同年及后续条目，地方影响仍须另证。 日期相邻不等于因果已经证实。

证据的边界是：《宋史》为元末官修，本纪保存宋廷纪年与处置；战果、数字、地方执行及对方动机须以编年、会要、对方史料或地方材料互校。 西夏文书、河西—宁夏考古与宋夏关系研究 可用于继续检验这一结论。

\noindent\eventanchor{SYD-1085-XIA-065}{冬十月甲子，夏國遣使進助山陵馬}\textbf{宋与西夏。}
从 《宋史》卷017〈元豐八年〉；西夏文书、碑刻及《辽史》或《金史》同年条（互校） 的 1085年（宋元豐八年、西夏大安11年；公元换算据两国年号对照） 条目看，〔冬十月甲子，夏國遣使進助山陵馬〕使 宋与西夏 的一个具体行动进入叙事；题名保留的城、路、水道或机构名称，是目前可以落地的空间线索。

把本项放回事件链，能够看到的前提为：本纪记录使节、贡赐或往来，不等于稳定统属关系。 而它所打开或收紧的下一步是：可与对方编年和交通、文书材料核对其后续落实。

《宋史》为元末官修，本纪保存宋廷纪年与处置；战果、数字、地方执行及对方动机须以编年、会要、对方史料或地方材料互校。 因此，西夏文书、河西—宁夏考古与宋夏关系研究 只能扩展问题，不能代替未保存的细节。

\noindent\eventanchor{SYD-1085-XIA-083}{以夏國主母卒，遣使弔祭}\textbf{宋与西夏。}
1085年（宋元豐八年、西夏大安11年；公元换算据两国年号对照） 的记载将 宋与西夏 与〔以夏國主母卒，遣使弔祭〕相连。此处可确认的是这项被记录的行动；材料未将地点细化到城址，不能用中枢所在地替它补足。

它承接的条件是：本纪记录使节、贡赐或往来，不等于稳定统属关系。 随后可追踪的变化为：可与对方编年和交通、文书材料核对其后续落实。

关于可说范围，须保留：《宋史》为元末官修，本纪保存宋廷纪年与处置；战果、数字、地方执行及对方动机须以编年、会要、对方史料或地方材料互校。 进一步讨论可参照 西夏文书、河西—宁夏考古与宋夏关系研究

\subsection{1086年（宋元祐元年、西夏天仪治平1年；公元换算据两国年号对照）}
\noindent\eventanchor{SYD-1086-XIA-035}{庚午，禁邊民與夏人爲市}\textbf{宋与西夏。}
〔庚午，禁邊民與夏人爲市〕见于 《宋史》卷017〈元祐元年〉；西夏文书、碑刻及《辽史》或《金史》同年条（互校） 的 1086年（宋元祐元年、西夏天仪治平1年；公元换算据两国年号对照） 记录。宋与西夏 的选择因此有了可回查的时间位置；题名保留的城、路、水道或机构名称，是目前可以落地的空间线索。

前一层压力可由以下线索辨认：本纪年条记录政令或机构处置；实施背景须参照制度材料。 这一处置留下的后续条件是：可在同年及后续政令中追踪；条文不单独证明地方执行。

证据的边界是：《宋史》为元末官修，本纪保存宋廷纪年与处置；战果、数字、地方执行及对方动机须以编年、会要、对方史料或地方材料互校。 西夏文书、河西—宁夏考古与宋夏关系研究 可用于继续检验这一结论。

\noindent\eventanchor{SYD-1086-XIA-066}{庚午，夏國遣使賀坤成節}\textbf{宋与西夏。}
1086年（宋元祐元年、西夏天仪治平1年；公元换算据两国年号对照），《宋史》卷017〈元祐元年〉；西夏文书、碑刻及《辽史》或《金史》同年条（互校） 所见的〔庚午，夏國遣使賀坤成節〕把 宋与西夏 的处置留在可复核的年序中；材料未将地点细化到城址，不能用中枢所在地替它补足。

本纪记录使节、贡赐或往来，不等于稳定统属关系。 记录所见的结果则是：可与对方编年和交通、文书材料核对其后续落实。 日期相邻不等于因果已经证实。

《宋史》为元末官修，本纪保存宋廷纪年与处置；战果、数字、地方执行及对方动机须以编年、会要、对方史料或地方材料互校。 因此，西夏文书、河西—宁夏考古与宋夏关系研究 只能扩展问题，不能代替未保存的细节。

\subsection{1086年}
\noindent\eventanchor{XIA-1086-EVENT-066}{李乾顺即位}\textbf{西夏。}
西夏 在 1086年 的材料痕迹是〔李乾顺即位〕，出处为 《宋史》〈夏国传〉及同年本纪；《辽史》《金史》或《元史》相关条；西夏文文书、碑刻与城址材料（据事核）。材料未将地点细化到城址，不能用中枢所在地替它补足。

把本项放回事件链，能够看到的前提为：西夏与金、宋的边境与册封关系中前任统治的结束与宫廷、部族或官僚的权力安排相互交织，促成“李乾顺即位”这一继承节点。 而它所打开或收紧的下一步是：继承并不自动消除既有矛盾；“李乾顺即位”后的任官、军权和对外策略需由后续条目追踪。

关于可说范围，须保留：西夏无存世连续国史，汉文叙事多出自对手；西夏文文书的释读、定年和地域代表性须逐项说明。 进一步讨论可参照 西夏文文书、河西—宁夏考古与宋夏金蒙关系研究

\subsection{1087 年（共享年序桥接）}
\noindent\eventanchor{SY-1087-CHRONOLOGY-BRIDGE}{【C级年序桥接】保留1087 年的多政权并行检索入口；本条不主张所列政权在当年发生直接互动，具体事件须另以单项材料入账。}\textbf{宋、辽与西夏（年序桥接）。}
从 《宋史》；《辽史》；《续资治通鉴长编》；西夏文书与碑刻（据事核） 的 1087 年（共享年序桥接） 条目看，〔【C级年序桥接】保留1087 年的多政权并行检索入口；本条不主张所列政权在当年发生直接互动，具体事件须另以单项材料入账。〕使 宋、辽与西夏（年序桥接） 的一个具体行动进入叙事；题名保留的城、路、水道或机构名称，是目前可以落地的空间线索。

它承接的条件是：相邻已入账事件之间缺少该年度的共享年序入口，易使跨政权材料被单一政权叙事遮蔽。 随后可追踪的变化为：保留该年度的检索与补账位置；后续仅在获得可核单项材料时拆分为具体政治、财政、军事、社会或文化事件。

证据的边界是：本条只确认该年位于可追踪的共享年序中；正史、地方文书和外部材料的保存密度不一，不能由桥接标签推断行动、统属、疆界或社会影响。 宋辽夏边境、财政与文书材料的并行年序研究 可用于继续检验这一结论。

\subsection{1087年（宋元祐二年、西夏天仪治平2年；公元换算据两国年号对照）}
\noindent\eventanchor{SYD-1087-XIA-036}{二年春正月乙丑，封秉常子乾順爲夏國主}\textbf{宋与西夏。}
1087年（宋元祐二年、西夏天仪治平2年；公元换算据两国年号对照） 的记载将 宋与西夏 与〔二年春正月乙丑，封秉常子乾順爲夏國主〕相连。此处可确认的是这项被记录的行动；材料未将地点细化到城址，不能用中枢所在地替它补足。

前一层压力可由以下线索辨认：本纪年条提供日期和中枢可见行动；前因不据单条补定。 这一处置留下的后续条件是：可回查同年及后续条目，地方影响仍须另证。

《宋史》为元末官修，本纪保存宋廷纪年与处置；战果、数字、地方执行及对方动机须以编年、会要、对方史料或地方材料互校。 因此，西夏文书、河西—宁夏考古与宋夏关系研究 只能扩展问题，不能代替未保存的细节。

\noindent\eventanchor{SYD-1087-XIA-067}{癸巳，以夏國政亂主幼，強臣乙逋等擅權逆命，詔諸路帥臣嚴兵備之}\textbf{宋与西夏。}
〔癸巳，以夏國政亂主幼，強臣乙逋等擅權逆命，詔諸路帥臣嚴兵備之〕见于 《宋史》卷017〈元祐二年〉；西夏文书、碑刻及《辽史》或《金史》同年条（互校） 的 1087年（宋元祐二年、西夏天仪治平2年；公元换算据两国年号对照） 记录。宋与西夏 的选择因此有了可回查的时间位置；题名保留的城、路、水道或机构名称，是目前可以落地的空间线索。

本纪年条记录政令或机构处置；实施背景须参照制度材料。 记录所见的结果则是：可在同年及后续政令中追踪；条文不单独证明地方执行。 日期相邻不等于因果已经证实。

关于可说范围，须保留：《宋史》为元末官修，本纪保存宋廷纪年与处置；战果、数字、地方执行及对方动机须以编年、会要、对方史料或地方材料互校。 进一步讨论可参照 西夏文书、河西—宁夏考古与宋夏关系研究

\noindent\eventanchor{SYD-1087-XIA-084}{辛-{丑}-，涇原言夏人寇三川諸砦，官軍敗之}\textbf{宋与西夏。}
1087年（宋元祐二年、西夏天仪治平2年；公元换算据两国年号对照），《宋史》卷017〈元祐二年〉；西夏文书、碑刻及《辽史》或《金史》同年条（互校） 所见的〔辛-{丑}-，涇原言夏人寇三川諸砦，官軍敗之〕把 宋与西夏 的处置留在可复核的年序中；题名保留的城、路、水道或机构名称，是目前可以落地的空间线索。

把本项放回事件链，能够看到的前提为：本纪从朝廷视角记录军政行动；出兵与战果需分辨。 而它所打开或收紧的下一步是：后续路线、控制与损失应与对方记录和遗址材料复核。

证据的边界是：《宋史》为元末官修，本纪保存宋廷纪年与处置；战果、数字、地方执行及对方动机须以编年、会要、对方史料或地方材料互校。 西夏文书、河西—宁夏考古与宋夏关系研究 可用于继续检验这一结论。

\subsection{1088 年（共享年序桥接）}
\noindent\eventanchor{SY-1088-CHRONOLOGY-BRIDGE}{【C级年序桥接】保留1088 年的多政权并行检索入口；本条不主张所列政权在当年发生直接互动，具体事件须另以单项材料入账。}\textbf{宋、辽与西夏（年序桥接）。}
宋、辽与西夏（年序桥接） 在 1088 年（共享年序桥接） 的材料痕迹是〔【C级年序桥接】保留1088 年的多政权并行检索入口；本条不主张所列政权在当年发生直接互动，具体事件须另以单项材料入账。〕，出处为 《宋史》；《辽史》；《续资治通鉴长编》；西夏文书与碑刻（据事核）。题名保留的城、路、水道或机构名称，是目前可以落地的空间线索。

它承接的条件是：相邻已入账事件之间缺少该年度的共享年序入口，易使跨政权材料被单一政权叙事遮蔽。 随后可追踪的变化为：保留该年度的检索与补账位置；后续仅在获得可核单项材料时拆分为具体政治、财政、军事、社会或文化事件。

本条只确认该年位于可追踪的共享年序中；正史、地方文书和外部材料的保存密度不一，不能由桥接标签推断行动、统属、疆界或社会影响。 因此，宋辽夏边境、财政与文书材料的并行年序研究 只能扩展问题，不能代替未保存的细节。

\subsection{1088年（宋元祐三年、西夏天仪治平3年；公元换算据两国年号对照）}
\noindent\eventanchor{SYD-1088-XIA-037}{辛-{丑}-，夏人寇塞門砦}\textbf{宋与西夏。}
从 《宋史》卷017〈元祐三年〉；西夏文书、碑刻及《辽史》或《金史》同年条（互校） 的 1088年（宋元祐三年、西夏天仪治平3年；公元换算据两国年号对照） 条目看，〔辛-{丑}-，夏人寇塞門砦〕使 宋与西夏 的一个具体行动进入叙事；材料未将地点细化到城址，不能用中枢所在地替它补足。

前一层压力可由以下线索辨认：本纪从朝廷视角记录军政行动；出兵与战果需分辨。 这一处置留下的后续条件是：后续路线、控制与损失应与对方记录和遗址材料复核。

关于可说范围，须保留：《宋史》为元末官修，本纪保存宋廷纪年与处置；战果、数字、地方执行及对方动机须以编年、会要、对方史料或地方材料互校。 进一步讨论可参照 西夏文书、河西—宁夏考古与宋夏关系研究

\noindent\eventanchor{SYD-1088-XIA-068}{乙亥，夏人寇德靜砦，將官張誠等敗之}\textbf{宋与西夏。}
1088年（宋元祐三年、西夏天仪治平3年；公元换算据两国年号对照） 的记载将 宋与西夏 与〔乙亥，夏人寇德靜砦，將官張誠等敗之〕相连。此处可确认的是这项被记录的行动；材料未将地点细化到城址，不能用中枢所在地替它补足。

本纪从朝廷视角记录军政行动；出兵与战果需分辨。 记录所见的结果则是：后续路线、控制与损失应与对方记录和遗址材料复核。 日期相邻不等于因果已经证实。

证据的边界是：《宋史》为元末官修，本纪保存宋廷纪年与处置；战果、数字、地方执行及对方动机须以编年、会要、对方史料或地方材料互校。 西夏文书、河西—宁夏考古与宋夏关系研究 可用于继续检验这一结论。

\subsection{1089 年（共享年序桥接）}
\noindent\eventanchor{SY-1089-CHRONOLOGY-BRIDGE}{【C级年序桥接】保留1089 年的多政权并行检索入口；本条不主张所列政权在当年发生直接互动，具体事件须另以单项材料入账。}\textbf{宋、辽与西夏（年序桥接）。}
〔【C级年序桥接】保留1089 年的多政权并行检索入口；本条不主张所列政权在当年发生直接互动，具体事件须另以单项材料入账。〕见于 《宋史》；《辽史》；《续资治通鉴长编》；西夏文书与碑刻（据事核） 的 1089 年（共享年序桥接） 记录。宋、辽与西夏（年序桥接） 的选择因此有了可回查的时间位置；题名保留的城、路、水道或机构名称，是目前可以落地的空间线索。

把本项放回事件链，能够看到的前提为：相邻已入账事件之间缺少该年度的共享年序入口，易使跨政权材料被单一政权叙事遮蔽。 而它所打开或收紧的下一步是：保留该年度的检索与补账位置；后续仅在获得可核单项材料时拆分为具体政治、财政、军事、社会或文化事件。

本条只确认该年位于可追踪的共享年序中；正史、地方文书和外部材料的保存密度不一，不能由桥接标签推断行动、统属、疆界或社会影响。 因此，宋辽夏边境、财政与文书材料的并行年序研究 只能扩展问题，不能代替未保存的细节。

\subsection{1089年（宋元祐四年、辽大安5年；公元换算据两国年号对照）}
\noindent\eventanchor{SYD-1089-LIAO-071}{丁-{丑}-，遼國使蕭寅等來賀坤成節，曲宴垂拱殿}\textbf{宋与辽。}
1089年（宋元祐四年、辽大安5年；公元换算据两国年号对照），《宋史》卷017〈元祐四年〉；《辽史》本纪同年条（互校） 所见的〔丁-{丑}-，遼國使蕭寅等來賀坤成節，曲宴垂拱殿〕把 宋与辽 的处置留在可复核的年序中；材料未将地点细化到城址，不能用中枢所在地替它补足。

它承接的条件是：本纪年条提供日期和中枢可见行动；前因不据单条补定。 随后可追踪的变化为：可回查同年及后续条目，地方影响仍须另证。

关于可说范围，须保留：《宋史》为元末官修，本纪保存宋廷纪年与处置；战果、数字、地方执行及对方动机须以编年、会要、对方史料或地方材料互校。 进一步讨论可参照 宋辽边政、使节记录与契丹碑刻研究

\noindent\eventanchor{SYD-1089-LIAO-137}{四年春正月壬申朔，不受朝，群臣及遼使詣東上閣門、內東門拜表賀}\textbf{宋与辽。}
宋与辽 在 1089年（宋元祐四年、辽大安5年；公元换算据两国年号对照） 的材料痕迹是〔四年春正月壬申朔，不受朝，群臣及遼使詣東上閣門、內東門拜表賀〕，出处为 《宋史》卷017〈元祐四年〉；《辽史》本纪同年条（互校）。材料未将地点细化到城址，不能用中枢所在地替它补足。

前一层压力可由以下线索辨认：本纪年条提供日期和中枢可见行动；前因不据单条补定。 这一处置留下的后续条件是：可回查同年及后续条目，地方影响仍须另证。

证据的边界是：《宋史》为元末官修，本纪保存宋廷纪年与处置；战果、数字、地方执行及对方动机须以编年、会要、对方史料或地方材料互校。 宋辽边政、使节记录与契丹碑刻研究 可用于继续检验这一结论。

\subsection{1089年（宋元祐四年、西夏天仪治平4年；公元换算据两国年号对照）}
\noindent\eventanchor{SYD-1089-XIA-038}{甲申，以夏人通好，詔邊將毋生事}\textbf{宋与西夏。}
从 《宋史》卷017〈元祐四年〉；西夏文书、碑刻及《辽史》或《金史》同年条（互校） 的 1089年（宋元祐四年、西夏天仪治平4年；公元换算据两国年号对照） 条目看，〔甲申，以夏人通好，詔邊將毋生事〕使 宋与西夏 的一个具体行动进入叙事；材料未将地点细化到城址，不能用中枢所在地替它补足。

本纪年条记录政令或机构处置；实施背景须参照制度材料。 记录所见的结果则是：可在同年及后续政令中追踪；条文不单独证明地方执行。 日期相邻不等于因果已经证实。

《宋史》为元末官修，本纪保存宋廷纪年与处置；战果、数字、地方执行及对方动机须以编年、会要、对方史料或地方材料互校。 因此，西夏文书、河西—宁夏考古与宋夏关系研究 只能扩展问题，不能代替未保存的细节。

\noindent\eventanchor{SYD-1089-XIA-069}{是歲，夏國、邈黎、大食、麻囉拔國入貢}\textbf{宋与西夏。}
1089年（宋元祐四年、西夏天仪治平4年；公元换算据两国年号对照） 的记载将 宋与西夏 与〔是歲，夏國、邈黎、大食、麻囉拔國入貢〕相连。此处可确认的是这项被记录的行动；材料未将地点细化到城址，不能用中枢所在地替它补足。

把本项放回事件链，能够看到的前提为：本纪记录使节、贡赐或往来，不等于稳定统属关系。 而它所打开或收紧的下一步是：可与对方编年和交通、文书材料核对其后续落实。

关于可说范围，须保留：《宋史》为元末官修，本纪保存宋廷纪年与处置；战果、数字、地方执行及对方动机须以编年、会要、对方史料或地方材料互校。 进一步讨论可参照 西夏文书、河西—宁夏考古与宋夏关系研究

\subsection{1090 年（共享年序桥接）}
\noindent\eventanchor{SY-1090-CHRONOLOGY-BRIDGE}{【C级年序桥接】保留1090 年的多政权并行检索入口；本条不主张所列政权在当年发生直接互动，具体事件须另以单项材料入账。}\textbf{宋、辽与西夏（年序桥接）。}
〔【C级年序桥接】保留1090 年的多政权并行检索入口；本条不主张所列政权在当年发生直接互动，具体事件须另以单项材料入账。〕见于 《宋史》；《辽史》；《续资治通鉴长编》；西夏文书与碑刻（据事核） 的 1090 年（共享年序桥接） 记录。宋、辽与西夏（年序桥接） 的选择因此有了可回查的时间位置；题名保留的城、路、水道或机构名称，是目前可以落地的空间线索。

它承接的条件是：相邻已入账事件之间缺少该年度的共享年序入口，易使跨政权材料被单一政权叙事遮蔽。 随后可追踪的变化为：保留该年度的检索与补账位置；后续仅在获得可核单项材料时拆分为具体政治、财政、军事、社会或文化事件。

证据的边界是：本条只确认该年位于可追踪的共享年序中；正史、地方文书和外部材料的保存密度不一，不能由桥接标签推断行动、统属、疆界或社会影响。 宋辽夏边境、财政与文书材料的并行年序研究 可用于继续检验这一结论。

\subsection{1090年（宋元祐五年、西夏天祐民安1年；公元换算据两国年号对照）}
\noindent\eventanchor{SYD-1090-XIA-039}{乙酉，夏人來議分畫疆界}\textbf{宋与西夏。}
1090年（宋元祐五年、西夏天祐民安1年；公元换算据两国年号对照），《宋史》卷017〈元祐五年〉；西夏文书、碑刻及《辽史》或《金史》同年条（互校） 所见的〔乙酉，夏人來議分畫疆界〕把 宋与西夏 的处置留在可复核的年序中；材料未将地点细化到城址，不能用中枢所在地替它补足。

前一层压力可由以下线索辨认：本纪年条提供日期和中枢可见行动；前因不据单条补定。 这一处置留下的后续条件是：可回查同年及后续条目，地方影响仍须另证。

《宋史》为元末官修，本纪保存宋廷纪年与处置；战果、数字、地方执行及对方动机须以编年、会要、对方史料或地方材料互校。 因此，西夏文书、河西—宁夏考古与宋夏关系研究 只能扩展问题，不能代替未保存的细节。

\subsection{1091 年（共享年序桥接）}
\noindent\eventanchor{SY-1091-CHRONOLOGY-BRIDGE}{【C级年序桥接】保留1091 年的多政权并行检索入口；本条不主张所列政权在当年发生直接互动，具体事件须另以单项材料入账。}\textbf{宋、辽与西夏（年序桥接）。}
宋、辽与西夏（年序桥接） 在 1091 年（共享年序桥接） 的材料痕迹是〔【C级年序桥接】保留1091 年的多政权并行检索入口；本条不主张所列政权在当年发生直接互动，具体事件须另以单项材料入账。〕，出处为 《宋史》；《辽史》；《续资治通鉴长编》；西夏文书与碑刻（据事核）。题名保留的城、路、水道或机构名称，是目前可以落地的空间线索。

相邻已入账事件之间缺少该年度的共享年序入口，易使跨政权材料被单一政权叙事遮蔽。 记录所见的结果则是：保留该年度的检索与补账位置；后续仅在获得可核单项材料时拆分为具体政治、财政、军事、社会或文化事件。 日期相邻不等于因果已经证实。

关于可说范围，须保留：本条只确认该年位于可追踪的共享年序中；正史、地方文书和外部材料的保存密度不一，不能由桥接标签推断行动、统属、疆界或社会影响。 进一步讨论可参照 宋辽夏边境、财政与文书材料的并行年序研究

\subsection{1091年（宋元祐六年、辽大安7年；公元换算据两国年号对照）}
\noindent\eventanchor{SYD-1091-LIAO-072}{六年春正月辛酉朔，不受朝，群臣及遼使詣東上閣門、內東門拜表賀}\textbf{宋与辽。}
从 《宋史》卷017〈元祐六年〉；《辽史》本纪同年条（互校） 的 1091年（宋元祐六年、辽大安7年；公元换算据两国年号对照） 条目看，〔六年春正月辛酉朔，不受朝，群臣及遼使詣東上閣門、內東門拜表賀〕使 宋与辽 的一个具体行动进入叙事；材料未将地点细化到城址，不能用中枢所在地替它补足。

把本项放回事件链，能够看到的前提为：本纪年条提供日期和中枢可见行动；前因不据单条补定。 而它所打开或收紧的下一步是：可回查同年及后续条目，地方影响仍须另证。

证据的边界是：《宋史》为元末官修，本纪保存宋廷纪年与处置；战果、数字、地方执行及对方动机须以编年、会要、对方史料或地方材料互校。 宋辽边政、使节记录与契丹碑刻研究 可用于继续检验这一结论。

\subsection{1091年（宋元祐六年、西夏天祐民安2年；公元换算据两国年号对照）}
\noindent\eventanchor{SYD-1091-XIA-040}{癸-{丑}-，詔鄜延路都監李儀等以違旨夜出兵入界，與夏人戰死，不贈官，餘官降等}\textbf{宋与西夏。}
1091年（宋元祐六年、西夏天祐民安2年；公元换算据两国年号对照） 的记载将 宋与西夏 与〔癸-{丑}-，詔鄜延路都監李儀等以違旨夜出兵入界，與夏人戰死，不贈官，餘官降等〕相连。此处可确认的是这项被记录的行动；题名保留的城、路、水道或机构名称，是目前可以落地的空间线索。

它承接的条件是：本纪年条记录政令或机构处置；实施背景须参照制度材料。 随后可追踪的变化为：可在同年及后续政令中追踪；条文不单独证明地方执行。

《宋史》为元末官修，本纪保存宋廷纪年与处置；战果、数字、地方执行及对方动机须以编年、会要、对方史料或地方材料互校。 因此，西夏文书、河西—宁夏考古与宋夏关系研究 只能扩展问题，不能代替未保存的细节。

\noindent\eventanchor{SYD-1091-XIA-070}{九月丁亥，夏人寇麟、府二州}\textbf{宋与西夏。}
〔九月丁亥，夏人寇麟、府二州〕见于 《宋史》卷017〈元祐六年〉；西夏文书、碑刻及《辽史》或《金史》同年条（互校） 的 1091年（宋元祐六年、西夏天祐民安2年；公元换算据两国年号对照） 记录。宋与西夏 的选择因此有了可回查的时间位置；题名保留的城、路、水道或机构名称，是目前可以落地的空间线索。

前一层压力可由以下线索辨认：本纪从朝廷视角记录军政行动；出兵与战果需分辨。 这一处置留下的后续条件是：后续路线、控制与损失应与对方记录和遗址材料复核。

关于可说范围，须保留：《宋史》为元末官修，本纪保存宋廷纪年与处置；战果、数字、地方执行及对方动机须以编年、会要、对方史料或地方材料互校。 进一步讨论可参照 西夏文书、河西—宁夏考古与宋夏关系研究

\noindent\eventanchor{SYD-1091-XIA-085}{夏人寇熙河蘭岷、鄜延路}\textbf{宋与西夏。}
1091年（宋元祐六年、西夏天祐民安2年；公元换算据两国年号对照），《宋史》卷017〈元祐六年〉；西夏文书、碑刻及《辽史》或《金史》同年条（互校） 所见的〔夏人寇熙河蘭岷、鄜延路〕把 宋与西夏 的处置留在可复核的年序中；题名保留的城、路、水道或机构名称，是目前可以落地的空间线索。

本纪保留灾害或工程的报告地点与朝廷处置；灾情范围需另证。 记录所见的结果则是：可与地方文书、河工和环境材料比较，不将单点记录外推为全域因果。 日期相邻不等于因果已经证实。

证据的边界是：《宋史》为元末官修，本纪保存宋廷纪年与处置；战果、数字、地方执行及对方动机须以编年、会要、对方史料或地方材料互校。 西夏文书、河西—宁夏考古与宋夏关系研究 可用于继续检验这一结论。

\subsection{1092年（宋元祐七年、辽大安8年；公元换算据两国年号对照）}
\noindent\eventanchor{SYD-1092-LIAO-073}{七年春正月甲辰，以遼使耶律迪卒，輟朝一日}\textbf{宋与辽。}
宋与辽 在 1092年（宋元祐七年、辽大安8年；公元换算据两国年号对照） 的材料痕迹是〔七年春正月甲辰，以遼使耶律迪卒，輟朝一日〕，出处为 《宋史》卷017〈元祐七年〉；《辽史》本纪同年条（互校）。材料未将地点细化到城址，不能用中枢所在地替它补足。

把本项放回事件链，能够看到的前提为：本纪年条提供日期和中枢可见行动；前因不据单条补定。 而它所打开或收紧的下一步是：可回查同年及后续条目，地方影响仍须另证。

《宋史》为元末官修，本纪保存宋廷纪年与处置；战果、数字、地方执行及对方动机须以编年、会要、对方史料或地方材料互校。 因此，宋辽边政、使节记录与契丹碑刻研究 只能扩展问题，不能代替未保存的细节。

\subsection{1093年（宋元祐八年、辽大安9年；公元换算据两国年号对照）}
\noindent\eventanchor{SYD-1093-LIAO-074}{丁巳，遼人遣使來弔祭}\textbf{宋与辽。}
从 《宋史》卷017〈元祐八年〉；《辽史》本纪同年条（互校） 的 1093年（宋元祐八年、辽大安9年；公元换算据两国年号对照） 条目看，〔丁巳，遼人遣使來弔祭〕使 宋与辽 的一个具体行动进入叙事；材料未将地点细化到城址，不能用中枢所在地替它补足。

它承接的条件是：本纪记录使节、贡赐或往来，不等于稳定统属关系。 随后可追踪的变化为：可与对方编年和交通、文书材料核对其后续落实。

关于可说范围，须保留：《宋史》为元末官修，本纪保存宋廷纪年与处置；战果、数字、地方执行及对方动机须以编年、会要、对方史料或地方材料互校。 进一步讨论可参照 宋辽边政、使节记录与契丹碑刻研究

\subsection{1093年（宋元祐八年、西夏天祐民安4年；公元换算据两国年号对照）}
\noindent\eventanchor{SYD-1093-XIA-041}{夏四月丁未朔，夏人來謝罪，願以蘭州易塞門砦，不許}\textbf{宋与西夏。}
1093年（宋元祐八年、西夏天祐民安4年；公元换算据两国年号对照） 的记载将 宋与西夏 与〔夏四月丁未朔，夏人來謝罪，願以蘭州易塞門砦，不許〕相连。此处可确认的是这项被记录的行动；题名保留的城、路、水道或机构名称，是目前可以落地的空间线索。

前一层压力可由以下线索辨认：本纪年条提供日期和中枢可见行动；前因不据单条补定。 这一处置留下的后续条件是：可回查同年及后续条目，地方影响仍须另证。

证据的边界是：《宋史》为元末官修，本纪保存宋廷纪年与处置；战果、数字、地方执行及对方动机须以编年、会要、对方史料或地方材料互校。 西夏文书、河西—宁夏考古与宋夏关系研究 可用于继续检验这一结论。

\subsection{1095 年（共享年序桥接）}
\noindent\eventanchor{SY-1095-CHRONOLOGY-BRIDGE}{【C级年序桥接】保留1095 年的多政权并行检索入口；本条不主张所列政权在当年发生直接互动，具体事件须另以单项材料入账。}\textbf{宋、辽与西夏（年序桥接）。}
〔【C级年序桥接】保留1095 年的多政权并行检索入口；本条不主张所列政权在当年发生直接互动，具体事件须另以单项材料入账。〕见于 《宋史》；《辽史》；《续资治通鉴长编》；西夏文书与碑刻（据事核） 的 1095 年（共享年序桥接） 记录。宋、辽与西夏（年序桥接） 的选择因此有了可回查的时间位置；题名保留的城、路、水道或机构名称，是目前可以落地的空间线索。

相邻已入账事件之间缺少该年度的共享年序入口，易使跨政权材料被单一政权叙事遮蔽。 记录所见的结果则是：保留该年度的检索与补账位置；后续仅在获得可核单项材料时拆分为具体政治、财政、军事、社会或文化事件。 日期相邻不等于因果已经证实。

本条只确认该年位于可追踪的共享年序中；正史、地方文书和外部材料的保存密度不一，不能由桥接标签推断行动、统属、疆界或社会影响。 因此，宋辽夏边境、财政与文书材料的并行年序研究 只能扩展问题，不能代替未保存的细节。

\subsection{1096 年（共享年序桥接）}
\noindent\eventanchor{SY-1096-CHRONOLOGY-BRIDGE}{【C级年序桥接】保留1096 年的多政权并行检索入口；本条不主张所列政权在当年发生直接互动，具体事件须另以单项材料入账。}\textbf{宋、辽与西夏（年序桥接）。}
1096 年（共享年序桥接），《宋史》；《辽史》；《续资治通鉴长编》；西夏文书与碑刻（据事核） 所见的〔【C级年序桥接】保留1096 年的多政权并行检索入口；本条不主张所列政权在当年发生直接互动，具体事件须另以单项材料入账。〕把 宋、辽与西夏（年序桥接） 的处置留在可复核的年序中；题名保留的城、路、水道或机构名称，是目前可以落地的空间线索。

把本项放回事件链，能够看到的前提为：相邻已入账事件之间缺少该年度的共享年序入口，易使跨政权材料被单一政权叙事遮蔽。 而它所打开或收紧的下一步是：保留该年度的检索与补账位置；后续仅在获得可核单项材料时拆分为具体政治、财政、军事、社会或文化事件。

关于可说范围，须保留：本条只确认该年位于可追踪的共享年序中；正史、地方文书和外部材料的保存密度不一，不能由桥接标签推断行动、统属、疆界或社会影响。 进一步讨论可参照 宋辽夏边境、财政与文书材料的并行年序研究

\subsection{1096年（宋紹聖三年、西夏天祐民安7年；公元换算据两国年号对照）}
\noindent\eventanchor{SYD-1096-XIA-042}{八月辛酉，夏人寇寧順砦}\textbf{宋与西夏。}
宋与西夏 在 1096年（宋紹聖三年、西夏天祐民安7年；公元换算据两国年号对照） 的材料痕迹是〔八月辛酉，夏人寇寧順砦〕，出处为 《宋史》卷018〈紹聖三年〉；西夏文书、碑刻及《辽史》或《金史》同年条（互校）。材料未将地点细化到城址，不能用中枢所在地替它补足。

它承接的条件是：本纪从朝廷视角记录军政行动；出兵与战果需分辨。 随后可追踪的变化为：后续路线、控制与损失应与对方记录和遗址材料复核。

证据的边界是：《宋史》为元末官修，本纪保存宋廷纪年与处置；战果、数字、地方执行及对方动机须以编年、会要、对方史料或地方材料互校。 西夏文书、河西—宁夏考古与宋夏关系研究 可用于继续检验这一结论。

\noindent\eventanchor{SYD-1096-XIA-071}{壬戌，夏人寇鄜、延，陷金明砦}\textbf{宋与西夏。}
从 《宋史》卷018〈紹聖三年〉；西夏文书、碑刻及《辽史》或《金史》同年条（互校） 的 1096年（宋紹聖三年、西夏天祐民安7年；公元换算据两国年号对照） 条目看，〔壬戌，夏人寇鄜、延，陷金明砦〕使 宋与西夏 的一个具体行动进入叙事；材料未将地点细化到城址，不能用中枢所在地替它补足。

前一层压力可由以下线索辨认：本纪从朝廷视角记录军政行动；出兵与战果需分辨。 这一处置留下的后续条件是：后续路线、控制与损失应与对方记录和遗址材料复核。

《宋史》为元末官修，本纪保存宋廷纪年与处置；战果、数字、地方执行及对方动机须以编年、会要、对方史料或地方材料互校。 因此，西夏文书、河西—宁夏考古与宋夏关系研究 只能扩展问题，不能代替未保存的细节。

\subsection{1097年（宋紹聖四年、西夏天祐民安8年；公元换算据两国年号对照）}
\noindent\eventanchor{SYD-1097-XIA-043}{庚午，夏人大至葭盧城下，知石州張構等擊走之}\textbf{宋与西夏。}
1097年（宋紹聖四年、西夏天祐民安8年；公元换算据两国年号对照） 的记载将 宋与西夏 与〔庚午，夏人大至葭盧城下，知石州張構等擊走之〕相连。此处可确认的是这项被记录的行动；题名保留的城、路、水道或机构名称，是目前可以落地的空间线索。

本纪从朝廷视角记录军政行动；出兵与战果需分辨。 记录所见的结果则是：后续路线、控制与损失应与对方记录和遗址材料复核。 日期相邻不等于因果已经证实。

关于可说范围，须保留：《宋史》为元末官修，本纪保存宋廷纪年与处置；战果、数字、地方执行及对方动机须以编年、会要、对方史料或地方材料互校。 进一步讨论可参照 西夏文书、河西—宁夏考古与宋夏关系研究

\noindent\eventanchor{SYD-1097-XIA-072}{庚子，知保安軍李沂伐夏國，破洪州}\textbf{宋与西夏。}
〔庚子，知保安軍李沂伐夏國，破洪州〕见于 《宋史》卷018〈紹聖四年〉；西夏文书、碑刻及《辽史》或《金史》同年条（互校） 的 1097年（宋紹聖四年、西夏天祐民安8年；公元换算据两国年号对照） 记录。宋与西夏 的选择因此有了可回查的时间位置；题名保留的城、路、水道或机构名称，是目前可以落地的空间线索。

把本项放回事件链，能够看到的前提为：本纪从朝廷视角记录军政行动；出兵与战果需分辨。 而它所打开或收紧的下一步是：后续路线、控制与损失应与对方记录和遗址材料复核。

证据的边界是：《宋史》为元末官修，本纪保存宋廷纪年与处置；战果、数字、地方执行及对方动机须以编年、会要、对方史料或地方材料互校。 西夏文书、河西—宁夏考古与宋夏关系研究 可用于继续检验这一结论。

\noindent\eventanchor{SYD-1097-XIA-086}{癸亥，于闐來貢，黑汗王攻夏人三州，遣其子以聞}\textbf{宋与西夏。}
1097年（宋紹聖四年、西夏天祐民安8年；公元换算据两国年号对照），《宋史》卷018〈紹聖四年〉；西夏文书、碑刻及《辽史》或《金史》同年条（互校） 所见的〔癸亥，于闐來貢，黑汗王攻夏人三州，遣其子以聞〕把 宋与西夏 的处置留在可复核的年序中；题名保留的城、路、水道或机构名称，是目前可以落地的空间线索。

它承接的条件是：本纪从朝廷视角记录军政行动；出兵与战果需分辨。 随后可追踪的变化为：后续路线、控制与损失应与对方记录和遗址材料复核。

《宋史》为元末官修，本纪保存宋廷纪年与处置；战果、数字、地方执行及对方动机须以编年、会要、对方史料或地方材料互校。 因此，西夏文书、河西—宁夏考古与宋夏关系研究 只能扩展问题，不能代替未保存的细节。

\noindent\eventanchor{SYD-1097-XIA-094}{甲午，涇原路鈐轄王文振敗夏人於沒煙峽}\textbf{宋与西夏。}
宋与西夏 在 1097年（宋紹聖四年、西夏天祐民安8年；公元换算据两国年号对照） 的材料痕迹是〔甲午，涇原路鈐轄王文振敗夏人於沒煙峽〕，出处为 《宋史》卷018〈紹聖四年〉；西夏文书、碑刻及《辽史》或《金史》同年条（互校）。题名保留的城、路、水道或机构名称，是目前可以落地的空间线索。

前一层压力可由以下线索辨认：本纪从朝廷视角记录军政行动；出兵与战果需分辨。 这一处置留下的后续条件是：后续路线、控制与损失应与对方记录和遗址材料复核。

关于可说范围，须保留：《宋史》为元末官修，本纪保存宋廷纪年与处置；战果、数字、地方执行及对方动机须以编年、会要、对方史料或地方材料互校。 进一步讨论可参照 西夏文书、河西—宁夏考古与宋夏关系研究

\noindent\eventanchor{SYD-1097-XIA-099}{壬寅，環慶鈐轄張存入鹽州，俘戮甚衆，及還，夏人追襲之，復多亡失}\textbf{宋与西夏。}
从 《宋史》卷018〈紹聖四年〉；西夏文书、碑刻及《辽史》或《金史》同年条（互校） 的 1097年（宋紹聖四年、西夏天祐民安8年；公元换算据两国年号对照） 条目看，〔壬寅，環慶鈐轄張存入鹽州，俘戮甚衆，及還，夏人追襲之，復多亡失〕使 宋与西夏 的一个具体行动进入叙事；题名保留的城、路、水道或机构名称，是目前可以落地的空间线索。

本纪年条记录政令或机构处置；实施背景须参照制度材料。 记录所见的结果则是：可在同年及后续政令中追踪；条文不单独证明地方执行。 日期相邻不等于因果已经证实。

证据的边界是：《宋史》为元末官修，本纪保存宋廷纪年与处置；战果、数字、地方执行及对方动机须以编年、会要、对方史料或地方材料互校。 西夏文书、河西—宁夏考古与宋夏关系研究 可用于继续检验这一结论。

\noindent\eventanchor{SYD-1097-XIA-103}{三月壬戌，夏人犯麟州神堂堡，出兵討之，及進築胡山砦}\textbf{宋与西夏。}
1097年（宋紹聖四年、西夏天祐民安8年；公元换算据两国年号对照） 的记载将 宋与西夏 与〔三月壬戌，夏人犯麟州神堂堡，出兵討之，及進築胡山砦〕相连。此处可确认的是这项被记录的行动；题名保留的城、路、水道或机构名称，是目前可以落地的空间线索。

把本项放回事件链，能够看到的前提为：本纪年条提供日期和中枢可见行动；前因不据单条补定。 而它所打开或收紧的下一步是：可回查同年及后续条目，地方影响仍须另证。

《宋史》为元末官修，本纪保存宋廷纪年与处置；战果、数字、地方执行及对方动机须以编年、会要、对方史料或地方材料互校。 因此，西夏文书、河西—宁夏考古与宋夏关系研究 只能扩展问题，不能代替未保存的细节。

\noindent\eventanchor{SYD-1097-XIA-107}{辛巳，西上閣門使折克行破夏人于長波川，斬首二千餘級，獲牛馬倍之}\textbf{宋与西夏。}
〔辛巳，西上閣門使折克行破夏人于長波川，斬首二千餘級，獲牛馬倍之〕见于 《宋史》卷018〈紹聖四年〉；西夏文书、碑刻及《辽史》或《金史》同年条（互校） 的 1097年（宋紹聖四年、西夏天祐民安8年；公元换算据两国年号对照） 记录。宋与西夏 的选择因此有了可回查的时间位置；题名保留的城、路、水道或机构名称，是目前可以落地的空间线索。

它承接的条件是：本纪从朝廷视角记录军政行动；出兵与战果需分辨。 随后可追踪的变化为：后续路线、控制与损失应与对方记录和遗址材料复核。

关于可说范围，须保留：《宋史》为元末官修，本纪保存宋廷纪年与处置；战果、数字、地方执行及对方动机须以编年、会要、对方史料或地方材料互校。 进一步讨论可参照 西夏文书、河西—宁夏考古与宋夏关系研究

\subsection{1098 年（共享年序桥接）}
\noindent\eventanchor{SY-1098-CHRONOLOGY-BRIDGE}{【C级年序桥接】保留1098 年的多政权并行检索入口；本条不主张所列政权在当年发生直接互动，具体事件须另以单项材料入账。}\textbf{宋、辽与西夏（年序桥接）。}
1098 年（共享年序桥接），《宋史》；《辽史》；《续资治通鉴长编》；西夏文书与碑刻（据事核） 所见的〔【C级年序桥接】保留1098 年的多政权并行检索入口；本条不主张所列政权在当年发生直接互动，具体事件须另以单项材料入账。〕把 宋、辽与西夏（年序桥接） 的处置留在可复核的年序中；题名保留的城、路、水道或机构名称，是目前可以落地的空间线索。

前一层压力可由以下线索辨认：相邻已入账事件之间缺少该年度的共享年序入口，易使跨政权材料被单一政权叙事遮蔽。 这一处置留下的后续条件是：保留该年度的检索与补账位置；后续仅在获得可核单项材料时拆分为具体政治、财政、军事、社会或文化事件。

证据的边界是：本条只确认该年位于可追踪的共享年序中；正史、地方文书和外部材料的保存密度不一，不能由桥接标签推断行动、统属、疆界或社会影响。 宋辽夏边境、财政与文书材料的并行年序研究 可用于继续检验这一结论。

\subsection{1098年（宋元符元年、西夏永安1年；公元换算据两国年号对照）}
\noindent\eventanchor{SYD-1098-XIA-044}{涇原路禽夏國統軍嵬名阿埋等，高麗、瞎征、西南蕃張氏、羅氏、程氏入貢}\textbf{宋与西夏。}
宋与西夏 在 1098年（宋元符元年、西夏永安1年；公元换算据两国年号对照） 的材料痕迹是〔涇原路禽夏國統軍嵬名阿埋等，高麗、瞎征、西南蕃張氏、羅氏、程氏入貢〕，出处为 《宋史》卷018〈元符元年〉；西夏文书、碑刻及《辽史》或《金史》同年条（互校）。题名保留的城、路、水道或机构名称，是目前可以落地的空间线索。

本纪记录使节、贡赐或往来，不等于稳定统属关系。 记录所见的结果则是：可与对方编年和交通、文书材料核对其后续落实。 日期相邻不等于因果已经证实。

《宋史》为元末官修，本纪保存宋廷纪年与处置；战果、数字、地方执行及对方动机须以编年、会要、对方史料或地方材料互校。 因此，西夏文书、河西—宁夏考古与宋夏关系研究 只能扩展问题，不能代替未保存的细节。

\subsection{1099 年（共享年序桥接）}
\noindent\eventanchor{SY-1099-CHRONOLOGY-BRIDGE}{【C级年序桥接】保留1099 年的多政权并行检索入口；本条不主张所列政权在当年发生直接互动，具体事件须另以单项材料入账。}\textbf{宋、辽与西夏（年序桥接）。}
从 《宋史》；《辽史》；《续资治通鉴长编》；西夏文书与碑刻（据事核） 的 1099 年（共享年序桥接） 条目看，〔【C级年序桥接】保留1099 年的多政权并行检索入口；本条不主张所列政权在当年发生直接互动，具体事件须另以单项材料入账。〕使 宋、辽与西夏（年序桥接） 的一个具体行动进入叙事；题名保留的城、路、水道或机构名称，是目前可以落地的空间线索。

把本项放回事件链，能够看到的前提为：相邻已入账事件之间缺少该年度的共享年序入口，易使跨政权材料被单一政权叙事遮蔽。 而它所打开或收紧的下一步是：保留该年度的检索与补账位置；后续仅在获得可核单项材料时拆分为具体政治、财政、军事、社会或文化事件。

关于可说范围，须保留：本条只确认该年位于可追踪的共享年序中；正史、地方文书和外部材料的保存密度不一，不能由桥接标签推断行动、统属、疆界或社会影响。 进一步讨论可参照 宋辽夏边境、财政与文书材料的并行年序研究

\subsection{1099年（宋元符二年、辽寿昌5年；公元换算据两国年号对照）}
\noindent\eventanchor{SYD-1099-LIAO-075}{二年春正月甲辰朔，-{御}-大慶殿，以雪罷朝，群臣及遼使詣東上閣門拜表賀}\textbf{宋与辽。}
1099年（宋元符二年、辽寿昌5年；公元换算据两国年号对照） 的记载将 宋与辽 与〔二年春正月甲辰朔，-{御}-大慶殿，以雪罷朝，群臣及遼使詣東上閣門拜表賀〕相连。此处可确认的是这项被记录的行动；材料未将地点细化到城址，不能用中枢所在地替它补足。

它承接的条件是：本纪年条记录政令或机构处置；实施背景须参照制度材料。 随后可追踪的变化为：可在同年及后续政令中追踪；条文不单独证明地方执行。

证据的边界是：《宋史》为元末官修，本纪保存宋廷纪年与处置；战果、数字、地方执行及对方动机须以编年、会要、对方史料或地方材料互校。 宋辽边政、使节记录与契丹碑刻研究 可用于继续检验这一结论。

\noindent\eventanchor{SYD-1099-LIAO-138}{遣中書舍人郭知章報聘於遼}\textbf{宋与辽。}
〔遣中書舍人郭知章報聘於遼〕见于 《宋史》卷018〈元符二年〉；《辽史》本纪同年条（互校） 的 1099年（宋元符二年、辽寿昌5年；公元换算据两国年号对照） 记录。宋与辽 的选择因此有了可回查的时间位置；材料未将地点细化到城址，不能用中枢所在地替它补足。

前一层压力可由以下线索辨认：本纪年条提供日期和中枢可见行动；前因不据单条补定。 这一处置留下的后续条件是：可回查同年及后续条目，地方影响仍须另证。

《宋史》为元末官修，本纪保存宋廷纪年与处置；战果、数字、地方执行及对方动机须以编年、会要、对方史料或地方材料互校。 因此，宋辽边政、使节记录与契丹碑刻研究 只能扩展问题，不能代替未保存的细节。

\subsection{1099年（宋元符二年、西夏永安2年；公元换算据两国年号对照）}
\noindent\eventanchor{SYD-1099-XIA-045}{甲申，夏人以國母卒，遣使告哀，且謝罪，卻其使不納}\textbf{宋与西夏。}
1099年（宋元符二年、西夏永安2年；公元换算据两国年号对照），《宋史》卷018〈元符二年〉；西夏文书、碑刻及《辽史》或《金史》同年条（互校） 所见的〔甲申，夏人以國母卒，遣使告哀，且謝罪，卻其使不納〕把 宋与西夏 的处置留在可复核的年序中；材料未将地点细化到城址，不能用中枢所在地替它补足。

本纪记录使节、贡赐或往来，不等于稳定统属关系。 记录所见的结果则是：可与对方编年和交通、文书材料核对其后续落实。 日期相邻不等于因果已经证实。

关于可说范围，须保留：《宋史》为元末官修，本纪保存宋廷纪年与处置；战果、数字、地方执行及对方动机须以编年、会要、对方史料或地方材料互校。 进一步讨论可参照 西夏文书、河西—宁夏考古与宋夏关系研究

\noindent\eventanchor{SYD-1099-XIA-073}{九月庚子朔，夏人來謝罪}\textbf{宋与西夏。}
宋与西夏 在 1099年（宋元符二年、西夏永安2年；公元换算据两国年号对照） 的材料痕迹是〔九月庚子朔，夏人來謝罪〕，出处为 《宋史》卷018〈元符二年〉；西夏文书、碑刻及《辽史》或《金史》同年条（互校）。材料未将地点细化到城址，不能用中枢所在地替它补足。

把本项放回事件链，能够看到的前提为：本纪年条提供日期和中枢可见行动；前因不据单条补定。 而它所打开或收紧的下一步是：可回查同年及后续条目，地方影响仍须另证。

证据的边界是：《宋史》为元末官修，本纪保存宋廷纪年与处置；战果、数字、地方执行及对方动机须以编年、会要、对方史料或地方材料互校。 西夏文书、河西—宁夏考古与宋夏关系研究 可用于继续检验这一结论。

\noindent\eventanchor{SYD-1099-XIA-087}{戊子，鄜延鈐轄劉安敗夏人於神堆}\textbf{宋与西夏。}
从 《宋史》卷018〈元符二年〉；西夏文书、碑刻及《辽史》或《金史》同年条（互校） 的 1099年（宋元符二年、西夏永安2年；公元换算据两国年号对照） 条目看，〔戊子，鄜延鈐轄劉安敗夏人於神堆〕使 宋与西夏 的一个具体行动进入叙事；材料未将地点细化到城址，不能用中枢所在地替它补足。

它承接的条件是：本纪从朝廷视角记录军政行动；出兵与战果需分辨。 随后可追踪的变化为：后续路线、控制与损失应与对方记录和遗址材料复核。

《宋史》为元末官修，本纪保存宋廷纪年与处置；战果、数字、地方执行及对方动机须以编年、会要、对方史料或地方材料互校。 因此，西夏文书、河西—宁夏考古与宋夏关系研究 只能扩展问题，不能代替未保存的细节。

\subsection{1100年（宋元符三年、辽寿昌6年；公元换算据两国年号对照）}
\noindent\eventanchor{SYD-1100-LIAO-076}{丙午，遣董敦逸賀遼主生辰，呂仲甫賀正旦}\textbf{宋与辽。}
1100年（宋元符三年、辽寿昌6年；公元换算据两国年号对照） 的记载将 宋与辽 与〔丙午，遣董敦逸賀遼主生辰，呂仲甫賀正旦〕相连。此处可确认的是这项被记录的行动；材料未将地点细化到城址，不能用中枢所在地替它补足。

前一层压力可由以下线索辨认：本纪年条提供日期和中枢可见行动；前因不据单条补定。 这一处置留下的后续条件是：可回查同年及后续条目，地方影响仍须另证。

关于可说范围，须保留：《宋史》为元末官修，本纪保存宋廷纪年与处置；战果、数字、地方执行及对方动机须以编年、会要、对方史料或地方材料互校。 进一步讨论可参照 宋辽边政、使节记录与契丹碑刻研究

\noindent\eventanchor{SYD-1100-LIAO-139}{丙寅，遼遣蕭穆來賀即位}\textbf{宋与辽。}
〔丙寅，遼遣蕭穆來賀即位〕见于 《宋史》卷019〈元符三年〉；《辽史》本纪同年条（互校） 的 1100年（宋元符三年、辽寿昌6年；公元换算据两国年号对照） 记录。宋与辽 的选择因此有了可回查的时间位置；材料未将地点细化到城址，不能用中枢所在地替它补足。

本纪年条提供日期和中枢可见行动；前因不据单条补定。 记录所见的结果则是：可回查同年及后续条目，地方影响仍须另证。 日期相邻不等于因果已经证实。

证据的边界是：《宋史》为元末官修，本纪保存宋廷纪年与处置；战果、数字、地方执行及对方动机须以编年、会要、对方史料或地方材料互校。 宋辽边政、使节记录与契丹碑刻研究 可用于继续检验这一结论。

\subsection{1101年}
\noindent\eventanchor{LIAO-1101-EVENT-042}{辽道宗去世，天祚帝即位}\textbf{辽。}
1101年，《辽史》相关本纪及志；《宋史》辽国传、同年本纪；宋使行录、契丹／汉文碑刻与考古材料（据事核） 所见的〔辽道宗去世，天祚帝即位〕把 辽 的处置留在可复核的年序中；材料未将地点细化到城址，不能用中枢所在地替它补足。

把本项放回事件链，能够看到的前提为：辽中期继承及辽宋夏关系中前任统治的结束与宫廷、部族或官僚的权力安排相互交织，促成“辽道宗去世，天祚帝即位”这一继承节点。 而它所打开或收紧的下一步是：继承并不自动消除既有矛盾；“辽道宗去世，天祚帝即位”后的任官、军权和对外策略需由后续条目追踪。

《辽史》为元末官修且契丹本方连续文书有限；宋使与碑刻的观察位置不同，战果、族称及制度覆盖范围不得互推。 因此，辽代契丹国家、都城、部族与跨境关系研究 只能扩展问题，不能代替未保存的细节。

\subsection{1101年（宋建中靖國元年、辽乾统1年；公元换算据两国年号对照）}
\noindent\eventanchor{SYD-1101-LIAO-077}{是歲，遼人來獻遺留物}\textbf{宋与辽。}
宋与辽 在 1101年（宋建中靖國元年、辽乾统1年；公元换算据两国年号对照） 的材料痕迹是〔是歲，遼人來獻遺留物〕，出处为 《宋史》卷019〈建中靖國元年〉；《辽史》本纪同年条（互校）。材料未将地点细化到城址，不能用中枢所在地替它补足。

它承接的条件是：本纪年条提供日期和中枢可见行动；前因不据单条补定。 随后可追踪的变化为：可回查同年及后续条目，地方影响仍须另证。

关于可说范围，须保留：《宋史》为元末官修，本纪保存宋廷纪年与处置；战果、数字、地方执行及对方动机须以编年、会要、对方史料或地方材料互校。 进一步讨论可参照 宋辽边政、使节记录与契丹碑刻研究

\noindent\eventanchor{SYD-1101-LIAO-140}{乙丑，遼使蕭恭來告其主洪基殂，遣謝瓘、上官均等往弔祭，黃寔賀其孫延禧立}\textbf{宋与辽。}
从 《宋史》卷019〈建中靖國元年〉；《辽史》本纪同年条（互校） 的 1101年（宋建中靖國元年、辽乾统1年；公元换算据两国年号对照） 条目看，〔乙丑，遼使蕭恭來告其主洪基殂，遣謝瓘、上官均等往弔祭，黃寔賀其孫延禧立〕使 宋与辽 的一个具体行动进入叙事；材料未将地点细化到城址，不能用中枢所在地替它补足。

前一层压力可由以下线索辨认：本纪年条提供日期和中枢可见行动；前因不据单条补定。 这一处置留下的后续条件是：可回查同年及后续条目，地方影响仍须另证。

证据的边界是：《宋史》为元末官修，本纪保存宋廷纪年与处置；战果、数字、地方执行及对方动机须以编年、会要、对方史料或地方材料互校。 宋辽边政、使节记录与契丹碑刻研究 可用于继续检验这一结论。

\subsection{1104年（宋崇寧三年、西夏贞观4年；公元换算据两国年号对照）}
\noindent\eventanchor{SYD-1104-XIA-046}{戊午，夏人入涇原，圍平夏城，寇鎮戎軍}\textbf{宋与西夏。}
1104年（宋崇寧三年、西夏贞观4年；公元换算据两国年号对照） 的记载将 宋与西夏 与〔戊午，夏人入涇原，圍平夏城，寇鎮戎軍〕相连。此处可确认的是这项被记录的行动；题名保留的城、路、水道或机构名称，是目前可以落地的空间线索。

本纪从朝廷视角记录军政行动；出兵与战果需分辨。 记录所见的结果则是：后续路线、控制与损失应与对方记录和遗址材料复核。 日期相邻不等于因果已经证实。

《宋史》为元末官修，本纪保存宋廷纪年与处置；战果、数字、地方执行及对方动机须以编年、会要、对方史料或地方材料互校。 因此，西夏文书、河西—宁夏考古与宋夏关系研究 只能扩展问题，不能代替未保存的细节。

\subsection{1105 年（共享年序桥接）}
\noindent\eventanchor{SY-1105-CHRONOLOGY-BRIDGE}{【C级年序桥接】保留1105 年的多政权并行检索入口；本条不主张所列政权在当年发生直接互动，具体事件须另以单项材料入账。}\textbf{宋、辽与西夏（年序桥接）。}
〔【C级年序桥接】保留1105 年的多政权并行检索入口；本条不主张所列政权在当年发生直接互动，具体事件须另以单项材料入账。〕见于 《宋史》；《辽史》；《续资治通鉴长编》；西夏文书与碑刻（据事核） 的 1105 年（共享年序桥接） 记录。宋、辽与西夏（年序桥接） 的选择因此有了可回查的时间位置；题名保留的城、路、水道或机构名称，是目前可以落地的空间线索。

把本项放回事件链，能够看到的前提为：相邻已入账事件之间缺少该年度的共享年序入口，易使跨政权材料被单一政权叙事遮蔽。 而它所打开或收紧的下一步是：保留该年度的检索与补账位置；后续仅在获得可核单项材料时拆分为具体政治、财政、军事、社会或文化事件。

关于可说范围，须保留：本条只确认该年位于可追踪的共享年序中；正史、地方文书和外部材料的保存密度不一，不能由桥接标签推断行动、统属、疆界或社会影响。 进一步讨论可参照 宋辽夏边境、财政与文书材料的并行年序研究

\subsection{1105年（宋崇寧四年、辽乾统5年；公元换算据两国年号对照）}
\noindent\eventanchor{SYD-1105-LIAO-078}{壬子，遣林攄報聘於遼}\textbf{宋与辽。}
1105年（宋崇寧四年、辽乾统5年；公元换算据两国年号对照），《宋史》卷020〈崇寧四年〉；《辽史》本纪同年条（互校） 所见的〔壬子，遣林攄報聘於遼〕把 宋与辽 的处置留在可复核的年序中；材料未将地点细化到城址，不能用中枢所在地替它补足。

它承接的条件是：本纪年条提供日期和中枢可见行动；前因不据单条补定。 随后可追踪的变化为：可回查同年及后续条目，地方影响仍须另证。

证据的边界是：《宋史》为元末官修，本纪保存宋廷纪年与处置；战果、数字、地方执行及对方动机须以编年、会要、对方史料或地方材料互校。 宋辽边政、使节记录与契丹碑刻研究 可用于继续检验这一结论。

\subsection{1105年（宋崇寧四年、西夏贞观5年；公元换算据两国年号对照）}
\noindent\eventanchor{SYD-1105-XIA-047}{己丑，夏人寇順寧砦，鄜延第二副將劉延慶擊破之}\textbf{宋与西夏。}
宋与西夏 在 1105年（宋崇寧四年、西夏贞观5年；公元换算据两国年号对照） 的材料痕迹是〔己丑，夏人寇順寧砦，鄜延第二副將劉延慶擊破之〕，出处为 《宋史》卷020〈崇寧四年〉；西夏文书、碑刻及《辽史》或《金史》同年条（互校）。材料未将地点细化到城址，不能用中枢所在地替它补足。

前一层压力可由以下线索辨认：本纪从朝廷视角记录军政行动；出兵与战果需分辨。 这一处置留下的后续条件是：后续路线、控制与损失应与对方记录和遗址材料复核。

《宋史》为元末官修，本纪保存宋廷纪年与处置；战果、数字、地方执行及对方动机须以编年、会要、对方史料或地方材料互校。 因此，西夏文书、河西—宁夏考古与宋夏关系研究 只能扩展问题，不能代替未保存的细节。

\subsection{1108 年（共享年序桥接）}
\noindent\eventanchor{SY-1108-CHRONOLOGY-BRIDGE}{【C级年序桥接】保留1108 年的多政权并行检索入口；本条不主张所列政权在当年发生直接互动，具体事件须另以单项材料入账。}\textbf{宋、辽与西夏（年序桥接）。}
从 《宋史》；《辽史》；《续资治通鉴长编》；西夏文书与碑刻（据事核） 的 1108 年（共享年序桥接） 条目看，〔【C级年序桥接】保留1108 年的多政权并行检索入口；本条不主张所列政权在当年发生直接互动，具体事件须另以单项材料入账。〕使 宋、辽与西夏（年序桥接） 的一个具体行动进入叙事；题名保留的城、路、水道或机构名称，是目前可以落地的空间线索。

相邻已入账事件之间缺少该年度的共享年序入口，易使跨政权材料被单一政权叙事遮蔽。 记录所见的结果则是：保留该年度的检索与补账位置；后续仅在获得可核单项材料时拆分为具体政治、财政、军事、社会或文化事件。 日期相邻不等于因果已经证实。

关于可说范围，须保留：本条只确认该年位于可追踪的共享年序中；正史、地方文书和外部材料的保存密度不一，不能由桥接标签推断行动、统属、疆界或社会影响。 进一步讨论可参照 宋辽夏边境、财政与文书材料的并行年序研究

\subsection{1109 年（共享年序桥接）}
\noindent\eventanchor{SY-1109-CHRONOLOGY-BRIDGE}{【C级年序桥接】保留1109 年的多政权并行检索入口；本条不主张所列政权在当年发生直接互动，具体事件须另以单项材料入账。}\textbf{宋、辽与西夏（年序桥接）。}
1109 年（共享年序桥接） 的记载将 宋、辽与西夏（年序桥接） 与〔【C级年序桥接】保留1109 年的多政权并行检索入口；本条不主张所列政权在当年发生直接互动，具体事件须另以单项材料入账。〕相连。此处可确认的是这项被记录的行动；题名保留的城、路、水道或机构名称，是目前可以落地的空间线索。

把本项放回事件链，能够看到的前提为：相邻已入账事件之间缺少该年度的共享年序入口，易使跨政权材料被单一政权叙事遮蔽。 而它所打开或收紧的下一步是：保留该年度的检索与补账位置；后续仅在获得可核单项材料时拆分为具体政治、财政、军事、社会或文化事件。

证据的边界是：本条只确认该年位于可追踪的共享年序中；正史、地方文书和外部材料的保存密度不一，不能由桥接标签推断行动、统属、疆界或社会影响。 宋辽夏边境、财政与文书材料的并行年序研究 可用于继续检验这一结论。

\subsection{1109年（宋大觀三年、西夏贞观9年；公元换算据两国年号对照）}
\noindent\eventanchor{SYD-1109-XIA-048}{闍婆、占城、夏國入貢}\textbf{宋与西夏。}
〔闍婆、占城、夏國入貢〕见于 《宋史》卷020〈大觀三年〉；西夏文书、碑刻及《辽史》或《金史》同年条（互校） 的 1109年（宋大觀三年、西夏贞观9年；公元换算据两国年号对照） 记录。宋与西夏 的选择因此有了可回查的时间位置；题名保留的城、路、水道或机构名称，是目前可以落地的空间线索。

它承接的条件是：本纪记录使节、贡赐或往来，不等于稳定统属关系。 随后可追踪的变化为：可与对方编年和交通、文书材料核对其后续落实。

《宋史》为元末官修，本纪保存宋廷纪年与处置；战果、数字、地方执行及对方动机须以编年、会要、对方史料或地方材料互校。 因此，西夏文书、河西—宁夏考古与宋夏关系研究 只能扩展问题，不能代替未保存的细节。

\subsection{1110 年（共享年序桥接）}
\noindent\eventanchor{SY-1110-CHRONOLOGY-BRIDGE}{【C级年序桥接】保留1110 年的多政权并行检索入口；本条不主张所列政权在当年发生直接互动，具体事件须另以单项材料入账。}\textbf{宋、辽与西夏（年序桥接）。}
1110 年（共享年序桥接），《宋史》；《辽史》；《续资治通鉴长编》；西夏文书与碑刻（据事核） 所见的〔【C级年序桥接】保留1110 年的多政权并行检索入口；本条不主张所列政权在当年发生直接互动，具体事件须另以单项材料入账。〕把 宋、辽与西夏（年序桥接） 的处置留在可复核的年序中；题名保留的城、路、水道或机构名称，是目前可以落地的空间线索。

前一层压力可由以下线索辨认：相邻已入账事件之间缺少该年度的共享年序入口，易使跨政权材料被单一政权叙事遮蔽。 这一处置留下的后续条件是：保留该年度的检索与补账位置；后续仅在获得可核单项材料时拆分为具体政治、财政、军事、社会或文化事件。

关于可说范围，须保留：本条只确认该年位于可追踪的共享年序中；正史、地方文书和外部材料的保存密度不一，不能由桥接标签推断行动、统属、疆界或社会影响。 进一步讨论可参照 宋辽夏边境、财政与文书材料的并行年序研究

\subsection{1111 年（共享年序桥接）}
\noindent\eventanchor{SY-1111-CHRONOLOGY-BRIDGE}{【C级年序桥接】保留1111 年的多政权并行检索入口；本条不主张所列政权在当年发生直接互动，具体事件须另以单项材料入账。}\textbf{宋、辽与西夏（年序桥接）。}
宋、辽与西夏（年序桥接） 在 1111 年（共享年序桥接） 的材料痕迹是〔【C级年序桥接】保留1111 年的多政权并行检索入口；本条不主张所列政权在当年发生直接互动，具体事件须另以单项材料入账。〕，出处为 《宋史》；《辽史》；《续资治通鉴长编》；西夏文书与碑刻（据事核）。题名保留的城、路、水道或机构名称，是目前可以落地的空间线索。

相邻已入账事件之间缺少该年度的共享年序入口，易使跨政权材料被单一政权叙事遮蔽。 记录所见的结果则是：保留该年度的检索与补账位置；后续仅在获得可核单项材料时拆分为具体政治、财政、军事、社会或文化事件。 日期相邻不等于因果已经证实。

证据的边界是：本条只确认该年位于可追踪的共享年序中；正史、地方文书和外部材料的保存密度不一，不能由桥接标签推断行动、统属、疆界或社会影响。 宋辽夏边境、财政与文书材料的并行年序研究 可用于继续检验这一结论。

\subsection{1111年（宋政和元年、辽天庆1年；公元换算据两国年号对照）}
\noindent\eventanchor{SYD-1111-LIAO-079}{是月，鄭允中、童貫使遼，以李良嗣來，良嗣獻取燕之策，詔賜姓趙}\textbf{宋与辽。}
从 《宋史》卷020〈政和元年〉；《辽史》本纪同年条（互校） 的 1111年（宋政和元年、辽天庆1年；公元换算据两国年号对照） 条目看，〔是月，鄭允中、童貫使遼，以李良嗣來，良嗣獻取燕之策，詔賜姓趙〕使 宋与辽 的一个具体行动进入叙事；题名保留的城、路、水道或机构名称，是目前可以落地的空间线索。

把本项放回事件链，能够看到的前提为：本纪年条记录政令或机构处置；实施背景须参照制度材料。 而它所打开或收紧的下一步是：可在同年及后续政令中追踪；条文不单独证明地方执行。

《宋史》为元末官修，本纪保存宋廷纪年与处置；战果、数字、地方执行及对方动机须以编年、会要、对方史料或地方材料互校。 因此，宋辽边政、使节记录与契丹碑刻研究 只能扩展问题，不能代替未保存的细节。

\subsection{1112 年（共享年序桥接）}
\noindent\eventanchor{SY-1112-CHRONOLOGY-BRIDGE}{【C级年序桥接】保留1112 年的多政权并行检索入口；本条不主张所列政权在当年发生直接互动，具体事件须另以单项材料入账。}\textbf{宋、辽与西夏（年序桥接）。}
1112 年（共享年序桥接） 的记载将 宋、辽与西夏（年序桥接） 与〔【C级年序桥接】保留1112 年的多政权并行检索入口；本条不主张所列政权在当年发生直接互动，具体事件须另以单项材料入账。〕相连。此处可确认的是这项被记录的行动；题名保留的城、路、水道或机构名称，是目前可以落地的空间线索。

它承接的条件是：相邻已入账事件之间缺少该年度的共享年序入口，易使跨政权材料被单一政权叙事遮蔽。 随后可追踪的变化为：保留该年度的检索与补账位置；后续仅在获得可核单项材料时拆分为具体政治、财政、军事、社会或文化事件。

关于可说范围，须保留：本条只确认该年位于可追踪的共享年序中；正史、地方文书和外部材料的保存密度不一，不能由桥接标签推断行动、统属、疆界或社会影响。 进一步讨论可参照 宋辽夏边境、财政与文书材料的并行年序研究

\subsection{1113 年（共享年序桥接）}
\noindent\eventanchor{SY-1113-CHRONOLOGY-BRIDGE}{【C级年序桥接】保留1113 年的多政权并行检索入口；本条不主张所列政权在当年发生直接互动，具体事件须另以单项材料入账。}\textbf{宋、辽与西夏（年序桥接）。}
〔【C级年序桥接】保留1113 年的多政权并行检索入口；本条不主张所列政权在当年发生直接互动，具体事件须另以单项材料入账。〕见于 《宋史》；《辽史》；《续资治通鉴长编》；西夏文书与碑刻（据事核） 的 1113 年（共享年序桥接） 记录。宋、辽与西夏（年序桥接） 的选择因此有了可回查的时间位置；题名保留的城、路、水道或机构名称，是目前可以落地的空间线索。

前一层压力可由以下线索辨认：相邻已入账事件之间缺少该年度的共享年序入口，易使跨政权材料被单一政权叙事遮蔽。 这一处置留下的后续条件是：保留该年度的检索与补账位置；后续仅在获得可核单项材料时拆分为具体政治、财政、军事、社会或文化事件。

证据的边界是：本条只确认该年位于可追踪的共享年序中；正史、地方文书和外部材料的保存密度不一，不能由桥接标签推断行动、统属、疆界或社会影响。 宋辽夏边境、财政与文书材料的并行年序研究 可用于继续检验这一结论。

\subsection{1113年（宋政和三年、辽天庆3年；公元换算据两国年号对照）}
\noindent\eventanchor{SYD-1113-LIAO-080}{甲午，以遼、女真相持，詔河北治邊防}\textbf{宋与辽。}
1113年（宋政和三年、辽天庆3年；公元换算据两国年号对照），《宋史》卷021〈政和三年〉；《辽史》本纪同年条（互校） 所见的〔甲午，以遼、女真相持，詔河北治邊防〕把 宋与辽 的处置留在可复核的年序中；题名保留的城、路、水道或机构名称，是目前可以落地的空间线索。

本纪年条记录政令或机构处置；实施背景须参照制度材料。 记录所见的结果则是：可在同年及后续政令中追踪；条文不单独证明地方执行。 日期相邻不等于因果已经证实。

《宋史》为元末官修，本纪保存宋廷纪年与处置；战果、数字、地方执行及对方动机须以编年、会要、对方史料或地方材料互校。 因此，宋辽边政、使节记录与契丹碑刻研究 只能扩展问题，不能代替未保存的细节。

\subsection{1114年}
\noindent\eventanchor{LIAO-1114-EVENT-043}{完颜阿骨打起兵反辽，女真军迅速取得初步胜利}\textbf{辽与女真。}
辽与女真 在 1114年 的材料痕迹是〔完颜阿骨打起兵反辽，女真军迅速取得初步胜利〕，出处为 《辽史》相关本纪及志；《宋史》辽国传、同年本纪；宋使行录、契丹／汉文碑刻与考古材料（据事核）。材料未将地点细化到城址，不能用中枢所在地替它补足。

把本项放回事件链，能够看到的前提为：“完颜阿骨打起兵反辽，女真军迅速取得初步胜利”发生在女真兴起与辽末多方角逐的具体压力与机会之中，相关行动者的选择不能脱离此前事件链解释。 而它所打开或收紧的下一步是：“完颜阿骨打起兵反辽，女真军迅速取得初步胜利”为后续政治、边防或资源配置留下条件；其社会影响与实际覆盖范围仍须以异类材料限定。

关于可说范围，须保留：《辽史》为元末官修且契丹本方连续文书有限；宋使与碑刻的观察位置不同，战果、族称及制度覆盖范围不得互推。 进一步讨论可参照 辽代契丹国家、都城、部族与跨境关系研究

\subsection{1115年}
\noindent\eventanchor{LIAO-1115-EVENT-044}{金建国，辽失去对女真关系的主导权}\textbf{辽与金。}
从 《辽史》相关本纪及志；《宋史》辽国传、同年本纪；宋使行录、契丹／汉文碑刻与考古材料（据事核） 的 1115年 条目看，〔金建国，辽失去对女真关系的主导权〕使 辽与金 的一个具体行动进入叙事；题名保留的城、路、水道或机构名称，是目前可以落地的空间线索。

它承接的条件是：“金建国，辽失去对女真关系的主导权”发生在女真兴起与辽末多方角逐的具体压力与机会之中，相关行动者的选择不能脱离此前事件链解释。 随后可追踪的变化为：“金建国，辽失去对女真关系的主导权”为后续政治、边防或资源配置留下条件；其社会影响与实际覆盖范围仍须以异类材料限定。

证据的边界是：《辽史》为元末官修且契丹本方连续文书有限；宋使与碑刻的观察位置不同，战果、族称及制度覆盖范围不得互推。 辽代契丹国家、都城、部族与跨境关系研究 可用于继续检验这一结论。

\subsection{1115年（金收国元年、宋政和五年；西夏年号、干支与月日待核；公元换算据各方本纪年号对照）}
\noindent\eventanchor{XIA-1115-EVENT-067}{金建国改变西夏北方外交环境}\textbf{西夏辽金。}
1115年（金收国元年、宋政和五年；西夏年号、干支与月日待核；公元换算据各方本纪年号对照） 的记载将 西夏辽金 与〔金建国改变西夏北方外交环境〕相连。此处可确认的是这项被记录的行动；材料未将地点细化到城址，不能用中枢所在地替它补足。

前一层压力可由以下线索辨认：“金建国改变西夏北方外交环境”发生在西夏与金、宋的边境与册封关系的具体压力与机会之中，相关行动者的选择不能脱离此前事件链解释。 这一处置留下的后续条件是：“金建国改变西夏北方外交环境”为后续政治、边防或资源配置留下条件；其社会影响与实际覆盖范围仍须以异类材料限定。

西夏无存世连续国史，汉文叙事多出自对手；西夏文文书的释读、定年和地域代表性须逐项说明。 因此，西夏文文书、河西—宁夏考古与宋夏金蒙关系研究 只能扩展问题，不能代替未保存的细节。

\subsection{1116 年（共享年序桥接）}
\noindent\eventanchor{SY-1116-CHRONOLOGY-BRIDGE}{【C级年序桥接】保留1116 年的多政权并行检索入口；本条不主张所列政权在当年发生直接互动，具体事件须另以单项材料入账。}\textbf{宋、辽与金（年序桥接）。}
〔【C级年序桥接】保留1116 年的多政权并行检索入口；本条不主张所列政权在当年发生直接互动，具体事件须另以单项材料入账。〕见于 《宋史》；《辽史》；《金史》；《续资治通鉴长编》相关年条（据事核） 的 1116 年（共享年序桥接） 记录。宋、辽与金（年序桥接） 的选择因此有了可回查的时间位置；题名保留的城、路、水道或机构名称，是目前可以落地的空间线索。

相邻已入账事件之间缺少该年度的共享年序入口，易使跨政权材料被单一政权叙事遮蔽。 记录所见的结果则是：保留该年度的检索与补账位置；后续仅在获得可核单项材料时拆分为具体政治、财政、军事、社会或文化事件。 日期相邻不等于因果已经证实。

关于可说范围，须保留：本条只确认该年位于可追踪的共享年序中；正史、地方文书和外部材料的保存密度不一，不能由桥接标签推断行动、统属、疆界或社会影响。 进一步讨论可参照 辽金转换与北宋北边情报、城防、人口流动研究

\subsection{1116年（宋政和六年、西夏雍宁3年；公元换算据两国年号对照）}
\noindent\eventanchor{SYD-1116-XIA-049}{高麗、占城、大食、真臘、大理、夏國入貢，茂州夷郅永壽內附}\textbf{宋与西夏。}
1116年（宋政和六年、西夏雍宁3年；公元换算据两国年号对照），《宋史》卷021〈政和六年〉；西夏文书、碑刻及《辽史》或《金史》同年条（互校） 所见的〔高麗、占城、大食、真臘、大理、夏國入貢，茂州夷郅永壽內附〕把 宋与西夏 的处置留在可复核的年序中；题名保留的城、路、水道或机构名称，是目前可以落地的空间线索。

把本项放回事件链，能够看到的前提为：本纪记录使节、贡赐或往来，不等于稳定统属关系。 而它所打开或收紧的下一步是：可与对方编年和交通、文书材料核对其后续落实。

证据的边界是：《宋史》为元末官修，本纪保存宋廷纪年与处置；战果、数字、地方执行及对方动机须以编年、会要、对方史料或地方材料互校。 西夏文书、河西—宁夏考古与宋夏关系研究 可用于继续检验这一结论。

\subsection{1118年（宋重和元年、辽天庆8年；公元换算据两国年号对照）}
\noindent\eventanchor{SYD-1118-LIAO-081}{庚午，遣武義大夫馬政由海道使女真，約夾攻遼}\textbf{宋与辽。}
宋与辽 在 1118年（宋重和元年、辽天庆8年；公元换算据两国年号对照） 的材料痕迹是〔庚午，遣武義大夫馬政由海道使女真，約夾攻遼〕，出处为 《宋史》卷021〈重和元年〉；《辽史》本纪同年条（互校）。题名保留的城、路、水道或机构名称，是目前可以落地的空间线索。

它承接的条件是：本纪从朝廷视角记录军政行动；出兵与战果需分辨。 随后可追踪的变化为：后续路线、控制与损失应与对方记录和遗址材料复核。

《宋史》为元末官修，本纪保存宋廷纪年与处置；战果、数字、地方执行及对方动机须以编年、会要、对方史料或地方材料互校。 因此，宋辽边政、使节记录与契丹碑刻研究 只能扩展问题，不能代替未保存的细节。

\subsection{1119年（宋宣和元年、西夏元德1年；公元换算据两国年号对照）}
\noindent\eventanchor{SYD-1119-XIA-050}{庚寅，童貫以鄜延、環慶兵大破夏人，平其三城}\textbf{宋与西夏。}
从 《宋史》卷022〈宣和元年〉；西夏文书、碑刻及《辽史》或《金史》同年条（互校） 的 1119年（宋宣和元年、西夏元德1年；公元换算据两国年号对照） 条目看，〔庚寅，童貫以鄜延、環慶兵大破夏人，平其三城〕使 宋与西夏 的一个具体行动进入叙事；题名保留的城、路、水道或机构名称，是目前可以落地的空间线索。

前一层压力可由以下线索辨认：本纪从朝廷视角记录军政行动；出兵与战果需分辨。 这一处置留下的后续条件是：后续路线、控制与损失应与对方记录和遗址材料复核。

关于可说范围，须保留：《宋史》为元末官修，本纪保存宋廷纪年与处置；战果、数字、地方执行及对方动机须以编年、会要、对方史料或地方材料互校。 进一步讨论可参照 西夏文书、河西—宁夏考古与宋夏关系研究

\noindent\eventanchor{SYD-1119-XIA-074}{己亥，夏國遣使納款，詔六路罷兵}\textbf{宋与西夏。}
1119年（宋宣和元年、西夏元德1年；公元换算据两国年号对照） 的记载将 宋与西夏 与〔己亥，夏國遣使納款，詔六路罷兵〕相连。此处可确认的是这项被记录的行动；题名保留的城、路、水道或机构名称，是目前可以落地的空间线索。

本纪年条记录政令或机构处置；实施背景须参照制度材料。 记录所见的结果则是：可在同年及后续政令中追踪；条文不单独证明地方执行。 日期相邻不等于因果已经证实。

证据的边界是：《宋史》为元末官修，本纪保存宋廷纪年与处置；战果、数字、地方执行及对方动机须以编年、会要、对方史料或地方材料互校。 西夏文书、河西—宁夏考古与宋夏关系研究 可用于继续检验这一结论。

\noindent\eventanchor{SYD-1119-XIA-088}{童貫遣知熙州劉法出師攻統安城，夏人伏兵擊之，法敗歿，震武軍受圍}\textbf{宋与西夏。}
〔童貫遣知熙州劉法出師攻統安城，夏人伏兵擊之，法敗歿，震武軍受圍〕见于 《宋史》卷022〈宣和元年〉；西夏文书、碑刻及《辽史》或《金史》同年条（互校） 的 1119年（宋宣和元年、西夏元德1年；公元换算据两国年号对照） 记录。宋与西夏 的选择因此有了可回查的时间位置；题名保留的城、路、水道或机构名称，是目前可以落地的空间线索。

把本项放回事件链，能够看到的前提为：本纪保留灾害或工程的报告地点与朝廷处置；灾情范围需另证。 而它所打开或收紧的下一步是：可与地方文书、河工和环境材料比较，不将单点记录外推为全域因果。

《宋史》为元末官修，本纪保存宋廷纪年与处置；战果、数字、地方执行及对方动机须以编年、会要、对方史料或地方材料互校。 因此，西夏文书、河西—宁夏考古与宋夏关系研究 只能扩展问题，不能代替未保存的细节。

\subsection{1120年（辽天庆十年、金天辅四年、宋宣和二年；干支、月日待核；公元换算据各方本纪年号对照）}
\noindent\eventanchor{LIAO-1120-EVENT-045}{宋金结盟共同攻辽，使辽面临南北夹击}\textbf{辽宋金。}
1120年（辽天庆十年、金天辅四年、宋宣和二年；干支、月日待核；公元换算据各方本纪年号对照），《辽史》相关本纪及志；《宋史》辽国传、同年本纪；宋使行录、契丹／汉文碑刻与考古材料（据事核） 所见的〔宋金结盟共同攻辽，使辽面临南北夹击〕把 辽宋金 的处置留在可复核的年序中；材料未将地点细化到城址，不能用中枢所在地替它补足。

它承接的条件是：女真兴起与辽末多方角逐里的军事消耗、边境供给与使节往来构成谈判条件，促成“宋金结盟共同攻辽，使辽面临南北夹击”所涉的协议或名义关系。 随后可追踪的变化为：“宋金结盟共同攻辽，使辽面临南北夹击”重排了贡赐、使节或停战的制度接口；条款是否执行及边地实际控制不能由盟约文字直接推定。

关于可说范围，须保留：《辽史》为元末官修且契丹本方连续文书有限；宋使与碑刻的观察位置不同，战果、族称及制度覆盖范围不得互推。 进一步讨论可参照 辽代契丹国家、都城、部族与跨境关系研究

\subsection{1122年}
\noindent\eventanchor{LIAO-1122-EVENT-046}{辽南京燕京被金军攻取}\textbf{辽与金。}
辽与金 在 1122年 的材料痕迹是〔辽南京燕京被金军攻取〕，出处为 《辽史》相关本纪及志；《宋史》辽国传、同年本纪；宋使行录、契丹／汉文碑刻与考古材料（据事核）。题名保留的城、路、水道或机构名称，是目前可以落地的空间线索。

前一层压力可由以下线索辨认：女真兴起与辽末多方角逐中既有的联盟、交通与补给压力，为“辽南京燕京被金军攻取”提供了直接的军事或动员背景。 这一处置留下的后续条件是：“辽南京燕京被金军攻取”改变了相关城镇、道路或政权间的军事选择；伤亡、俘获和控制范围仅可按各方材料分别确认。

证据的边界是：《辽史》为元末官修且契丹本方连续文书有限；宋使与碑刻的观察位置不同，战果、族称及制度覆盖范围不得互推。 辽代契丹国家、都城、部族与跨境关系研究 可用于继续检验这一结论。

\subsection{1122年（宋宣和四年、辽德兴1年；公元换算据两国年号对照）}
\noindent\eventanchor{SYD-1122-LIAO-082}{丙子，遼人立燕王淳為帝}\textbf{宋与辽。}
从 《宋史》卷022〈宣和四年〉；《辽史》本纪同年条（互校） 的 1122年（宋宣和四年、辽德兴1年；公元换算据两国年号对照） 条目看，〔丙子，遼人立燕王淳為帝〕使 宋与辽 的一个具体行动进入叙事；题名保留的城、路、水道或机构名称，是目前可以落地的空间线索。

本纪年条提供日期和中枢可见行动；前因不据单条补定。 记录所见的结果则是：可回查同年及后续条目，地方影响仍须另证。 日期相邻不等于因果已经证实。

《宋史》为元末官修，本纪保存宋廷纪年与处置；战果、数字、地方执行及对方动机须以编年、会要、对方史料或地方材料互校。 因此，宋辽边政、使节记录与契丹碑刻研究 只能扩展问题，不能代替未保存的细节。

\subsection{1123年（宋宣和五年、辽德兴2年；公元换算据两国年号对照）}
\noindent\eventanchor{SYD-1123-LIAO-083}{丙戌，遼人張覺以平州來附}\textbf{宋与辽。}
1123年（宋宣和五年、辽德兴2年；公元换算据两国年号对照） 的记载将 宋与辽 与〔丙戌，遼人張覺以平州來附〕相连。此处可确认的是这项被记录的行动；题名保留的城、路、水道或机构名称，是目前可以落地的空间线索。

把本项放回事件链，能够看到的前提为：本纪年条提供日期和中枢可见行动；前因不据单条补定。 而它所打开或收紧的下一步是：可回查同年及后续条目，地方影响仍须另证。

关于可说范围，须保留：《宋史》为元末官修，本纪保存宋廷纪年与处置；战果、数字、地方执行及对方动机须以编年、会要、对方史料或地方材料互校。 进一步讨论可参照 宋辽边政、使节记录与契丹碑刻研究

\subsection{1123年}
\noindent\eventanchor{XIA-1123-EVENT-068}{金在灭辽过程中与西夏建立新的外交关系}\textbf{西夏与金。}
〔金在灭辽过程中与西夏建立新的外交关系〕见于 《宋史》〈夏国传〉及同年本纪；《辽史》《金史》或《元史》相关条；西夏文文书、碑刻与城址材料（据事核） 的 1123年 记录。西夏与金 的选择因此有了可回查的时间位置；题名保留的城、路、水道或机构名称，是目前可以落地的空间线索。

它承接的条件是：“金在灭辽过程中与西夏建立新的外交关系”发生在西夏与金、宋的边境与册封关系的具体压力与机会之中，相关行动者的选择不能脱离此前事件链解释。 随后可追踪的变化为：“金在灭辽过程中与西夏建立新的外交关系”为后续政治、边防或资源配置留下条件；其社会影响与实际覆盖范围仍须以异类材料限定。

证据的边界是：西夏无存世连续国史，汉文叙事多出自对手；西夏文文书的释读、定年和地域代表性须逐项说明。 西夏文文书、河西—宁夏考古与宋夏金蒙关系研究 可用于继续检验这一结论。

\subsection{1124年}
\noindent\eventanchor{LIAO-1124-EVENT-047}{耶律大石率部西行，西辽形成的条件开始出现}\textbf{辽与西辽。}
1124年，《辽史》相关本纪及志；《宋史》辽国传、同年本纪；宋使行录、契丹／汉文碑刻与考古材料（据事核） 所见的〔耶律大石率部西行，西辽形成的条件开始出现〕把 辽与西辽 的处置留在可复核的年序中；材料未将地点细化到城址，不能用中枢所在地替它补足。

前一层压力可由以下线索辨认：“耶律大石率部西行，西辽形成的条件开始出现”发生在女真兴起与辽末多方角逐的具体压力与机会之中，相关行动者的选择不能脱离此前事件链解释。 这一处置留下的后续条件是：“耶律大石率部西行，西辽形成的条件开始出现”为后续政治、边防或资源配置留下条件；其社会影响与实际覆盖范围仍须以异类材料限定。

《辽史》为元末官修且契丹本方连续文书有限；宋使与碑刻的观察位置不同，战果、族称及制度覆盖范围不得互推。 因此，辽代契丹国家、都城、部族与跨境关系研究 只能扩展问题，不能代替未保存的细节。

\subsection{1124年（宋宣和六年、西夏元德6年；公元换算据两国年号对照）}
\noindent\eventanchor{SYD-1124-XIA-051}{夏國、高麗、于闐、羅殿入貢}\textbf{宋与西夏。}
宋与西夏 在 1124年（宋宣和六年、西夏元德6年；公元换算据两国年号对照） 的材料痕迹是〔夏國、高麗、于闐、羅殿入貢〕，出处为 《宋史》卷022〈宣和六年〉；西夏文书、碑刻及《辽史》或《金史》同年条（互校）。材料未将地点细化到城址，不能用中枢所在地替它补足。

本纪记录使节、贡赐或往来，不等于稳定统属关系。 记录所见的结果则是：可与对方编年和交通、文书材料核对其后续落实。 日期相邻不等于因果已经证实。

关于可说范围，须保留：《宋史》为元末官修，本纪保存宋廷纪年与处置；战果、数字、地方执行及对方动机须以编年、会要、对方史料或地方材料互校。 进一步讨论可参照 西夏文书、河西—宁夏考古与宋夏关系研究

\subsection{1125年（辽保大五年、金天会三年、宋宣和七年；干支、月日待核；公元换算据各方本纪年号对照）}
\noindent\eventanchor{LIAO-1125-EVENT-048}{天祚帝被金军俘获，辽朝统治终结}\textbf{辽与金。}
从 《辽史》相关本纪及志；《宋史》辽国传、同年本纪；宋使行录、契丹／汉文碑刻与考古材料（据事核） 的 1125年（辽保大五年、金天会三年、宋宣和七年；干支、月日待核；公元换算据各方本纪年号对照） 条目看，〔天祚帝被金军俘获，辽朝统治终结〕使 辽与金 的一个具体行动进入叙事；材料未将地点细化到城址，不能用中枢所在地替它补足。

把本项放回事件链，能够看到的前提为：“天祚帝被金军俘获，辽朝统治终结”发生在女真兴起与辽末多方角逐的具体压力与机会之中，相关行动者的选择不能脱离此前事件链解释。 而它所打开或收紧的下一步是：“天祚帝被金军俘获，辽朝统治终结”为后续政治、边防或资源配置留下条件；其社会影响与实际覆盖范围仍须以异类材料限定。

证据的边界是：《辽史》为元末官修且契丹本方连续文书有限；宋使与碑刻的观察位置不同，战果、族称及制度覆盖范围不得互推。 辽代契丹国家、都城、部族与跨境关系研究 可用于继续检验这一结论。

\subsection{1125年（宋宣和七年、辽延庆2年；公元换算据两国年号对照）}
\noindent\eventanchor{SYD-1125-LIAO-084}{壬辰，金人以擒遼主，遣李孝和等來告慶}\textbf{宋与辽。}
1125年（宋宣和七年、辽延庆2年；公元换算据两国年号对照） 的记载将 宋与辽 与〔壬辰，金人以擒遼主，遣李孝和等來告慶〕相连。此处可确认的是这项被记录的行动；材料未将地点细化到城址，不能用中枢所在地替它补足。

它承接的条件是：本纪年条提供日期和中枢可见行动；前因不据单条补定。 随后可追踪的变化为：可回查同年及后续条目，地方影响仍须另证。

《宋史》为元末官修，本纪保存宋廷纪年与处置；战果、数字、地方执行及对方动机须以编年、会要、对方史料或地方材料互校。 因此，宋辽边政、使节记录与契丹碑刻研究 只能扩展问题，不能代替未保存的细节。

\subsection{1126年（宋靖康元年、辽延庆3年；公元换算据两国年号对照）}
\noindent\eventanchor{SYD-1126-LIAO-085}{遼故將小䩴䩮攻陷麟州建寧砦，知砦楊震死之}\textbf{宋与辽。}
〔遼故將小䩴䩮攻陷麟州建寧砦，知砦楊震死之〕见于 《宋史》卷023〈靖康元年〉；《辽史》本纪同年条（互校） 的 1126年（宋靖康元年、辽延庆3年；公元换算据两国年号对照） 记录。宋与辽 的选择因此有了可回查的时间位置；题名保留的城、路、水道或机构名称，是目前可以落地的空间线索。

前一层压力可由以下线索辨认：本纪保留灾害或工程的报告地点与朝廷处置；灾情范围需另证。 这一处置留下的后续条件是：可与地方文书、河工和环境材料比较，不将单点记录外推为全域因果。

关于可说范围，须保留：《宋史》为元末官修，本纪保存宋廷纪年与处置；战果、数字、地方执行及对方动机须以编年、会要、对方史料或地方材料互校。 进一步讨论可参照 宋辽边政、使节记录与契丹碑刻研究

\subsection{1126年（宋靖康元年、西夏元德8年；公元换算据两国年号对照）}
\noindent\eventanchor{SYD-1126-XIA-052}{十一月丙寅，夏人陷懷德軍，知軍事劉銓、通判杜翊世死之}\textbf{宋与西夏。}
1126年（宋靖康元年、西夏元德8年；公元换算据两国年号对照），《宋史》卷023〈靖康元年〉；西夏文书、碑刻及《辽史》或《金史》同年条（互校） 所见的〔十一月丙寅，夏人陷懷德軍，知軍事劉銓、通判杜翊世死之〕把 宋与西夏 的处置留在可复核的年序中；材料未将地点细化到城址，不能用中枢所在地替它补足。

本纪年条提供日期和中枢可见行动；前因不据单条补定。 记录所见的结果则是：可回查同年及后续条目，地方影响仍须另证。 日期相邻不等于因果已经证实。

证据的边界是：《宋史》为元末官修，本纪保存宋廷纪年与处置；战果、数字、地方执行及对方动机须以编年、会要、对方史料或地方材料互校。 西夏文书、河西—宁夏考古与宋夏关系研究 可用于继续检验这一结论。

\noindent\eventanchor{SYD-1126-XIA-075}{夏四月戊戌，夏人陷震威城，攝知城事朱昭死之}\textbf{宋与西夏。}
宋与西夏 在 1126年（宋靖康元年、西夏元德8年；公元换算据两国年号对照） 的材料痕迹是〔夏四月戊戌，夏人陷震威城，攝知城事朱昭死之〕，出处为 《宋史》卷023〈靖康元年〉；西夏文书、碑刻及《辽史》或《金史》同年条（互校）。题名保留的城、路、水道或机构名称，是目前可以落地的空间线索。

把本项放回事件链，能够看到的前提为：本纪保留灾害或工程的报告地点与朝廷处置；灾情范围需另证。 而它所打开或收紧的下一步是：可与地方文书、河工和环境材料比较，不将单点记录外推为全域因果。

《宋史》为元末官修，本纪保存宋廷纪年与处置；战果、数字、地方执行及对方动机须以编年、会要、对方史料或地方材料互校。 因此，西夏文书、河西—宁夏考古与宋夏关系研究 只能扩展问题，不能代替未保存的细节。

\subsection{1139年}
\noindent\eventanchor{XIA-1139-EVENT-069}{李仁孝即位}\textbf{西夏。}
从 《宋史》〈夏国传〉及同年本纪；《辽史》《金史》或《元史》相关条；西夏文文书、碑刻与城址材料（据事核） 的 1139年 条目看，〔李仁孝即位〕使 西夏 的一个具体行动进入叙事；材料未将地点细化到城址，不能用中枢所在地替它补足。

它承接的条件是：西夏与金、宋的边境与册封关系中前任统治的结束与宫廷、部族或官僚的权力安排相互交织，促成“李仁孝即位”这一继承节点。 随后可追踪的变化为：继承并不自动消除既有矛盾；“李仁孝即位”后的任官、军权和对外策略需由后续条目追踪。

关于可说范围，须保留：西夏无存世连续国史，汉文叙事多出自对手；西夏文文书的释读、定年和地域代表性须逐项说明。 进一步讨论可参照 西夏文文书、河西—宁夏考古与宋夏金蒙关系研究

\subsection{1144年}
\noindent\eventanchor{XIA-1144-EVENT-070}{西夏与金以册封、贡赐和边境安排维持和平}\textbf{西夏与金。}
1144年 的记载将 西夏与金 与〔西夏与金以册封、贡赐和边境安排维持和平〕相连。此处可确认的是这项被记录的行动；材料未将地点细化到城址，不能用中枢所在地替它补足。

前一层压力可由以下线索辨认：西夏与金、宋的边境与册封关系里的军事消耗、边境供给与使节往来构成谈判条件，促成“西夏与金以册封、贡赐和边境安排维持和平”所涉的协议或名义关系。 这一处置留下的后续条件是：“西夏与金以册封、贡赐和边境安排维持和平”重排了贡赐、使节或停战的制度接口；条款是否执行及边地实际控制不能由盟约文字直接推定。

证据的边界是：西夏无存世连续国史，汉文叙事多出自对手；西夏文文书的释读、定年和地域代表性须逐项说明。 西夏文文书、河西—宁夏考古与宋夏金蒙关系研究 可用于继续检验这一结论。

\subsection{1193年}
\noindent\eventanchor{XIA-1193-EVENT-071}{李纯佑即位}\textbf{西夏。}
〔李纯佑即位〕见于 《宋史》〈夏国传〉及同年本纪；《辽史》《金史》或《元史》相关条；西夏文文书、碑刻与城址材料（据事核） 的 1193年 记录。西夏 的选择因此有了可回查的时间位置；材料未将地点细化到城址，不能用中枢所在地替它补足。

西夏与金、宋的边境与册封关系中前任统治的结束与宫廷、部族或官僚的权力安排相互交织，促成“李纯佑即位”这一继承节点。 记录所见的结果则是：继承并不自动消除既有矛盾；“李纯佑即位”后的任官、军权和对外策略需由后续条目追踪。 日期相邻不等于因果已经证实。

西夏无存世连续国史，汉文叙事多出自对手；西夏文文书的释读、定年和地域代表性须逐项说明。 因此，西夏文文书、河西—宁夏考古与宋夏金蒙关系研究 只能扩展问题，不能代替未保存的细节。

\subsection{1206年}
\noindent\eventanchor{XIA-1206-EVENT-072}{成吉思汗统一蒙古诸部，西夏北方环境突变}\textbf{西夏与蒙古。}
1206年，《宋史》〈夏国传〉及同年本纪；《辽史》《金史》或《元史》相关条；西夏文文书、碑刻与城址材料（据事核） 所见的〔成吉思汗统一蒙古诸部，西夏北方环境突变〕把 西夏与蒙古 的处置留在可复核的年序中；材料未将地点细化到城址，不能用中枢所在地替它补足。

把本项放回事件链，能够看到的前提为：“成吉思汗统一蒙古诸部，西夏北方环境突变”发生在蒙古扩张下的西夏存亡的具体压力与机会之中，相关行动者的选择不能脱离此前事件链解释。 而它所打开或收紧的下一步是：“成吉思汗统一蒙古诸部，西夏北方环境突变”为后续政治、边防或资源配置留下条件；其社会影响与实际覆盖范围仍须以异类材料限定。

关于可说范围，须保留：西夏无存世连续国史，汉文叙事多出自对手；西夏文文书的释读、定年和地域代表性须逐项说明。 进一步讨论可参照 西夏文文书、河西—宁夏考古与宋夏金蒙关系研究

\subsection{1209年}
\noindent\eventanchor{XIA-1209-EVENT-073}{蒙古军攻西夏，西夏接受和议与联姻安排}\textbf{西夏与蒙古。}
西夏与蒙古 在 1209年 的材料痕迹是〔蒙古军攻西夏，西夏接受和议与联姻安排〕，出处为 《宋史》〈夏国传〉及同年本纪；《辽史》《金史》或《元史》相关条；西夏文文书、碑刻与城址材料（据事核）。材料未将地点细化到城址，不能用中枢所在地替它补足。

它承接的条件是：蒙古扩张下的西夏存亡里的军事消耗、边境供给与使节往来构成谈判条件，促成“蒙古军攻西夏，西夏接受和议与联姻安排”所涉的协议或名义关系。 随后可追踪的变化为：“蒙古军攻西夏，西夏接受和议与联姻安排”重排了贡赐、使节或停战的制度接口；条款是否执行及边地实际控制不能由盟约文字直接推定。

证据的边界是：西夏无存世连续国史，汉文叙事多出自对手；西夏文文书的释读、定年和地域代表性须逐项说明。 西夏文文书、河西—宁夏考古与宋夏金蒙关系研究 可用于继续检验这一结论。

\subsection{1217年}
\noindent\eventanchor{XIA-1217-EVENT-074}{西夏对蒙古军事要求的配合出现紧张}\textbf{西夏与蒙古。}
从 《宋史》〈夏国传〉及同年本纪；《辽史》《金史》或《元史》相关条；西夏文文书、碑刻与城址材料（据事核） 的 1217年 条目看，〔西夏对蒙古军事要求的配合出现紧张〕使 西夏与蒙古 的一个具体行动进入叙事；材料未将地点细化到城址，不能用中枢所在地替它补足。

前一层压力可由以下线索辨认：“西夏对蒙古军事要求的配合出现紧张”发生在蒙古扩张下的西夏存亡的具体压力与机会之中，相关行动者的选择不能脱离此前事件链解释。 这一处置留下的后续条件是：“西夏对蒙古军事要求的配合出现紧张”为后续政治、边防或资源配置留下条件；其社会影响与实际覆盖范围仍须以异类材料限定。

西夏无存世连续国史，汉文叙事多出自对手；西夏文文书的释读、定年和地域代表性须逐项说明。 因此，西夏文文书、河西—宁夏考古与宋夏金蒙关系研究 只能扩展问题，不能代替未保存的细节。

\subsection{1223年}
\noindent\eventanchor{XIA-1223-EVENT-075}{李德旺即位，西夏王室更替发生在蒙古压力下}\textbf{西夏。}
1223年 的记载将 西夏 与〔李德旺即位，西夏王室更替发生在蒙古压力下〕相连。此处可确认的是这项被记录的行动；材料未将地点细化到城址，不能用中枢所在地替它补足。

蒙古扩张下的西夏存亡中前任统治的结束与宫廷、部族或官僚的权力安排相互交织，促成“李德旺即位，西夏王室更替发生在蒙古压力下”这一继承节点。 记录所见的结果则是：继承并不自动消除既有矛盾；“李德旺即位，西夏王室更替发生在蒙古压力下”后的任官、军权和对外策略需由后续条目追踪。 日期相邻不等于因果已经证实。

关于可说范围，须保留：西夏无存世连续国史，汉文叙事多出自对手；西夏文文书的释读、定年和地域代表性须逐项说明。 进一步讨论可参照 西夏文文书、河西—宁夏考古与宋夏金蒙关系研究

\subsection{1226年}
\noindent\eventanchor{XIA-1226-EVENT-076}{蒙古军大举攻入西夏}\textbf{西夏与蒙古。}
〔蒙古军大举攻入西夏〕见于 《宋史》〈夏国传〉及同年本纪；《辽史》《金史》或《元史》相关条；西夏文文书、碑刻与城址材料（据事核） 的 1226年 记录。西夏与蒙古 的选择因此有了可回查的时间位置；材料未将地点细化到城址，不能用中枢所在地替它补足。

把本项放回事件链，能够看到的前提为：蒙古扩张下的西夏存亡中既有的联盟、交通与补给压力，为“蒙古军大举攻入西夏”提供了直接的军事或动员背景。 而它所打开或收紧的下一步是：“蒙古军大举攻入西夏”改变了相关城镇、道路或政权间的军事选择；伤亡、俘获和控制范围仅可按各方材料分别确认。

证据的边界是：西夏无存世连续国史，汉文叙事多出自对手；西夏文文书的释读、定年和地域代表性须逐项说明。 西夏文文书、河西—宁夏考古与宋夏金蒙关系研究 可用于继续检验这一结论。

\subsection{1227年}
\noindent\eventanchor{XIA-1227-EVENT-077}{西夏灭亡，蒙古接收河西和宁夏区域}\textbf{西夏与蒙古。}
1227年，《宋史》〈夏国传〉及同年本纪；《辽史》《金史》或《元史》相关条；西夏文文书、碑刻与城址材料（据事核） 所见的〔西夏灭亡，蒙古接收河西和宁夏区域〕把 西夏与蒙古 的处置留在可复核的年序中；题名保留的城、路、水道或机构名称，是目前可以落地的空间线索。

它承接的条件是：蒙古扩张下的西夏存亡中既有的联盟、交通与补给压力，为“西夏灭亡，蒙古接收河西和宁夏区域”提供了直接的军事或动员背景。 随后可追踪的变化为：“西夏灭亡，蒙古接收河西和宁夏区域”改变了相关城镇、道路或政权间的军事选择；伤亡、俘获和控制范围仅可按各方材料分别确认。

西夏无存世连续国史，汉文叙事多出自对手；西夏文文书的释读、定年和地域代表性须逐项说明。 因此，西夏文文书、河西—宁夏考古与宋夏金蒙关系研究 只能扩展问题，不能代替未保存的细节。
